% =====================================================================
%  GENERATED; DO NOT EDIT BY HAND.  Built by
%    python3 ebook-method/tydskland-efter-hegel/tr_parts.py build --write
%  from ebook-method/tydskland-efter-hegel/tr_template.tex and the translators' parts
%  (texts/hoeffding/tydskland-efter-hegel/.parts/tr/*.texfrag).  To change the English, edit the part
%  (or this template) and rebuild.
%
%  Harald Høffding, PHILOSOPHIEN I TYDSKLAND EFTER HEGEL (Kjøbenhavn: Gyldendal, 1872): English
%  translation, "Philosophy in Germany after Hegel".  Translated 2026-10-09/10 from transcription.tex
%  (the text of record, collated to OPEN 0; see its header) by translator agents working from one brief
%  and one glossary (ebook-method/tydskland-efter-hegel/TRANSLATOR-BRIEF.md, GLOSSARY.md), each part checked
%  by script against the Danish (tr_parts.py build): every page marker once and in order, footnotes,
%  counter resets and paragraph breaks per printed page, the Greek character for character.
%  The margin numbers [N] are the pages of the 1872 printing.
%  TRANSCRIPTION REPAIRED WHILE TRANSLATING (TRANSLATION-PLAYBOOK §6c): the translators' 200 odd spots were
%  read at the image and the transcription mended through ebook-method/tydskland-efter-hegel/patches.py (POST2,
%  POSTMIS: 113 faults mended, 60 misprints marked, a stray copy of p. 139 removed from p. 143); the English follows
%  the repaired text, with a "% print:" line at the top of each page that has a newly marked misprint.
%
%  CONVENTIONS.  Quotations Høffding gives in German, Latin or Greek stay in the original; what he gives
%  in Danish (including his own Danish renderings of German authors) is translated from his Danish.
%  Titles of books and journals as printed.  Letterspacing (\emph in the Danish) -> \emph.  Footnotes at
%  the same anchor, marked *), **) per printed page as in the print.  Misprints the transcription marks
%  PRINTED AS IS are translated in their evident sense, with a "% print:" comment at the site.
%  Terms: GLOSSARY.md (also in the book's RESUME-NOTES).  American spelling.
% =====================================================================
\documentclass[12pt,a4paper,openany]{book}
\usepackage[utf8]{inputenc}
\usepackage[T1]{fontenc}
\usepackage{textalpha}
\usepackage{libertinus}
\usepackage{libertinust1math}
\usepackage[english]{babel}
\usepackage[protrusion=true,expansion=true]{microtype}
\usepackage{setspace}
\usepackage[margin=1.2in]{geometry}
\usepackage{fancyhdr}
\setlength{\headheight}{28pt}
\usepackage{hyperref}
\hypersetup{pdftitle={Philosophy in Germany after Hegel (1872)}, pdfauthor={Harald Høffding}, hidelinks}
\pagestyle{fancy}
\fancyhf{}
\fancyhead[LE,RO]{\thepage}
\fancyhead[RE]{\textit{Harald Høffding: Philosophy in Germany after Hegel}}
\fancyhead[LO]{\textit{1872}}
\fancypagestyle{plain}{\fancyhf{}\fancyhead[LE,RO]{\thepage}}
\setstretch{1.3}
\usepackage{marginnote}
\usepackage{graphicx}
\DeclareUnicodeCharacter{0254}{\reflectbox{c}}   % ɔ (should not occur in the English)
\renewcommand{\thefootnote}{%
  \ifcase\value{footnote}\or *)\or **)\or ***)\else \arabic{footnote})\fi}
\newcommand{\tocline}[2]{\noindent #1\dotfill\ #2\par}
\newcommand{\opage}[1]{\marginnote{\footnotesize\textit{[#1]}}}
\newcommand{\chap}[1]{\clearpage\addcontentsline{toc}{chapter}{#1}\vspace*{2em}\begin{center}{\large #1}\end{center}\medskip}
\newcommand{\chaphere}[1]{\addcontentsline{toc}{chapter}{#1}\begin{center}{\large\bfseries #1}\end{center}}
\pdfstringdefDisableCommands{\renewcommand{\footnote}[2][]{}\renewcommand{\opage}[1]{}}

\begin{document}

\begin{titlepage}
\begin{center}
\vspace*{3em}
{\LARGE\bfseries Philosophy in Germany}\\[0.8em]
{\large \emph{after Hegel.}}\\[2em]
By\\[1em]
{\large\bfseries Dr. H. Høffding.}\\[2em]
{\small Translated from the Danish}\\[0.4em]
\vfill
{\small\bfseries Copenhagen.}\\[0.4em]
\emph{Published by the Gyldendal Bookshop (F. Hegel).}\\[0.4em]
{\small \emph{I. Cohen's Printing Office.}}\\[0.4em]
{\small 1872.}
\end{center}
\end{titlepage}
\thispagestyle{empty}
\vspace*{4em}
\begin{center}\small
``Auch wir können noch an \emph{wirkende}, aber wir können\\
nicht mehr an \emph{hexende} Ideen glauben.''
\end{center}
\begin{flushright}\footnotesize
H. Lotze: Geschichte der Aesthetik in Deutschland p. 415.
\end{flushright}
\clearpage
\thispagestyle{empty}
\begin{center}{\large\bfseries \emph{Contents.}}\end{center}
\begin{center}\rule{3em}{0.4pt}\end{center}
\hfill{\small Page}\par
\tocline{\textbf{Chapter One:} \emph{Introduction}}{1}
{\small 1) Historical Orientation p. 1. 2) Characterization of Hegel's System p. 7. 3) Main Problems of Post-Hegelian Philosophy p. 14.}\par\smallskip
\tocline{\textbf{Chapter Two:} \emph{The Critical Consciousness}}{17}
{\small 1) The Controversy over Individual Immortality p. 17. 2) Strauss p. 29. 3) Bruno Bauer p. 40. 4) Ruge p. 44. 5) Feuerbach p. 49. 6) Max Stirner p. 88. 7) Das Verstandesthum und das Individuum p. 92.}\par\smallskip
\tocline{\textbf{Chapter Three:} \emph{Modern Materialism}}{94}
{\small 1) Empirical Materialism (Vogt, Moleschott, Büchner) p. 99. 2) Ethical-Speculative Materialism (Czolbe) p. 109.}\par\smallskip
\tocline{\textbf{Chapter Four:} \emph{Speculative Theism}}{117}
{\small 1) Schelling's Later Doctrine p. 121. 2) The Younger Fichte p. 149. 3) Weisse p. 181.}\par\smallskip
\tocline{\textbf{Chapter Five:} \emph{The Idealism of Reality}}{220}
{\small 1) Trendelenburg p. 222. 2) Fechner p. 246. 3) Hartmann p. 268. 4) Lotze p. 286.}\par\smallskip
\tocline{\textbf{Chapter Six:} \emph{Characterization and Results}}{319}
\begin{center}\rule{3em}{0.4pt}\end{center}
\clearpage


% --- p. 1 ---
\setcounter{footnote}{0}%
\chap{CHAPTER ONE}
\begin{center}\rule{2em}{0.4pt}\end{center}
\opage{1}


1. It might seem to be a task as thankless as it is
difficult to give an account of post-Hegelian philosophy
in Germany. For this period does not carry the
interest that is called forth by great epoch-making
advances and breakthroughs; it is essentially a time of
epigones, following the magnificent philosophical
production of the preceding half-century. And on the
other hand, post-Hegelian philosophizing lies so near
to us that a truly historical and objective account
must be difficult, if not impossible.

These misgivings vanish, however, on the consideration
that it is necessary for everyone who strives
after a philosophical world-view to arrive at
clarity as to which decisive problems are stirring
in the thought of the age, and in which direction their treatment
and solution are going. However deep and inward the
enrichment that is won by constantly returning to
the writings of the great masters, just so important and indispensable
is it, on the other hand, to get one's bearings with
regard to the thinkers of one's own time, even if they should not
be counted in the class of the masters. And when it is
this interest that spurs the inquiry, then
the purely objective and historical point of view also falls away
% --- p. 2 ---
\setcounter{footnote}{0}%
\opage{2}and gives place to a more properly philosophical and critical
mode of consideration.

The point of departure is marked by the Hegelian philosophy,
which by its universal character reached so deeply
into the whole spiritual life of the time, and was at the same time itself the result
of an idealistic direction that runs through the whole of
modern philosophy. Yet Hegelianism is only our
relative point of departure. It does not form, as has often
been thought, a nodal point in the history of modern philosophy:
that already lies in its being the most decided
idealism. The nodal point must rather be sought in the
philosopher in whom the idealistic and the realistic
lines of development intersect and for a time become one
on the basis of a principle lying deeper. This philosopher,
the central thinker of modern times, is Immanuel
Kant. This can be shown in brief with regard both
to his course of development, his principle and his method.

After Kant had long been occupied partly with studies in natural science
and partly with the study of the Leibnitz-Wolfian
idealistic philosophy, it was, by his own
declaration, David Hume, the representative of thoroughgoing empiricism,
who woke him from his ``dogmatic slumber.''
And that he now saw his task as finding the point of unity of the
two opposed directions is seen, among other things, from the fact that in
his famous work, ``Kritik der reinen Vernunft,'' after
every main section he tries to show how his doctrine
lets both Leibnitz and Locke come into their own. ---
What, then, is the philosophical principle by which Kant
reaches this superior standpoint? This principle
must be sought in the very way in which he conceives the
essence of philosophical knowledge. He does not, like
realism, take his standpoint in the external, in the objects,
which were supposed to act on man from without and stamp knowledge
into him, and just as little does he, idealistically, let
the human spirit develop knowledge out of its own interior.
Thought, to use Bacon's image, is neither the ant,
% --- p. 3 ---
\setcounter{footnote}{0}%
\opage{3}which merely gathers and heaps up the material, nor the spider,
which spins its web out of itself, but the bee, which
extracts the juice from the flowers and in its own peculiar way
digests and transforms it. It is the inner reciprocal relation of thought
to things in knowledge that Kant directs his eye upon;
he seeks to determine the essence of the two factors and the laws
of their interaction. The knowledge of these laws
Kant calls transcendental knowledge. This
lies at the ground of every other knowledge, whether
this is of a metaphysical or an empirical nature.
--- From this principle follows a corresponding method.
This is at once inductive and deductive. It proceeds
from the fact of knowledge and the elements contained
in it, and to that extent the inquiry becomes inductive
and analytic; but since knowledge is not an
external fact, but inherent in itself, its own principle,
on which all its forms and all its content depend,
the development becomes synthetic and deductive. This
unity shows itself in the combination of the critical
and the dialectical in Kant's investigations. Kant distinguishes
critically, but he does so in order to let the oppositions
come forward and come into their own, so as
through this to reach once more the underlying unity.
It is this that has been called the teasing element in Kant, that
one is not allowed to rest content with the sharp
distinctions, but is always pointed beyond them anew.
The dialectic breaks forth precisely in the carrying through of the critique.


Yet Kant left tasks enough for his successors.
He had indicated in principle the goal that was to
be reached, but in reality he had not subdued
the contending powers. In particular, by the distinction
between the subjective forms of knowledge and the thing in
and for itself, he had opened the way to a dualism that threatened
the very great goal to which he had so energetically pointed,
the inner unity of idealism and realism. One
% --- p. 4 ---
\setcounter{footnote}{0}%
\opage{4}saw, moreover, soon enough that this distinction rested on a
self-contradiction. For on the presupposition that our forms of knowledge
have a merely subjective significance, the assumption
of a thing in and for itself is impossible: no determinations of thought
--- and so not the concepts of being and
cause either --- can, consistently, be applied to it;
it can neither \emph{be} nor be the \emph{cause} of the sensation
of an external object. Without the concept of ``Ding
an sich'' one could not, as Jakobi rightly remarked,
get into Kant's system, and with it one could
not remain there. It was thus a necessity lying in the nature
of the case that the realistic element in Kant's doctrine
should first be separated out. To overcome the dualism
one had to hold to what is clear and certain in itself,
thought itself and the ultimate source of its forms, the original
unity of self-consciousness, what Kant called
``the transcendental apperception of self-consciousness.'' Over against
this rich and luminous world of forms stood ``das Ding an
sich'' as an indeterminate, fog-like mass, whose connection
with the former could not be shown. Fichte constructed
the world out of subjectivity and set it as the latter's
task constantly to overcome anew the existence thus arisen.
The objective thus still remains with him
as a negative principle, which in eternal unrest
is posited in order to be sublated again. The limits of this subjective idealism
Schelling oversteps. All science rests,
he says, on the presupposition that the absolutely
ideal is also in itself the absolutely real. The Absolute
is pure identity, where no differences and
oppositions any longer hold. This is absolute
idealism: the ideal is all, because all is Idea. The
Schellingian philosophy of nature would prove the identity of nature
with the world of Ideas, and sets as its goal ``the perfect
presentation of the intellectual world in the laws and forms of the phenomenal
world, and so, conversely, the
comprehension of these laws and forms from out of the
% --- p. 5 ---
% print: af Sein (misprint: af for auf; translation odd spot)
\setcounter{footnote}{0}%
\opage{5}intellectual world.'' Schelling's knowledge of nature thus aims
solely at the ideal significance of phenomena, whereas the
task of explaining the real connection of things did not
exist for him. He saw a new period about
to begin, in which humanity was to be led from the earlier imperfect
knowledge to the beholding of the Absolute,
the one true reality. This prophetic beholding
Hegel brought under the form of the concept, in that by his
dialectical method he developed the content of reality in
its unity with the essence of thought (of the Idea). Hereby even
the tinge of subjectivism that still
clung to Schelling's constructions vanished, and the idealistic
consequences of Kant's doctrine were developed to their
highest point.

While here, then, the realistic element in Kant's
doctrine was gradually quite eliminated, other
thinkers nonetheless took their point of departure precisely in this element.
Already as a young student and Fichte's hearer in Jena,
Herbart worked out a critique of Schelling's first
writings, in which he disputes the legitimacy of transferring
the form of thought, of knowing, the self-developing, active unity,
to the content and object of knowing. For the content of knowing
is absolute being, but this is incompatible
with the concept of positing and producing itself; it
is one with absolute rest and stillness, whereas the I
is an eternal vortex streaming forward and back; therefore
the two must be held well apart from each other. In this critique
Herbart's later doctrine is already contained. Knowing is to
withdraw into itself, so that the content can freely
present itself. The content stems from experience, and
thought has only to examine whether the concepts of experience
are free of contradictions. Now since this is far from being the case,
the given must be so treated that through it one
can penetrate to the true essence of the matter: ``so
viel Schein, so viel Hindeutung af Sein.'' Thought thus
gets here only a formal significance; by its
% --- p. 6 ---
\setcounter{footnote}{0}%
\opage{6}critique it is to separate out the absolute position, the true being,
that is unconsciously contained in sensuous sensation.
--- Schopenhauer agrees with Herbart insofar as
in his critique of the Kantian philosophy he declares that
Kant's error lay in the way in which he derived ``das
Ding an sich,'' not in his acknowledging it
at all. Hereby he came into so decided an opposition
to speculative idealism that it even degenerated
into downright fanaticism, in that he thought he could only
derive what he called speculative philosophy's
misunderstanding of Kant from a disregard of truth
for personal and egoistic ends. The thing in and
for itself, however, does not remain with Schopenhauer as
an indeterminate something; on the contrary, this foundation of all
phenomena, of the whole of nature, is for him one with
what we find in our own interior as will. But the will
is different from and absolutely independent of knowledge,
which is only a phenomenon of the brain, a fleeting illusion,
whereas the will is the only original and real.
``Das Ding an sich'' has thus here become one and all,
just as, conversely, in Hegel thought, the concept, was one and
all. When Schopenhauer says that he has carried through
what Kant had left unfinished, Hegel could from
his standpoint say the same.

We see, then, that after Kant the ways part again,
after they had for a short time united in one nodal point.
But the idealistic direction, both on account of
its own essence and of the character of the age, became
the all-dominating one in the generation following Kant.
The beginning of the century was idealistic and romantic;
people reveled in the great ideas that had dawned
upon consciousness during the struggle of modern times
against the powers of the past. And speculative idealism
commanded precisely the forms in which the interests thus awakened
could find their expression and come to clarity
about themselves. The great significance of Hegelianism for the history of culture
% --- p. 7 ---
\setcounter{footnote}{0}%
\opage{7}rested precisely on its striving to lead
all essential oppositions back to a comprehensive,
all-reconciling unity. Only after the dissolution of the Hegelian system
do Herbart and Schopenhauer acquire a more far-reaching
significance, so that, in particular around the
first of these men, a firmly closed school forms itself
(whose organ is the ``Zeitschrift für exakte Philosophie'').
But this does not prevent their place in the history of philosophy
from having to be assigned them in a series that develops out of Kant's doctrine
in parallel with speculative idealism.

The philosophical phenomena with which we shall be occupied
in this work stand in a particular relation to
both of these series. The development within the period
we shall treat (1830---1870) is conditioned from the first by
its relation to Hegelian speculation; but during the latter's
critique and dissolution the antipodes of speculation
(Herbart and Schopenhauer) naturally make their influence felt,
just as do the other elements that it had pushed
back, and that precisely now develop with astonishing
force (exact empirical science). And amid all
this one constantly seeks back to Kant, as the one from
whom the whole movement proceeded, in order to get one's bearings
in him with regard to all the important main problems.

2. After this general historical consideration,
by which we have determined the position of our period within
the larger division of the history of philosophy to which
it belongs, we must now characterize more closely our particular
point of departure, the Hegelian philosophy, whose dominion
over minds was at its height precisely at the death of its founder (1831).

Hegel from the very outset set his view in opposition
to two directions prevailing at the beginning of the century.
According to the one, man could
come into relation with the infinite only through presentiment or immediate feeling.
But in this way a fog is cast
over the life and the manifoldness of thought and existence,
% --- p. 8 ---
\setcounter{footnote}{0}%
\opage{8}and one contents oneself with ``the indeterminate enjoyment
of the indeterminate Deity.'' The other direction does
indeed maintain the significance of thought, but conceives the Absolute as
immediate identity, in which all oppositions have
vanished, and which therefore can only be grasped in an intellectual
intuition with artistic genius. But this too
is unscientific. The Absolute cannot be presented
or grasped immediately (``wie aus der Pistole geschossen'');
as undifferentiated darkness it would contain
no explanation of existence. Immediate
feeling and intuition must be raised to real, speculative
knowledge by a scientific method and by virtue
of a scientific principle.

The very demand for a method implies that the
immediate is not left in peace, but is subjected to a test,
submitted to reflection. The negating
force of reflection penetrates, distinguishing, into the immediate unity;
the moments are set against one another; there arises strife and
opposition instead of peace and repose. But precisely through this
the true essence of the unity comes to light. Negation
is not only the dissolving but also the uniting
power. The individual moments set loose do not acquire immediate
independence; that would merely have posited an
immediate multiplicity instead of an immediate unity.
But the individual moments exclude and negate one another,
and through this system of negations a higher form of the
original unity is reached, won through the completion of reflection.
The method is thus dialectical, i.e., it
lets the content itself develop in all its fullness through
contradiction, by the driving-forward force of negation.

In applying negation, thought adds nothing
of its own. Negation is immanent, is contained in
the essence itself. The true is not, as Spinoza would have it,
the absolutely positive substance. Spinoza is right that
every determination is a negation, but he overlooks
that every negation is at the same time a determination, and that
% --- p. 9 ---
\setcounter{footnote}{0}%
\opage{9}only through contradiction is the concrete, real
content reached. The substantial unity must, in order to be
real, be articulated, and the articulation takes place by virtue of
the immanent negation that lets the differences come forth
out of the unity. Hereby substance is developed to freedom,
to subjectivity.

From this it follows that in the dialectical method
the movement of thought coincides with that of the content. The absolute
subjectivity develops itself, comes to itself
through the speculative development of concepts. In
dialectical speculation, then, spirit has regained the
lost repose in the infinite to which skepticism and criticism
had put an end, and which the wings of the philosophy of feeling and intuition
had not been strong enough
to reach again. Thought now grasps with freedom and consciousness
what the fathers gave themselves up to with immediate faith;
this advance ``the world-spirit has succeeded''
in making now.

Concrete existence is as such logical
existence. For it is determinateness, quality, that
marks concrete existence; determinateness is the relation
to, the likeness with, itself; but this identity
is abstraction, hence thought. Conversely, the pure concept,
which relates only to itself, is immediate identity,
i.e., being. Being is the entirely abstract, immediate
relation to itself, the abstract moment of the concept,
which precisely by its absolutely abstract character is subject
to negation and thereby opens the dialectical
movement. By this movement is gained the further
development of what already lay in itself (an sich) in the point of departure,
the original unity of thinking and being.
For Hegel, therefore, the content of logic is the true
reality. ``Only in its concept,'' it says in the introduction
to the Logic, ``does something have actuality; insofar as
it is different from its concept, it ceases to be
actual and is a non-being (ein Nichtiges); to this
% --- p. 10 ---
% print: Aandige et (misprint: et for er)
\setcounter{footnote}{0}%
\opage{10}non-being belong palpability and sensuous
being-outside-itself.'' What we do not come upon, then, in
the logical-dialectical deduction is only slag
that does not belong to the rich ore of reality. --- Immediate
being, by its absolutely abstract character, passes over
into its opposite, and so we are led constantly
from one moment to another. But the externality
between these moments vanishes when it is seen that
the determination of being by its other is properly its determination
by itself, that the negation is essentially an
inner one. Essence is being at once sublated and affirmed;
it has its opposite in seeming, but
itself posits the seeming in order through it to assert itself as
essence, as the ground in existence, as the inner in the
outer. Absolute actuality is the unity of these
oppositions. And this actuality is the Absolute
itself, which in its inner necessity, which becomes transparent
to thought, opens itself in a kingdom of freedom and inwardness
and thereby proves itself to be subjectivity,
to be spirit. From this highest point of view it holds,
then: only the spiritual is the actual. Science
is the reality of spirit, the kingdom it builds itself in its
own element.

Logic, according to this, ought properly to be the only
science. What it could not develop could,
as a ``Nichtiges,'' not become an object of science at all.
Yet Hegel's urge for knowledge is here stronger
than his consistency; it could not content itself
with the logical shadow-world. But what else can
there be besides absolute actuality, and how
can the transition to it be made from logic?

We stand here at the Achilles' heel of the Hegelian system,
the transition from logic to the philosophy of nature. The difficulty
shows itself in the double turn
Hegel himself gives the matter. At the close of the Logic it is
said that the Idea, which is absolutely sure of itself and
% --- p. 11 ---
\setcounter{footnote}{0}%
\opage{11}rests in itself, freely releases itself and transposes itself into the form of externality, as ``äusserliche Idee,'' in order then, through the ascending forms of nature and spirit, to
raise itself again to the inwardness of the concept. But in the philosophy of nature the transition is described as a falling away of the Idea from
itself. Nature is the sphere of pure externality and
as such unfit to hold fast the concept, which is
fulfilled in nature only quite abstractly. The concept is in
nature, but is not nature itself, whereas in logic
concept and reality were one. The metamorphoses
through which the Idea fights its way step by step
to freedom through this its state of estrangement therefore
concern not nature but only the concept in
itself, which is the true interior of natural things. The dialectical
process of the concept and the real movement of nature have nothing
to do with each other.

A similar relation we find in the philosophy of spirit:
an opposition between the dialectical development of the concept of spirit
and the real life of spirit. It is a main Hegelian
proposition that what is rational is also actual,
what is actual also rational. The Idea is
not so powerless that it only ought to be, and is not
actual. Hegel himself complains that these propositions
have been misunderstood and urges that their explanation be sought
in his Logic. Now according to the fundamental thought of that work,
as we have shown, the Idea itself, the fully developed concept,
is the true, the only actuality, whereas the contingent,
the mere phenomenon, has no claim to the name of
actuality. It is clear, then, that ``actuality''
here means the ideal, the rational actuality; but
it is then also clear that Hegel overlooks the essential
problem that must present itself here, namely the question
of the relation in which this ideal actuality
stands to external, factual reality. The proposition of the
unity of reason and actuality is therefore ambiguous,
but precisely in its ambiguity characteristic of Hegel and
% --- p. 12 ---
\setcounter{footnote}{0}%
\opage{12}his age. In particular it contains the explanation of
how he, who in his youth professed
thoroughly radical views in religious as well as in
political respects, could in his mature manhood, without denying
his formal principle, appear as the champion of orthodox
dogma and conservative politics,
--- and how, again, the so-called left side of
his school, by consistently holding fast to this same
principle, could arrive at entirely negative and revolutionary
results both in the religious and in the social
domain.

The Idea of the good is determined in the Logic as a merely
subjective demand or ought, which constantly has reality
outside itself and thereby always hinders its own
absolute realization, for which reason it must find its fulfillment
in the Idea of the true, by virtue of which the concept
and reality are one, so that the Idea thus does not need
subjective activity in order to attain reality.
The individual subject as such then vanishes
before the eternal reality of the absolute Idea. In consequence
of this, the philosophy of right then teaches that ethical
spirit has reality in and for itself and does not depend on
individuals. The ethical is in itself the universal, which does indeed
need the individual's consciousness of it and the subjective passion
thereby awakened, but in such a way that the individual
is only one of the vanishing moments that the Idea in
its development posits and sublates again as single forms
of its reality. Ethical striving, free
personal life, historical development, thus acquire
no true significance. The essence of the Idea in its absolute
autarky is and remains the only reality. Subjective
spirit is wholly absorbed in objective spirit,
i.e., in the state and the communities embraced by it, which
are the expression of the ethical Idea.

Yet above objective spirit there rises still
a highest sphere, absolute spirit, which in itself
% --- p. 13 ---
\setcounter{footnote}{0}%
\opage{13}comprehends all earlier stages. Here the highest
reconciliation takes place through art, religion and science,
and we meet here, in the relation between the two last powers,
a peculiar application of the proposition of the unity of the rational
and the actual. The religious standpoint here stands
for actuality. For speculative
knowledge it is indeed only a finite standpoint, since it moves
in the immediate forms of representation and feeling, encumbered with
the individual; but yet in the matter itself there is to be
no opposition. What was formerly
believed as mystery is now to be revealed to
spirit in the universal and necessary form of thought. And
under this change of form the content is preserved undistorted.
For both philosophy and religion have the infinite
as their content, only that in religion, which precisely in this
points beyond itself, it appears under finite
forms. These are the principles of the speculative reconciliation
of science and church doctrine. ---

With inner necessity the account of the fundamental thoughts of the
Hegelian system became at the same time a critique, in that
the contradiction between the goal that speculative idealism
set itself and the result it really reached came to light, at the
decisive points, in involuntary
admissions. The contradiction concerns the very fundamental problem
of the relation between reason and reality. But yet
it must not be overlooked what great merits speculative
idealism has earned with respect to this same problem.
Kant had been led by his sharp critique to
a separation of the two elements, so that thought
with its forms stood on the one side, things with
their real essence on the other; thought
thereby became purely formal and subjective, the essence of things
the remote, the formless and indeterminate. Speculative
idealism, by immersing itself with enthusiasm and energy
in the essence of thought, has demonstrated that this ideally includes
the essence of objects within itself; seen from the ideal side,
% --- p. 14 ---
% print: Forskjellem (misprint: for Forskjellen)
\setcounter{footnote}{0}%
\opage{14}the opposition between the subjective and the objective is
sublated. Thought bears within itself the whole wealth of existence
as its possible content. The great error of speculative
idealism was now this: to confuse possibility with
actuality, logical actuality with real actuality.
Therefore it thought it could reach the manifoldness of concrete
existence by a purely logical path. When, therefore, at
the dissolution of speculative philosophy, the problem of
the relation between reason and reality presents itself
anew, reason here means not merely the subjective
form of thought, but the whole manifold and content-rich
world of self-consciousness and the Idea, and reality not
the formless ``thing in and for itself,'' but existence revealed in its concrete
and individual forms. Speculation
had reduced the difference between these elements
to a mere difference of form, which can be dialectically posited
in order to be sublated again; real actuality became
only the imperfect form of ideal actuality.
In an age whose whole direction was idealistic,
this magnificent attempt to carry through the
ideal as the only true reality was bound to awaken enthusiasm
and adherence; but when the interest of reality
awakens, the result must prove to rest on an illusion.

3. The question of \emph{the relation between reason
and reality} can be taken from two different
sides. Reason, in opposition to reality, means
the universal, the law for the being and acting of individual beings,
the purpose for, and the value of, their real
development. Thus the question becomes that of \emph{the relation
between the universal and the individual}.
But as the expression of the universal law and of purpose, reason
stands, secondly, as the ideal over against the
real as what exists in material reality. The
second particular form of the general problem then becomes
the question of \emph{the relation between the spiritual
and the material}.

% --- p. 15 ---
\setcounter{footnote}{0}%
\opage{15}Speculative idealism regarded, as we have
seen, individuals as vanishing moments, in the last instance insignificant, in the development of the universal Idea.

As soon as criticism awoke, it threw itself upon this
point. First an end is made of all ambiguity by
drawing, \emph{pantheistically}, the last consequences
of the concept of the Idea as the only reality. But
soon the consciousness of the significance of the single and the individual
presses forward again, and one seeks to
claim for it a reality that goes further than being a
merely dialectical point of passage. This endeavor
in turn leads in a double direction. On the one hand one sees singularity
only in finite individuality, and as this
point of view is carried through, everything universal vanishes step by step
and is in the end reduced to being a property
or predicate of the individual. This standpoint is
the \emph{anthropological}, in its extreme the \emph{egoistic}.
On the other hand singularity is also seen as the concentration by
which the universal grasps itself, and which is therefore a
necessary condition for the reality of the universal. One then
seeks to develop the speculative concept of the Absolute
into a fuller and more concrete idea of God, in which
the moment of personality comes into its own. This
endeavor thus goes in a \emph{theistic} direction.

By its volatilization of the individual, speculative
idealism had violated the real as a whole. Efficient
causes are set aside for the higher, often
mystical ``necessity of the Idea,'' and material stuff dissolved
before the force of speculative construction as
ice before the sun's rays. But the following age
developed more and more in a realistic direction,
and the reaction against idealism rose to such a degree
that one saw in the Idea only a figment of the brain and in the material
the only reality. Over against this \emph{materialism},
however, there steps another direction, which just as much as
it, but with more sobriety and scientific rigor,
% --- p. 16 ---
\setcounter{footnote}{0}%
\opage{16}recognizes the necessity of the real foundation and
the mechanical forms of action, but which in all this only
sees forms and means for the realization of higher, significant purposes
and Ideas, a direction which we
can therefore rightly call \emph{the Idealism of Reality}.


The struggle over the first problem is waged especially in the two
first decades after Hegel's death (the thirties and forties),
the struggle over the second problem in the two following decades
(the fifties and sixties). Yet it will be most natural
not to follow the order of time, but to put together the kindred
answers to both problems. We shall then
get four main groups: 1) the critical consciousness; 2)
modern materialism; 3) speculative theism; 4)
the Idealism of Reality, of which the first two can be described
as negative in opposition to the last two as
positive. For the critical consciousness develops through
the dissolution of speculative idealism and leads only
indirectly and in the form of a demand to a real,
new content; and materialism, which takes its principle
outside thought and knowledge, must consistently come
to stand in a skeptical and hostile relation to philosophy.
The last two groups, on the other hand, make real attempts
to develop a philosophical world-view.

Over against the magnificent productivity that German
philosophy unfolded after Kant, with an inspired
union of many-sidedness and profundity, had first broken the
new paths for thinking, there may seem to be something
insignificant and merely provisional, something dependent and
unfinished, in the character of post-Hegelian philosophy. We
meet here, as already remarked, no thinkers of the
first rank, no epoch-making, world-historical discoveries
in the world of thought. But one must consider
the task that is set. Speculative idealism
could quickly complete its magnificent systems once
it had first found the fundamental thought by the path of intuition;
% --- p. 17 ---
\setcounter{footnote}{0}%
\opage{17}the method was the abstractly dialectical one, and the movement
took place in the pure ethereal region of thought, which was then confused
with reality itself. Post-Hegelian philosophy is
scarcely so agreed on anything as on its sharp critique of
this abstract idealism. It wants to deal with
reality as reality, but this is not grasped
in one day by thought. Here one must go forward step
by step. Neither the right of thought nor that of reality
may be violated; the points must be found where both meet,
the law of this agreement found, and what is thus
won gradually gathered into a comprehensive total view.
This work is seemingly less brilliant
than that of speculation; but it demands no less
effort and far more resignation. Speculative
philosophy gave itself up confidently to the consistency of its fundamental thought,
as the poet to the power of imagination; speculation
is the lyric of philosophy. Where, on the other hand, it is a matter
of grasping reality as reality and yet asserting
the right of the truly ideal, there flight does not go so easily,
and therefore not always so high either. Yet here too
the eternal urge toward freedom and truth that
is implanted in the human spirit has not left itself without witness,
but has shown itself in clear and deep thoughts,
borne by noble personalities.

\begin{center}\rule{3em}{0.4pt}\end{center}
\chaphere{CHAPTER TWO.}
\begin{center}\rule{2em}{0.4pt}\end{center}


I. \emph{The Controversy over Individual Immortality}.
--- If the Hegelian philosophy had really
``reconciled the ideal with reality, our demands
with what we possess, our wishes with what has been
% --- p. 18 ---
\setcounter{footnote}{0}%
\opage{18}attained''\footnote[1]{J. L. \emph{Heiberg}: Om Philosophiens Betydning for den nuværende Tid. p. 48.}, then its fundamental thought of the unity of reason
and reality could not stand in pure generality
--- for then it would itself be one of the demands
that were now to be satisfied. It had at least
to be capable of giving a definite answer to
the decisive questions which real life leads us to
raise. The system therefore first had to pass the test
of being confronted with the religious and social problems.
But already with regard to the first point that
was made a subject of debate, the question
whether the Hegelian philosophy taught the immortality of the soul,
the indeterminateness of the system with respect to the concrete
and individual showed itself clearly. Hegel had with his
whole intensity immersed himself in his fundamental thought in and for
itself, without making the concrete problems the subject
of separate investigations. Therefore this very
first discussion could lead to a split within his
school into tendencies, each of which believed it interpreted
Hegel in the right way.

Even before Hegel's death his doctrine had been reproached
with destroying the belief in immortality. And on the
other side Ludwig Feuerbach, in his anonymous work
``Gedanken über Tod und Unsterblichkeit'' (1830), had, starting from
the Hegelian system itself, pantheistically denied individual
immortality. But the controversy proper broke out
only when Fr. Richter published his work ``Die Lehre von
den letzten Dingen'' (1833), in which with conscientious
exactness and thoroughness he showed that the consequence
of Hegel's doctrine was the abolition of the old doctrine of
a beyond and another life. Two men who had already
come forward with criticism of the Hegelian system,
C. H. Weisse and I. H. Fichte (son of the famous
I. G. Fichte), declared themselves in agreement with Richter insofar
as they denied the possibility of maintaining
% --- p. 19 ---
\setcounter{footnote}{0}%
\opage{19}individual immortality on Hegel's presuppositions, whereas
they held fast to the possibility of establishing it in
another way. Against these interpretations of Hegel there then
came forward the man who at this time had the greatest
standing among the adherents of the school, Carl Friederich Göschel,
to save the orthodoxy of the system on this point,
with his treatise ``Von den Beweisen für die Unsterblichkeit
der menschlichen Seele im Lichte der spekulativen
Philosophie'' (1835)\footnote[1]{Concerning the younger Fichte's, Weisse's and Göschel's writings in this matter, one may compare \emph{Poul Møller's} interesting treatise ``Om Muligheden af Beviser for Menneskets Udødelighed''. Efterladte Skrifter, 5th volume.}. With Göschel's view
we begin our account, in order then to pass on to
the various dissenters.

Göschel's treatise has its model in the work
with which Hegel was occupied just before his death, on the proofs
of the existence of God. Hegel sought to show how
the three main forms of these proofs point back
to one pervading thought, which is one with the fundamental thought of speculative
philosophy about the relation between
the infinite and the finite, --- if one only breaks through
the external shell of the understanding by which
they are covered in the popular presentation. Something similar
Göschel now attempts with the proofs of the immortality of the
soul. The so-called metaphysical proof, which
corresponds to the cosmological proof of the existence of God, starts
from the soul as an absolutely simple and immaterial being
and concludes from this presupposition that it cannot
undergo any dissolution. It concludes from external
contingency to inner essential being. This
proof points beyond itself. It is by virtue of the unity of
consciousness that the individual is the absolutely simple
being. As consciousness, man stands over against the
rest of existence. Conscious mental life is the purpose
at which life aims; consciousness is to develop
% --- p. 20 ---
\setcounter{footnote}{0}%
\opage{20}in such a way that it takes up the whole of existence into its inwardness.
As the bearer of an eternal purpose, conscious
mental life must itself be eternal. Yet this teleological
proof does not lead us beyond the opposition between subject and
object, selfhood and otherness. Only when the subject has
developed so far that it has not merely overcome
corporeality and appropriated the objective content of
existence, but has also at last found and made itself one
with its own and the All's eternal ground and source in absolute
subjectivity, only then is the individual consciousness
developed into spirit, and the circle is closed, since the final
grounding has been found for the contingent individuality
with which the beginning was made. At this summit
the soul returns to its origin with a richly developed personal
and eternal content, and its immortality can therefore
be derived from its own essence, so that here we get a
proof we may call the ontological.

Göschel does not content himself with this speculative
recasting of the vulgar proofs, but seeks to establish
that Hegel's doctrine, by its whole character, must lead
to maintaining the eternal reality of the individual\footnote[1]{Jahrbücher für wissenschaftliche Kritik. 1834.}. He then refers
partly to the logical determination of the relation between
the universal and the individual, partly to the relation between
spirit and nature, and partly, finally, to the conception of the idea of God.
--- The Hegelian logic posits an inner relation
between the universal and the particular. The universal in
and for itself is barren and abstract; to be the true
universal it must develop itself into concretion, posit its
inner differences as its moments and thus become
totality, unity, singularity. Conversely, the particular is
in and for itself atomistic and abstract; it is seen in its
truth only when it is seen as that which includes in itself and
expresses universal determinations: the true being of the particular
is to be the universal. On these dialectical
% --- p. 21 ---
% print: en guddommelige (misprint: for guddommelig)
\setcounter{footnote}{0}%
\opage{21}determinations Göschel now believes he can build a proof
of the eternal validity of personality. The individual, he says,
proves itself to be the highest, the truth of the universal,
because it has taken the universal up into itself; therefore
the universal is no longer an alien or hostile
power before which the individual must perish.
--- In nature it is certainly the case that the species,
the universal, is maintained through the arising and
passing away of individuals; but from the debt which the individual here
must pay to the species by its death, man is freed,
since as spirit he has gone beyond the opposition
of species and individual, the universal and the particular. In
the sphere of spirit the essence of the species is freely taken up and
thereby overcome by personality. --- Vulgar
pantheism certainly had to let the individual
perish, because it conceives the universal only as objective
and impersonal; but speculative logic develops
the idea of God as absolute subjectivity, in which the individual
can only be rejuvenated, not perish. God is the highest consciousness,
the eternal thought, which upholds and preserves
the differences, the individual determinations, that it thinks.
The human spirit is a divine thought, and
God's thought cannot perish but must be eternal, as
God's essence is. If individuals were not preserved,
the memory of the absolute spirit would be of death and annihilation,
and what then would become of its own life and
immortal power? ---

The younger \emph{Fichte} (``Die Idee der Persönlichkeit
und der individuellen Fortdauer'') acknowledges it as
Hegel's great merit to have shown how far the
way of a priori dialectical deduction can lead; but this
merit one would only obscure by carrying it
over into a domain where it does not hold. The personal
and individual cannot be known by way of dialectical
construction; to reach it, experience must be
reinstated in its rights over against abstract
% --- p. 22 ---
\setcounter{footnote}{0}%
\opage{22}speculation. Göschel therefore errs when he finds the
personality of God contained in Hegel's doctrine. It is the Hegelian
expression ``overreaching subjectivity'' that has
given occasion to this misunderstanding, which falls away
on closer examination of the way in which Hegel
develops the concept of subjectivity. God is for Hegel not
merely substance or dead identity, but the living
process which eternally posits and sublates the oppositions,
is constantly an infinite other, but through this constantly
continues to be itself. It never exhausts itself in
any single self-formation, but reaches out beyond every
particular form. And precisely because God, on Hegel's conception,
is the absolute process in its eternal unrest, he lacks,
according to I. H. Fichte, the kernel and center of absolute
personality. Hegel's merit is to have shown conceptually
the infinite life and movement of the Godhead;
but the other main side he lets recede in
return, namely the eternal rest, the still clarity of
identity, the bond of the self-comprehending unity,
by which the All-active becomes for itself the
absolute One and Blessed. The concentration, ``das unverrückbare
Selbst'', is lacking. Therefore God in Hegel's
system must work his way up through the lower stages in order
to reach the determination of absolute spirit. God becomes
result, not principle.

If this is how matters stand with the divine
personality, then it cannot go otherwise either
with the human. If the Absolute is the eternal
process whose moments have no abiding significance
but are posited and sublated to infinity, then the consistent
Hegelian doctrine with regard to the individual
can only be that the individual is necessary as a dialectical
moment, but that the single individual as such
cannot be anything substantial and enduring. The problem
of individual immortality cannot really
arise at all on Hegel's presuppositions; it
% --- p. 23 ---
\setcounter{footnote}{0}%
\opage{23}first presents itself, and can first find its answer,
when one takes one's stand on experience, physiological,
psychological, ethical-religious experience.
How the younger Fichte thought this
proof could be carried out we shall not dwell on here\footnote[1]{In Poul Møller (loc. cit.) one will find a more detailed account.}, since
it stands in close connection with the whole rest of his
doctrine, which we shall later come to occupy ourselves
with. It need only be emphasized that by his sharp
characterization of Hegel's system he contributed much to clarifying
the points of view, even if his own attempts to tread
the way of ``experience'' still go rather far in the direction
of the fantastic.

C. H. Weisse too denied that Hegel's philosophy
could maintain the validity of personality. He acknowledged
the great and epoch-making significance of the Hegelian logic,
but ascribed to it only formal validity, since it
does not reach the reality of nature and spirit themselves. Indeed,
even as a formal science it is imperfect, since it
abstracts from space and time, these fundamental forms of all
reality. Thereby it has come about that nature came
to stand as pure externality, since space and time
could not be seen in their essential significance for the Idea.
If this significance is acknowledged, then nature meets thought
as the free, rich reality which cannot
be exhausted by logical construction but must be made
the object of a philosophical empiricism. As regards the philosophy of spirit,
\emph{Weisse}, whose first
studies concerned chiefly classical antiquity, whereby there
was awakened in him a lively and clear consciousness of the
worth and significance of the beautiful in personal
life as well, maintains that Hegel wrongly set science as
the highest stage. The idea of truth, which points to
universal reason, must enter into an inward unity with
% --- p. 24 ---
\setcounter{footnote}{0}%
\opage{24}the idea of beauty, which shows us the transfiguration of individual matter
by the infinite, in order then finally, through
ethical-religious experience, to be merged in the Idea of the good.
By this way alone is the highest expression for the idea of God
as absolute personality reached. Therefore Weisse also maintains
(``Die philosophische Geheimlehre über die Unsterblichkeit
des menschlichen Individuums'') that man
can win the conviction of his immortality only by the personal,
ethical-religious way. According to the Hegelian
system the single individual has played out its role
when the world-spirit has passed through it, and the world-spirit
then hurries forward to new, richer and more perfect forms.
But when the true, the beautiful and the good are taken up and worked
into personal life, this gives it an eternal
and infinite significance and takes it up, reborn, into
the divine being. What has thus been
transfigured cannot perish again. But on the
other side it also follows of itself that only those individuals
who are reborn in this way live eternally; to those
who have not reached this stage of spirit there belongs
no immortality. ---

It thus turns out, as the result of Weisse's
and the younger Fichte's investigations, that the Hegelian system,
strictly speaking, does not reach at all the point where the
concrete, individual problems arise. Both of them also point,
though still in an indeterminate way, to the
power which speculation had violated or at least overlooked,
namely experience. Nevertheless these inquirers believed
they could hold fast to the Hegelian method at least
within logic and metaphysics; at the transition to
the concrete spheres experience was then to step in
as a supplement. The Hegelians rightly objected
to this that it is an impossibility, since the dialectical
method is precisely the self-movement of the content. Göschel
maintained in his work ``Der Monismus des Gedankens''
that the claim that reality was more than logical
% --- p. 25 ---
\setcounter{footnote}{0}%
\opage{25}more than rational and had something peculiar to
itself outside the logical form, led to dualism and
materialism, insofar as the material is postulated
from outside as a supplement to logic\footnote[1]{Cf. \emph{Göschel's} statement in ``Jahrbücher für wissenschaftliche Kritik'' 1835. p. 714: ``Es ist die feinste Potenz des Materialismus, wenn sie meinen, dass im Leben mehr enthalten sei, als das Denken mit seinen Momenten, und zur Fülle des Lebens noch etwas Anderes hinzutrete als das Denken; vielmehr ist das, was mehr, was voller, dichter zu sein scheint, ein Mangel etlicher Kategorien.''}. --- Only later
did criticism advance to taking hold of the dialectical
method itself, namely when Trendelenburg in his ``Logische Untersuchungen'' (1840), which we shall come
to later, carried through the polemic in such a way that he penetrated
to the innermost marrow of the system. For the present it is, as
already remarked, more practical and concrete starting points
that criticism chooses for itself. ---

Göschel agreed with his two opponents just discussed
in maintaining individual immortality, but he claimed,
what they denied, that it could be derived from Hegel's
principles. Against both parties, then, stand Feuerbach and
Richter, who deny immortality by virtue of these same
Hegelian principles.

According to \emph{Feuerbach} the modern age suffers from a pernicious
hollowness and emptiness, because, especially in Protestantism,
it has set personality on the throne and lets
everything revolve around it. Hence the modern doctrine of immortality.
The ancients immersed themselves with childlike inwardness
and joy of life in the real world, and it did not occur
to them to desire any continuance beyond it;
insofar as they had the notion of another life, it was
as a shadow-life. The old Christendom and the
medieval Church likewise had really no belief in immortality
in the modern sense; they believed in a judgment of the world,
in heaven and hell, but the single
% --- p. 26 ---
% print: -Kjærlighedsgjerning (misprint: a stray hyphen at the line start)
\setcounter{footnote}{0}%
\opage{26}personality was for them so wholly absorbed in the great purposes and
exemplars that they had no special interest in
its purely individual continuance. It is only modern,
selfish reflection that cannot let go of its dear
I but even transforms death into a shadow, a vanishing thing
before individuality, in whose enjoyment one
wallows. In modern religion the concern is no
longer God but the individual; God is the periphery,
the individual himself the center.

Death must be reinstated in its rights, acknowledged
to be more than apparent death. For death has, in
the first place, a great ethical significance. True action
is the self-surrender of love in opposition to the
egoism of natural man. The truly moral
and religious man places his essence in what stands
above him; his true I is the infinite object
to which he surrenders his finite I. This free
work of love exists in nature as the necessity of
death; the question is put to man
whether he will make this necessity into the
highest freedom through the self-renunciation of love;
only thus, only before death, can death be overcome.

In a metaphysical respect death is an expression of
the reality of the infinite: in death the infinite
being breaks through the finite limits within which
individual existence held it. Death thus comes
not from lack and want but from wealth and fullness. ---
In a psychological respect it holds that in self-consciousness
a distinction must be made between consciousness itself and
its object. What is conscious, the object, is the single
individual; consciousness itself is absolutely universal, is
in all men equal to itself. The various persons
know themselves in consciousness, just as colors are seen in
light. As little as light itself is one with a single
color, so little is consciousness in itself one with the single
person. Death, now, is the fact in which
% --- p. 27 ---
\setcounter{footnote}{0}%
\opage{27}spirit, consciousness, proves its freedom over against the
single objects; as colors change because they are not
one with light, so persons change because they are not
consciousness itself but only its objects, or
because they are only subjects, not subjectivity. Yet for all that
they are not without significance. Consciousness, absolute
subjectivity, is an eternal memory in which everything is preserved
that has been a member of the great whole of consciousness.
That man is immortal means that he is a being
of worth and significance.

What Feuerbach had set forth in flaming
words and with epigrammatic brevity, Fr. Richter developed
at all possible length. His work is of greatest
interest for the guileless directness with which
he is conscious of pleading the cause of the good when he annihilates
the doctrine of a beyond, in his eyes so pernicious.
He reproaches Hegel and his disciples for not
having drawn the full consequences of their doctrine.
Through the Hegelian system with its thoroughgoing immanence
a new world has, in Richter's opinion, arisen;
but the existing order of things has nevertheless been
left standing untouched. This lies perhaps in the
one-sidedly theoretical character of Hegelianism; for Hegel,
as for Aristotle, theory is the highest. But one
should not only think absolutely, but also act absolutely,
live speculatively. Religion is not a particular stage of consciousness,
but, as the innermost absorption of spirit, includes
all stages of consciousness, science too, within
itself. In science the religious disposition reaches knowledge
of itself and of reality and through this passes
over into practical action. --- From this
religious standpoint Fr. Richter carries through the critique
of the notions of a beyond. He conceives the
infirmity of the modern age in the same way as Feuerbach.
The condition of a real ethical striving is that one
becomes conscious that this world is a real world, closed
% --- p. 28 ---
\setcounter{footnote}{0}%
\opage{28}in itself, and that one surrenders oneself to the universal
purpose as an unconditional servant. In this consists the true
eternal life. Criticism therefore does not touch the true content
of the dogma of immortality; it only sees as present
what the dogma pushes out into a future,
otherworldly existence. The motive for assuming such an existence is
either idle, unbounded wishes without justification
or an equally unethical demand for reward for the good
in another life. Good action has its reward and its
purpose in itself. By death the single man is reminded
that his individual life is not the world's final
end, that the species outlives him, and thus he lives
not for himself alone but for the species. In the certainty
of his death man has the first occasion for
distinguishing himself from the animal. --- In an anthropological respect
the belief in immortality is founded on the old dualistic
conception of the relation between soul and body. The soul
is regarded as dwelling in the body and leaving it at
death, which is the separation of the two. But true anthropology
sees man as a whole, as a harmonious
unity from the cradle to the grave. ---

The one question of individual immortality
necessarily had to draw a multitude of others connected
with it in its train; so deeply does it reach into
the whole world-view. Thus we have seen
how it led to discussing the conception of the idea of God,
the essence of knowledge (speculation --- experience)
and the problem of immanence and transcendence. Even
this first debate therefore laid the ground for a split
within the Hegelian school, whose opposed parties
Strauss later (Streitschriften. 3rd Heft. p. 95) designated
as Right and Left, after the parties in the French
assemblies. If one asks whether the conservative Right
or the radical Left most nearly expresses the own opinion of the school's
founder, then no other answer can be given
than that Hegel himself, personally, no doubt believed himself in
% --- p. 29 ---
\setcounter{footnote}{0}%
\opage{29}agreement with the doctrine of the Church, but that the philosophical
consequence of his system goes in the opposite direction.
The full consequence works its way forward step by step
and is seldom or never grasped in its whole extent
by the first originator of the thought. Hegel would, as
Ruge aptly remarks (Deutsche Jahrbücher. 1841.
p. 142), hardly have regarded the neo-Hegelian development
of his doctrine with more favorable eyes than Kant regarded Fichte's
development of the fundamental thoughts in ``Kritik der reinen Vernunft.''
--- We now pass on to a still more burning
problem of the philosophy of religion, through which that opposition
between the two tendencies especially comes forward.

2. \emph{David Friederich Strauss} (b. 1808). ---
Speculative philosophy had placed itself in a relation
to positive religion similar to that which mysticism,
guided by its practical-personal interest, takes up. Just as
the mystic believes he acknowledges sacred history
when he sees it as an inner history, taking place in his
own soul during its struggling and laboring to give birth to
God within itself and by the way of suffering to die away in him, so
speculation believes it has explained historical
Christianity when it has shown the inner relation
between the infinite and the finite and the necessity
of the latter's being merged in the former. Hegel too had
received the seed of his philosophy of religion through the study
of Master Eckardt. We have already seen its fundamental thought,
but shall now consider its application
with regard to the center of the Church's doctrine, Christology.

The finite world is, according to Hegel, an appearance
in external reality of the inner differences and
moments of the Absolute. The finite spirit is indeed in itself one with
the Godhead, and originally it also rests securely in
this unity, which immediately fills consciousness. But
step by step the self-consciousness of the finite spirit grows,
and with it also its isolation. The absolute content
vanishes more and more from it and becomes an
% --- p. 30 ---
% print: Methaphysiske (misprint: so printed)
\setcounter{footnote}{0}%
\opage{30}alien thing. The period of dissolution of classical antiquity is the
point in time when the infinite content has completely vanished,
so that empty self-consciousness remains behind
with proud resignation before the collapsed world
or with negative joy that everything is shown in its nothingness.
This standpoint Hegel designates as the unhappy
self-consciousness: the complete finitude of self-consciousness
in the midst of its apparent infinity. Finitude weighs with all its weight, with all its pain
and misery, on the spirits that have lost the eternal
powers and do not know how they are to call them to life
again. The only rescue is that the human spirit
recognizes that the finitude under which it groans
is not alien to the essence of the Godhead, but is one
of its forms. God himself is not outside finitude,
but has entered into it. Distress and pain are
then only the expression of the infinity of spirit, of its incommensurability
with finite existence, through which
it must nevertheless pass in order to develop its own essence.
The finite spirit now no longer sees itself as
isolated; self-consciousness is again united with infinite
substance. This is the Idea of the God-man, the idea of
the unity of the divine and the human nature,
which dawns upon consciousness.

This construction, in which the metaphysical and
the historical fuse together, still needs an important
psychological supplement in order to lead to the goal\footnote[1]{The metaphysical construction is developed in particular in the philosophy of religion and points back to the logic, the historical especially in the phenomenology, the history of philosophy and the philosophy of history. --- Already in Schelling, moreover, the basic features of these constructions occur. (``Vorlesungen über das akademische Studium''. 8th and 9th lecture).}. Not
only speculatively, for thinking consciousness, must the idea
of the unity of the divine and the human
reveal itself; it must also come forward for immediate,
% --- p. 31 ---
\setcounter{footnote}{0}%
\opage{31}popular consciousness. From this point of view
that progress of the world-spirit shows itself in the phenomenon
in this way, that it becomes the faith of the world that the Godhead has
come into being as a self-consciousness, as a single man,
that this divinity is seen, felt, heard by
believing consciousness. What is in itself the property of the human spirit
then comes forward as something external and beyond,
which stands over against it. We thus come back here to the
Hegelian doctrine already touched on earlier, that the content of religion
and of philosophy is the same, only the form
different. The unity of the divine and the
human nature appears, considered speculatively, as
Idea, considered religiously, as the historical Christ.
This difference of form is not to impair the content, although
Hegel holds fast that the form of the concept is the higher, to
which one must advance from the historical form, or
the form of representation.

In this speculative Christology a large part of
Hegel's disciples saw a reconciliation of faith and knowledge, and theology
celebrated ``the new era'' with enthusiasm.

Then in the year 1835 appeared Strauss's ``Leben Jesu, kritisch
bearbeitet''. The year 1835 was, as has rightly
been said, in theology what 1848 was in politics. ---
Strauss here came to the result that the Gospel
history is a mythical history which developed in
the consciousness of the earliest congregation as the expression of the great
thought that had dawned upon it of the unity of the divine
and the human. By the concept of myth
the unclear and ambiguous in the Hegelian Christology is removed;
it furnishes the necessary middle term for explaining
how the Idea can, for consciousness, assume the
shape of a single individual. This shape is only a
vanishing form, since the universal Idea cannot absolutely
reveal itself in the individual phenomenon, but
needs the whole species as its medium. All the predicates which the Church's doctrine ascribes to Christ must therefore
% --- p. 32 ---
\setcounter{footnote}{0}%
\opage{32}be transferred to humanity as a whole; they also contradict
one another if they are really to belong to a single individual.
But humanity is God become man,
son of a visible mother and an invisible father, of nature
and spirit; it is the worker of miracles, insofar as spirit
in the course of history ever more completely masters
nature; it is the one who dies, rises again and ascends into heaven,
insofar as from the sublation of its immediate,
natural condition a higher spiritual life constantly proceeds.

When speculative Christology deduced the
positive dogma, this deduction really only aimed
at the psychological and historical necessity of
the illusion in which the universal appears as an individual
single thing. Hegel's orthodox disciples naturally
conceived the matter differently. Göschel emphasized against
Strauss that if one lets the Idea have reality only
in the whole species, one sinks back to Kant's standpoint,
where the Idea has only subjective validity: for if the Idea
is nowhere completely realized, then it exists
only in the subjective summing up of its moments scattered over
the whole species, and the Hegelian fundamental thought
of the unity of thinking and being is sublated. ---
Strauss answers to this (Streitschriften. 3rd Heft. p. 99 f.)
that this objection could also be directed against the assumption
of the realization of the Idea in a single individual: for since
this realization cannot be absolute in any
single moment of his life, it would exist only
for the thinking and remembering summing up of all
the moments of his life. It is Göschel himself who falls back
to the Kantian standpoint, since he does not credit
the Idea with the power to permeate multiplicity without
losing its unity. Göschel demands the general concept
realized outside the sum of the individuals that
belong under it, like the scholastic who asked for
fruit to eat, but fruit \emph{in general}, neither
apples, pears nor any particular fruit.

% --- p. 33 ---
\setcounter{footnote}{0}%
\opage{33}Strauss was attacked not only with dialectical weapons.
His book had stirred up commotion in the theological camp,
and a confused crowd pressed forward to answer or
at least to revile the giver of offense. The well-known champion
of orthodoxy, Hengstenberg, with his ``Evangelische Kirchenzeitung'',
was not the last. Strauss has, according to Hengstenberg,
``a heart like Leviathan, as hard as the
nether millstone'', and belongs to those who are enthusiastic
for the spirit of the abyss, for the great beast to whom
it was given to speak great things and blasphemies etc.
To this Strauss, who moreover valued
Hengstenberg's orthodox zealotry more than the groping and
half-hearted procedure of the usual apologetics, replied
with a dignified statement, inspired in all its calm, in which
he maintains the right of free science and rejects all
irrelevant objections drawn from morality
(Streitschriften. 3rd Heft. p. 22 ff.).

We must especially emphasize that Strauss defends not only
the freedom of philosophical thinking but also that of historical
inquiry. In this too Strauss marks
an advance in relation to Hegel. Hegel on the
whole did not care for historical criticism, any more than Goethe did. Both
feared that the great figures of antiquity would lose their
luster when criticism took hold of them with its irreverent
hands. Hegel saw in criticism a relapse into
Kant's subjectivism and demanded that one should immerse
oneself in the historical content without troubling about the
critical difficulties. He mocks Schleiermacher's
critical investigations of Plato's dialogues and Niebuhr's
of the earliest Roman history. This is closely
connected with Hegel's whole position toward reality.
He shows it at once too much and too
little respect. True reverence for reality
consists in a serious testing of how far the representation
of the real also corresponds to the real
in itself. Just as little as the natural scientist accepts
% --- p. 34 ---
\setcounter{footnote}{0}%
\opage{34}the immediate sense-image as real nature,
just as little can the historical inquirer stop at what
immediate faith holds fast as historical truth.
This principle, which springs from the very essence of scientific
inquiry, Strauss was the first to apply to a domain
that had previously been kept closed to it. How confused
a conception people still had of the significance of criticism
can be seen from the demand made by an
ardent Hegelian who at that time belonged to the extreme Right
but later went over to the extreme Left, Bruno Bauer.
Criticism, in his opinion, is not to be merely skeptical
and negative, but is through negation to re-erect
its object in its whole extent and leave nothing
behind that is not positively justified. Only then is
it true, positive criticism. But --- Strauss answers --- the critic
must surely at least presuppose the possibility
that there is something in his object which cannot be justified,
since otherwise he would have had no motive
for taking up the critical standpoint. That possibility
can be kept entirely away only if one starts
from the assumption that the absolutely divine character of the object
makes every corruption or addition impossible; but
this is in itself a purely arbitrary assumption.

From another point of view, on the other hand, Strauss could rightly
be reproached that his criticism was not positive.
He stops the investigation at the content of the Gospel
writings and believed he had reached the end when
he had shown what thought must lie at the
root of the mythical clothing in the individual narratives.
The next step was to investigate the origin of the writings
themselves, their connection with the whole
age and with the spiritual movements that crossed
one another in it. Only thereby, by being set into
the whole historical context, does that content find
its full explanation. Strauss's teacher, the famous founder of the Tübingen school,
Ferd. Chr. Baur
(† 1861), struck out on this way. In this domain he stands beside
% --- p. 35 ---
\setcounter{footnote}{0}%
\opage{35}F. A. Wolf and Niebuhr in the domain of classical philology.
It is true that in the Tübingen school too, and surely
not without reason, uncritical after-effects
of Hegelianism have been found. But that proves only that a principle
correct in itself can receive a one-sided application on account
of the circumstances under which it first appeared; it proves
nothing against the truth of the principle. ---

Strauss designates his general standpoint as
absolute immanence. His whole world-view
comes forward especially in his work ``Den christelige Troeslære
i dens historiske Udvikling og i dens Kamp med
den moderne Videnskab'' (1840---41). Here it is above all
the conception of the idea of God that matters. Hegel's
doctrine of the subjectivity of the Absolute had been understood by his orthodox
disciples in such a way that the old dogma of the personality and
transcendence of God was directly
grounded in it. On this point, as in the doctrine of immortality
and in Christology, they therefore considered their
master to be in full agreement with the doctrine of the Church.
But Strauss shows that by his doctrine of absolute
subjectivity Hegel only meant to insist that the Absolute
cannot be conceived as dead, resting substance,
but as the infinite life which is its own product and
develops itself through an infinity of differences and
oppositions. Spinoza did not comprehend the Absolute as
eternal self-determination; he showed how everything
flowed back into substance, but not how
it flows out of it again, --- only the venous system,
not the arterial system, not the whole infinite circulation of life,
by which substance through its other closes
together with itself and so becomes person and spirit.
This conception, according to Strauss, avoids two dangerous
rocks. It does not set the Godhead, in an external way,
finished over against the world, like vulgar theism;
nor does it bind God's essence to the development of the world in such a way
that he should need finite existence
% --- p. 36 ---
\setcounter{footnote}{0}%
\opage{36}and the human spirit in order to be himself.
Not by man, who after all is not by himself, but
through and in man, whom it has itself posited, does
the Absolute come to self-consciousness and personality.
God's personality must not be thought as a single personality,
but as an all-personality; instead of personifying
the Absolute one must comprehend it as that which to
infinity personifies itself.

The single personifications of the Absolute are
only transient moments. If they became fixed, that would
be a sign of the impotence of absolute spirit. The
substantial is not in the individuals, but beyond them in
absolute spirit, to which they are related as changing
accidents. This holds with regard to the whole of
existence developing itself in time. Time is a
finite form of representation, which needs external points of support
for intuition in order to be able to maintain itself. Thereby
the finite differences are fixed, which in themselves possess no reality.
The unity in the difference of the moments of time is eternity;
by thinking time one deprives it of beginning
and end, its represented points of support, and it sinks
together of itself into the indifference of eternity as
its ground, from which it just as well proceeds to
infinity. Eternity and time are related as substance
and its accidents.

It follows from this that Strauss's pantheism cannot
avoid falling back at last into the Spinozistic
doctrine of substance. Individual life is a mere phenomenon,
behind which thought seeks the true, which is the
undifferentiated eternity of substance. It is only representation
that can still ascribe real reality to finitude,
and representation is, according to Strauss, an inadequate
form of knowledge.

But a still more awkward circumstance shows itself.
There always remains an opposition between
substance and the individual consciousness in which
% --- p. 37 ---
\setcounter{footnote}{0}%
\opage{37}it comes to itself. Substance is even to
lie \emph{beyond} single beings --- and yet Strauss's
last word in his dogmatics is that the otherworldly \emph{under}
all shapes is the enemy which speculative
criticism is to combat. The standpoint of immanence is
thus not carried through. There is still a step to take
--- but it goes in quite another direction than toward
Spinoza. Self-consciousness must sublate every transcendent
relation in order to become wholly clear and transparent to itself;
it must remove every alien power as a
play of shadows. Along this way we are also led by the further
development of the ``critical consciousness'' in Bruno Bauer
and Feuerbach.

Strauss's sharp critical thought has thus loosened
the unity between the moments that in Hegel's doctrine were
in intimate connection. His whole critical tendency brought
him into opposition to the dialectical method and to
all-mediating speculation, an opposition that could only
be sharpened by the way in which the orthodox
Hegelians misused the doctrine of dialectical ``contradiction''
to abolish the difference between ``Widersinn''
and ``Widerspruch'' (Streitschriften. 3rd Heft. p. 23. 113).

The opposition between substance and individuals
proved to stand in close connection with the psychological
opposition between thought and representation. Consideration
of this opposition will throw a decisive
light on the Straussian and thereby on the whole speculative
philosophy of religion. --- Strauss counts it as
the merit of Christianity to have brought the unity of
the divine and the human to the consciousness of the human race.
This Idea is therefore what is abiding in
Christianity. The old world was essentially dualistic;
the peoples felt the Godhead as a distant and alien
power, which only in the earliest time, in the mythical
golden age, went about in friendly fashion among men. Against
this dualism Christianity came forward with its doctrine
% --- p. 38 ---
\setcounter{footnote}{0}%
\opage{38}that God had become man. Thus both worlds were
united. But the union was accomplished only
at a single point, as an exception, a miracle. The
miraculous, hence dualistic, way of looking at things remained
dominant. Christianity is thus the divine's
becoming this-worldly on the presupposition of its otherworldliness,
the beginning of the modern system of immanence,
but still on the old soil of the transcendent.
It now depends on whether one will regard it as the
peculiarity of Christianity that the immanence of the divine
is restricted to a single point, which is then
separated by an absolute, insurmountable chasm from the rest of
humanity, or whether the true essence of Christianity might not
be preserved if one developed the seed of immanence
it contains with the help of the results of human
development. If one holds fast to the narrower,
stricter and therefore essentially dualistic conception of
Christianity, then one comes into an irreconcilable conflict
with human development. If, on the other hand, one enters
into a more conciliatory conception, then one can say that
Christianity is constantly developing further from its original
starting point. We should then have an analogy
in the history of political freedom. The Greek states acknowledged
freedom, but only as the property of a certain number of citizens,
in opposition to slaves and metics. They
thus wanted to know of freedom only at this point.
But historical development has led to extending
it to all, in accordance with the awakening consciousness of the
share of all in the same human nature. (Zwei
friedliche Blätter. 1839). --- With this way of looking at things
Strauss's treatment of the dogmas agrees. He
grasps each single dogma at its origin and follows it
through the whole course of historical development. Through
this process the dogmas break their finite form
and are purified to the purity of the concept. The form of representation,
in which religion has clothed the infinite content,
% --- p. 39 ---
\setcounter{footnote}{0}%
\opage{39}is not merely an outer shell, but causes a pervasive
dualism in the whole view. Hegel is
therefore wrong to regard this difference of form as inessential.
From the imperfect one must advance to the perfect,
from faith to knowledge.

In the philosophical world-view that replaces
the positive doctrine of the Church, consciousness finds the same
satisfaction as the believer in his dogmas, and the
difficulties and contradictions which these bring with them fall
away there. But will this liberation then be able to fall to
the lot of all, or will it always be restricted to a few
individuals? This question Strauss leaves unanswered;
he sees in the opposition between those who know and the
non-philosophizing a chasm between two classes of
human society which will perhaps never be filled.
He only emphasizes that it is of no use to cover over
this chasm, since it is there once and for all. ``So let the
believer let the knower, and the latter the former, calmly go his
way; we let them keep their faith, so they must also
let us keep our philosophy; and if the excessively
pious should succeed in excluding us from their
Church, we shall count it a gain; of false
attempts at union there has been enough, only a sharpening of the oppositions
can lead further.'' (Troeslære. Danish transl.
I. p. 309).

Religion is thus only an imperfect form of
philosophy. The religious really receives only a negative
determination; it is that which is to be sublated, translated
into another language, a child's babbling in relation to the
rational speech of the grown man. This points back to a psychological
presupposition from which Strauss starts, as does
the whole speculative philosophy of religion,
whose final consequence he draws, --- the presupposition,
namely, that religion has its origin in the drive for knowledge
and is therefore essentially an intellectual phenomenon.
If this presupposition is correct, then it will
% --- p. 40 ---
\setcounter{footnote}{0}%
\opage{40}hardly be possible to raise any decisive objection against Strauss;
then all religion falls away with positive religion.
But if this presupposition is correct, then it also becomes
difficult to understand the chasm that is to go on existing
between those who know and those who believe. How can
Strauss doubt that this chasm will cease, when the two elements
it separates are essentially akin, are moments
in the same development? Can one suppose that an
immeasurable barrier should have been set between two
stages of knowledge, and must one not rather ask whether
the refusal of religious consciousness to dissolve itself into
speculative knowledge might not be grounded in the fact that it
does not find its whole essence expressed in knowledge, but
misses essential elements, elements that make up precisely
what is religious in religion?

Thus on the one side the ease with which Strauss
can show the dissolution of dogma into speculative concept,
on the other side the resistance which, he admits,
religious consciousness offers to making this
transition, --- this opposition awakens doubt whether
the psychological presupposition from which he starts
is correct. Strauss has not entered upon a real testing
of it. He has inaugurated a new epoch with
regard to the historical foundation of the philosophy of religion,
but he has overlooked the psychological foundation, which
is still more important, since the historical is not
intelligible without it. On this point too we are referred to
Feuerbach, who broke through the barriers that still
held Strauss captive.

3. \emph{Bruno Bauer} (b. 1809) is a living testimony
to the ambiguity in the Hegelian principle of the unity of reason
and reality, a principle that can just as well
be reactionary as revolutionary, just as well
conservative as radical. Thus Bruno Bauer could,
from the extreme Right, the most far-reaching orthodoxy, suddenly
throw himself over into the opposite extreme. In his critique
% --- p. 41 ---
\setcounter{footnote}{0}%
\opage{41}of Strauss's ``Leben Jesu he reproached Strauss, as already
mentioned, for not carrying through completely
the comprehension of the objective; criticism stops at a
certain point and does not prove itself to be ``the invincible
power of self-consciousness'', which lets nothing remain standing
as an impenetrable object, but negates
everything objective in order to recognize itself in it. In Bauer's
opinion the true, ``positive'' criticism justifies the object
in its entirety by comprehending it. --- But
a few years later he attacked Strauss from quite another side.

Strauss --- says Bruno Bauer --- has in his criticism
remained standing at an essentially theological standpoint.
He believes his task accomplished when he has shown the
mythical character of the Gospel narratives, and then
thinks that these myths formed themselves in the earliest congregation
under the influence of tradition. But here he points
back to something mysterious. Every individual work must
surely have proceeded from an individual thinking and acting
and cannot grow out of a general consciousness which
with mystical power inspires the single persons, without any one of them
being able by himself to be called the originator of the narrative.
No, criticism must be carried through in such a way that one gets a
single, definite author to hold on to. The evangelists
are to be compared with Homer and Hesiod, who, as
Herodotus says, made the Greeks their gods.

Just as Strauss in his historical criticism lets the
mystical power overwhelm real self-consciousness,
so his whole view aims at maintaining
the Absolute as objective substance. It is thereby that
he essentially stands at a religious or theological
standpoint. For the religious standpoint is that
where self-consciousness still has its content separated from
itself, whereby the content is made finite. At its summit
religious consciousness is absolutely alien to
itself. In nature religion the gods are still regarded
as particular, single beings, and as such are akin
% --- p. 42 ---
\setcounter{footnote}{0}%
\opage{42}to men; therefore the alienness is not carried through,
but state, family life and art receive immediate
religious significance. ``The chains the human
spirit bore in the service of these religions were wound about with
flowers; splendidly and festively adorned like a sacrificial animal,
man brought himself as a sacrifice to his religious powers;
the chains themselves deceived him as to the harshness
of his service.'' But the flowers withered, the ancient
world of nature collapsed, and ``the vampire of spiritual
abstraction'' accomplished its work. Only the I, self-consciousness
in its whole infinity, remained behind amid
the ruins of the perished world, as the abyss into which everything
had sunk. But the I could not yet see through
itself; its infinity was still an alien being
for it, an absolute, almighty Godhead which stood threatening
over against it. This is the unveiled mystery of Christianity.
Self-consciousness has in the Gospel narratives
to do with itself, but that is to say, with the most frightful
parody of itself, a Moloch that demands every
natural and real relation as a sacrifice. But since it
is nevertheless the image of self-consciousness, it is understood how
it has been able to hold humanity captive so long, to impose
such terrible burdens on it. The human race has
been brought up in bondage under its own likeness, in order
that it might the more thoroughly win freedom and establish
itself in it\footnote[1]{Kritik der evangelischen Geschichte der Synoptiker. 3ter Band.}.

Pantheism, which sets an objective substance as
the object of consciousness, has not yet reached this
freedom, is still within bondage. The opposition
must be carried to the end, where only atheism is the right
name. Religious consciousness must be completely negated.
It is altogether false when Strauss thinks that philosophy
is to take the place of religion, and that it can
do so because they both aim at the same absolute
% --- p. 43 ---
\setcounter{footnote}{0}%
\opage{43}content. Religious consciousness is indeed the seed of
perfect self-consciousness, but this seed is, precisely
because it is only a seed, a false and pernicious power.
Those in whom reason is not developed have indeed the disposition
to true humanity, but precisely because it is only a
disposition, they are not human beings but inhuman beings. Therefore
no settlement with religious consciousness can be thought
of. Truths cease to be truth
when they are translated into the form of religion. There is no
truth in religion. Since religious self-consciousness
is not free and does not see through its own essence, it still relates
itself in an arbitrary and contingent way
to its object; it is the obscure drive that posits
the object of religion, a drive which on the one side
is an anxiety before the dark, uncomprehended power that weighs
on consciousness, on the other side is the passionate
wish that wants to see its fantasies realized\footnote[1]{Bruno Bauer in ``Deutsche Jahrbücher.'' 1843. p. 82---95.}.

When, then, criticism has been carried through to the point where
the dark power is altogether annihilated, every resistance,
every object removed, what then remains? Bruno
Bauer answers: pure criticism. Nothing has truth except
critical consciousness itself. Every object is posited
only to be negated again; insofar as it is acknowledged, it ceases
to be true. He who seeks satisfaction
in anything, he who wants to be enthusiastic about anything, must go
back to the demonic power of religion in order to let himself
be robbed of his freedom. The expression atheism no longer fits,
since it still includes a regard to the object
that is negated, and therefore does not designate absolutely free
self-consciousness. State, family, art and science are dissolved
in this titanic movement, in which his brother
Edgar Bauer took upon himself in particular the critique of the state\footnote[2]{On him see ``Nyt dansk Maanedsskrift''. 1871---72. p. 223.}.

The movement is here carried with tearing speed to its
% --- p. 44 ---
\setcounter{footnote}{0}%
\opage{44}extreme limit. It is a fanaticism of the concept that
appears here, first with the demand for absolute comprehension, for
seeing through the depths of existence so that not a single
obscurity remains, --- and when this fails, with
an equally absolute devotion to the opposite striving, to
let everything be annihilated by the self-consciousness in which
it would not be resolved. The Hegelian panlogism is
here carried through in such a way that it has returned to its
starting point, the invisible point where thinking and
being are one, because only thinking is. This invisible
point is what Bauer designates as pure criticism; but
whereas Hegel started from this point in order to conquer
the whole world, Bruno Bauer has arrived at it by
letting everything sink together into nothing. He ends in nothing
because he begins with nothing.

4. \emph{Arnold Ruge} (b. 1802). --- Just as Hegel in
the philosophy of religion had applied his principle
of the unity of reason and reality in such a way that it
meant the reconciliation of dogma with speculative thought,
since, as Strauss expressed it, the dogmas came through the
dialectical purgatory into the heaven of the Idea with skin
and hair, so he applied the same proposition in the philosophy of right
in such a way that the state of the Restoration period,
this expression of broken promises, was thereby
established as ``rational''. Here as little as there did
Hegel apply any real critique of what lay before him. Since he was
convinced that the Idea is not so impotent that it should only
be a subjective ideal which has not won for itself
any form of reality, but that in the state we must see
the world the moral Idea has created for itself, his concern
was only to show reason immediately in the
historically developed forms of state.

Yet Hegel lived to see, the year before his death, a
reality he did not find rational, namely the July Revolution,
and his last days were embittered by a quarrel into which he
came on that occasion with one of his disciples, the philosopher
% --- p. 45 ---
\setcounter{footnote}{0}%
\opage{45}of right Gans, who wanted to draw liberal consequences
from his doctrine\footnote[1]{Gans began his ``Vorlesungen über die Geschichte der letzten fünfzig Jahre'' by reproaching German historiography with turning exclusively toward the past and setting aside the history of its own time as ``Tagesgeschihte''.}. But it was only some ten
years later that the speculative view of history
went through a crisis similar to the one which, in the philosophy of religion,
had broken out with Strauss and
Bruno Bauer.

Arnold Ruge published, from the year 1838 on, a journal
by the name of ``Hallische Jahrbücher'', which originally belonged
to what was called the Hegelian Center. It appeared
with reverence for the positive and, with
Hegel, considered the difference between religion and philosophy a
mere difference of form. But year by year it became more radical.
First Strauss's, then Bruno Bauer's and Feuerbach's
views in the philosophy of religion became the dominant ones
in the Jahrbücher. Ruge separates himself from the Center and the Right
of the Hegelian school by reproaching them for their
position as inactive spectators of the world-historical
process. ``There is only one atheism, doubt
of the spirit of history.'' Hereupon followed an ever sharper
critique of the Prussian political conditions, which at last
compelled Ruge to move to Dresden (1841), from where he
now published the journal as ``Deutsche Jahrbücher.

This opens with a preface in which he charges the Hegelians
with having a finished system, a
presupposed dogmatics, and finding the Idea without further ado in the
political status quo. Against these three powers, philosophical,
theological and political orthodoxy, he
wants to open a war of liberation. And if it is asked under what
banner this struggle is to be waged, Ruge answers:
``There is only one sign in the heavens which lights and warms
the spiritual world, and the name of this star is
% --- p. 46 ---
\setcounter{footnote}{0}%
\opage{46}philosophy.'' But then philosophy cannot be conceived in
Hegel's way, who thought that its calling was only to comprehend
the preceding reality. Hegel, says Ruge, thereby himself
acknowledges after all that thought is a new and higher
stage, since it judges reality. And since thought
is thus the last, it is, which Hegel overlooked, also
the new, an anticipation of the coming reality.
In the logic Hegel teaches that the negation of the preceding
leads over to a new stage; but in world history
he will not admit this. And yet the movement
in world history rests on human consciousness
making inward to itself the content it has won, which is precisely
what happens through philosophical critique. Thereby one advances
to a new task, which sets an ``ought'' in place of
an ``is'', and thus the right carrying through
of Hegel's dialectic leads to recognizing also the truth of
Kant's doctrine of the autonomy of the will and Fichte's doctrine of
the self-determining I. The philosopher is an apostle
of the future, who theoretically accomplishes the movement of spirit
that is to be accomplished in practice in the life of the state. It
is no reproach against theory that it is revolutionary:
``nur umwälzende Gedanken sind Gedanken.''

Critique thus seeks the concept of the preceding reality
in order to derive from it the concept of the following reality.
But Hegel has confused the metaphysical
and the historical concept, and has thereby misjudged
the significance of critique. It is the metaphysical concept
that drives critique forward, since it proves
not to have its complete expression in the historical.
Thus the Idea of the state is not realized in a single
definite state. Hegel, on the other hand, immediately raised the
historically existing to metaphysical dignity. In
this immediate transition the opposition is overlooked, and
with it critique. Only through critique, on the
other side, does the way lead back again from the concept to reality.
``The true connection of the concept and
% --- p. 47 ---
\setcounter{footnote}{0}%
\opage{47}reality is not the apotheosis of existence into concept, but
the incarnation of the concept into existence.''

If it is asked how this incarnation of the concept
takes place, then Ruge's determination of the essence of religion
is to be brought out. The Hegelian philosophy and
the German spirit as a whole have, he says, concealed the practical
significance of theory. Theory becomes practical when
it passes over into enthusiasm and resolve, i.e., when
it becomes religion. For religion in itself is nothing
other than the practical pathos for the ideal, for truth,
faithfulness to the Idea and self-sacrificing exercise of
it. Religion is praxis: it everywhere brings with it,
on account of the resistance of finite existences to the Idea, its
peculiar struggle and the sacrifice of earthly interests.
This is the true religion, the religion of the this-worldly,
in opposition to the dualistic religions,
which, by making man alien to himself
and to existence, only hinder him in the exercise of freedom
instead of furthering it. Never has anything other than freedom
been religion. But now its form
and content become richer than before. The whole ideal world of humanity
must go into the crucible of the heart, in order to
come forth from it again as a glowing stream and form a new world.

Mere theory, on the other hand, which does not deepen itself in itself
into religious pathos, does not get beyond the
so-called liberalism which never attacks the very
root of consciousness. It believes that it only matters
to reform the political forms, while for the rest it
lets everyone think as he will and be saved
by his faith. This indifferentism shows the superficiality
of liberalism. It does not see that what matters
is a reform of consciousness, and therefore never gets
further than mere wishes for freedom. With the reform of consciousness
the reform of the world would be given.
When this reform is carried out, liberalism becomes
democratism; every alien power is resolved into
% --- p. 48 ---
\setcounter{footnote}{0}%
\opage{48}free self-consciousness, and in the state every opposition vanishes,
since the Church is resolved into the school in order through
the education of the people to absorb all rabble, and since the military and judicial systems are directly administered by the self-organizing
people. ---

At this point in its development the journal was
banned by the Saxon government (1843). Ruge later,
as a member of the Frankfurt Parliament in 1848, gained the
experience that there was still far to go before the realization
of the reform of consciousness. The demand Ruge
had made makes him what he wanted the philosopher
to be, a prophet of the future. For that demand
for the reform of consciousness is in its generality the task of
all history, and it is a grandiose idealism
to assert it immediately without having positive
starting points for its realization. Therefore
Ruge is not a practical politician, but a noble disciple
of Plato's Republic\footnote[1]{Part of the section on Ruge has, with additions and abridgments, been printed in ``Nyt dansk Maanedsskrift''. 1871---1872.}.

In a philosophical respect what is interesting
about Ruge is that he has maintained the practical
significance of philosophy. He has maintained the necessity that the Idea
should not merely become an object of speculation, but be appropriated
by living personal consciousness. Thereby he has
overstepped the limits of the Hegelian psychology and philosophy of religion,
within which heart and imagination
could not receive such a positive significance. Indeed, he
abandons, by his determination of the relation of thought to
historical reality, the very fundamental thought of Hegel
of the unity of thinking and being. If critique is a
necessary intermediate determination between the concept and reality,
then that unity is not immediate and absolute.
Through the historical critique and the critique in the philosophy of religion
we thus catch sight of new metaphysical consequences.

% --- p. 49 ---
\setcounter{footnote}{0}%
\opage{49}5. \emph{Ludwig Feuerbach} (b. 1804). --- Each of
the phenomena with which we have so far been occupied appeared
at once with a peculiar, firmly marked character.
And insofar as in Bruno Bauer we were
witness to a swing from one extreme to the other,
this swing took place in an instant, with restless haste. It is otherwise
with the thinker at whom we now stand, and who is the truly
typical representative of the process of self-dissolution
of speculative idealism, but in such a way that he at the same time points out
the positive elements for a new development. Feuerbach
runs through the whole scale from the extreme right to the
extreme left, from absolute idealism to the boundary of
materialism, in such a way that each single step is clearly
and emphatically marked, whereby the transitions become
gradual, indeed almost imperceptible to the superficial glance.
He is filled with a fermenting unrest which has followed
him from his earliest youth, when as an orthodox theologian he
took delight in the speculative reconciliation of faith
and knowledge, right up to the sensualist subjectivism in which
he ends. Ever dissatisfied with what is given, with what has
so far been attained, he ever seeks new forms and expressions.
He himself relates how little desire he had, on the occasion
of the publication of his collected works, to return
to writings which for him marked a stage left behind
--- and every one of his writings did so as soon
as he had put the last hand to it\footnote[1]{``Jeder fertigen Schrift sagte ich für immer Adieu; jede hatte mir nur meine Fehler und Mängel zu Bewusstsein gebracht und daher nichts andres in mir zurückgelassen, als das dringende Verlangen, ihr Andenken durch eine neue Schrift auszulöschen. Und nun wurde mir auf einmal zugemuthet, meinen unzufriedenen, schriftwidrigen, unbiblischen Geist auf alle meine längst meinem Sinn entschwundenen Schriften zu richten. Welche Zumuthung! Wider den Strom des Lebens soll ich schwimmen? wider den Lauf der Natur statt vorwärts, rückwärts gehen?'' (Werke. 1ster Band. p. V---VI).}. This unrest,
this striving after ever renewed advance, gives him
% --- p. 50 ---
\setcounter{footnote}{0}%
\opage{50}the character of a genuine thinker who passionately wants to
fight his way forward to the sources of truth; but there is something
Tantalean in this striving of his, in that he never succeeds
in reaching a standpoint that satisfies him, and
which he can then develop and carry through completely.
He is right when he characterizes himself as
a ``dissatisfied spirit.'' --- The restless, never satisfied
yearning gives his language a peculiar beauty;
in a series of turns of phrase, antitheses, images and comparisons
he seeks to exhaust what he cannot
hold fast in a sharply closed conceptual form. In this he finds
the unity of his work as an author expressed. ``You
thought as a philosopher,'' he says to himself, ``but you
did not write as a philosopher; you always transformed the being of thought,
when you uttered it, into a being of flesh and
blood. You knew that philosophy as such, mere
reason, pure thought, is nothing for man,
that one can convince man of a truth only
when one makes it, from a being of reason, into a being
that is like man, a sensuous being\dots{} Everywhere
you have joined the abstract to the concrete, the
non-sensuous to the sensuous, the logical to the anthropological.''
(Werke I. p. VII.---IX.). The development consists,
then, in this: that what at first was only a form,
a clothing, at last becomes the essence itself. The linguistic
expression thus enters into an ever more inward unity
with the content: ``Formerly you thought, and said it
too, if not literally, then in fact: the true must
be present, real, sensuous, intuitable, human;
now you say, consistently, the reverse: only the
real, the sensuous, the human is the true.''

In Feuerbach's philosophical development we can
distinguish three different stages, if we pass over
what lies before his public appearance
(his orthodox-theological standpoint). These three
% --- p. 51 ---
\setcounter{footnote}{0}%
\opage{51}stages we characterize in general before we go on
to occupy ourselves more closely with each single one of them.

The first stage is pantheistic. The movement
here goes absolutely in the direction of the universal, the objective,
the Idea. Everything personal and individual is sought to be sublated
as imperfection. By being united with subjectivity
the purity of the Idea is violated; it no longer comes
to hold in and for itself. But of the objective reality of the Idea
there is no doubt; the Idea is, on the contrary, that which alone truly
is. As that which has validity in and for itself, it can
be apprehended and grasped only by theoretical knowledge.
Theory is the truth. Religion and theology
miss the truth because they place themselves at the practical
standpoint, where finite, personal motives
hold, and not the matter in and for itself. Self-love,
the drive for blessedness, is already here designated as the
origin of religion, and a corrective is thereby applied
against the speculative philosophy of religion, which regarded
religion only as an imperfect form of theory. ---
To this first stage (1830---1838) belong the following
writings: Gedanken über Tod und Unsterblichkeit; Geschichte
der neuern Philosophie von Baco bis Spinoza;
Darstellung und Kritik der Leibnitzischen Philosophie;
Pierre Bayle.

The real moment of which Feuerbach had thus become
aware in his investigations in the philosophy of religion
gained increased significance for him through the zealous
study of the natural sciences to which he now devoted
himself. He thereby became aware of the necessity and authority
of sensuous experience, and gained an inward
conviction of the reality of sensuousness, which
was bound to overturn the validity of the Idea as the one true
being. From a being as objective substance the
Idea was now referred to subjective being as a concept of the species.
Man develops universal concepts by virtue of the consciousness
of the species, whose essence they express. The Idea
% --- p. 52 ---
% print: Wesen der Christenthums (misprint: der for des)
\setcounter{footnote}{0}%
\opage{52}is no longer a being of reason, but being for reason,
reason itself. As such the universal, the Idea,
still has validity over against finite empirical experience. Reason
is seen as man's true I, which gives the laws under
which drive and feeling are to bend. But just as the
objective in general is only a reflection of the subjective,
so theology is anthropology, and religious
representations are productions of man's overflowing,
insatiable yearning for the satisfaction of the life of feeling.
--- This stage is represented above all by
Feuerbach's chief work ``Wesen der Christenthums'' (1841).

The third stage is introduced by the declaration (Philosophie
der Zukunft. 1843) that whereas he had formerly
still put the sensuous and natural in a subordinate
place, as that which was to subordinate itself to reason,
to the knowledge of the universal, he now
declares: only the sensuous is the real. Not reason,
but the body in its totality is the true I.
The universal is only an illusion; the individual drive,
in its sympathy as well as in its egoism, is the true
motive. --- In ``Wesen des Glaubens im Sinne Luthers''
and ``Wesen der Religion'' religion is explained as having sprung
from the egoistic drive for blessedness, an explanation
whose extreme carrying-through is given in the work
``Theogonie''. Finally, in his last work,
``Gottheit, Freiheit, Unsterblichkeit'' (1866), Feuerbach goes right to the boundary
of materialism, in that he dwells with particular relish
on the proposition: ``der Mensch ist, was er isst.''
What holds him back is the conviction of the
inner activity of the human being. Whereas he formerly
regarded the psychological motives from which religion,
in his opinion, sprang as pernicious,
he now himself builds his ethics on them. The standpoint
which at his first stage he designated as the
practical one he has now himself occupied, after having driven
religion out of it.

% --- p. 53 ---
\setcounter{footnote}{0}%
\opage{53}After this preliminary survey we can designate as the
general law of movement which governs the whole
the law of subjectivization, i.e.\ the law that every objective
being which is posited as the object of spirit is, precisely thereby,
spirit's own work, which therefore, when spirit
becomes conscious of itself, is taken back again, in that the positing
force reflects itself into itself and recognizes its
object as having sprung from its own essence. This
law we shall now find applied in the particular.

a. Feuerbach's first standpoint is purely metaphysical.
We have already mentioned his ``Gedanken über Tod
und Unsterblichkeit'', where everything personal is sublated in the
universal, over against which it could only egoistically want
to assert itself. This work can be designated as a
rehabilitation of death: death is the expression of the
nothingness of the individual, is the absolute victory of universal
being. --- The same speculative-objective standpoint is
pervasive in his historical writings from this period.
In his ``Geschichte der neuern Philosophie'' he praises
Spinoza and shows, through the development of the ontological
proof, that God's existence is not to be separated from
his essence; as God's essence, so also his
existence falls within reason, which is itself the highest
being. It is the Hegelian logic that is still
held fast.

The break with this panlogical view is introduced
by intensive occupation with real
actuality, and first in the domain of the philosophy of religion.
Feuerbach immersed himself with zeal in individual
religious phenomena, in religious psychology;
hence the wealth of characteristic and individual
traits which is found in his writings. It was this individual element
that interested him, whereas the speculative
philosophy of religion contented itself with a survey of
the general, objective content of the dogmas, without regard
to the psychological soil from which it had
% --- p. 54 ---
% print: usyndig (misprint: usyn-|dig, for usynlig)
\setcounter{footnote}{0}%
\opage{54}arisen. Feuerbach soon saw that it was more than a
difference of form that separated religion and philosophy. He
declared it to be the πρῶτον ψεῦδος of speculative philosophy to regard the drive for knowledge as the psychological
source of religion. In his work on Leibnitz he refers
religion quite generally to the practical domain,
because it owes its origin and its power to practical
motives, hope and fear. Philosophy is not,
as Hegel thought, to seek to find out the rational concept
of the content of religion; dogma is incomprehensible because
it is irrational; the only task is to explain
how it has arisen. In the treatise on Bayle
we find the further development of this.

Theology, it is said here, has the miracle for its principle,
in that it proceeds from an arbitrary, unconditioned will
as the cause of existence. The question of the ground of
things therefore does not exist at all in the theological domain;
it is disposed of once and for all with an answer
which is no answer; the will is the asylum ignorantiæ
by which all science is turned away. The miracle cancels all reason
and experience. In the miracle one sees only the two extremes;
one sees, e.g., wine where one saw water before;
but the transition, the actual event itself, is invisible
and can be experienced only with the eyes of faith. The real
miracle therefore has no advantage over the invented one; for what is
essential in both, the supernatural causality, is put
in by the apprehending subject and springs solely from
its personal need. Faith in the miracle is the essence of
the miracle. --- In contrast to this, philosophy,
as the representative of the idea of science, proceeds from the matter
in itself and seeks to find the right grounding from the
nature of the matter. The principle of philosophy is reason, the mother of
lawfulness and necessity.

When the question is the origin of Christianity,
theology derives it from a supernatural
act which, passing over all conditions, acts
% --- p. 55 ---
\setcounter{footnote}{0}%
\opage{55}immediately at this particular point in history. But
that is to say, in other words, that theology does not recognize
this question as a problem at all. Philosophy,
on the other hand, first seeks to determine the nature of religion
in general and the general laws under which
it appears in the human spirit and in history. These
laws we find in the representations springing from the nature
of religion, on whose diversity the difference between
the particular religions rests. The content of a religion can
very well be irreligious, i.e.\ stand in conflict with the true
essence of religion. The other factor philosophy must take into
account is the time when Christianity arose. Under the
unhappy conditions of the age, when everything that had
borne classical antiquity was perishing, the mind could see through
the nothingness of all finite circumstances and freely rise
above them; amid the general moral corruption
a pure, quiet, inward heart, precisely through its opposition
to its surroundings, could not but be thrust back into itself,
so as to find in the inner world what the outer denied.
Only in such a time could the idea of morality be grasped in
its purity --- and that is the only positive thing in Christianity,
when one keeps its final purpose in view.

What Feuerbach here calls the positive in religion
is not, however, what makes it a ``positive religion.''
In the difference between true and false religion the content
is decisive. This content is the concept of God,
the way in which God is thought and is an object for consciousness.
But the concept of God is essentially the object of philosophy
and morality, which religion makes the object
of a cult. What makes religion religion,
its difference from morality and philosophy, religion
in itself, is therefore not yet good, holy and divine;
it is so only when its content is a divine one.
Holiness is the peculiar religious category.
But holiness is not a primitive concept;
the holy must first be recognized as rational and good
% --- p. 56 ---
\setcounter{footnote}{0}%
\opage{56}before it can have validity. Only the true is holy,
but the holy as such is not yet true. When
the holy is separated from this, its justifying ground,
and developed for itself, religion becomes positive, ``a
religion secularized into a church and bounded by its formulas of faith,''
``a concern distinct from reason and morality.''
The cause of the church becomes the cause of religion; when
man only believes what the church teaches, and performs
the holy acts which it enjoins, then he is
saved. The Catholic Church is the only successful definition
of a positive, ecclesiastical religion; it has most
faithfully, most consistently realized its essence, but
precisely it has also expressed most sharply the indifference
of the religious or ecclesiastical interest toward the interests
of morality. Positive religion regards
God as a single being toward whom one has special
duties, which fill the heart with fear and anguish and
let the humane, essential and universal relations of duty recede
into the background. Positive religion necessarily becomes the
affair of passion. Anguish at the loss of salvation drives to faith.
The craving of the drive for self-preservation and for blessedness overwhelms
reason and subdues the capacity for pure and unbiased love of truth.
Faith, it is true, renounces reason and what is higher
in man, but not what is lower, the egoistic motives.
The positively religious therefore stands in an external
relation to the ethical. Virtue is made not the principle
but a side issue: faith, it is said, makes one blessed, but
faith cannot be without works --- just as the Epicureans
said: pleasure makes one happy, but pleasure is
not without virtue.

As virtue is subordinated to faith, i.e.\ to the drive for salvation,
so the good is subordinated to the divine will,
represented as an arbitrary, personal will. The good has
for philosophy its ground in itself, but for theology
the good is good only because God wills it. If one objects
to this that God wills only the good, and that his
% --- p. 57 ---
\setcounter{footnote}{0}%
\opage{57}will is therefore no arbitrary will, then one has
precisely thereby recognized that the Idea of the good is
the eternal a priori which determines the will and first makes
it good. Only by virtue of the Idea of the good can one think
the good as a predicate. It therefore comes to the same thing
whether one makes God, as the good, or the good in
itself the principle of the ethical; God is here only a name,
the essence is the concept of the good. Yet this name ought to be left out,
so that the good shall not merely be seen as a perhaps
accidental predicate of a presupposed subject, but can
be seen in its absolute independence. Kant's categorical
imperative was the manifesto in which ethics proclaimed
its freedom and independence to the world. If Kant
still reached only the bare necessity of duty, and did not get
beyond the elementary doctrine of ethics, Fichte rose
to ideas so exalted that positive religion has
nothing to set beside them. Fichte's ideas are, it is true,
severe, almost superhuman, but ethics is not
pedagogy either; it does not accommodate itself to
human weaknesses and needs. The moral
spirit has vanished where the Idea of the good is not thought
free of all personality and loved and realized
for its own sake. ---

Feuerbach's chief endeavor here aims at
asserting the validity of the Idea. Although he apprehends
it not merely as an abstract metaphysical Idea but in a more concrete
form, as the ethical Idea, which is the unity of ethics and
metaphysics, of will and knowledge, he nevertheless rejects
the intrusion of any subjective element. The method
rests on the law of subjectivization stated earlier:
he seeks to discover the original working of consciousness
by which everything else is grounded. Here he then finds that it
is the ethical Idea by virtue of which all will is good,
whether it is thought as divine or as human
will; consequently, he concludes, the ethical Idea is
what is original, the good will what is derived. And as
% --- p. 58 ---
\setcounter{footnote}{0}%
\opage{58}regards the Absolute, the motive then falls away for
thinking of a special subject distinct from the Idea.
It is only man who, as a sensuous being, does not
coincide with the Idea, but precisely for that reason the Idea appears
to him as the ethical law. In this law is contained
the true religion, while positive religion
has its ground in the fact that the human drive for blessedness
will not subordinate itself to the conditions of the Idea.

Through the weight Feuerbach has here laid on the
real moment, in both an ethical and a religious respect,
his judgment on speculative philosophy was bound to be modified.
He still praises it as Hegel's greatness
that he investigated the logical and metaphysical fundamental concepts
in and for themselves, without presupposing a subject.
But he nevertheless declares that the time of systematic speculation
is past; what is needed is a critical-genetic philosophy.
Immanent speculation is not able to develop
real actuality; thought must work this over
critically, so as then --- genetically --- to show its origin
and through that its essence.

But that Feuerbach did not rest content with the view
he had developed in ``Pierre Bayle'' --- the
reason for this must be sought in the indeterminacy with which the relation of the ``Idea''
to the empirical is conceived, particularly in the respect of the philosophy of
religion. Either the Idea is what is original ---
and Feuerbach must presuppose that here --- but how
will he then explain the empirical distortions
of its essence which he finds in the positive religions?
Or else the empirical is what is first, but then
a change must take place in the determination of the ideal,
since this is no longer valid by itself but is conditioned
by the empirical.

These questions furnish the motive for Feuerbach's subsequent
works, in which he enters upon a new stage.

b. The transition to this is formed by some smaller
treatises in which the fundamental thoughts already come
% --- p. 59 ---
\setcounter{footnote}{0}%
\opage{59}forward that are dominant in Feuerbach's chief work,
``Wesen des Christenthums.'' Of these treatises we single out
in particular ``Kritiken des modernen Afterchristenthums''
and ``Philosophie und Christenthum.'' None of man's
essential faculties --- so the line of thought runs here
--- can go beyond itself and grasp something that is wholly
foreign to it. Thinking cannot grasp a thought
that is not its own: precisely by grasping it, it would
indeed prove that it was in reality its own
thought. It is the same with feeling. If
feeling were able to go beyond human
thought, it would also have to go beyond the essence of man,
and so be no human feeling. The philosophy of religion
must therefore recognize and treat religion
as psychology. The opposition between philosophy
and positive religion rests on the opposition between
thinking and the heart on the one side and imagination
and the life of feeling on the other. The heart is the natural,
healthy life of feeling that makes itself one with reason.
The life of feeling, on the other hand, is the sick, supernatural heart, which
disdains determination and limitation, and for which
thinking is only a finite activity, whereas
it takes refuge in imagination, which abandons itself to
infinity without letting itself be governed by rational
and natural laws. The life of feeling and imagination, in faith
in an otherworldly, heavenly infinity, lift man out
beyond all finite relations, but thereby also dissolve all
the bonds of love and humanity.

In ``Wesen des Christenthums'' these thoughts are developed
further in a deep inner unity and coherence. --- Man
is not driven, like the animal, blindly and immediately by
the essence of the species; rather, in consciousness this becomes an object
for man. Since the essence as such is infinite,
consciousness must also be so. And yet
consciousness is won only through the relation to an object, hence
through a limitation. But this object has no
% --- p. 60 ---
\setcounter{footnote}{0}%
\opage{60}significance in and for itself, apart from the force that posits
it. Since the essence cannot go beyond itself,
the object to which the subject stands in an essential
and necessary relation can be nothing other than the subject's
own essence, objectively expressed. The object is
the subject's manifest I; its power over man is
the power of his own essence. The power of the object of reason is
reason's own power; the beloved object rules
with love's own power. He who thinks the
infinite thinks and confirms the infinity of thought;
he who feels the infinite feels and confirms the infinity
of feeling.

Feuerbach states more precisely the relation of these principles
to earlier views in the philosophy of religion, in particular
Schleiermacher's and Hegel's. Schleiermacher, this
``last theologian of Christianity,'' had made feeling
the religious organ. But if religion is a matter of feeling,
then the content of religion becomes indifferent; it becomes
\emph{religious} content only by being the object of feeling,
and the essence of feeling itself then properly becomes the
divine. It was only theological prejudice when
Schleiermacher thought he could find a way from feeling
to the objective doctrine of faith. --- Hegel, on the other hand, determined
religion as the stage of consciousness at which man's
knowledge of God became one with God's knowledge of himself.
But here again it holds: the organ is one with the object,
the means is the essence itself. The knowledge of God
is man's self-knowledge; man's knowledge of
God is man's knowledge of himself. The mystical
relation is a relation to oneself.

The result, then, is: God is man's manifest
inner self, his expressed self, the public confession of
his secrets of love. Yet this holds only
for philosophical consciousness. Religious consciousness
is still only man's \emph{indirect} self-consciousness,
i.e.\ in religion man becomes conscious of his own
% --- p. 61 ---
% print: er virkeligt, (misprint: er for et)
\setcounter{footnote}{0}%
\opage{61}essence as of a foreign essence, first places
his essence outside himself before he finds it in himself.
Historical progress in the history of religion consists
in this foreignness being sublated, in self-consciousness
recognizing itself more and more in its object. Every
advance in a religious respect is therefore a deeper
self-knowledge. Yet it is peculiar to
every positive religion, even if, like Christianity, it
makes man's own essence its content, that this
content stands as something external, as something distinct from self-consciousness
itself. This imperfection is first removed when
one passes over from the sphere of religion into that of philosophy.

It is admitted, to be sure, that the divine attributes,
on account of the imperfection of human knowledge,
are borrowed from the human essence and
transferred to the Godhead; but it is nevertheless held fast
that the divine subject itself is not thereby humanized.
But that is a distinction which religious consciousness
does not make and cannot make. For it God
is a real, living being, not that which is in itself without determination
and lies behind the attributes. And besides ---
where is the boundary? If the attributes are anthropomorphisms,
then the subject is an anthropomorphism too. If
love, goodness and personality are anthropomorphisms,
then existence itself is also an anthropomorphic
determination. Rather, the case is the reverse; it
is the attributes, the predicates, that have validity in
themselves, and the subject is only the existing predicate.
The real atheist is he for whom the divine
predicates, love, wisdom, justice,
are nothing, not he who merely abolishes the subject
of these predicates. The subject shares the fate of the predicates,
not the reverse.

It follows of itself that man, who transfers
essential predicates from himself to the Godhead and
thereby first properly produces it, cannot go
% --- p. 62 ---
\setcounter{footnote}{0}%
\opage{62}on ascribing the same predicates to himself, at
least not in the same degree and significance. He must
necessarily deny in himself what he posits in his
God. Hence the remarkable fact that the more human
a being the Godhead becomes, the greater becomes
the difference between God and man. Man plunders
himself to enrich his God: the more God comes
to resemble man, the less man comes
to resemble God. Man denies his knowledge
in order to see in God the omniscient one, his blessedness in order to see in
God the absolutely blessed one, his good nature in order to see in God
the absolutely holy one.

This, however, is still only the one side of the essence of religious
consciousness, which one could call the
religious arterial system or the religious force of repulsion.
Man here thrusts his own essence away from himself and makes
it an objective being. The other side is that the force of repulsion
is sublated by the opposite force, attraction,
that over against the arterial system there stands a venous system,
in which the blood flows back to the heart. Religious consciousness
does, it is true, raise God infinitely high above
itself, but it nevertheless sees itself as the essential
object of the divine activity. In the object
the essence lies expressed; as that which has man
for the object and end of its working, the Godhead is
thus after all a human being. The truth of religion
consists in its setting man in relation
to his own essence; its untruth and limit consist
in this essence being apprehended as a foreign, indeed opposed,
being --- herein lies the source of religious fanaticism,
``the metaphysical principle of the bloody human sacrifices.''

If religion springs from the essence of man
and yet sets up an opposition between God and man,
then the germ of this opposition, with which
religion begins, must be demonstrable in man's own
% print: Spise (misprint)
essence. This germ lies in the psychological
% --- p. 63 ---
\setcounter{footnote}{0}%
\opage{63}opposition between understanding and heart. The understanding is untouched
by the sufferings of the heart, has no drives and passions,
no needs, and therefore also no
defects and weaknesses, as the heart has. While
the heart turns toward the special, the individual,
the understanding is the representative of the universal, the
pure, passionless light of intelligence. God as the metaphysically
and morally perfect being is therefore the objectified
essence of the understanding. But over against this man feels
himself as the finite and sinful one. The feelings of the heart do
not come into their own. Therefore the heart represents to itself,
on the foundation of that abstract being of the understanding, which is only
the mathematical point of religion, God as infinite
love. Thereby man is reconciled with God.
The deepest view of divine love
even lets God himself become man. The incarnation is
the sensuous revelation of God's human nature;
God felt the human need.

But this self-doubling is not a harmless illusion.
This shows itself in the fact that, although the Godhead is seen as
love, love is nevertheless not its essential determination,
not valid in and for itself, but is attached to
a presupposed divine subject, which as it were lurks in
the background of love and is still something by itself apart
from love --- ``a demonic being whose personality, separable
from love and really separated from it,
delights in the blood of heretics and unbelievers, the phantom of religious
fanaticism.'' Whence, then, come these pernicious
elements, this untruth in the essence of religion?

The relation to God as to another being is partly
a natural, involuntary and unconscious relation, partly a
conscious and reflected one. The first has its root in the very
origin of religion and rests on its essential
standpoint, the practical. Man's relation to God
is one with the relation to his own salvation: God is himself
salvation realized and the unlimited power
% --- p. 64 ---
\setcounter{footnote}{0}%
\opage{64}to realize man's salvation. God as the object of religion
--- and only as such is he God
--- is the object of the life of feeling, not of reason, an object
for practice, not for untroubled theory, for the heart's distress,
not for freedom of thought. Reason, the universal,
acts on practical man only through the representation
of a personal being. In religion the light of spirit
is refracted in the medium of imagination and the life of feeling,
whereby one and the same being is made intuitable as
doubled. About what lies behind this practical consciousness
religion does not concern itself. Immediate
religious consciousness rests with childlike security in
this unity of the divine and human essence.
But the more reflection awakens, the more effort is spent
on sublating the unity, because one surmises that it
might in its consequences lead to sublating the religious
object as such. The difference which at
the immediate standpoint is only the refraction of one
and the same light is now fixed into a dogmatic opposition.
This endeavor runs through the so-called
proofs of the existence of God. They seek to establish that
religion is right when it posits man's essence
as something foreign and objective. God's existence is then
apprehended here as a special determination that is ascribed to the
religious object, and by which this is placed outside consciousness.
God is not merely to be in faith, in feeling,
in the inwardness of the life of feeling, but to have a being distinct
from them and independent of them. But a real being, distinct
from thought and independent of it, can only be a sensuous
being. Only by ascribing such a being to God can
one maintain his existence ``outside us.''
Theology
seeks to avoid this by declaring that this ``outside'' is not
to be taken in its literal sense. But in what sense,
then? A spiritual being is only a being in thought,
not at all outside us. Reflection therefore leads
necessarily to atheism, in the sense that God's existence
% --- p. 65 ---
\setcounter{footnote}{0}%
\opage{65}is negated as an external, empirical being independent of
thought. But the truly divine is, as we have already
seen, not this bare and mere existence, but
is contained in ``the divine attributes.'' If theology
will not take this step, there is nothing else
for it to do than to go back to the immediate
religious stage of consciousness, which does not distinguish between inner
and outer. By the power of imagination, which is the intellectual faculty
corresponding to the life of feeling, an external existence is posited
which is nevertheless not external, a sensuous being which
is nevertheless not sensuous. This hovering between the sensuous
and the non-sensuous is precisely the essence of the power of
imagination. Only imagination preserves from atheism.

The anguish and passion with which man, as
reflection awakens, holds fast the objective existence of the religious
object stand in close connection
with the fact that it is not only what is true, universal and justified
in man's essence that is made into God, but
also what is individual, arbitrary and finite. This
two-sided essence of religion Feuerbach finds expressed above all
in Christianity's doctrine of the relation between
faith and love. Love makes man one
with God and determines God as a human God.
Faith separates man from God, in that it posits him
as distinct from and raised above the humane determinations of essence.
This brings it about that love is limited,
cannot move freely; God is love, it is true, but love
is not God. Love has God within itself, but
faith has him outside itself, and in
Christianity love is made dependent on faith. In that God is separated
from man, one man is separated from another,
for God is only man's mystical concept of the species; his
separation from man is thus the dissolution of the common
bond. Faith is in its essence exclusive: this definite thing
you shall believe if you would not lose your salvation.
The believer has God for him, the unbeliever has him
% --- p. 66 ---
% print: under hvilken (misprint: hvilken for hvilket (Billede))
\setcounter{footnote}{0}%
\opage{66}against him; faith is goodness, unbelief wickedness and obstinacy.
The commandment that one shall love one's enemies does not hold
toward God's enemies. Therefore there lies in faith
an evil principle: ``The flames of hell are only the sparks
of the annihilating glance, glowing with wrath, which
faith casts on the unbelievers.'' In the essence of love,
on the other hand, lies universality. True love
is sufficient to itself and demands no special title or
authority. It recognizes virtue in sin itself, truth
in error itself. In love the essence of the species is realized
in a subjective way, through disposition.
It is therefore in itself one with reason, which is the objective
reality of the species. In opposition to reason and to the
heart, which is akin to it in essence, stands the life of feeling, the
self-feeling of individuality as such. Christianity
is the contradiction between heart and life of feeling, because it is
the contradiction between love and faith: faith comes
from the life of feeling, love from the heart. The figure of Christ
is the image under which the unity of the species pressed itself
upon the consciousness of the people. The essence of the species, which manifests itself
in the energy of love, overcomes every
particularity, every limit. Where, therefore, the consciousness
of the species as species arises, there Christ
vanishes, without his true essence perishing.

In this relation between faith and love
Feuerbach sees the clearest expression of the necessity of
going beyond the religious. If it is established that theology
is anthropology, that God's essence is man's
own essence, then it is necessary to let it come
to the turning point at which this clear self-consciousness replaces
the earlier religious consciousness, foreign to itself.
And that not merely for theoretical reasons. For
the self-doubling in religious consciousness is not a
harmless illusion. In the first place the life of feeling, the
real psychological factor and source of religion, is the
mind at odds with nature and therefore yearning beyond nature,
% --- p. 67 ---
\setcounter{footnote}{0}%
\opage{67}which can see its wishes satisfied only in a miraculous way.
In opposition to this it must be maintained that
nature belongs to the essence of man and has true,
real validity and significance in itself. Next,
the ethical must be released from its limits, no longer be derived
from faith. The right, the true and the good always have
the ground of their holiness in themselves. Ethical relations
are in and for themselves religious relations. Life is altogether
of a divine nature and needs no
special consecration. ---

What in this chief work of Feuerbach's particularly attracts
attention is the way in which
religion is apprehended. In that it is traced back to the life of feeling,
the source of unlimited wishes, and is determined as the
purely practical standpoint, where fear and hope are the
leading motives, the opposition to Hegel's intellectualist
apprehension of religion stands out clearly. Feuerbach
has himself declared that his doctrine is wrongly regarded
as a further development of the Hegelian;
it arose precisely out of opposition to it. What
for Hegel was the secondary and merely formal, the heart
and feeling, into which the objective content of religion
descends in order to become man's property, that
is for Feuerbach precisely the essential and primitive.
He stands closer, on the other hand, to Schleiermacher, who made
feeling the religious organ; and while Hegel polemicized
against Schleiermacher because the latter took so subjective
a point of departure, it is, as already stated,
Feuerbach's objection to Schleiermacher that he is not
subjective enough, did not carry through his subjective
point of departure absolutely.

The advance in relation to the treatise on Bayle
consists precisely in the more determinate application of the subjective
point of view to positive religion. This is seen,
e.g., in the treatment of the concept of miracle. Whereas
in ``Pierre Bayle'' he establishes the nothingness of the miracle by
% --- p. 68 ---
\setcounter{footnote}{0}%
\opage{68}demonstrating its incongruity with the divine
essence, in ``Wesen des Christenthums'' he regards it
as the supranaturalist satisfaction of a human
wish, and applies to it a subjective-ethical, no
longer an objective-metaphysical, critique.

The individual life of feeling from which positive religion
springs is morbid and unjustified, because it
seeks to raise itself above the universal law of essence of reason. The
individual is thus still definitely subordinated to the universal.
While the life of feeling sacrifices the species to the individual,
reason sacrifices the individual to the species. But the universal
is no longer apprehended, as in ``Pierre Bayle,'' under
the form of the objective Idea. The Idea is now seen as man's
own essence, as the law of the species in man, albeit
above man. The universal keeps its reality, only
this is transferred from objective being to subjective.
Feuerbach declares expressly that in the critique
of religion it is necessary to rest on the
universal: ``only the species is able at once to abolish
and to replace religion.''

But his restless spirit drove him ever further. By
virtue of the method of subjectivization he necessarily had to
raise the question whether the \emph{subjective} being of the universal
too might not be an illusion, which
like theology dissolves into pure anthropology. Thereby
the way would then be paved for overcoming the opposition
that runs through his earlier writings,
the opposition between the universal and the individual.
This is the task whose solution Feuerbach seeks
at the third stage of his work as an author.

c. Feuerbach himself characterized the stages of his development
thus: ``God was my first thought, reason
my second, man my third and last thought''.
(Sämmtl. Werke. II p. 410). This last stage, however,
he has not yet quite reached in ``Wesen des Christenthums'',
although he there means to carry through the proposition
% --- p. 69 ---
\setcounter{footnote}{0}%
\opage{69}that theology is anthropology. For he there still upholds
reason, the consciousness of the universal essence of the species,
as that which is indeed in man, but yet \emph{above}
man. There is still a moment of transcendence
left in immanence. Feuerbach is to that extent still
an idealist, since in determining the essence of man he
proceeds from the consciousness of the universal and derives
the ethical principle from it.

Yet there is already to be found, in an essay in ``Deutsche
Jahrbücher'' for 1841 (included in the 2nd volume of his works):
``Ueber den Anfang der Philosophie, a proclamation
of realism. The beginning of philosophical and empirical
knowledge is here declared to be one and the same act.
Philosophy is not merely to appropriate the results of empirical inquiry;
it cannot arrive at reality except by
beginning with reality. It must see in empirical
activity a philosophical activity, in seeing a
thinking, in the organs of sense the organs of philosophy.
Psychology is not a special science, but the
first general science, which deduces the particular principles
from the various relations of the I. To the I,
then, belongs the body as well, which as the passivity of the I has
not merely a natural-historical or empirical, but also a speculative
and metaphysical significance.

What is here intimated is developed further in the work ``Grundsätze
der Philosophie der Zukunft'' (1843) in the following
way. The task of modern times was to develop
theology into anthropology. This task Protestantism solved
from its practical standpoint, in that it
considered only what God was as man and for
man; in a rational, theoretical way the task has been solved
in modern speculative philosophy, whose consummation is
the Hegelian philosophy. The essence of God was here developed
in such a way that it became one with the very essence of reason,
since the divine determinations of essence (``attributes'')
proved to be identical with those of reason, and in particular
% --- p. 70 ---
\setcounter{footnote}{0}%
\opage{70}the determinations of essence of speculative reason. Infinity,
necessity, immutability, eternity belong to
the eternal truths and laws of reason, and consequently
to reason itself. In reason the thinking and the
thought are one. Whereas, therefore, ordinary consciousness
distinguishes between God as thinking and as thought by
man, these two sides now coincide as
identical. The divine thinking is only a represented
thinking, which has its reality in the a priori knowledge
of speculative philosophy; and if one also
thought of God as the one who knew everything sensuous, then
this divine knowledge has found its realization
in the empirical sciences of modern times.

But in transforming God into
reason, modern philosophy let reason itself take on divine
attributes. Pantheism is a theological atheism, a
negation of theology, but within theology's own domain.
This shows itself above all in the abstraction of speculation,
which is the theoretical after-effect of medieval
asceticism. Hegel has therefore only proved the divinity
of the understanding in its opposition to sensuousness,
the world, man himself. If the beyond of theology has
become a here-and-now, then in return the
real world itself has become a beyond.

The new philosophy, the philosophy of the future, must therefore
stand in sharp opposition to speculation and
first and foremost correct the concept of being. Being is
no general concept that can be separated from existing
things themselves. Being is one with that which is.
When Hegel in the Phenomenology shows that, since every thing
is a ``this here'', the particular is immediately the
universal, he has proved nothing but the contradiction
between the word, which is universal, and the thing, which
is always particular. Beyond this contradiction thought,
which must always rely on the word, cannot get. But
the question of being is not merely theoretical; it is
% --- p. 71 ---
% print: Spørgsmaale (misprint: for Spørgsmaal)
\setcounter{footnote}{0}%
\opage{71}a practical question, ``eine Frage auf Tod und Leben''.
By its opposition to the universal word, being becomes
unutterable. Life, the secret of being, begins only
where words stop. If, therefore, unutterability is the
same as irrationality, then all existence is irrationality,
precisely because it is this determinate existence; --- but
it is not so: existence has in and for itself
meaning and reason, although it is unutterable.

The new philosophy thus takes the concrete particular in earnest,
which speculative philosophy believed it let
come into its own when it took up particularity as a
moment in the concept. Not as a moment of the concept, but
as a true and real being, and that is to say: as
a sensuous being, is the real to be the object of
knowledge. Truth, reality and sensuousness
are one. Only through the senses is an object \emph{given}. Being
can be apprehended only by being, by existence, by sensing,
intuition, feeling, love: ``das Sein ist also ein
Geheimniss der Anschauung, der Empfindung, der Liebe.''
Love is passion, and only passion is the mark
of existence. Abstract, passionless
thinking sublates the difference between being and not-being.
Whereas earlier philosophy said: what
cannot be thought, is not, the new philosophy says:
what cannot be loved, is not. Only what can be
loved and worshipped and can thus be an object of religion,
only that can also be the object of philosophy.

Everything can be perceived in a sensuous way, if not immediately,
then mediately, if not with the crude and
vulgar senses, then with the cultivated ones, if not
with the eyes of the chemist and the anatomist, then with those of the philosopher.
Nor is the sensuous that which lies
ready to hand, the thoughtless, which is understood
of itself. It is only through a long and difficult
cultivation that man is educated to the intuition of
the sensuous; immediately, the beings of imagination, of
% --- p. 72 ---
% print: væsenslige (misprint: for væsentlige?)
% print: Philophiens (misprint: for Philosophiens)
\setcounter{footnote}{0}%
\opage{72}representation and of the abstract understanding lie far
nearer to him. Intuition and thinking must determine
each other mutually. Intuition by itself has no principles,
and thinking in and for itself does not get beyond
logical unity with itself. The movement of thought
must be interrupted by sensuous intuition and can
therefore not, as Hegel would have it, close itself into ``a circle
of circles''. The symbol of the new philosophy is not
a circle, but an ellipse. Yet it is no alien
power that determines the movement of thought, for the laws
of reality are thought's own laws; thought does
not form a state within the state.

The new philosophy makes man its principle of knowledge,
not reason. Man is the reality
and subject of reason; it is man who thinks, not
reason itself. Whereas earlier philosophy said:
``only the rational is the true and real'', the new
philosophy says instead: ``only the human
is the true and real.'' But by the human one
must not think merely of the understanding. It is not
merely by thinking that man distinguishes himself from
the animal, but by his whole essence, of which thought is only
a single consequence. Since the human, including
nature as the basis of man, now becomes
the sole object of philosophy, anthropology becomes
the universal science.

But the individual man does not have the whole essence of man
in himself. Only in the community and
unity of one man with another man does the essence of man
reveal itself. Only the unity of I and Thou is God. It
is this truth that is expressed in the dogma of the Trinity,
the truth that no being in and for itself is true and
perfect, but that truth and perfection
are only the unity and connection of essential beings.
The last and highest principle of philosophy is therefore the
% --- p. 73 ---
\setcounter{footnote}{0}%
\opage{73}unity of man with man. From this unity spring
all essential relations and the principles of all the various
sciences.

Whereas earlier philosophy had a double
truth, a theoretical and a practical, in that it set
itself up as truth in and for itself in opposition to religion,
truth for man, the new
philosophy overcomes this opposition: it has not only
the same object as religion, but also the essence of religion
in itself, is itself the true religion. ---

``The new philosophy'' is, according to Feuerbach, to be
different in essence, toto genere different from the earlier one,
and the decisive opposition is to rest on its
recognizing sensuousness consciously, which speculation
had done only indirectly by seeing the real enclosed as a moment
in the highest principle. The dissolution
of absolute idealism had of necessity to
lead to a point where attention was drawn
to the positive elements that had been misjudged
and overlooked by speculation. Feuerbach represents
precisely the stage in the development where the negative momentum
brings the critic face to face with the positive
and real. But precisely his negative and critical line of thought,
which is itself to a high degree influenced by the theological
and speculative opponents it has fought, makes
it difficult for him to develop the new
elements in a positive way. He has rather beheld the new with enthusiastic,
almost tumultuous feeling than clearly and stringently
stated fruitful principles for its treatment,
rather set up a great demand than contributed to its
solution. Add to this that his restless, ``dissatisfied
spirit could not avoid self-contradictions, since it could not
bring itself to clarify logically what it grasped with mystical
impetus through the imagination. ``His philosophizing
remained divination. His ``
% --- p. 74 ---
% print: nn (misprint: for nu)
\setcounter{footnote}{0}%
\opage{74}sensuousness'' is a new form of absolute thinking, which
wholly disregards practical experience''\footnote[1]{\emph{Alb. Lange}: Geschichte des Materialismus. p. 284. 306.}.

The principle of knowledge is to be man himself in
the whole extent of his essence. Whereas earlier philosophy
sought the ultimate grounding of knowledge in the metaphysical
Idea of knowledge, which displays its essence through
man's psychologically conditioned knowledge, Feuerbach makes
psychological-human knowledge one with
metaphysical knowledge. The species in its multiplicity is
now to possess what was formerly concentrated in the Idea
of knowledge. But what reality can Feuerbach ascribe to the species
as such? If it cannot gather its multiplicity
into a unity, knowledge becomes only an aggregate,
whose fragments exist for the representation of each individual,
without the whole being able even, as Kant would have it, to
become a guiding rule and purpose. In reality there remains
only the knowledge of each individual man. And according to
Feuerbach the individual man knows only what he
has apprehended through sensing. The new philosophy,
he says, is with respect to its basis nothing other than
sensation raised to consciousness, the heart that has been
brought to understanding. The standpoint thus becomes complete
sensualism. Feuerbach says that he will lead
thinking back to unity with experience; but it
becomes in reality a complete dependence on
experience. When he says that everything can be apprehended, at least mediately,
by sensing, if not by the
crude, then by the developed sense, he thereby intimates
an intermediate determination which should have prevented
the sensualistic consequences; for what distinguishes
the cultivated sense from the immediate one is precisely
the spiritual work, the energy of thought, whose results
deposit themselves immediately in sensuous intuition itself.

The sensuous is to be the sole object of
% --- p. 75 ---
% print: ikke han (misprint: han for kan)
\setcounter{footnote}{0}%
\opage{75}knowledge. But the sensuous is this particular and individual thing,
which is unutterable because the word, as the universal,
cannot grasp it in its individual determinateness; ---
and yet it is to possess reason and validity despite
this unutterability. --- This seems bound to make
the ``new philosophy'' altogether impossible. For it was
precisely to raise sensation to consciousness, to bring
the heart to understanding, and thus to convert the particular into the
universal. When Feuerbach maintains that the individual is
rational, although this reason cannot be uttered or
comprehended, he posits the paradox. His determinations
here form a complete parallel to S. Kierkegaard's
determinations of ``the single individual'' and ``the God-man''.
Just as Kierkegaard, on account of the impossibility
of comprehending, declares faith to be a passion, so
for Feuerbach too passion is the mark of
existence: only that is which is an object of passion.
The truth can be grasped only with man's whole
essence, not with thought alone. In an aphorism from
the same time as the work we are here concerned with,
we read: ``Science, it is said, does not solve the riddle
of life. --- Granted; but what follows from that? That you
should run over to faith? That would be going from
the featherbed into the straw. No, but that you should go over to life,
to praxis. The doubts that theory cannot solve,
praxis solves for you''. (Werke II p. 411). The new philosophy
thus itself acknowledges that it cannot give satisfaction
in theoretical respects.

In practical respects the new philosophy has made
the advance of overcoming the opposition between the
scientific and the religious standpoint; it is itself
religion, since it aims only at what is the
object of passionate sensation and love.
Yet a difficulty arises here similar to the one in
theoretical respects. Not the individual man, but
only the union of I and Thou is to express the
% --- p. 76 ---
\setcounter{footnote}{0}%
\opage{76}essence of man. What is unclear is whether man seeks this union
because the universal essence of the species leads him beyond
his individuality, or because his individuality, passionate and egoistic,
finds its satisfaction in the Thou. ``My principle,'' says
Feuerbach in an aphorism, ``is ego and alter ego, egoism
and communism; for the two are inseparable.'' For love
must presuppose the egoism of others as
justified. But all this implies is that the individual and
the universal \emph{are to} be united, not how they can
be one. In ``Wesen des Christenthums'' the universal was
still sole ruler and egoism the root of all evil;
now it has taken egoism up as coordinate with itself, without the tension
ceasing. In practical as well as in theoretical respects
there always remains a sting of contradiction
that drives the restless thinker further forward.

Feuerbach soon modified his conception of religion
in accordance with ``Grundsätze der Philosophie
der Zukunft'', and in such a way that through this
the very principles stated in that essay were developed
further and more sharply. ``Das Wesen des Glaubens im Sinne
Luthers'' (1844) is, according to his own declaration in the preface
to his collected works, the first work in which
the abstract being of reason no longer haunts
his head, but in which the sensuous, real essence of nature and of humanity
has come into its own, and
in which, therefore, what he recognizes as essential
in religion is at the same time what is essential for him
himself. What is peculiar about this work is that
it places the religious at once lower and higher
than ``Wesen des Christenthums''. There he still saw
the universal as an essential factor in religion, although
it was overwhelmed and distorted by the individual need
of the heart. But the essay on Luther determines
God as man's satisfied drive for happiness,
as the being that realizes human wishes,
whose highest potency is the wish for blessedness. It is set down
% --- p. 77 ---
\setcounter{footnote}{0}%
\opage{77}as Luther's peculiarity that he did not, with the
old doctrine of the Church, lay the main weight on the objective
content of dogma, but throughout stressed the subjective
moment: ``Alles was wir im Glauben erzählen, ist
für uns geschehen und kommet uns heim,'' says Luther.
And Feuerbach adds: ``Luther erst hat das Geheimniss
des christlichen Glaubens ausgeplaudert''. Faith in Luther's
sense consists in faith in God as a being
that essentially stands in relation to man. Luther does indeed make
man nothing before God; his doctrine is
a hymn to God, a lampoon on man. But
God could not be everything for man if man
himself were not nothing. Luther is inhuman toward
man only because his God is absolutely human. The God
whom Luther removes absolutely from man is a God who
himself wills the highest that man can will for himself.
Faith is therefore in the last instance love of oneself.
But even so faith is not condemned, as in ``Wesen
des Christenthums'', as absolute egoism. It is indeed
still maintained that only the species can be the standard
of the good, if there is to be an absolute, not a
relative and egoistic, standard; but on the other hand
it is stated that when love is built on faith, which
in Feuerbach's view means the same as that
love of others is built on love of
oneself, this can have the right meaning that
I cannot do good to others if I myself have
nothing good, that I must therefore possess in order to be able to
impart.

From this one sees, among other things, that it has been wrong to
think that it was opposition to the doctrine of religion
that later led Feuerbach to embrace egoism
as the sole moral principle. He sees self-love more and more
as the only active force in religion, and it is the awakening sympathy with this
pathological element in human nature that
% --- p. 78 ---
\setcounter{footnote}{0}%
\opage{78}alters his ethical view. One can trace in Feuerbach's
development, running in parallel, a bringing forward of the pathological
factor as the source of religion and of egoism
as the principle of morality, since, as he himself puts it,
what is essential in religion more and more becomes what is
essential for him himself.

The earlier writings in the philosophy of religion concerned themselves
exclusively with Christianity, the highest
form of religion, where the object of faith was a being raised absolutely
above nature. The question now arises
how the religious consciousness has been able
to develop to such a height, how it
has acquired the force and the means to posit a supernatural
being outside itself. This question Feuerbach answers
in the essay ``das Wesen der Religion'' (1845).
He here traces the essence of religion back to the Schleiermacherian
determination, the feeling of dependence. That
on which man is and feels himself dependent is originally
nothing other than nature, which is therefore the
proper object of religion. Man feels that he does not owe
his existence to himself, and so honors as God what
he feels himself dependent on. As high a value as he sets
on his existence, so high a value he sets on God.
And it is nature that is man's source of life; every
good gift comes not from above, but from below, from the depths
of nature. The higher presupposes the lower, not
the reverse, and the higher a being is, the more it
presupposes. The most composite, most dependent being is
the highest. The individualized being is a more
divine being than the one without individuality; but
it is bound to its conditions, ceases with them,
just as it begins with them. The foundation of man,
eternal nature, persists amid an infinite arising
and passing away; it alone cannot be derived from anything higher.
God regarded as an objectively real being is himself nothing
other than the essence of nature, which by the aid of abstraction
% --- p. 79 ---
\setcounter{footnote}{0}%
\opage{79}is separated from real, sensuous nature, and
by the aid of imagination transformed into a human
or humanlike being. God therefore presupposes
nature, as the abstract presupposes the concrete.
To want to derive nature from God is to want to derive the
real essence of nature from its abstract essence, existing only in
thought. The reason why such a deduction is constantly demanded
and attempted is that thought involuntarily
makes the abstract and universal, which stands nearer to it
and is for it the higher, into the presupposition of
the concrete and real, although in reality the relation
is the reverse.

However high religious faith may raise itself above
nature, nature nevertheless remains its proper
object. On account of his unfulfilled wishes man
feels himself dependent, and the necessary power of nature,
which everywhere meets him, now accommodating,
now as a barrier, is then apprehended as the free, arbitrary
power that grants or denies the object of the wish.
Man puts his own life of feeling and of will into
nature. This is the first stage, nature religion proper.
When man begins to separate his own
essence from nature, to see his true life not as an
immediate life of nature, but as a human life of society,
when he thus raises himself above nature in his understanding
and his will, then he also transfers the higher
intellectual and moral attributes to the Godhead
and thereby separates it more and more from nature.
This separation is carried through where, as in Christianity,
religious faith springs from the wish for an
unlimited blessedness, a wish that wholly leaps over
the natural barriers within which the wishes of the heathens
moved. The Greeks did not make the
divine, i.e.\ the possible being, the model and
standard for the real, but the real being the
standard for the possible. But the God of the Christians is
% --- p. 80 ---
\setcounter{footnote}{0}%
\opage{80}transcendent, like their wish. The religion of spirit makes
the opposite movement to nature religion. The latter
makes real beings into represented beings, into
beings for the imagination, since it is faith in
nature as a human being; the religion of spirit
makes represented beings, the objects of thought and imagination,
into real beings, since it is faith in the
human being as the being of nature.

Here Feuerbach's sympathy with the real foundation of
the positive religions comes clearly to light. He now no longer
makes the ethical and universal the basis in judging
them, but makes nature the standard and
demands that the limits set in it be respected.
He is thereby led over into a naturalism which in its further
development could lead back to his older pantheistic
view, if he did not, by the way
in which he declares spirit to be something merely derived, an abstraction
that has only subjective significance, rather approach
a complete materialism. What is interesting
in his last writings is now the conflict between this
materialistic tendency (a natural consequence of the
weight he laid more and more on the real) and the
energy with which he holds fast to his anthropological standpoint
and thereby to his original ideal point of departure.
From these writings\footnote[1]{Die Unsterblichkeitsfrage vom Standpunkte der Anthropologie (1846), Vorlesungen über das Wesen der Religion (1848), Theogonie (1857), Gottheit, Freiheit und Unsterblichkeit (1866).} we shall now bring out the main features,
partly of his general philosophical, partly in particular
of his ethical view and his view in the philosophy of religion, as
it takes shape at the close of his activity.


It is first the conception of the relation between the
universal and the individual that draws attention.
The existing being is always only the individual;
% --- p. 81 ---
\setcounter{footnote}{0}%
\opage{81}the kind is only a predicate. But thinking reverses the relation,
makes the subject a predicate, the predicate a
subject. Thought presents the universal that occurs
in all particulars as a thing by itself, and posits
as a continuum what is in reality something discrete,
life's infinite ``many times as an identical
``once''. The first step to wisdom is to recognize
this essential and insuperable opposition between
thinking and reality.

Since Feuerbach had earlier put self-consciousness in close
connection with the universal essence of the species, which through
it became an object for the individual, the conception
of consciousness naturally now had to be modified by the altered
conception of the universal. The secret of individuality,
he says, is the unity of I and not-I, and
this unity is also the ground of religion. Without the not-I
man would have no religion, since he would then
himself be God; without the I, or as an I that
did not distinguish itself from the not-I, he would be only a
plant or an animal. The deeper man goes into
himself, the more he sees the difference between nature and
man vanish, the more he recognizes that he
is only a conscious unconscious, a not-I that is I.
Therefore man is the very deepest and most profound
being. But man cannot endure or comprehend
his own depth, and therefore splits his essence into an I
without a not-I, which he calls God, and a not-I
without an I, which he calls nature.

When the universal is only as a predicate of the individual,
the conscious only as a predicate of the unconscious,
it might seem as if Feuerbach had already
overstepped the limit of materialism. And yet
he protests against ``the extravagant and transcendent
materialism'' that goes beyond the essence of man in order
to explain it from purely external causes (something Feuerbach
undeniably approaches himself in his proposition: ``der Mensch
% --- p. 82 ---
\setcounter{footnote}{0}%
\opage{82}ist, was er isst''). He wants to hold fast to ``an Archimedean
point between materialism and spiritualism'', an ``immanent
materialism'' that keeps within the anthropological,
and he maintains an immediate unity of
soul and body which lets nothing lie between them and
leaves no room for any separation or opposition
between material and immaterial. This unity he finds
concentrated in the brain as punctum saliens, where
the body is spirit, spirit body. This indissoluble
immediacy is the unassailable property of life, indeed
is the very essence of life in opposition to thinking and
knowing. In vain does idealism seek, from its theoretical
standpoint, to solve the problem of the reality of objective
being, while it rejects sensation
as something merely subjective. Sensation is indeed subjective,
but its ground is objective, and the identity of
subject and object reveals itself to us in the feeling of
well-being. To that extent religion is right when it
derives the reality of the world not from the understanding, but from
the will.

With these determinations Feuerbach has further
developed what we designated above as his paradox:
the unutterability and incomprehensibility of being. Being is conceived
abstractly and individualistically, essence becomes an invisible
point that cannot be specified in theory, and the opposition
between thought and life leaves only a practical,
a religious, solution open. Therefore Feuerbach ends
by saying: ``Meine Philosophie ist keine Philosophie.''

But it is only indirectly that he can be regarded as
the progenitor of modern materialism. In his energetic
holding fast to the anthropological point of view against
transcendent materialism there lies hidden an idealistic
principle which he never gives up. On this rests
also the decisive opposition between Feuerbach
and his contemporary, Auguste Comte in France, with whom
he otherwise has much in common, notably in his view
% --- p. 83 ---
\setcounter{footnote}{0}%
\opage{83}of the stages of development of the human spirit. Parallel with Comte's
three stages: theology --- metaphysics --- positive knowledge
run Feuerbach's three: religion --- speculation (pantheism)
--- anthropology. But Comte's ``positive knowledge is
precisely what Feuerbach calls transcendent, since it
takes its point of departure outside man and explains
man's essence purely from external causes\footnote[1]{Cf. on Comte ``Nyt dansk Maanedsskrift''. 2nd vol. p. 200---208.}.

In ``Theogonie'' Feuerbach occupies himself anew
with the investigation of the origin of religion and carries through,
in a multiplicity of turns and examples,
taken ``from the sources of classical, Hebrew and Christian
antiquity'', the doctrine already set forth earlier, that
man's relation to God is his relation to himself,
not in the sense in which this was conceived in ``Wesen
des Christenthums'', as man's relation to his universal
essence, to the Idea of the species, but as a relation to
his own self-love. Religion is no longer indirect
self-consciousness, but indirect self-love.
The wish is the theogonic principle: it is the expression of a
lack, a barrier, a suffering, but a vehement, a revolutionary
expression, a willing that this lack and barrier
shall not be. In and with the wish there therefore arises the notion
of a God, i.e.\ a being that is raised above this
barrier and can raise man above it. If one
would say that not the wish for happiness, but the wish
to become good is the ground of religion, then
Feuerbach answers that religion with good right grounds virtue
precisely on blessedness. Man cannot be good if
he is not happy. Purely theoretically one can make
happiness dependent on virtue, but in life things
go otherwise, where it is feeling, needy, desiring
beings that appear. He, therefore, who would
make men better must first make them happier.
% --- p. 84 ---
\setcounter{footnote}{0}%
\opage{84}Religion is itself the oldest form of culture and has in itself
the same purpose as culture, but wants to attain it without
using the means.

Here it becomes evident that Feuerbach is in principle at one
with religion, as he conceives it. Just as
religion does, in his opinion, so he himself
now builds the ethical wholly on egoism\footnote[1]{Feuerbach's course of development expresses itself in the following propositions, each superseding the one before: a) ``Wesen des Christenthums'': the corruption of Christianity shows itself in its building love on faith. b) An aphorism (contemporary with ``Phil. d. Zuk.''): love of others must recognize their egoism as justified. ``Wesen des Glaubens'': it can have a correct meaning to derive love from faith (love of others from love of oneself). c) ``Theogonie'': no virtue without happiness; therefore religion is right, and egoism is the principle of morality.}, and in
his vehement polemic against every ethics that does not proceed
from it, he can then, from his standpoint, rely
on religion as factual testimony for his doctrine,
just as he did above on the question of
the reality of objective being. How wonderfully has
the position not shifted in the course of Feuerbach's development!
He began by attacking religion from the standpoint
of philosophy; then he fought philosophy
because it was only the abstract counter-image of religion;
and at last he relies on religion in his
polemic against philosophy.

In his ``Pierre Bayle'' Feuerbach had spoken with enthusiasm
of Kant's and Fichte's ethics. Now he rejects
it with contempt. Duty cannot be made ground
and cause; it is a consequence, an effect, a drive that
asserts its right against the lust for domination of another drive.
The bad conscience is the angry shadow of
a drive that another and stronger drive has violently
put to death or threatened with it. The essential drive,
% --- p. 85 ---
\setcounter{footnote}{0}%
\opage{85}which underlies all other drives, is the drive for happiness,
with which the will is one. To will means
not to want to suffer. Right and morality are not thereby abolished;
they are only set on their natural feet.
In his very egoism man has the criterion of what
is right or wrong. The thief himself wants his own
property to be respected, the murderer his life to
be sacred. The contradiction between the criminal's will
with regard to himself and his conduct toward
others is the inner ground of conscience, of the consciousness
of right and wrong. I cannot demand that my property,
my egoism, be recognized if I do not recognize
that of others. And in any case the egoism of others will
surely keep my egoism within the proper limits; even
if it is a matter of indifference to me whether I am good or bad, it
is not a matter of indifference to the others. The good is what
corresponds to the egoism of all men, the evil what
satisfies the egoism of some individuals at the
expense of others.

In this last proposition, however, a principle is intimated
that leads us once again beyond the morality of egoism. ``All
men'' is, namely, only an atomistic expression for
what Feuerbach earlier called ``the species''. With this
his utterances in a practical-social direction also agree,
which can by no means be grounded on the principle
of mere egoism. Occasionally, for example, he pleads the cause of the proletarians.
He shows that it is false when it is said
that they have nothing and therefore seek to obtain everything; they have,
on the contrary, a very great deal, namely human self-respect,
the drive to cultivation, and the desire to work. What they demand is that
the value of labor shall be recognized, that the worker shall not
be regarded as a mere means, but as an independent
end that has validity in and for itself. But
Feuerbach's practice comes even more into conflict with his
theory when, against one-sided political views
which hold that everything depends on political and social
% --- p. 86 ---
\setcounter{footnote}{0}%
\opage{86}reforms and improvements in man's material conditions,
or, as Feuerbach puts it, that the evil sits not
in the head or heart of humanity, but in its
stomach, when against such views he maintains
that there are, after all, many diseases of the stomach that have
their ground in the head, and declares that he for his part
has made it his particular task to fathom and to cure
the diseases of the head and heart of humanity.
Likewise he declares that he who is still a slave of
prejudices has no right to free himself from
material and political oppression. That was also the reason
why he abstained from direct participation in the Revolution
of 1848; for it sprang from faith in
ideas and institutions from which the consciousness of the human race
must be freed before there can be any question of
practical and political liberation. Feuerbach here demonstrates
what is moreover the motive of his whole activity
with all its changing standpoints, an interest in
truth and spiritual freedom which can by no means be derived
from egoism.

Thus even at this last point,
to which we can follow Feuerbach, we find ourselves placed between
two-sided impressions. This remarkable figure
does not let itself be caught in any determinate form; he is always
in contradiction with himself and therefore always engaged
in developing, never finished. He wrestles
with the problems, but never says the last word and
really ends by declaring their solution impossible.
Having originally set out from absolute idealism,
he has caught sight of the significance of realism and now cannot
find the necessary intermediate determinations, but gives himself
over to it with immediate enthusiasm and without wholly
giving up the presuppositions of his former standpoint.
When he says, e.g., that the particular is unutterable,
it is Hegel's own proposition from the Phenomenology
that he adopts, only that he declares it to be the positive
% --- p. 87 ---
\setcounter{footnote}{0}%
\opage{87}expression of true being, which Hegel saw precisely as something
negative, as an expression of not-being. Therefore
Feuerbach has only got as far as becoming the prophet and mystic
of realism.

Feuerbach's chief merit lies in the domain of the philosophy of religion and of psychology. Here his
expositions are of epoch-making significance for the explanation
of the origin of religions in the human spirit. But
there is nevertheless this defect in Feuerbach's psychological investigations,
that they do not lead beyond the concrete and
empirical; his strength is what he calls dramatic
psychology, i.e.\ psychology in close connection
with the objects in which the human soul reveals
itself in its totality. A deeper grounding of these
concrete movements, which Feuerbach for the rest depicts with
such great mastery, he cuts himself off from giving
by the very fact that he makes psychology or anthropology itself
the fundamental science. The anthropological
standpoint does indeed prevent him from stopping at a pure sensualism
or materialism, but only by making his
doctrine wholly subjectivistic in the last instance. To this
he has come through the method of subjectivization: reason
is the subject of God, and man is the subject
of reason. But what significance these subjective forms
have in themselves and in existence as a whole is not
investigated.

This empirical and subjectivistic individualism of Feuerbach's
prevents him from gaining a really historical
view. The historical rests on the interaction
between the universal, the objective purposes and ideas, and
the concrete, individual beings who seek to appropriate
them and to realize them in themselves and outside themselves. But
one of these factors is lacking in Feuerbach. He
is content with the psychological How and does not concern
himself with the metaphysical and historical What.
In this he forms a complete opposition to Hegel,
% --- p. 88 ---
\setcounter{footnote}{0}%
\opage{88}who in history takes account only of the universal ideas,
while the acting individual and the whole psychological
consideration are indifferent and worthless to him. Yet
Hegel's view stands nearer to history to this extent,
that it is able to appropriate history's content in its
inner richness and diversity, whereas Feuerbach is bound
to one and the same circle of concepts, since the psychological
factors are essentially the same at every
point of history. Thus for him all positive
religions become at bottom one and the same being, whose
inner differences become insignificant, because,
in view of the pathological element in the essence of religion, he
overlooks the other important elements in it, which point beyond
the purely individual drive of the subject and show that
through these forms of spirit too something higher
works its way forward. Thus he lets neither the rational
element nor the element of resignation in religion
come into its own. And this has a harmful influence not only
on his conception of the historical religions,
but also on his conception of the very essence of man
as this displays itself through them.

Feuerbach has himself demonstrated the truth of his individualistic
principle by stamping his view and
his writings with the peculiar mark of his own individuality.
To agree with him one would therefore really have
to resemble him in psychological respects as well, and it is
not surprising if objective criticism must raise
many objections against him. But still we part
from him with the conviction that we have been occupied
with one of the few really original thinkers after
Hegel. ---

6. \emph{Max Stirner}. --- Feuerbach had driven the opposition
to speculative philosophy so far that he
at last declared the true philosophy to be no philosophy
at all. And yet there is more speculation in him
than he himself will admit. What speculative thought
% --- p. 89 ---
\setcounter{footnote}{0}%
\opage{89}seeks at all times, the unity in the oppositions of existence,
that is also Feuerbach's task, although he
sometimes in paradoxical fashion cuts off the possibility of
its solution. And even when he does this, he expressly holds fast
to the terms of the opposition, if in one-sided
form; it is only the expression of their unity and connection
that he despairs of finding.

If formal consistency could be the standard
in judging philosophical views, Feuerbach contains
contradictions enough. And these contradictions
appear all the more glaring the more he takes occasion, from the real
and supposed contradictions he discovers in earlier
views, for passionate, often scornful
attacks, and the more he himself believes he has opened
the entrance to a wholly new knowledge, ``toto genere''
different from the earlier one. To that extent it is a just
nemesis when the weapon he himself so passionately
turned against others is also turned against himself,
when he, who believed himself the most radical revolutionary,
now, like the parties succeeding one another in
the French Revolution, is himself outstripped by still more
extreme views, by which he is placed outright
alongside the old world-view he believed he
had overthrown. But even so the real
philosophical movement stops with Feuerbach. The formal
consistency, which at last appears even merely in the form of
experiment, cannot deserve the name of philosophical
thinking, although it has its philosophical interest to
see how absolute idealism at last ends in
pure formalism.

The work ``Der Einzige und sein Eigenthum'' (1845)
by the pseudonymous Max Stirner (Dr. Schmidt, known as an author
on social economy, who died in Berlin at
the beginning of the sixties) was for the 19th century
what Helvetius's book ``de l’esprit'' was for the
18th: the most ruthless and unreserved proclamation
% --- p. 90 ---
% print: vort egen (misprint: egen for eget)
% print: betvinge (misprint: at lost before betvinge)
\setcounter{footnote}{0}%
\opage{90}of egoism as the only real and justified motive
for action. But whereas Helvetius relies
on materialism, Stirner is an idealist to the extent that he
carries the Feuerbachian method of subjectivization
beyond the point where Feuerbach himself stopped\footnote[1]{Stirner does indeed direct his polemic chiefly against ``Wesen des Christenthums'', but his objections also strike those writings of Feuerbach that belong to his third stage, and some of which had already appeared when Stirner's book came out.}.

The child is a realist, the youth an idealist and the man
an egoist. It goes the same way with the history of humanity
as a whole. The ancients saw the truth
in the world and let their own I bow to the conditions of the world,
which, however, their striving aimed
to subdue and dissolve. The consummation of this dissolution is
Christianity, where the world is a nothing, spirit the
only truth. In modern times the movement aims
at dissolving the truth of spirit. As the world was for the
ancients, so spirit is for the moderns ``eine Wahrheit,
hinter deren Unwahrheit sie zu kommen suchen und
endlich wirklick kommen.'' This dissolution has approached
its consummation in the latest philosophy, where
God is recognized as man's own essence. But
also only approached its consummation. For even Feuerbach
and Bruno Bauer still stand at bottom on the
theological standpoint. For what is the use of
Feuerbach's saying that the beyond is only our own essence,
when he then makes our own essence into God and sets
it in such opposition to us that it becomes a
new beyond, a new heaven within us? We are not
that which is in us, any more than we are that which is outside
us. Feuerbach splits us into an essential and an inessential
I and thereby pushes us back into dualism. God
remains God, even if one drives him out of
heaven and shuts him up in the human breast.
% --- p. 91 ---
\setcounter{footnote}{0}%
\opage{91}Feuerbach therefore really represents the last convulsive
effort to hold fast to Christianity;
he is what he himself said Schleiermacher was,
the last theologian of Christianity.

The last step, the complete sublation of transcendence,
has not yet been taken so long as a universal is still recognized
over which the individual does not dispose as
his property. Everything holy is something alien. If
you believe in the truth, you are of course the servant of the truth,
do not belong to yourself. The universal purposes lay claim
to me for their own sake, are thus egoists;
but is my egoism not just as justified? Let us
then reverse the relation and make the universal our
servant! Political liberalism works for freedom
and right, social liberalism (communism) for
society; but freedom, right and society are only chimeras.
I do everything for my own sake. Only I
myself determine what shall belong to me. Help yourself!
War of all against all is the watchword. Those who talk of
love are possessed, have a screw loose. For love
makes me dependent, I am possessed as much
as I possess; the object of love cannot
be absolutely one with me.

When every authority both within and without,
in heaven and on earth, has been felled, and ``the Unique One, by
transforming everything into ``his property'', has himself become the
only reality, it is understandable that he can hurl out
a proposition such as this: ``Ich begehe getrost die Sünde
wider den heiligen Geist!'' The movement is consummated; ``the
Unique One'' has seen through the untruth of spirit.

Therefore a critique hardly seems possible here either;
for when the truth is his property, he cannot
be bound by any universal principle. With good reason
Stirner's book has been called the most radical ever
written. It is radical also in the sense
% --- p. 92 ---
\setcounter{footnote}{0}%
\opage{92}that it seems to want to tear up truth by the roots, so
that truth shall not be used as a weapon against it.

Yet Max Stirner humorously gives himself a practical
denial by putting the following dedication at the front of the book:
``Meinem Liebchen, Marie Dähnhardt.'' It
shows that even if the egoistic standpoint can formally
be carried through in theory, it is nevertheless impossible in practice.

A theoretical exposure Stirner next gives himself
by constantly presenting egoism as a new morality,
a new truth, which is to replace the old. That looks
strange at the very moment when all truth
is to have been rooted out! But that is also why even Stirner could
still manage to be counted among the Right.

7. ``\emph{Das Verstandesthum und das Individuum}''.
(1846). The author of this anonymous work was,
according to some, a clergyman from Köthen by the name of Karl Schmidt,
who wished to show to what results the ``critical
consciousness'' must finally come once it
has given itself over to the demonic power of formal consistency.
After showing, by extracts from the writings of the various critics,
with what rushing speed Strauss,
Feuerbach, the brothers Bauer and Max Stirner replace and
overthrow one another in turn, the critical movement is taken up
at the last-named author, who himself sets up
an ideal at the moment when he believes he has overthrown
all ideals. A world of egoists is now to
be the perfect world. But then Stirner is no
longer ``der Einzige''; he is one of many, not the
absolute individual. Everything that ends in --- ist is
not individual.

Against Max Stirner's ``Einzige'' is then set the individual,
the absolute individual, which is free of every remnant of the rule of the understanding.
But the individual finds itself existing in a
double world, a physical and a psychical, a world of nature
and a world of spirit. The first question is then
how the individual stands toward these worlds.

% --- p. 93 ---
% print: Tilværelsen uendelige (misprint: for Tilværelsens)
\setcounter{footnote}{0}%
\opage{93}In reality only the particular exists. Nature
as a whole is a figment of thought, which arises from
natural scientists seeing every particular in connection and bringing
it under universal laws; thereby they become dogmatists.
``The Individual'', on the other hand, does not think the particulars, the atoms;
thereby it would make them universal; --- it intuits
them, stares at them, and by this staring
at the particular lets the world, the whole and the connection
pass away.

Spirit is the faded, the anemic individual,
the dull gray in gray instead of the motley
multiplicity of reality. In ``the realm of spirit'' only the
universal powers count, not the individual. ``The Individual'', on the other hand, sees all the institutions and forms of the human world only as
the crystallized representations of men, their fixed
ideas; therefore they have no validity for it. Yet
it goes along with these forms, that is to say, plays with
them, enjoys them, as long as it pleases. The dissolution of spirit
means the dissolution of all earnestness of life, all passions,
all enthusiasm, all pathos.

But what, then, is the individual itself? It cannot be uttered,
since language can express only the universal. The individual
can be designated only by gesture, by pointing. The individual
is incommensurable, has no other measure than
itself. The individual is nothing and is everything; it is a Proteus
that constantly changes shape. It constantly breaks
with its past and yet makes no progress. In such
and countless other turns one could characterize
``the Individual'', and yet one would not come
nearer to the matter. The characterization must constantly be revoked
and corrected --- and at last ``the Individual''
nonetheless mockingly breaks off the whole endeavor with the words: ``Ich
bin ich selbst allein!'' ---

German speculation, which began with the I
as a monad, out of which the infinite content of existence
developed, thus ends in the I as the absolute
% --- p. 94 ---
\setcounter{footnote}{0}%
\opage{94}atom, which makes every determination, every content impossible.
But the final consequences here are, to be sure,
of such a nature that they could be reached only by way of
thought-experiment.

\chaphere{CHAPTER THREE.}


The reasons for the emergence of materialism in our time,
and for the great following it has found, are of various
kinds, partly negative, partly positive. The magnificent
philosophical enthusiasm that had carried the movement
beyond its true limits was bound to
bring with it partly a slackening, partly a reaction, which
turned skeptically and polemically against every properly
philosophical attempt. Philosophy here only shares the fate
of the other spiritual powers, which in the middle of the
century seem to lie dormant after the vigorous life of
its beginning. The age threw itself with all its
energy into the currents of outward life and for the time being turned
its back on the theoretical and aesthetic ideals with
which the preceding generation had been inspired.
In its recklessness it even seems to turn its back on
the ethical ideals as well. In any case it does not, like its predecessor,
satisfy the religious need
by the way of inner, free thinking and experience; here too
it wants to be ``practical,'' and from this it is understood that the
most outward literalism shares dominion in the consciousness of the multitude
with materialism.

But philosophy itself, too, had through the self-dissolution of absolute
idealism come to such an
extreme that materialism could seem the only possibility
that still stood open. Step by step the
% --- p. 95 ---
% print: af hvilken Tanken (misprint: hvilken for hvilket (Substrat))
% print: Tryllekreds og. (misprint: a point after og, for nothing)
\setcounter{footnote}{0}%
\opage{95}universal, the Idea, had lost ground, and the real particular had
stepped into the foreground. At last this stood as an indissoluble
atom, while thought was dissolved into a nothing. What
was there then to do but to take one's point of departure in
the real atoms and in their composition to find a
substitute for the lost Idea?

The materialists often appeal to Feuerbach. He
had sought to show that man in religion and speculation
only carries the forms of his own essence over into the objective,
forms it after himself. Therefore, the materialists conclude,
it cannot be thought, reason, that
rules in existence. Thought arises in man only
as the product of material causes. Feuerbach held fast
to human thought as the only possible
standpoint that could not in turn be derived from without; but
materialism, which he therefore describes as transcendent
and extravagant, breaks this magic circle and
presupposes matter as the factual substrate from which thought
is derived. The materialists appeal to Strauss as well,
in that from his pantheism they eliminate the ideal
factor, so that only the material substance remains,
which bears everything and underlies everything.

Yet the materialists will not put up with being derived
from the earlier philosophical movement. Their doctrine
is to be the necessary result of the magnificent upswing and progress
of exact empirical science in
our century. Trusting in this, they make short work
of the whole of earlier philosophy. So much must be granted,
that the dominion of empirical science and of practical interests
is perhaps the most important cause of
the emergence of materialism, just as most of
its spokesmen are natural scientists, even if they have acquired
a quasi-philosophical education. But it is another
matter whether materialism really is the true consequence
of exact empirical science, or whether
the connection between the two is not merely apparent and
% --- p. 96 ---
\setcounter{footnote}{0}%
\opage{96}accidental. This the following account must decide. In
any case it is only the great results of the intervening scientific
development, not distinctive
and new thoughts, that produce any contrast between the
materialism of the 18th and that of the 19th century. ``If
materialism in our time did not gain a peculiar
interest through the excessive false enthusiasm for a
truly great circle of scientific education, whose
results and analogies it misuses in the most frivolous
manner, we would not need to occupy
ourselves with it here''\footnote[1]{H. \emph{Lotze}: Medicinische Psychologie. p. 30.}.

Nevertheless it is of great theoretical and practical
interest to get one's bearings with regard to modern
materialism. When on all sides people agree on
the recognition that philosophy must appropriate the
real foundation, as exact empirical science
demonstrates it, that it must once more descend from the heaven
of the Idea to the real world, then modern
materialism only pronounces the one-sided consequence of this,
in that it lets philosophy be absorbed into the
supposed consequences of the special sciences. Materialism is the irrational,
unphilosophical restoration of reality;
it represents the extreme left in the efforts toward a
philosophy of the future. --- Materialism has at the same time a
popular, democratic tendency, in contrast to the
spiritual aristocracy of earlier science. This striving
after the popular goes so far that Büchner in the preface
to his book ``Kraft und Stoff'' as good as declares
that philosophical expositions that cannot be understood
by every educated person are not worth the printer's ink; philosophy
is to be ``geistiges Gemeingut''\footnote[2]{Büchner, to be sure, cannot himself have fulfilled this demand, since immediately afterwards he reproaches his opponents with not having understood, for \opage{97}lack of scientific insight, the inner core of his view. To be consistent he should have accused them of lack of education.}. Carl Vogt
% --- p. 97 ---
\setcounter{footnote}{0}%
\opage{97}travels about as a missionary of materialism and gives
popular lectures, even appearing in the churches for lack of other
premises. And this democratic direction of materialism
stands in connection with its
whole practical tendency. It regards itself as the possessor
of the solution of the social and political problems,
in that it means to take hold of the matter from below and thereby accomplish
what the highest ethical and religious ideas have not
been able to do. To the philosophical interest
the interest of cultural history is here closely joined. ---

With regard to materialism in our century it is
worth noticing that the two most important
clashes in the materialist controversy took place between
parties who all stood on the ground of the special sciences
but disagreed about how the final consequences
were to be drawn. This shows at once how matters stand
with materialism's claim to be the true consequence
of the special sciences. --- The famous chemist
Justus Liebig had in his ``Chemiske Breve''
declared that natural science, far from cutting off the way
to the knowledge of the highest ideas, was on the contrary itself
a link in the chain that leads to them. Jakob Moleschott,
whose well-known proposition ``Ohne Phosphor kein
Gedanke!'' Liebig had particularly attacked, answered him in
his book ``Der Kreislauf des Lebens'' (1852), where he
recklessly held fast to the materialist doctrines.
The same year in which this book appeared, another controversy
had already broken out, between Rudolph Wagner in Göttingen and
Carl Vogt in Geneva. Yet this reached its high point only some years
later, when Wagner at a meeting of natural scientists
in Göttingen set forth as the content of his confession of faith
an ethereal soul-substance that did not die with
% --- p. 98 ---
\setcounter{footnote}{0}%
\opage{98}the body, a faith in which he saw the biblical doctrine of the soul
expressed, but which he thought he could hold fast only
by absolutely separating his faith from his knowledge. He
therefore declares himself an adherent of a collier's faith as soon as
one goes beyond the domain of exact science:
``In Sachen des Glaubens liebe ich den schlichten, einfachen
Köhlerglauben am meisten.'' Against this
Vogt came forward with his crushing polemic ``Köhlerglaube und
Wissenschaft'' (1855). Vogt is a spirit who in his own
domain has not a little in common with Feuerbach. Here
is the same impetus, the same reckless cynicism in thought
and expression. But as far as thinking proper is concerned,
the difference is considerable indeed; beside Feuerbach's
fine, penetrating analyses, Vogt's expositions look
rather bearish. What is most
appealing in him is his hearty humor, which did not
leave him even after he, as ``Reichsregent'' in 1848,
suffered shipwreck of his ideal dreams. Before that time,
he said later, he had indeed doubted the perfection of man,
but had nonetheless regarded him as the most perfect
animal; but in Frankfurt he had lost
this faith too.

What made Moleschott's and Vogt's position relatively
easy was that their opponents were not clear
about the form in which the cause of the ideal
is to be defended. Liebig speaks in a vague and confused
way of knowledge and revelation, so that one does not
see whether it is the free knowledge of the ideal or
the positive doctrine of the Church that he wants to assert. Naturally
Moleschott at once avails himself of this to present him
as the apologist of a medieval dogmatics. Still
more unclear is Rudolph Wagner. His ``collier's faith''
proves on closer inspection to be a whole scholastic
system, composed of biblical doctrine, from which he
in particular asserts the doctrine of the descent of all men
from Adam and Eve as necessary to salvation, and of
% --- p. 99 ---
\setcounter{footnote}{0}%
\opage{99}the fantastic assumption of an immortal soul-substance
that moves the fibers of the brain as the musician moves those of the pianoforte,
and propagates itself by division from parents to children. It
is no wonder that not only a materialist like Carl
Vogt, but also natural scientists like Virchow and philosophers
like Lotze, Zeller and Weisse had to protest
against such a doctrine. But this unclarity has been
to the detriment of what the controversy yielded.

Besides Moleschott and Vogt we must further single out
two principal materialist writers, Louis Büchner
and Heinrich Czolbe. What gives them particular interest
is that in materialism --- strange as it sounds
--- they really seek the satisfaction of a spiritual need.
They are led by a philosophical and ethical interest. In particular,
Czolbe stands in contrast to the other materialists
in that he comes to materialism by way of
deduction, a contrast that is well suited
to orient us with regard to the true peculiarity of modern
materialism; for which reason we will divide our
account into two parts, empirical and
speculative materialism.

I. \emph{Empirical Materialism}. --- None of
the materialists takes his stand so resolutely on the standpoint of the bare special science,
and from it pronounces absolute
judgments on the great problems of existence, as \emph{Carl}
\emph{Vogt}. He himself describes his standpoint as a physiological and natural-historical one.
Physiology has solved its task when
it shows which is the organ of thought, how it
comes into function, and on which conditions this function
depends. The natural history of the human race has nothing
to do with the essence of man as this can be determined
conceptually, but examines the millions
of human beings who are spread over the earth, and investigates
their peculiarities, and that back to a
time ``when man has left us hardly more traces
of his existence than any other wild animal that
% --- p. 100 ---
% print: Wissenschafft (misprint: so printed)
\setcounter{footnote}{0}%
\opage{100}inhabited the same region.'' If one then asks about the determination
of man's relation to the animals, only the
anatomical and physiological characteristics come into consideration;
all other considerations are indifferent\footnote[1]{See, besides ``Köhlerglaube und Wissenschafft'' (1855), the lectures ``über den Menschen, seine Stellung in der Schöpfung und in der Geschichte der Erde'' (1863).}.

On such a standpoint the special science must
necessarily place itself if it is to be able to carry through
its exact method. This always presupposes a certain abstraction;
the phenomenon must be isolated, drawn out of its
living, many-sided reality, in order to be checked
from a single definite side. Here, at the beginning of the special science,
we therefore also come up at once against its limit.
If it is only the anatomical characteristics that come
into consideration, then it is also precisely only as an anatomical
being that man has been considered, and one cannot
without further ado make this consideration the absolute one.
But against such a limitation Carl
Vogt now protests most decidedly; he wants immediately
to make physiology the absolute science.

Geoffroy St. Hilaire and Quatrefages held that one
must, in the natural-historical consideration of man, also
take account of the intellectual, the moral and
the religious, which make up precisely what is distinctively human,
and, relying on this, they set up a separate
human kingdom in contrast to the animal and plant kingdoms.
Vogt is not content to reject this from a purely
natural-historical point of view. He undertakes to show that
in the spiritual elements that the French researchers
emphasize there lies nothing peculiar to man.
When one ascribes religiousness to man, this is supposed
to mean only that man forms certain notions
of the phenomena whose ground he cannot grasp,
while the animal, owing to its lesser spiritual
% --- p. 101 ---
\setcounter{footnote}{0}%
\opage{101}endowment, does not feel prompted at all to think about
the ground of such phenomena. Yet in this way Vogt would
easily come to set a greater opposition between man
and animal than agrees with his view; he
therefore explains the actual origin of the religious from
fear of the unknown, which is regarded as something
supernatural, and this fear, which in man
becomes fear of God and produces religion, the animal too
is supposed to possess. Thus the dog
and others of our ``intelligent domestic animals'' feel fear of ghosts
just as well as the Breton or the Basque. Nor
is the moral supposed to be anything peculiar to man,
since the disposition to it already shows itself in the animals,
especially those that live in societies; moreover the concepts
of good and evil are relative, depending on differing
culture, on the needs of society and on the mutual
relations of individuals. The psychology of materialism thus reduces
itself to what man has in common with the animal.

Into Vogt's physiological and anatomical investigations
we cannot go deeply here. It is sufficient to
bring out the way in which he seeks to state his
results precisely. The brain, it is said, is the organ of all the various
so-called functions of the soul; these functions
are bound to certain parts and places in the brain and can
be exercised only by this organ. Even if, though reluctantly, he
must grant Virchow that it is impossible to explain
the fact of consciousness in terms of natural science, he
nonetheless thinks that one will be able to succeed in pointing out
with determinateness the ganglion cells by whose stimulation this or
that special sensation is aroused. That consciousness
cannot come about without ganglion cells can be proved;
\emph{how} it arises in them cannot be said; but
it stands firm that function and organ depend on each other
and are indissolubly bound up with each other. The function
must therefore also be dependent on the material
state of the organ and suffer under its changes.

% --- p. 102 ---
\setcounter{footnote}{0}%
\opage{102}It is clear that here only a relation of dependence,
of conditioning, has been demonstrated. Because two moments condition
each other, the one need not be produced by
the other. There may well be a qualitative difference
present (which naturally must not be confused
with absolute heterogeneity). The relation between the homogeneous
and the heterogeneous in the two moments one will
find by bringing out not only the points where the conditioning
is absolute, but also those where it proves
to be relative, each moment acting in its own peculiar
way. The qualitative difference that remains
then points toward a third, more comprehensive element
from which both those moments spring. Here, then,
when we apply this to the example before us,
the insufficiency of the psychological consideration
becomes evident just as much as that of the physiological. Only a metaphysical
investigation can fix the relation between
spirit and matter. The materialists can interpret
the results of physiology in their favor only by presupposing
as given that spirit is merely the product of matter.
It is their tragic fate that they, who shun
metaphysics like the plague, only manage to start from the
purely crass and crude metaphysics. The metaphysical
simply cannot be excluded.

The polemic of Vogt and those of his opinion is in the highest
degree influenced by the very opponents they combat.
Their hereditary enemy is dualism, which posits an external separation,
an alien relation between soul and body,
and this enemy they discover in every conception that
will not with them derive spirit from matter. Thus
Rudolph Wagner has produced much confusion by
his doctrine of the soul-substance and of the special organ of faith.
He is an opponent after the materialists'
own heart, for he takes things just as externally and immediately
as they do themselves. Against such views
they are in the right; but in the intoxication of victory they let themselves be carried
% --- p. 103 ---
% print: Erkjendeslære (misprint: Er-|kjendeslære, for Erkjendelseslære)
\setcounter{footnote}{0}%
\opage{103}further and think that by the proof of the conditioning of spirit
by the material they have shown the latter to be the complete
cause of the former. ---

While Vogt goes straight from zoology to the decision
of the main problems, we find in \emph{Moleschott} attempts
to give a deeper grounding through a kind of theory of knowledge
and metaphysics. In him one notices the influence
of Spinoza, Hegel and Feuerbach from his
earlier studies. Philosophy is, in his conception,
the spiritual expression of the sum of observations that
sensuous perception has gained at a given time. To
philosophize is to think. Thinking consists in the working up
of sensuous observation. Philosophy and
experience are mutually to merge into each other, so that intuition
is thought and the understanding intuits with consciousness.
One seems to hear a Schellingian here. But
it soon becomes apparent that the unity of the ideal and the real
element in knowledge is one-sided: the right of thought derives,
as Moleschott puts it, only from the grace of the
senses. Thinking arises from observation, and observation
is perception of the impressions of moving matter
on our nerves and through them on the brain. The
thinking human being is therefore the sum of his senses,
consciousness a property of matter. What is known
is only the properties of things, but the sum of
these properties revealed to the senses makes up
the essence of things, just as the sum of the senses made up thought.
The tree in and for itself is not different from the tree
as it is perceived; for the tree is only insofar as
it is perceived. Without a relation to the eye, into which it sends
rays of light, the tree does not exist. All being is a being through
properties, and every property consists in a relation.

These propositions contain several problems,
which Moleschott, however, passes over without more ado. In particular
he does not notice how he himself has pressed in a
sting of skepticism that must make all materialism
% --- p. 104 ---
% print: der. Han (misprint: hans|der.: Hæn- lost at the line turn (Hænder))
\setcounter{footnote}{0}%
\opage{104}impossible. If the external is only by being mirrored in the
internal, what assurance can one then have of the reality of external
actuality, let alone that one should be able to
posit the movements of the external as the cause of the internal?
Moleschott answers that, in consequence of Feuerbach's doctrine,
the concept of force, and with it every concept of activity,
must be understood only as an anthropomorphism without objective
significance. But then the same must hold of
the concept of matter, the concept on which Moleschott builds
his whole doctrine.

The doctrine of the circulation of matter is, namely, according to Moleschott,
the new world-view of which the Encyclopedists
in the 18th century were only the enthusiastic
forerunners; it is the hinge on which the philosophy of the present
turns. Modern chemistry has weighed the tooth of
time. It has discovered that matter is immortal, never
perishes, but is only converted into other and yet other
forms. Between the organic and the inorganic world,
between plants and animals, matter circulates in unending motion.
E.g.: ``The miner who in the Wetterau or in
Estremadura digs for phosphate of lime seeks more
than gold; he digs for wheat, for human beings.
The miner raises the treasure of the spirit that the farmer puts
into circulation, giving the wheel of time its first push.
The miner, who in the sweat of his brow and at the risk of his life
wins his livelihood by struggle, does not know whether perhaps
the matter for the finest head is not slipping through his
hands. He perhaps sets centuries in motion by
his hidden labor.'' Therefore, says Moleschott,
one does not name the circulation of matter without a certain awe. For
just as trade is the soul of intercourse, so
the eternal circulation of matter is the soul of the world. On the relation
to this circulation the peculiarity of the different beings depends.
Matter underlies everything, goes
through everything. What then is the general concept of matter?
It is found by abstracting from all the properties
% --- p. 105 ---
\setcounter{footnote}{0}%
\opage{105}by which one matter differs from another. There
then remain three common and decisive properties:
weight, the filling of space, and force or the capacity for
motion. Force is thus only a property of
matter. This holds, as we have already seen, also
of the highest force, consciousness. Matter governs everything,
man too: do spices not promote digestion,
do onion and vanilla not arouse sensuous drives, do
wine, coffee and tea not rule the ``mood of the brain''?

In this determination of the concept Moleschott has not
carried the abstraction through. For it is clear that of
the three properties in which, according to his doctrine, the essence of matter
is to express itself, two, weight and the filling of space,
fall under the third, the general determination
of force. Modern physics also leads to
conceiving matter as concentrated force. But
this points us in a quite different direction from materialism.

From his doctrine of the eternal circulation of matter
Moleschott now derives peculiar practical consequences.
Just as man as a whole is determined from first to last
by the influence of matter, so too the will is
only the necessary expression of a state in the brain called forth by
external influences. The concepts of good
and evil and the concept of right arise from the needs of society.
The right to punish is grounded in the need
for self-preservation. From this unconditional determinism
Moleschott derives a mood of mildness and reconciliation
toward the criminal, in accordance with Mme. de Staël's words:
tout comprendre, ce serait tout pardonner, words that
in his opinion ought to stand at the head of the gospel of modern
times, just as the commandment of love of one's
neighbor is the core of Christian morality. Yet
Moleschott himself forgets this ``golden word'' in the face of the
Hessian Chamber of Deputies, which to his indignation rejected
a motion for the abolition of the death penalty. Another
% --- p. 106 ---
\setcounter{footnote}{0}%
\opage{106}materialist writer, H. Czolbe, has seen this
inconsistency and declares that Moleschott's ``soft heart''
has here run away with him. Yet this does not bring Moleschott
out of the optimistic joy in existence
that he derives from his materialist determinism:
``With matter life circulates through all parts of the world,
with life thought, with thought the will that is good by natural necessity.
In spite of all its evils the earth is and remains
a paradise after all.''

Although individual beings merely emerge in the
unending circulation of matter only to sink back into it again,
man's activity nonetheless acquires a special significance in that,
guided by chemical, physical and physiological
knowledge, it can economize on matter by leading it along
paths and through combinations where it can exert the highest effect
by the shortest route. By the right
mixture of carbonic acid, ammonia etc. we work toward
the highest development of humanity. The natural scientist is
the Prometheus of our time, chemistry the central science.
The solution of the great social questions of our time lies in
the hand of the natural scientist; for poverty is immediately
only a lack of matter, and since the doctrine of the circulation of matter
assures us that there is enough matter in nature, science
must some day get so far that it teaches a
distribution of matter by which poverty becomes impossible.

What interests us in this ethics, built on ``ein stoffliches
Glaubensbekenntniss,'' is to see
how materialism here, as on so many other
points, goes beyond itself in seeking to draw its
final consequences. Moleschott involuntarily comes
to set ``a holy task, a purpose that is to be realized
by means of the laws according to which the movements of matter
are governed.'' But here it is then thought and
law that are the governing factor, and the absolute relation of dependence
on matter is broken. Moleschott will no doubt
reply to this that the activity of thought is only a side track,
% --- p. 107 ---
% print: tilsidt (misprint: for tilsidst)
\setcounter{footnote}{0}%
\opage{107}onto which the movement of matter as it were turns, in order to
return again with greater intensity. But already
in this lies a self-surrender of materialism. For the concept of matter,
under materialist presuppositions, makes impossible
a circulation in any other sense than the mere
external succession of various mutually alternating
forms. As soon as a purpose is to be developed through
this circulation, a concentrated unity has been posited which
is different from matter as such, and to which matter
stands in a serving relation. ---

The most popular of the materialist writers
is Louis Büchner, whose book ``Kraft und Stoff. Empirisch-naturphilosophische
Studien'' has been translated into
English, French and Spanish and has gone through some ten
German and French editions. Nevertheless this book is
very superficial, but it interests by Büchner's lively,
though often rather naive, enthusiasm for the cause.

We find here first a striving after a sharp conception
of the nature of experience. Büchner teaches that
scientific research at last comes up against a natural
limit, which lies in man's very faculty of knowledge.
We can form no notion of, nor
have any knowledge of, the Absolute, i.e. of what goes
beyond the sensuous world. All our knowledge and representing
is only relative and rests on a mutual comparison
of things. Therefore it does not attain to an
explanation of the inner relation between spirit and matter either.
Büchner declares repeatedly that what he wants to demonstrate
is only the necessary and indissoluble relation between
spirit and matter, force and matter, not the how
of this relation. But although the relation between spirit
and matter cannot, in his opinion, be developed positively from the standpoint
of natural science, it can nonetheless
be determined negatively, that is, empirical inquiry cuts off
the possibility of a transcendent, dualistic relation, in that
it shows matter to be the necessary condition of spirit.

% --- p. 108 ---
% print: paralelle (misprint: so printed)
\setcounter{footnote}{0}%
\opage{108}These propositions correctly characterize the standpoint of empirical
science. By holding fast to them Büchner would
not have become a materialist either. But unfortunately he contradicts
them in the same breath in which he puts them
forward. Already in the preface he finds in the shipwreck of speculative
philosophy of nature the clearest proof that
the world is not the realization of a thought, but a
complex of things and facts. In a similarly
dogmatic manner he calls the relation between spirit and
nature (which he himself has declared inexplicable) a
causal relation, and understands by this that it is matter
that is the cause, spirit that is the effect. But
the unclarity comes out most strongly in the relation between force
and matter in general. He starts from the Moleschottian
proposition: without matter no force, without force no
matter, and like Moleschott he understands this proposition more precisely
to mean that force is a property of matter.
The untenability of this assertion we have seen above.
When Büchner became aware of the physicists', in particular
Helmholtz's, investigations into the conservation of force,
he inserted into the 5th edition of his book a chapter on
``the immortality of force,'' in which the circulation of force is placed
as a necessary correlate beside the circulation of matter,
so that ``both together from eternity and to eternity
form the sum of phenomena that we call the world.''
He does not consider that he has hereby passed over
into a complete dualism between two parallel circulations,
but calmly returns to his presupposition: ``In matter
dwell all natural and spiritual forces; matter is the primal ground
of all being.'' And yet he holds fast that all force
is immaterial: ``We have defined force as a property
of matter and have seen that the two are inseparable; yet
the two lie \emph{according to the concept} far apart from each other,
indeed in a certain sense negate each other. At least we do not
know how one would define spirit and force otherwise
than as something immaterial, which in itself excludes
matter and is opposed to it.'' --- But how can
% --- p. 109 ---
\setcounter{footnote}{0}%
\opage{109}something that in itself excludes matter be a property
of it?

Just as the immaterial force is nonetheless to be a
property of matter, so the immaterial
soul is to be a product of matter. Büchner refers
here in particular to successive development. Ideas arise
through man's perceiving what is common to things in
the external world, forming from it an ideal shape,
and finally attributing to this the predicate of the true, the beautiful
or the good. But this constantly presupposes the main thing:
the very infinitizing, idealizing faculty
of forming such a shape. By presupposing
it one presupposes precisely what was to be explained. Because
art borrows its single features from nature, the ideal
does not therefore rest merely on a combination of these single
features. Because, e.g., the Japanese consider black teeth a
beauty, while it seems to them abominable to
have white teeth ``like a dog,'' the
aesthetic ideal is not therefore something purely accidental and external.

In Büchner the different directions are thus refracted.
He belongs to materialism by principle, but involuntarily
the true conception stirs in him, and he
thus comes into contradiction with himself. If
this refraction and self-contradiction is now thought through to the end, while
materialism is nonetheless held fast, then one comes
to a point where one no longer starts from experience,
but from the inner conviction of the truth of materialism,
by which conviction the
pronouncements of experience must then be governed.

2. \emph{Ethical-Speculative Materialism}. ---
This name seems to contain a self-contradiction. Materialism
seems to be so closely bound to empirical inquiry
that the latter cannot be overstepped without materialism too
being given up. Yet the very section on
empirical materialism must already have taught us to see the incorrectness
of this generally widespread opinion. It turned
% --- p. 110 ---
% print: og Lydphænomener (misprint: og doubled over the line turn (og|og))
\setcounter{footnote}{0}%
\opage{110}out that the materialists could appeal only wrongly
to the authority of empirical science. On the one hand empirical science does not
engage in dogmatic views,
and on the other hand, when all is said and done, exact empirical science seems
to lead in an idealistic rather than a materialistic
direction. This will appear especially from its treatment
of the concept of matter and from its doctrine of sense perception,
and will show itself in the 5th chapter of this book
from various sides. It is therefore a testimony to
a rare freedom from prejudice when the last of materialism's
champions whom we shall bring forward here, \emph{Heinrich Czolbe}
(a physician, like Büchner), declares it a proof of
ignorance to assert that modern natural science
leads to materialism. It is his conviction that
when one, like Vogt and Moleschott, accepts modern
physiology's doctrine of sensing, one has no
other way out than to throw oneself into the arms of some
speculative system or other. Since materialism has so far
combated theology and speculative philosophy,
the most important task therefore now remains, namely to
combat modern natural science\footnote[1]{``Die Entstehung des Selbstbewusstseins''. Berlin 1856. --- ``Die Elemente der Psychologie vom Standpunkte des Materialismus'' in the 26th volume of ``Zeitschrift für Philosophie und philosophische Kritik.''}.

Modern physics and physiology teach that the sense organs
are affected from without in a purely mechanical way, but
that these physical effects are converted in the organism itself,
in a manner differing according to the peculiarity of the individual sense organs,
into subjective sensation.
As long as one adheres to this doctrine, one has, according to Czolbe,
not got the better of the supersensible. The transition
from the physical effects of impact and pressure to phenomena of light and
of sound stands as something inexplicable and leads
to the assumption of a peculiar subjective principle in
% --- p. 111 ---
\setcounter{footnote}{0}%
\opage{111}the sentient being. The specific energy of the sensory nerves
is a completely mystical entity. By virtue of
the principle of materialism there cannot be such a gulf between
the physical agent and the psychical sensation.
The sensory nerves can only be passive substrates, and
the sense organs have only the significance of reinforcing
the motion before it reaches the nerves; and it must
be essentially the same motion that propagates itself from without
right through without any qualitative change.
But on what, then, does the difference between the
different senses rest? This question, which in
Czolbe's opinion one usually cuts short instead of answering,
he answers by setting up the concept of intensive quantity,
i.e. a quantity that is measured not by number and volume
(extensively) but by speed and strength. The difference
between the different senses now rests on this,
that not only do the different motions (sound, light
etc.) propagate themselves with different speeds, but
also the nerves of the different senses (auditory and optic nerves
etc.) have different elasticity, so that the motions
from without are propagated in different ways. Just as
intensive quantity takes the place of specific quality,
so specific elasticity takes the place of specific energy.
And just as the different speed of the motion is to explain
the different senses, so its
different strength calls forth the different feelings and drives:
too great a strength calls forth pain, too little need and
drive, a middle strength the feeling of pleasure.

All inner experiences (sensations, feelings etc.),
which are in this way explained purely mechanically, are comprised
under one common property, consciousness. This consists
in the direction of all inner experiences that runs back into
itself, by which they form an indissoluble unity. Consciousness
then presupposes an organ that can give the
motions propagated to it such a direction running back into
itself. Such an organ is the
% --- p. 112 ---
\setcounter{footnote}{0}%
\opage{112}brain, and the subject of consciousness (that which is conscious) is
therefore the motions that are communicated to the brain by the senses.
The unity of self-consciousness is expressed by these
motions circling beside one another; it is not
necessary to assume a centripetal direction of the motions
in the brain.

One need not be very familiar with physiology
and psychology to see how boldly Czolbe comes
forward against all views hitherto accepted. Czolbe is in
this respect very open-hearted. He states outright
that his doctrine conflicts with ordinary physiology
and is reactionary, since it defends a view that
natural scientists regard as antiquated. But he declares
at the same time that when, from the standpoint of materialism,
which excludes everything supersensible, one is to give an explanation
of sense-sensation and consciousness, one must
come to results like those now presented. The
deepest point of difference between materialist and
spiritualist philosophy lies, in his opinion,
in the different settlement of the question whether the
external agentia propagate themselves mechanically through the nerves
to the brain or whether they only touch the sense organ and in
it undergo a transformation. If he had attributed decisive
significance to empirical science, he would have had to
suspend the whole problem of materialism and spiritualism
in order first to come to a conviction about the
direction in which experience, exactly conceived, points.
But he starts from the principle of materialism and deduces
from it how the factual connection must
be. Therefore his philosophy of nature is as speculative as
Schelling's and Hegel's. The principle is carried through by a
series of hypotheses, of which the one brings the other
with it. And the most remarkable thing is that Czolbe, precisely
by carrying through his principle in this way, is led
beyond it. Sensing is explained by intensive quantity,
consciousness as rotating motions; but then
% --- p. 113 ---
\setcounter{footnote}{0}%
\opage{113}one must also conclude conversely: where there is intensive
quantity and rotating motion, there is also sensing
and consciousness. The whole of nature thus becomes ensouled.
This, says Czolbe, is only a real conception of
pantheism. When one conceives spirit as nature,
one must also conceive nature as spirit. Although
he thus sees an inner necessity in the reversal of materialism into idealism,
he is nonetheless himself surprised by
the consequences of his principle: ``Scheint es mir doch
selbst, dass ich durch die Konsekvenzen, zu denen das
Princip mich zwang, in eine märchenhafte Gedankenwelt
gerathen bin''. The point is that the illusion on which
materialism rests here comes to light.
Czolbe explains sensation and consciousness no more
than the physiology he rejects. What is decisive in
sensation and consciousness is the inner apprehension and
representation of the external real. But in Czolbe's doctrine
we constantly see one external real acting on
another; even consciousness is to be only a certain motion
of external parts in an external space. Here, then, no
transition is made from the external to the internal. Czolbe has himself
barred his way to it by positing an absolute external. There
is no absolute external except in the materialist's thought ---
and precisely because it is in his thought, it is after all not
absolutely external; without knowing it, he thinks thought
along with it when he thinks the external. Thus in Czolbe
the circular motion is at bottom nothing other than a symbol of
the inner unity of spirit. It is his need for intuitive
concepts that brings him to let this symbol
crystallize into a fixed finite form, which he then
regards as the thing itself.

For the principle of materialism, which he determines negatively
as the exclusion of everything supersensible from thinking,
he determines positively as springing from
``das geistige Bedürfniss nach Anschaulichkeit des Denkens.''
Already from this it can be seen that his point of departure is
% --- p. 114 ---
\setcounter{footnote}{0}%
\opage{114}a subjective one. But an intuitive explanation is for Czolbe
one with a mechanical explanation; he uses outright
``physikalisch'' and ``anschaulich'' as synonyms. Yet
it turns out, in the very concept of matter, that the explanations of physics
do not always fulfill the demands of intuitiveness,
at least in Czolbe's sense. He therefore demands
extended, bounded, impenetrable atoms,
and rejects the centers of force of physics as ``a supersensible
concept contradicting the principle of materialism.''
He further wants to explain the mutual attraction of bodies
by means of ``intuitive fundamental conditions,''
without employing mysticism --- in the following way: bodies
are separated, but are nonetheless all enclosed by space, and
attraction is the necessary intuitive intermediate determination
between connection and separation. This
explanation, whose merit Czolbe places especially in the fact that
it excludes force, that supersensible concept (``die
moderne Phrase: `kein Stoff ohne Kraft', ist eine Unmöchlichkeit,
ein Absurdum''), is a counterpart to the
speculative philosophy of nature's construction of matter.
Speculation constructs matter out of force, Czolbe
force (attraction) out of matter. Here the speculative
character of his materialism comes out clearly.

The principle of intuitiveness stands in Czolbe in connection
with a whole ethical world-view, in which,
as he himself is clearly aware, the ultimate motive
of his materialism lies. He repeatedly reminds us
that the principle of intuitiveness is akin to
the whole spiritual direction of Greek antiquity, and by his own statement
it was the naturalism of the poet Hölderlin, ensouled by enthusiasm for
Greek antiquity, that awakened in him the first
seeds of his world-view. What is harmful
in the present state of culture consists, in his
opinion, especially in the remnants of the supranaturalist
world-view. Transcendence, in particular the doctrine of
the personal God, he regards as a remnant of the egoism of
% --- p. 115 ---
\setcounter{footnote}{0}%
\opage{115}paganism, and he therefore regards materialism as
the completion of Christianity. The Greek contentment
with earthly life conflicts with egoism and presupposes
much self-mastery and resignation; it deserves
the name of piety far more than what is generally
so called. Czolbe finds his ethical and cultural-historical
fundamental view expressed in Strauss's
well-known words (in the essay on Julian the Apostate):
``The free, harmonious humanity of the Greeks and the
self-reliant manliness of the Romans is that toward which we are now
again in the process of working our way out of the long
Christian Middle Ages, enriched with the spiritual and
moral yield of that age.'' Dissatisfaction with the world of
phenomena is a moral weakness. But it
is not only ``der Blödsinn der Transscendenz'' that
Czolbe wants to combat; it is also the one-sided preponderance of the abstracting
understanding, not only in philosophy
but also in the exact sciences. He wants to bring about
harmony between intuition, feeling and understanding, and
thereby in the whole personality.

To such a degree is the ethical conception of the principle of materialism
the one that underlies it in Czolbe that he
in the last instance grounds it by an appeal to feeling.
The immense gulf between materialist and
spiritualist philosophy is produced by the different bent of heart,
conditioned by different experiences of life, in
aesthetic as well as in ethical respects. In reality
a credo ut intelligam always takes place; feeling
governs knowledge; a purely objective knowledge
does not exist. Therefore Czolbe harbors much sympathy for
the classic theological combatant of materialism,
Fabri; like him he sees in the results of knowledge
the expression of the direction of the will. It is a genuinely
theological proposition when Czolbe states that it is ``folly
in the superlative'' to want to philosophize about materialism
without an inward conviction of its truth!

% --- p. 116 ---
\setcounter{footnote}{0}%
\opage{116}Czolbe is a Stoic. For the Stoics too the
ethical-psychological motive was the main thing; they seek a doctrine
that can satisfy their practical tendency, and thereby come,
in a way similar to Czolbe, to a materialist
view. The self-contradiction of both rests on this,
that they postulate the necessity of materialism out of an
ideal principle. But in this they have revealed the secret
of materialism. Even Vogt and Moleschott at
bottom start from something ideal: the truth and right
of knowledge and the purpose of historical development. Materialism
not only turns, when consistently carried through,
into its own opposite, but even presupposes this
as its motive. We have dwelt at greater length on Czolbe
precisely because he shows this so clearly. Not only is he
the noblest and most consistent among the champions
of materialism, but he also has the clearest consciousness
that its significance in the present is predominantly
one of cultural history\footnote[1]{Czolbe on the whole comes forward with a modesty and resignation rare among the materialists. He does not begin, as they do, by condemning all earlier doctrines and institutions; here too he is more consistent than they, in that he recognizes what is historically given as necessary, while they wanted to be reformers. In inward conviction of the truth of his cause he is also assured that it will prevail when natural development leads to it. Therefore he is tolerant, which is likewise no materialist virtue, although consistently it would have to be one.}. Modern materialism represents
the most extreme reaction against absolute
idealism. It is a radical expression of the realistic
direction of the age, a memento of theoretical and practical
problems that demand the most serious reflection. And
if it should be the case, as a pessimistic
view might lead one to think, that we are heading toward a
new Middle Ages, ruled not, like the first, by
% --- p. 117 ---
\setcounter{footnote}{0}%
% print: Spiritualisme. men (misprint)
\opage{117}spiritualism, but by materialism, then we can in any
case console ourselves that this has already declared
its own untruth both in its principle and in its consequences.

\chaphere{CHAPTER FOUR.}


Faced with a grand development of thought such as absolute
idealism, whose one-sidedness has become evident,
different standpoints can be taken. One can reject it altogether
and, without entering into any negotiation whatsoever
with it, seek new points of departure for oneself. So
materialism, which with barbarous ignorance looks down
on the spiritual movement of the preceding age. One
can also work one's way out of it by dissolving the unity
to which it had led back the elements of existence,
and then, by way of formal consequences,
carrying each of these elements out into its abstract
extreme. So the critical consciousness. But in opposition
to these schools the endeavor is then natural
to correct the overcome standpoint as far as possible
by new elements which do not dissolve it, but
bring it nearer to truth and reality. This
endeavor saves historical continuity. One-sided
criticism cuts the thread and prevents the real
appropriation of the results of the development. For this
to be possible, a sympathetic dwelling on what came before
is necessarily presupposed.

It might well seem as if such an endeavor
were impossible with a system of thought like absolute
idealism. Precisely by its absolute character this
system seems unfit to enter as an element
% --- p. 118 ---
\setcounter{footnote}{0}%
\opage{118}into a more comprehensive connection. It would mediate
everything, but seems not itself capable of being mediated. All or
nothing, so it seems, must be the word here, as indeed the
critical consciousness, too, in the course of its development, at breakneck speed
converted Hegelianism's All into an absolute Nothing.

But it was precisely, as already remarked earlier,
the error of speculation that it posited as absolute what
is only relative. Speculation was right that the essence of
thought is not purely subjective and formal, cut off
from all inner connection with reality; thought
is an element of reality just as much as external reality,
which in its pure objectivity is also only an element.
By developing its own essence and content according to
their inner, necessary laws, thought therefore to that extent develops
the essence of reality. But precisely this \emph{to that extent}
speculation overlooked. It did not consider that
that development has only ideal validity, but does not exhaust
the whole essence of reality. Thereby speculation became
panlogism; the content of reality was seen
fully expressed in the content of thought. By an optical
illusion ideal and real reality were confused.
The ground of this illusion the inquirers we now pass
on to seek in the fact that speculation, in the course of its grand development,
forgot its original point of departure, the Kantian
critique of reason. Kant had precisely carried through and
held fast the distinction between thinking and being, the
subjective and the objective, and that so absolutely that, as
we have seen in the first chapter, the whole force of speculation
was needed to overcome it. But this distinction
rested on the irrefutable truth that human
knowledge must turn critically against itself and
not immediately regard its own forms as real
being. Only then is there also room for an influence of experience,
in principle, on the whole world-view.
These two sides, the abandonment of ``the monism of thought'' in the sense
of an immediate unity of thinking and being,
% --- p. 119 ---
\setcounter{footnote}{0}%
\opage{119}and the assertion of the real element, are closely connected
and point back to Kant\footnote[1]{Cf., with regard to the significance now again ascribed to Kant: C. H. \emph{Weisse's} academic address: ``In welchem Sinne die deutsche Philosophie jetzt wieder an Kant sich zu orientiren hat.'' Leipzig 1847. I. H. \emph{Fichte}: ``Bericht über meine philosophische Selbstbildung'' (Vermischte Schriften 1ster Bd. 1869). --- In Berlin \emph{Schelling} recommended the study of Kant's Kritik der reinen Vernunft as the foundation for understanding his lectures.}.

The question then arises what validity
and significance speculative knowledge can now
have. This question is answered differently by the inquirers
whom we shall bring forward as representatives of this
school. Their first development fell under the dominion of speculative
philosophy, and they then each seek in
their own way to hold fast what was justified in it. This endeavor
appears above all as a more concrete development
of the idea of God, which is why we have designated this school
as speculative theism. Although on
several points an influence of positive-dogmatic
notions shows itself in these thinkers, one would nevertheless do
them an injustice in thinking that underlying their endeavor
lay the tendency to assert the dogma of a personal
God at any price\footnote[2]{In Weisse's work ``Das philosophische Problem der Gegenwart. Sendschreiben an I. H. Fichte'' we read: ``Ich weiss wohl, dass Sie sich dieses Ansdrucks [der göttlichen Persönlichkeit, der Persönlichkeit Gottes], den die Gegner freilich jetzt als das Schibolet Ihrer Philosophie, so wie nicht minder der meinigen zu betrachten lieben, in Ihren früheren Schriften nur anbequemungsweise, und nicht ohne den Vorbehalt, dass er auf Gott nur eine uneigentliche Anwendung leide, bedient habe. Auch ich meinerseits habe, obwohl aus andern Gründen und in anderer Weise als Sie, mich stets dagegen verwahrt, dass man mir das Bekenntniss des Begriffs der göttlichen Persönlichkeit ohne Weiteres und mit der Voraussetzung unterlege, als sei mir dieser Begriff ein für sich selbst klarer, und bedürfe es über die Bedingungen seiner Gültigkeit nicht \opage{120}erst einer nähern wissenschaftlichen Verständigung . . . als sei es mir im Allen nur um die Persönlichkeit Gottes quand même zu thun.'' (p. 10).}. Their endeavor has its
% --- p. 120 ---
\setcounter{footnote}{0}%
\opage{120}natural ground in their opposition to speculative
panlogism and pantheism\footnote[1]{Only in the case of the younger Fichte could there, as we shall see, be reason to speak of an outright tendency.}.

By the way of dialectical speculation one can, in
their view, reach only the general horizon under which
the whole of reality appears, the highest, most
abstract laws of all existence. To attain a knowledge of reality
in the proper sense one must take a
step beyond that circle of universal laws and determinations.
This step consists at once in an appropriation
of concrete experience and in the further development of
the idea of God from universal reason to the consciously willing
and acting being which alone furnishes the grounding
of real reality. At this point speculation passes
over into philosophical empiricism and philosophy of religion.
There is thus in the line of thought of this school a double
movement, an ascending and a descending one: first
one seeks upward to the highest principle, and when this is
won and posited as real, one seeks to explain the given existence
from it.

It follows of itself that the point where the one
of these movements passes over into the other must also
become the most difficult point in these views. The movement
of the concept is to lead to its own cessation, activity
to pass over into receptivity. This turning point becomes a
standstill, and the movement forms not an ellipse but
an angle. This comes, as we have already had occasion to remark above (p. 24 f.),
from the fact that one still
in many ways stands on the speculative standpoint
one seeks to get beyond. Not without reason
has a dualism been found in thus supplementing
% --- p. 121 ---
\setcounter{footnote}{0}%
\opage{121}speculative thinking with elements that cannot organically
be developed from it. On such a stumbling block, as a rule,
the attempts to reconcile the old with
the new do run aground. Yet the significance of this school does not therefore become
slight, and above all, as will appear, in respect of the philosophy of religion.


I. Schelling's \emph{Later} Doctrine. --- It is one of the heroes
of the brilliant period of speculation that we first
come to deal with. Schelling (b. 1775,
d. 1854), who at the beginning of the century came forward with
the work of genius of his youth, in which the fundamental thought of absolute idealism
was sketched in bold and large strokes, appeared
toward the middle of the century with a doctrine that
was at once to hold fast absolute idealism and
to lead beyond its domain into an entirely new world, hitherto
only dimly sensed. He is an example of uncommon
vitality. Early mature (Magister in his 17th,
philosophical author in his 19th year), he also laid
early the foundation of the system that has secured him a place
in the history of philosophy. But despite this early completion
in essentials, his thought did not rest;
it worked at an ever further-reaching development of
the original principles. The later doctrine arose
through such an organic development, in which the guiding motives
can be shown at every point. And yet the result
of the later doctrine comes to form a complete
opposition to the earlier doctrine. What earlier
was the Only and the Absolute is later reduced to something provisional
and relative. We shall now see in brief how
Schelling came to the view we have to deal with here.

At every time --- says Schelling in one of his earliest
writings --- two main directions in philosophical
striving stand opposed to each other. The one takes its point of departure
in what is and in the real and seeks to explain from it
thought and the ideal; the other conversely proceeds
from the ideal and subjective and seeks to reach the
% --- p. 122 ---
\setcounter{footnote}{0}%
\opage{122}real. Neither of these directions can reach an absolute completion,
since the opposition always remains standing on account
of the one-sided point of departure. In the Absolute,
on the other hand, both opposites must be one, and
true philosophy must therefore take up the absolute standpoint,
the standpoint of reason, which is related impartially,
indifferently, to those opposites. Insofar as these
nevertheless come forth in the finite, their relation becomes
yet only quantitative difference, a preponderance of the ideal
over the real (spirit) or of the real over the
ideal (nature). But this quantitative difference is
only in the phenomenon; reason, the Absolute, the Godhead
is in itself absolute indifference. Only insofar as
this absolute standpoint is abandoned and a relation of reflection is posited
does the finite world have reality\footnote[1]{``Briefe über Dogmatismus und Kriticismus'' (1796). ``Darstellung meines Systems'' (1801).}.

But the question cannot be pushed back why
it is necessary at all to enter into this
relation of reflection instead of holding fast the absolute standpoint.
In the work ``Philosophie und Religion'' (1804) Schelling
seeks to answer it, and thereby brings forward the point
where the germ of his later view lies. He denies
that there is a continuous transition from the Absolute
to the finite. The origin of the world of sense is
thinkable only through a break, a falling-away from the Absolute, for
there is in God no positive ground for the coming-into-being
of real things. In I-hood lies the principle of the falling-away and
thereby of finitude. The falling-away, which is eternal like
the Absolute itself, consists namely in a being in itself
instead of a being in the Absolute. Thereby
absolute indifference is cancelled and the opposites come into force.

There is thus an irrational element in the finite,
in that being does not follow from the concept, the particular not from
the universal, but something independent of the concept must
% --- p. 123 ---
\setcounter{footnote}{0}%
\opage{123}be added. This irrational is something lower than the Absolute. On the other hand it is impossible to assume a higher
potency above the Absolute; by that the latter would indeed precisely
be reduced to something relative. The Absolute as such is
God\footnote[1]{Already in a letter to Hegel of 4 February 1795 Schelling says: ``Für uns sind die orthodoxen Begriffe von Gott nicht mehr. Wir reichen weiter noch als zum persönlichen Wesen.''}. ---

For a further development a double possibility now
presented itself. If absolute identity is held fast,
the finite can only become something negative, an
outer, untrue form for the inner, positive unity of the Absolute.
The carrying through of this line of thought would, as Weisse
has rightly remarked, have led to the Hegelian view.
But the other possibility is this, that that irrational
element, which is the principle of the finite development of freedom,
is more strongly accentuated, regarded from a more
positive point of view. Then it will become necessary, as it were,
to let absolute identity fade and to find (what
Schelling in the treatise just cited expressly denied
the possibility of) a fuller and more active expression
for the highest principle; for absolute identity does not permit,
as we have seen, any positive development of freedom
in the finite. This latter way Schelling struck
out on already in a treatise of 1809: ``Philosophische
Untersuchungen über das Wesen der menschlichen
Freiheit und die damit zusammenhängenden Gegenstände.''


The system of identity, it says here, has overcome
the opposition between thought and its object, between
reason and reality. But there now appears
a higher form of the same problem, in that the question is asked
about the relation between freedom and necessity.
The investigation of this finds a point of connection in the
earlier investigations through the important separation which the
% --- p. 124 ---
\setcounter{footnote}{0}%
\opage{124}doctrine of identity made between the being as existing and
the being as ground of existence. This separation
came forth in the philosophy of nature, Schelling conceiving
nature as a scale of quantitative differences or
potencies, of which the one always forms the basis for
the following. In this distinction between ground and existence
he sees the decisive difference between his
own doctrine and Spinoza's. The ground is not God regarded
absolutely, i.e. as existing; it is that in God
which is not God, nature in God, his prius (though not
in time or essence). All things are in God, not in his
proper, actual essence, but in the eternal ground
of his essence. And the opposition between ground and existence
is not dualism, for it presupposes a primal unity, the primal ground
or, in Jakob Böhme's expression, the unground,
which lies prior to all oppositions. The primal ground is
the absolute indifference which we already know from the
earlier doctrine, the Absolute as absolute (das schlechthin
betrachtete Absolute). But this indifference is thus
only the absolute point of departure, the universal presupposition
for the concrete development, which consists in the opposites
being posited and again overcome to identity through
love, in personal existence.

Earlier Schelling knew no decisive opposition
between indifference and identity (absence of difference
and unity). In the authentic exposition of 1801 reason is expressly
determined first (§ 1) as indifference
and then posited as one with identity (§ 9).
In a somewhat later work, which was first printed after
his death (Werke 1ste Abth. 5ter Bd), he distinguishes more precisely,
so that the Absolute is called God, but reason
God's reflex, which is related to God as indifference
to identity. But here too he expressly holds fast
the pantheistic standpoint and treats all problems
from it. The same is the case in the work ``Philosophie
und Religion''. When Schelling, therefore, in the
% --- p. 125 ---
\setcounter{footnote}{0}%
\opage{125}treatise on freedom would have his whole earlier doctrine regarded
as comprising only what he there calls the primal ground
and the ground (nature), this rests on
an illusion; his earlier standpoint was absolute and
closed in itself, but he now sought beyond it, sought an
existence, a reality, which he did not find within
its bounds.

While his earlier view has an antique character,
partly through the weight laid on nature, partly
in that art and the state are seen as the highest
spheres of spirit, he now throws himself into a medieval
theosophy, in which nevertheless, despite much that is purely scholastic,
one recognizes the ingenious inventor of the doctrine of identity in
many extraordinarily profound investigations. But his
assurance and feeling of victory have left him. He has
aptly been compared, at his earlier standpoint,
with the great Napoleon, who with the assurance of genius
compelled everything to bow before superior power.
And this image has been carried further by
comparing Schelling in his later doctrine with Napoleon's
attempt at constitutional government in the 100 Days.
Here is the same wavering and indeterminacy in the whole
bearing, signs that stage and actor no
longer suit each other. ``Instead of the arrogance
with which the teacher of identity played with the crowned
heads of science, there has come an uncertain, indeed anxious,
deference to the judging public''\footnote[1]{\emph{Erdmann}: Ueber Schelling, namentlich seine negative Philosophie. Halle 1857. p. 9.}. Schelling
made many attempts to publish his new
doctrine, but always drew back again. Even printed
works he had destroyed, having second thoughts at
the last moment before publication. But at the same time
he published pronouncements full of arrogance and self-confidence and of great bitterness against Hegel's doctrine, which had
% --- p. 126 ---
\setcounter{footnote}{0}%
\opage{126}outflown his and now, by its rich and grand systematic structure,
its ingenious execution and its great promises,
ruled over minds for several decades. Only
some thirty-odd years after the publication of the treatise
on freedom did he come forward in Berlin before a wider circle
with his new doctrine, and in authentic printed writings it appeared
only after his death.

In one of his few public pronouncements in that interval
he determines in general outline the opposition
between his earlier and later doctrine\footnote[1]{Preface to the German translation of a work by Cousin 1834 (Werke 1ster Abth. 10te Bd).}. Here appears
for the first time the famous separation of negative
and positive philosophy. Philosophy must begin with
what must necessarily be thought, what cannot not be
thought. That without which no knowledge is possible
(the a priori) is the negative in knowledge;
that by which real knowledge is brought about
is the positive. The absolute prius, even for the Godhead,
is contained in the science of reason, which fathoms
that which is itself, the highest that is (subject-object).
Yet this must not be misunderstood, as it is by
Hegel (``ein später Gekommener''), as if this science of reason
were to be purely abstract, a merely logical process.
This misunderstanding has only furnished one more proof
that it is impossible to reach reality starting from the
purely rational. Not being, as Hegel would have it, but that which
is, is the beginning of philosophy; and when philosophy,
as negative philosophy, has developed this according to
all the possibilities of its essence, a development that leads
us beyond the opposition between rationalism and empiricism,
then the question arises about that which is no
longer negative condition but which is positive cause.
But here a new principle is then needed, a new beginning:
``I do not want that which merely is; I want that
% --- p. 127 ---
\setcounter{footnote}{0}%
\opage{127}which is, which is or \emph{exists}''. In this sense
a great and, in the main, final change lies ahead
for philosophy, which on the one hand will furnish the
positive explanation of reality, while on the
other hand reason is acknowledged to be in possession
of the absolute prius, even for the Godhead.

While the polemic against Hegel in this pronouncement
aroused offense among the Hegelians, the further
hints won approval with many who had already beforehand
entered into opposition to the Hegelian system. Yet
these hints were still rather indeterminate; one could
not yet see what character the opposition between
negative and positive philosophy would take. Never, therefore,
as has rightly been said, has Schelling been praised
from so many different sides; everyone interpreted his pronouncement
according to his own opinion.

In the year 1841 he was called upon to lecture
in Berlin --- in order to counteract the negative
criticism which had developed out of Hegelianism. This
system, which in its time could pass for the philosophy of the established order,
was now in radical opposition, and Schelling's
``positive philosophy'' was to be placed in the breach. From a
theological chair in Berlin his appearance was hailed
as a turn of Providence. The first lecture,
which he published, is remarkable not only for the self-confidence
and authority with which he expresses himself\footnote[1]{An opponent said of this lecture that in Germany it was read like a speech from the throne: ``and the likeness was great enough; the speaker expressed himself with much dignity about himself, made great promises and avoided everything that could embarrass him.'' \emph{Noack}: Schelling und die Philosophie der Romantik. 2ter Bd. p. 466.}, but also
for a clearer determination of the relation between negative
and positive philosophy. He declares that he has not
come to tear down, but to build up. ``Nothing
shall be lost through me of what has been
% --- p. 128 ---
\setcounter{footnote}{0}%
\opage{128}acquired for genuine science since Kant; how should I,
moreover, give up the philosophy which I myself founded,
the invention of my youth? Not to set another
philosophy in its place, but to add to it a new science
which has hitherto been considered impossible, in order thereby
to secure it again on its true foundation, to give back
to it the firm footing which it had lost by going
beyond its natural bounds, in that something was made
into the whole which could only be a fragment of a
larger whole, --- that is the task and the purpose.'' Philosophy
(i.e. the Hegelian) assures us that it is religious in
its results, but this is disputed from the side of life.
This disproportion between science and life, which has
called forth so strong a reaction, must be removed; philosophy must
be brought out of its difficult position. And
this Schelling knows himself capable of, since he is in possession
``not of a philosophy that explains nothing, but of one
that furnishes information that has been longingly
desired and urgently demanded, --- a philosophy that
extends the common consciousness beyond its present
bounds''\footnote[1]{Erste Vorlesung in Berlin. Werke 2te Abth. 4ter Bd. p. 360 ff.}.

The opposition between negative and positive philosophy
thus rests on the fact that the latter is in its essence religious
philosophy. And if this is compared with the utterances in the preface
to Cousin, then the relation becomes this, that positive
philosophy, through the religious, finds the way to a
conception and explanation of real existence. In
the posthumous writings (Werke, 2te Abth.) we now find
its further execution.

A distinction must be made between what is and that which
is what is, the existing subject. The first
is the real, the second the reality of the real.
Abstracted from reality the real is only possible, is
% --- p. 129 ---
\setcounter{footnote}{0}%
\opage{129}the possible content of reality. The real is a
What (qvid), reality, existence a That (qvod). To
the infinite possibility of being corresponds the infinite
possibility of knowledge, and if one proceeds from reason
as the principle in which both possibilities are one,
their opposition being sublated, then it becomes the task of the science of reason,
the first science, to go through
the infinite potencies of the real, to exhaust
their whole fullness. In that this first science paves
the way for the knowledge of the true, positive reality
by separating out what has only the possibility of existence,
it is critical, and corresponds to Kant's Kritik der
reinen Vernunft; it thereby becomes at the same time negative, insofar as
it is defensive against what merely is
(that which is not what is); and in that it occupies itself
only with what must necessarily be thought,
it follows the law of pure logical necessity and is
rational and a priori philosophy. But this rational
must not be understood as opposed to the empirical.
For in order that the possibilities of what is may be exhausted,
it is necessary also to think them as being
and real, whereby they become an object of experience.
Thought goes into experience and follows it from potency
to potency. The limit of experience is here, as Kant
taught, also the limit of thought, because the essence of both is
one. Since one now did not hold fast that what really exists
(that which is what is) has nothing whatever
to do with this whole development, which draws itself through
the realms of thought, nature and history in order to exhaust
the infinite content of the possibilities, --- since one
(Hegel) confused subject and predicate, existence
and potency, one came to the doctrine that God must
go through the spheres of nature and history in order to
come to consciousness of himself, a pantheism which
in its consequences must necessarily lead to atheism.
The doctrine of identity was also pantheism, and this
% --- p. 130 ---
\setcounter{footnote}{0}%
\opage{130}is to that extent also necessary for negative philosophy,
which constantly moves within the Absolute,
the domain of the universal Idea, and cannot go out beyond it.
But this pantheism has only logical significance; for
negative philosophy does not yet raise the question
about what exists, but keeps within the
pure concept.

What possibilities are now contained in this pure
concept? What is is, first (as primum
cogitabile), pure capacity, potentia, the subject that is
in itself; but while this that may be (— A)
constantly goes out beyond itself, has its being outside
itself, the second moment, corresponding to it, is what is
as that which purely is (+ A), which is actus purus, pure
object. As the unity of these moments, as the highest
Idea (summum cogitabile), what is is finally subject-object (± A), unity of potentia and actus, the sum
of all possibilities. All three moments together are
the Absolute, but, mind you, always only as
possibility, as Idea. Here thought then comes up against its
last limit, in raising the question
what it is that is the Idea, is the highest thinkable.
But this question cannot yet be answered; it
cannot even come forward in a decisive sense at
this stage.

By seeking out the potencies in the Absolute and following
them in their possible functions we do not get outside
the Idea. That which may be will, when it is put into
function, spread itself in quantitative infinity like the
Greek ἄπειρον, without measure and bound. Now in that the
second moment, that which purely is, as actus and determinate
limit (πέρας), reacts against this and brings negation
into the infinite becoming, it produces fixed qualities
and individual forms (τὸ περαῖνον, which is at the same time
τὸ εἰδοποιοῦν). If finally the third, higher moment comes forth,
then we have the Aristotelian οὗ ἕνεκα, the unity of
% --- p. 131 ---
% print: deri er Anelse (misprint: er for en)
\setcounter{footnote}{0}%
\opage{131}form and matter. The relative unity of the three moments (about the absolute unity one can only ask
at the end of the negative science) is the soul, as
world-soul and individual soul. Through the soul
everything that is stands in connection with the Idea and that which is
the Idea. With these four principles the whole world of Ideas is
given. In that the one moment points to the other,
there lies in this a presentiment that there is a higher being
than mere possibility, being bound to the Idea. But
although every transition from potency to act is a willing,
and we thus have in the relation between the moments in the Absolute
a primal will, as all being is only through
a willing, a self-assertion (object = resistance), --- yet
this willing is still only, as was already
expressed in the treatise on freedom, a hunger and thirst
for existence, not existence itself.

If we are to reach the point where this infinite craving
is satisfied, we must leave the Absolute, the standpoint of what merely
is. The way away from this we find
by calling forth the last possibility that still remains:
the possibility of a being outside the Idea. This
possibility cannot be sought in the Idea itself. It is I. G.
Fichte's immortal merit to have shown how
the real world first arises by the I's stepping out from
immediate unity with the Absolute into
the relation of reflection, positing itself as self-consciousness and
thereby in opposition to an objective reality (cf.
above p. 122). The possibility of this lies in the fact that the soul,
which according to the foregoing is the relative unity of the potencies,
is this only insofar as it does not tear itself loose from
the Absolute. If it does not preserve the relation of unity to
this, but attempts to make the relative unity into an
absolute one --- and this happens by its giving up God (that
which is what is) and thinking, in its indwelling in the Idea, that it has
and is everything --- then there thereby arises a world outside
God, whose step-by-step development the science of reason must
% --- p. 132 ---
\setcounter{footnote}{0}%
\opage{132}follow as the last series of possibilities that lie
within its domain. Hereby first arises matter
in the physical sense, as vis inertiæ, and time and space
as forms of the being-outside-themselves of things (eccentricity,
displacement from the center). The world now comes
to stand as alien over against conscious spirit, which,
by an act that lies beyond the limit of consciousness,
and which it therefore cannot grasp at this stage,
is set in this opposition. The only way to get beyond this
by its own powers is to penetrate and
overcome the world by its knowledge and its action.

The ancients already expressed this in a profound way.
The figure of Prometheus, one of the primal thoughts that
work their way forth of themselves'', denotes the human spirit as
that which by its very spark of divinity turns away from
the Godhead and by its own ingenuity seeks to found a
new world without God. And as Prometheus is chained to
the rock, where his liver is constantly devoured by the eagle, and
can only be freed from his torment when a god comes forward
to release him, so the human spirit now wanders
the long way in search of the real God, whom it has
lost, but who is the final cause, the personal
bearer of the Ideas on the possession of which it has founded
its autonomy.

The I has not only to struggle with the material
world which surrounds it; it also finds itself enclosed
by ideal bounds, as a member of a spiritual whole, where
the particular is subordinated to the totality, the preceding
to the following. This total order cannot have
any arbitrary or accidental origin; it is a
law which springs from the essence of what is, the source
of all right and duty. It is Kant's imperishable merit
to have brought forward this law and shown its
independence and self-sufficiency. Only insofar could
his doctrine of the autonomy of the law of reason be misunderstood,
as one might here, by reason, think of the
% --- p. 133 ---
\setcounter{footnote}{0}%
\opage{133}subjective, human reason. For the law, as said, has
its ground in what is itself. But as law of reason
it may not be traced back immediately to God; that would
be to deny the reality of rational ethics and thereby of rational
philosophy. --- As the factual foundation
for the realization of the rational law the state
(juridical legislation) steps forward over against the individual, yet
not as a restriction of his freedom, but on the contrary
to make it possible. The state is in its essence
merely serving, a means for a free social life, in which
the personal virtues can be exercised. With the help of this
firm foundation the individual is to win the inner independence
which is at the same time the best protection against political
oppression. The development of the state does indeed consist in
ever more of the universal law of reason passing over into the reality
of the state; but the goal of history cannot be
sought in a perfect state. Such a one belongs only
in the world of the ideal, of possibilities\footnote[1]{Earlier (e.g. in ``Vorlesungen über die Methode des akademischen Studiums 1803'') Schelling extolled Plato's state and placed the ancient life of the state above the modern, because it did not, like the latter, have private life outside itself.}.

The universal, impersonal power of law and the state is a
discipline to which man is subjected and from whose pressure he
seeks to free himself. Man is caught in the contradiction
that he has willed himself in opposition to the
Absolute, but that this willing nevertheless leads precisely to dependence
on impersonal powers. There then comes a
turning point where the I bends back again from outward-turned striving
into itself, to the contemplative life. It
now resigns, gives up reaching the goal by the way
of action and seeks by immersion in the Idea to press forward to
the God it longs for. The first stage here is
mystical piety, where the soul strives to give up
itself, its own will (se désapproprier sa volonté.
% --- p. 134 ---
\setcounter{footnote}{0}%
\opage{134}Fénelon) in order to sink wholly into God, with disinterested
love. To this is joined, secondly, art,
in which the human spirit, in completely selfless producing,
makes itself like the Godhead. Finally spirit seeks,
thirdly, its purest form, the one most freed from the potential,
in rational science itself. It now beholds
God as the highest thought, the pure νόησις τῆς νοήσεως, the absolute subject-object. But it is still
only a purely negative determination under which God
is here conceived. God is indeed seen in his pure self-being,
but only as the last to which one comes, not as
the living cause from which one can proceed. For Aristotle
God is thus expressly the end, the final
purpose, which itself moves nothing, does not go beyond
itself, but is ἄπρακτος τὰς ἔξω πράξεις. Therefore it is a
misuse here to speak of absolute spirit. We are
still at the merely ideal God, bound to the Idea, who is indeed
summum cogitabile, but nevertheless belongs within the domain of negative
philosophy.

How then can we overstep the bounds of this
domain? Not by thinking. For this follows
only the law of pure logical necessity and can therefore
not go beyond the Idea, but seeks only to exhaust
all its possibilities. It must be in the will that
the demand arises to posit God no longer as
mere Idea, but above all thought and reason. The absolute
reconciliation is not reached in the rational sphere. The development
within this led us at its height to the
contemplative life (mysticism, art and science); but
from this man is constantly driven back to practical
life: there must be action. Now the whole vanity of the present
condition dawns on him, and he
seeks beyond the Idea to the living and acting God,
who is no longer lost in what is (the Idea), but
is the lord of being. God is thrust out of the Idea, sought and
grasped as personality. It is not reason, the
% --- p. 135 ---
\setcounter{footnote}{0}%
\opage{135}universal, but the individual himself who longs for salvation
and finds it in the real God.

Negative philosophy, in running through
the scale of possibilities, became at the same time an encyclopedia
of the sciences. Every single stage within this
development corresponds to one of the particular sciences. These
thus lie on the way from negative to positive philosophy. Negative philosophy is the universal
science, which has no particular object,
precisely because it comprises the whole potency of knowledge.
But in that it has now completed the development and exhausted
the possibilities, and in that the transition to positive philosophy
is made by God's being freed from the Idea and seen
as that which is what is, philosophy acquires a
different character. It now acquires its own determinate object, becomes a science, not merely the science.
For as the absolutely Idea-free, God is not a universal,
but a particular, a This (ἕν τι, τόδε τι ὄν), not being
something, but being absolutely. There exists no universal
at all. Only the particular \emph{is}; the universal is only
insofar as the absolute single being is the universal, has
it for its predicate. It is wrong to say that the Idea
is the true, the only reality. The Idea is only insofar as
the Ideal is. The Ideal is the cause of being (αἰτία τοῦ
εῖναι) for the Idea. If the absolutely immanent concept (the
absolute Idea) is the end point of negative philosophy, then
the absolutely transcendent concept (the existing
absolute single being) is the starting point of positive philosophy.
Negative philosophy was a priori science,
in that it followed the potencies of what is in order to
exhaust their essence; its outline was the realm of eternal possibilities,
of eternal truths. It was only a
thinking, not a properly knowing science, and kept
constantly a hypothetical character: ``if what is
is, then it must be so constituted.'' A priori, i.e. prior to being. The realm of potencies has now been
% --- p. 136 ---
\setcounter{footnote}{0}%
\opage{136}traversed, and thought is brought to a halt at the absolute act,
which it cannot comprehend, precisely because everything universal, all
possibility, is excluded. Thought can only come to it,
must relate itself purely as beholding.

There is no immanent transition from negative to
positive philosophy. The realm of the immanent has been traversed.
The absolutely existing, the Idea-free God, is,
as independent of every Idea, also independent of the highest Idea
of negative philosophy. Therefore positive
philosophy comes to stand only in an accidental connection
with the negative. He who feels no urge to
pass over to the positive will rest content in the
negative, and conversely one will be able to pass over
immediately to the positive without first having run through
the negative\footnote[1]{Werke, 2te Abth. 1ster Bd. p. 564.}. Negative philosophy comes,
through immanence, when it understands itself, to
the transcendent being. Positive philosophy begins
with the merely existing, pure transcendence,
and now seeks to make it immanent, to show the
qvid that is borne by this qvod. In other words, while
negative philosophy developed the essence of the Godhead
up to the point where the question arose about the existence
of this divine, positive philosophy begins
with the absolutely existing and shows that it
is God. Positive philosophy thus begins with
a being that lies above and prior to all thought (das
unvordenkliche Sein), and sees this being as the bearer
of the whole realm of the potencies. By the development of the potencies
on this basis God is known as the lord of being.
The name speculation belongs especially to positive
philosophy. We are here no longer in the domain of necessary
thinking, but what is now at stake is a free,
i.e. speculative, thinking. To speculate means to look
around for possibilities by which a certain aim can be
% --- p. 137 ---
\setcounter{footnote}{0}%
\opage{137}attained in science. These possibilities (major premises
in a hypothetical syllogism) are shown to be real a posteriori.
For the method takes up real experience at every point:
``If the necessarily existing is
God, then a, b, c are real; now according to experience
a, b, c really exist; consequently the necessarily existing
is really God.'' --- The first question now becomes
that of the possibility of the real world.

The first potency, that which may be, has a
constant tendency to spread itself, as it were to overflow.
And since the development of the potencies is not to be merely
a purely necessary one, like that by which the second potency,
that which purely is (the Son), leads the first back
to the unity (of the Spirit) in God himself, --- since God wants not
merely to be God for himself, but also to communicate his
essence, just as genius constantly needs to produce,
--- there arises the possibility that he sets the potencies
whose unity he is in tension, in an opposition
that is to be overcome through a successive process. Hereby
the unity is turned around, and the universe (unum versum)
arises. But as yet only in God. The possibility of a
world extra deum (not merely præter deum) lies in what
negative philosophy already taught in its way, that
consciousness (man), in which the outflowing
potencies return again to God and which is thus
the bond between the world and God, can loosen this
bond again, because it wants to be lord of the potencies, as
God is. What man forgets is that while
God is by the necessity of his essence lord of the potencies and
of being, man is so only in a relative sense,
on the presupposition of his immediate unity
with God. The moment that tempts man again
to set the potencies in tension is again that first
potency, the infinitely bearing and producing principle,
which in the imagination presents to him the possibility of a world.
Herein works at the same time the power which Greek
% --- p. 138 ---
\setcounter{footnote}{0}%
\opage{138}mythology depicts as Nemesis, the world-law that
wills that everything shall be clear and without ambiguity, that no
accidental relation may persist, that untried happiness
is inadmissible. In Christianity this principle is depicted
as Satan, who sins ἀπ’ἀρχῆς, i.e. is the soliciting
principle that stirs in man's will so that this
may be tried. Conceived according to its philosophical Idea,
Satan is therefore the proper primum movens in history,
without which the life of the world would stagnate. Against
this power the second potency (Logos, the Son) comes forward,
not to negate it absolutely, but to subject it
to itself, to make it into matter and basis, and thus
lead it back again into an inner relation to God. It
is this process which the philosophy of mythology and revelation
pursues, whereby it (the most important part of
positive philosophy) becomes the highest philosophy of history.

The presupposition from which both mythology and
revelation proceed, and without which they are altogether
inexplicable phenomena, is that there has been an original
real relation between man and God. A psychology
drawn from the states of consciousness now usual
is as little sufficient here as the mechanical
laws that hold under the present state of the universe
can be transferred to the original becoming, the first
living coming-into-being. Herein lies the reason why all earlier
attempts to explain mythology and revelation
have failed.

When man, by an act lying beyond all
history and all consciousness, set the potencies in tension,
the original, real relation to God was abolished\footnote[1]{Another consequence of this catastrophe has been considered in particular in negative philosophy, namely the coming-into-being of an external, material world. On the other hand negative philosophy could not, according to its essence, go into how ``the \opage{139}lord of being'' stands toward that catastrophe, and what, according to it, the relation between him and man becomes, or rather, it could only conceive this from the negative side.}, and
the infinite power now works as a force that threatens
% --- p. 139 ---
% print: høire Sphære (misprint: høire for høiere)
% print: med at søndersprænge (misprint: med doubled over the page turn (p. 138 ends "truer med"))
\setcounter{footnote}{0}%
\opage{139}to burst consciousness asunder if it is not pushed
back within fixed bounds. The limiting power,
victorious step by step, which in the human spirit
takes up the struggle with that demonic force, and which
thereby itself works its way forth to reality, is,
as already stated, the second divine potency
(the Son). This theogonic process is for the time being purely
mythological. It aims at materializing the
first, chaotic potency, reducing it to a servant, to
matter and foundation. Polytheism is the first result
of this process. For since true unity has been put out
of force, the first potency wanting to rule alone, while
the other potencies constantly strive to take
their place again, the different stages form themselves,
in this strife, for consciousness into different figures of gods,
nodal points on the way consciousness must
travel in order to find again what it has forfeited.

The first stage is the worship of the stars. Here
the properly mythological has not yet begun, insofar
as the acting potency still stands at the transition
between being center and periphery, just as in
the planets, too, the original tendency to be center
is still present, and, by being constantly pushed back,
becomes the ground of the constant motion. The next
stage is that where the potency is reduced to matter,
to periphery, which is mythologically signified by
female deities appearing beside the male
(Urania beside Uranus: Herod. I, 131; the
Persian, Babylonian and Arabian religion). As the philosophy
of mythology on the whole is a philosophy of nature,
only in a higher sphere (both present, each in its
way, a theogony), so this point in the
% --- p. 140 ---
\setcounter{footnote}{0}%
\opage{140}mythological development corresponds in particular to what must be thought
as the proper beginning of nature, where that which at first
was itself something that is, is reduced to ground and matter (mater—materia).
Soon the higher potency, too, appears (Dionysus:
Herod. III, 8) in its first annunciations, still
obscure and not understood. But the god himself is older than
his name, his presence in consciousness is older
than his complete realization in it. Now
consciousness is filled with the struggle between the two powers, the
material and the ideal god (Phoenician, Phrygian, Egyptian,
Indian mythology). With the material god
consciousness fears to lose God altogether; only through
pain and death do liberation and the victory of the ideal god come about.
Hence the oppositions in mythological consciousness:
grief and lamentation because the god is dead, joy because he
is born again. Greek mythology finally forms
at once the conclusion and the recapitulation of the
mythological process. In the history of the Greek world of gods
we have, in the succession of Uranus, Kronos and Zeus,
a reflection of the struggle between the ideal and material
principle, until at last the clear Olympian harmony
governs everything. But Greek mythology has, besides
this exoteric side, also an esoteric one, in the mysteries,
whose content is the emerging consciousness of mythology's
own essence. Demeter signifies the higher, ideal
consciousness turned toward God. She grieves at Persephone's
death, because with her (consciousness as
that which is subjected to the material, mythological
principle) she believes she has lost everything. The mysterious reconciliation
then comes about in that, instead of the god who perished,
she receives the abiding one, who can never perish. But
herein lies a moment of the future: the mystical Dionysus
(Ἴακχος) is the ever-coming one (Eleusis --- Advent).
Here then is the end of mythology, where it has not only
explained itself but also pointed beyond itself.

The transition from mythology to revelation does not take place
% --- p. 141 ---
% print: førte og staaer (misprint: førte for første)
\setcounter{footnote}{0}%
\opage{141}by a new substantial content being taken up. The difference
rests on the acting principle. Both mythology
and revelation presuppose, as already remarked earlier,
a real relation to God lying beyond all conscious thinking and poetic invention.
But this relation can either be
merely natural and unfree, subject to necessity, where
then, on account of blindness and entanglement, superstition
soon arises. So in the natural, wild-growing religions.
Or else it is not merely a natural potency
that acts, but man is freed from blind,
unfree religion by its being a personality that comes forth
as acting. Therefore the person and history of Christ
are the true content of Christianity, whereas the
Christian dogmas are only products of a wretched
philosophy\footnote[1]{Earlier (Methode des akad. Studiums 1803) Schelling had precisely praised the Church Fathers because with the help of philosophy they had known how to develop so profound a system as the doctrine of the Church out of the meager content of the biblical books, which stands far behind that of the Indian religious books.}. In paganism the process took place purely
within consciousness. A subjective reconciliation was here
sufficient, insofar as only the consequences of the separation from
God were to be overcome. Otherwise when it is a matter
of abolishing the cause itself, of completely breaking the
divine displeasure (the first potency); this could
only happen through a real historical act. In place of
the ecstatic history in mythology there now enters, in the incarnation,
real, external history.

But the true understanding of this historical fact
has not hitherto been reached, because it was not put
in connection with the whole past of the history of religion.
It was not seen that Christ was implicitly also what
acted in natural religion. Through the theogonic
process of mythology the second potency has subjected the
first to itself and stands as victor in the human spirit. The question
now is how it, the Son of Man, as it
% --- p. 142 ---
\setcounter{footnote}{0}%
\opage{142}may rightly be called, since human consciousness
has brought about its becoming under this form, will use
the power and authority it has acquired (μόρφη θεοῦ): whether it
will keep this for itself (ἁρπαγμὸν ἡγεῖσθαι Phil. II,
6) or use it to lead the world back to God
(παραδιδόναι τὴν βασιλείαν τῷ θεῷ καὶ πατρί. 1 Cor. XV, 24).
(Cf. the story of the temptation). The becoming-man, where the
second potency, instead of godlikeness, takes on the
form of a servant in order to suffer death, is the answer to this. In
the incarnation and the historical events that attach
to it, the inner, transcendent, divine
history has broken through the bounds of external history, and
this now acquires significance only insofar as it can prove
its connection with that. Now the third potency, too,
the Holy Spirit, comes to break through,
in that the first potency no longer chaotically and satanically
withdraws from the higher animating unity.

The law of the history of the Church's development is contained in
the relation between the three chief apostles: Peter stands
parallel to Moses, is the lawgiver, the representative of the stable;
Paul is an Elijah, the representative of freedom and movement;
John finally is for Christianity
what John the Baptist was for Judaism,
the apostle of the future. The Petrine church is therefore Catholicism
with its fixed, law-bound unity; the Pauline
is Protestantism with its radical principle of freedom;
free unity has its type in John, who
unites Peter's simplicity with Paul's dialectical sharpness.
For the goal that is to be reached is that church which is for the first time truly universal
(if church is still the right word here),
which is built up in the spirit alone and consists only in the complete
understanding of Christianity and its real
fusion with general science and
knowledge. Free religion is indeed initiated but not
realized by Christianity. Consciousness must,
in order to reach free religion, free itself from
% --- p. 143 ---
\setcounter{footnote}{0}%
\opage{143}Christianity just as much as from natural religion; both are in
themselves only sources of unfree knowledge. But this
free religion is reached through the very understanding of Christianity,
namely when one sees that this is at the same time the
highest science, which leads to the final historical
connection of everything. Philosophical religion or
the religion of the spirit is thus one with the philosophy of religion,
and is distinguished from the so-called religion of reason
in that it has mythology and revelation
as its presuppositions and is their real completion. ---

Schelling's later doctrine has a decidedly historical character.
Everything here is development toward a highest goal. To be sure,
in the doctrine of identity, too, existence was not a dead being;
everything was life, namely the infinite self-knowledge of the Absolute,
its eternal joining of itself as
object with itself as subject. But it was a movement
that was one with absolute rest; since the point of view
was taken only in the Absolute, every inhibiting
moment vanished; there were no bounds, no strife. Gradually,
however, Schelling became aware of an irrational
element which burst this immanent unity and
opened the view to a world of separation and unrest.
And the more he fixed his attention on this
element, which for him became one with the essence of will, the
more the Absolute stepped into shadow, and it was now regarded
as that which lay prior to the concrete development, the
undifferentiated unity of the opposites which could only at the
end of the development be reconciled to real identity.

But the decisive opposition between the older and
the later doctrine rests, however, on the movement now taking place
about two centers. Negative philosophy leads to
a point where transcendence is put into force, and
positive philosophy then begins with the transcendent
principle in order to lead back again to the immanent. The necessity
of the double center is sought in the opposition
between the universal and the particular. Negative philosophy
moves within the universal and therefore does
% --- p. 144 ---
\setcounter{footnote}{0}%
\opage{144}not reach reality, for only the particular exists. The opposition
here becomes so strong that only the will yearning for personal
redemption is supposed to be able to make the transition.
Everything real, thought is supposed to be able to know,
but not that the real is (quid, not
quod). The determination of existence has, by being separated (``thrust out'')
from the universal, certainly become irrational;
but with what right does Schelling ascribe reality to so enormous an
abstraction as the distinction between existence and the
existing? Apart from the existing being,
existence is only a figment of the imagination. The absolutely particular
in this sense can only be beheld, not thought, and
insofar as Schelling here speaks of a free thinking, this
becomes only a fantastic speculation.

Weary of the grand movement in the world of thought,
the thinker, driven by an obscure urge, seeks beyond
it. Absolute idealism now becomes only an introduction
to the highest mysteries. Thus the Neoplatonists
began precisely where Plato and Aristotle ended.
The highest principle, the final cause of everything, is reached,
according to Plotinus, only when the soul is awakened by a light that is
above reason to seek beyond reason and to turn
toward something higher. What Neoplatonism is in relation to
classical Greek philosophy, Neo-Schellingianism is
in relation to speculative idealism.

In other thinkers, too, whose development has led
them to a break with speculative idealism, we find
a similar relation. Ludwig Feuerbach rejected
speculation as an apotheosis of reason, of the
universal in opposition to the particular and sensuous, and
corrects, just as Schelling does, the concept of being thus:
``Being is one with that which is.'' The question of
being is for him no longer a logical question,
but is a matter of life and death, is a secret of intuition and of love.
When he says that existence
has meaning and reason in it, although it cannot be
% --- p. 145 ---
\setcounter{footnote}{0}%
\opage{145}uttered, this is quite reminiscent of Plotinus's words about the relation
to the highest: ἔχειν οὐ κωλυόμεθα κἂν μὴ λέγωμεν.
In writings that are roughly contemporary with Schelling's
lectures in Berlin and with Feuerbach's ``Philosophie
der Zukunft'', Søren Kierkegaard distinguishes in a corresponding
way between knowledge and reality (``Philosophiske
Smuler'', ``Uvidenskabelig Efterskrift''): what
reality is cannot be appropriated in knowledge: all knowledge
about reality is possibility; existence is always the
particular. Here too the concept of being is corrected in opposition
to the volatilization of speculative idealism. In
Kierkegaard the determination of existence acquires the meaning
of ethically conditioned reality; the particular that
is is the existing individual, who stands under the
ethical-religious law. In Feuerbach existence is given
to sense perception; the existing particular is the sensuous
object. In Schelling, finally, reality is
the object of speculative beholding, which alone can be grasped
when one seeks beyond the usual limit of consciousness;
the existing particular is here, as in Plotinus, the
figment of the imagination that appears to thought when it,
weary of working through the categories, sinks into an
ecstatic intuition called forth by the unsatisfied
``hunger for existence''\footnote[1]{Yet another parallel suggests itself: between Schelling and Schopenhauer. Thus \emph{Edv. v. Hartmann} has recently emphasized that Schelling's later doctrine, as a unity of negative and positive philosophy, has the merit of comprising Hegel's principle (the Idea) and Schopenhauer's (the will) as equally justified and equally indispensable sides of the one principle. Philosophie des Unbewussten. 3te Aufl. p. 766.}.

But that this opposition cannot be held fast is seen
in Schelling from the fact that positive philosophy is not
content with the mere ``that'' (qvod), but builds a
whole transcendent development on the basis of ``das unvordenkliche Sein''. The opposition thus holds only in the
% --- p. 146 ---
\setcounter{footnote}{0}%
\opage{146}starting point itself. And even here we are led beyond it;
for thought must after all know its own point of departure, must
touch the terminus a qvo of its own movement; in
the point of departure itself, too, thought and being thus go together,
and being does not become ``unvordenklich''.

This might now also seem to be Schelling's
meaning, to the extent that he indeed places the task of positive philosophy
in making the transcendent immanent. This
takes place through the mythological process, which is a process of consciousness,
even though consciousness here is in the hands of an obscure,
fatalistic power, is outside itself (ecstatic).
But at the end of this process Schelling
again establishes the transcendent, even in miraculous form: the
ecstatic history, so that reconciliation can be carried through,
comes forth in the incarnation as real history.
But the possibility and necessity of this are not shown.
Since the religious relation is from the first determined as a
relation of consciousness, why then can reconciliation not
be carried through by a process of consciousness? Especially since
already the mythological process, according to Schelling, must
show the possibility of this, in that it has indeed step by step
pushed back the obscure, dispersing and tempting principle.
One could then imagine the development carried further
by the transition a potentia ad actum which is now Schelling's
favorite determination, imagine ``the second potency''
completing its becoming and thereby also making possible
the reality of the third potency (the Spirit), so that unity
and reconciliation are reached without a new transcendence.
But then one would certainly have to give up deducing the
ecclesiastical doctrine of the Trinity and of the incarnation. It is in order
to satisfy in this respect that Schelling has here
abolished the rational relation between potentia and actus.

The whole absolute opposition between ``what'' and
``that'' shows, in the case of Schelling, as in that of Feuerbach and
Kierkegaard, that the presuppositions of speculative idealism
are still held fast. Otherwise
% --- p. 147 ---
% print: Karakteristik (misprint: for Karakteristisk)
\setcounter{footnote}{0}%
\opage{147}Schelling could not think that thought, by developing its own
possibilities, at the same time developed the essence of things (qvid), that
experience and thinking thus coincided. Negative
philosophy in itself stands entirely on the standpoint of speculation,
and in particular the philosophy of nature is retained
entirely from the earlier doctrine. Just as uncritically
as Schelling in that doctrine related himself to real experience,
he relates himself in positive philosophy to
mythology and revelation. His doctrine rests here on
a very uncertain historical foundation\footnote[1]{One will easily convince oneself of this by a glance at p. 139 f. above. It should be emphasized in particular that Schelling conceives the succession of Uranus, Kronos and Zeus as a struggle between mythological notions that displace one another, whereas the notion of the struggle of the gods has sprung only from attempts to explain the conflicting powers in nature or to picture to oneself the origin and history of the gods (the Olympian gods).}. Characteristic
in this respect is the remark that, since his philosophy of revelation
rests on an extension of philosophy
and of philosophical consciousness which was possible
only by going back to principles, but historical
criticism does not know this extension of consciousness,
he can take no account whatever of
this criticism\footnote[2]{Werke. 2te Abth. 4ter Bd. p. 231.}. In other words, he treats criticism
in the same way as theology does. The mythological development
itself he sees as an inner, metaphysical one,
which is not conditioned or touched by physical
or psychological influences.

Nevertheless Schelling has the great merit of having
asserted the necessity of a positive and sympathetic
understanding of the historical religions. His philosophy of religion
thereby forms an opposition to the speculative-idealist one.
Feuerbach, too, forms such an
opposition; but he stressed so one-sidedly the lower,
pathological element in religion that it became impossible
% --- p. 148 ---
\setcounter{footnote}{0}%
\opage{148}to explain how it could also contain a real
kernel of truth. Schelling holds fast the truth of reason
as foundation (prius) and can thereby also assert
the truth of the religious relation.

What is justified in Schelling's concept of a positive
philosophy is on the whole this, that the existing stands
independently over against comprehending thought and must be acknowledged
by it, not dissolved and volatilized. But
this positive cannot be reached by any other way than that of
exact empirical science and of historical criticism.
And when experience acquires independent significance,
the existence of things (quod) will also prove to be a
necessary presupposition for thought's being able to comprehend their
essence (qvid). Schelling's emphasis on the concept of existence
is more a prophetic annunciation than a real
beginning. He has beheld in imagination what is
the task for laborious work in reality.

Schelling's positive philosophy is not orthodox. It
excludes not only theology, but faith too, as
positive faith, cannot possibly acquire real significance, since
positive philosophy is indeed precisely to lead beyond mythology
and revelation to the free, philosophical religion
which is one with the philosophy of religion. ``The philosophy of revelation
is more radical than is generally
believed. From the theological side one thought at once to be able
to use Schelling's new system as a counterweight to the
negative criticism, and he was, as we have seen, not free
from coquetting with orthodoxy. But the old thinker
has not denied his past to the degree that was
assumed from the most opposite sides, that of negative criticism as well as
that of orthodoxy. Theology has therefore
here, as so often in relation to philosophical systems,
confined itself to accepting individual premises or
slogans, while it cautiously holds back from the consequences.
These consequences themselves are, moreover, expressed
only in brief hints; but already by these a principle is
% --- p. 149 ---
\setcounter{footnote}{0}%
\opage{149}established which furnishes the best corrective
against the great thinker's own uncritical doctrines.

2. \emph{The Younger Fichte}. --- This name, too,
recalls one of the greatest figures in the history of modern philosophy. The recollection lies all the
nearer, since the son himself declares that he found the first impulses
to his own inquiry in his father's thinking
and personality. His inquiry thereby received a stamp which
it never lost, even through the various phases it passed through.
--- Johann Gottlieb Fichte (1762---1814)
seized and developed, as was mentioned earlier (p. 4),
with restless energy a central point in the Kantian
philosophy, and thereby carried it beyond its original
limitation and unclarity. He was thereby led
to a subjective idealism, in which he also found a
foundation for his practical striving. Heroic labor
for the sake of labor itself, and with ascetic contempt for
every resting in what is given and attained, was Fichte's
ethical principle. But later he sought, beyond this
subjective and practical idealism, this eternal and never
quieted becoming and acting, to grasp an eternal being; beyond
the strict necessity of duty, to find the religious,
the blessed life, where the Absolute no longer
stands as a distant, unattainable end, but is immanent in
man. From his purely subjective starting point he believed
he had arrived at this summit, in that he no
longer regarded consciousness as absolute, but as the
phenomenal revelation of the Absolute.

It was from this later form of the father's doctrine
that the son, Immanuel Hermann Fichte (b. 1797), took his
starting point. In studying his father's posthumous writings
he found the theoretical expression of what in his
childhood had met him practically in the personal way of thinking
and life of both his parents\footnote[1]{Cf. ``J. G. Fichtes Leben und literarischer Briefwechsel'' \opage{150}2te Aufl. 1862. Von I. H. Fichte, --- a book which Feuerbach (in ``Pierre Bayle'') rightly held in high esteem.}. For the mother
% --- p. 150 ---
\setcounter{footnote}{0}%
\opage{150}had entered in the most inward way into the whole spirit of her husband,
and because of his early
death it was she through whom this spirit had to act upon
the son. ``To her influence,'' he says\footnote[1]{Bericht über meine philosophische Selbstbildung. Verm. Schr. 1869. 1ster Bd. p. 32.}, ``I owe
everything of higher awakening, of unshakable assurance, that has remained with me all my life
long,
even if not philosophically interpreted, still less philosophically
comprehended. But the drive, the inextinguishable need,
was thereby implanted in me to see that ethical-religious element
justified and brought to understanding for me also by
way of the concept.'' This statement contributes to the explanation
of the strong as well as the weak sides of
the younger Fichte. By virtue of the motive he has here
given for his philosophical inquiry (it was this that
led him over from philological studies to philosophy),
he calls himself a philosopher of personality; the conviction
of the reality of personality and of personal life
is what runs through the whole of his long scientific
activity. But this personal interest,
which is present from the very first and demands its satisfaction,
also hinders in many ways the free, unbiased
view, the pure, consistent development. However much
merit the younger Fichte has earned by
his opposition to one-sided pantheism and idealism,
and later to materialism, the effect of it is nevertheless weakened,
on the other hand, by the tendentious, by
the immediate appearance of the conscious aim,
and in the results themselves, too, much that is
unclear and downright fantastic thereby emerges.

As the need to clarify his ethical-religious conviction
was the positive motive for I. H. Fichte's
% --- p. 151 ---
\setcounter{footnote}{0}%
\opage{151}philosophizing, so he has further declared that the
opposition in which this conviction was bound to come
to the reigning system, the Hegelian, has been a
``constant goad driving him forward. The absolute
idealism, as Schelling and Hegel had developed
it, was bound to repel him in a twofold
respect. First by its pantheistic
result, according to which everything single and individual was something
untrue and something vanishing; --- the reality of personality
was, after all, a fundamental truth for him from the first. But secondly
also by the postulate of an absolute knowledge.
For Fichte it was a necessary demand that the
validity of knowledge must first be established before thought
could raise itself to the Absolute. The subjective standpoint,
which Schelling and Hegel believed they had overcome,
still had its justification for him. --- ``I must
confess,'' he says (op. cit. p. 53), ``that the involuntary
protest against it [against the results of Hegel's system],
but also the wish to do justice by the speculative
way to the new demands that thereby arose, have since
that time become the spur that really urged on
my philosophizing.''

The study of the history of philosophy led him\footnote[1]{``Beiträge zur Charakteristik der neuern Philosophie'' (1829. 2te Aufl. 1841). ``Ueber Gegensatz, Wendepunkt und Ziel heutiger Philosophie'' (1832).} to
distinguish between a twofold series of views which
stood in decisive opposition to each other. The
one is that which, like Schelling and Hegel, at once takes an
objective starting point and from there develops a doctrine of
being as objective substance. This direction rests,
in its speculative as well as in its mystical forms,
on the unity of thinking and being, and leads in its final
consequence to pantheism. The other direction is
the subjective-reflective, which is chiefly represented
% --- p. 152 ---
% print: methaphysiske (misprint: so printed)
\setcounter{footnote}{0}%
\opage{152}by Kant; its one-sidedness consists in its being unable to
raise itself above the limits of subjective consciousness, being unable
to reach an objective knowledge. The true union
and mediation between the two directions he finds in the
philosophical system's beginning with the theory of knowledge
and only through it raising itself to the objective standpoint.
In his chief metaphysical work, ``Grundzüge
zum Systeme der Philosophie,'' the first part is therefore the theory of knowledge:
``das Erkennen als Selbsterkennen'' (1833);
the second part is ontology (1836), ``the science of the forms of reality
and, since only the Absolute is the real,
of the Absolute's forms of reality''; the third part,
``die spekulative Theologie oder allgemeine Religionslehre''
(1846), forms the conclusion, in that from the concept of the world,
as ontology developed it in
developing the forms of reality, a way is sought back to the idea of God,
the final explanation and grounding of the given fact of the world.


Consciousness must, so the theory of knowledge begins,
begin with itself, with its immediately given essence.
Only step by step can it develop and educate
itself to truth. Consciousness is itself the beginning,
because it is what is immediately certain, and the theory of knowledge
is the history of the development of consciousness to
thinking and under thinking. From its first standpoint,
where it is still lost in intuition, in immediate
unity with the objective, consciousness works
its way forward to representation, thought and knowledge. When
it now, at its highest stage, reflects on its own
essence and, with penetrating reflection and dissolving
skepticism, has annihilated every immediate content,
there remains only its own essence itself as the
sole certainty. The whole of objectivity has vanished into subjectivity.
Knowledge has therefore shrunk from the broad
manifoldness of external reality into a narrow circle, into
the formal knowledge of itself. This is a necessary
% --- p. 153 ---
\setcounter{footnote}{0}%
\opage{153}point of transition to the free knowledge of the truth.
That it can be only a point of transition,
--- that consciousness cannot hold itself at this point,
lies in the fact that by making everything nothing it would also sublate
itself, since it would then become consciousness of
nothing. Consciousness can have reality only insofar as
it has an absolute content as its basis, indissoluble by
reflection. This content is \emph{absolute}, because it
is not, like the finite objective, posited by consciousness,
but is itself its root and foundation; it is an
absolute being, because it is not consciousness itself,
but is in consciousness. The Absolute is the
eternal fundamental condition of all consciousness. Consciousness is thus only something
relative, the image and revelation of an unconditioned being,
which must be as surely as there is
consciousness at all. From this it follows also that speculative
thought does not stand in opposition to intuition; as
consciousness in its ultimate ground is a sensing of
the Absolute, so thought, when it is to grasp
all the Absolute's forms of revelation, must constantly take elements of intuition
into itself. The highest form of knowledge
is therefore speculative-intuiting knowledge, and all
speculation in the last instance is theosophy.

The determination of the essence and concept of the Absolute
the theory of knowledge must leave to ontology. The theory of knowledge
can only establish that the concept of the Absolute
has reality; ontology has the task of developing the
eternal forms in which everything concretely real appears.
But the thinking by which this development is carried through
is no longer, as in Hegel, abstract
and void of content; on the contrary, it has through the theory of knowledge
gained certainty of the Absolute, and now seeks
only to develop more fully what lies in its essence.
From the more abstract and one-sided forms it seeks, by
way of the immanent dialectic, ever more concrete
determinations. At its summit, when it has
% --- p. 154 ---
\setcounter{footnote}{0}%
\opage{154}exhausted the categories, the system of the eternal forms of the world,
ontology becomes a doctrine of Ideas or speculative theology, in that
it ends in something real that is no longer afflicted with
the sting of contradiction, but contains the eternal
archetype of reality, to which all those forms of the world stand
in a serving relation. In the doctrine of Ideas the negative dialectic gives way
to the positive; for it is now real relations that are
at issue. The movement within the world of Ideas is an archetype of
the real world-process. When, therefore, the doctrine of Ideas leads to
the result that the Idea of spirit, of personality, is the
highest reality, this is a truth which
the real development of the world also in fact establishes, which, however,
it is no longer the business of metaphysics but of real philosophy
(the philosophy of nature and the philosophy of history)
to establish. Hegelian philosophy was bound to lead
to pantheism precisely because it did not make this separation
between formal and real philosophy. In Hegel
the dialectical development of the categories is the Absolute's own
self-development; therefore the Absolute is first complete,
first absolute spirit, at the end of the development.
But it must be strictly maintained that the ontological development of concepts
separates what in God is eternally one, what he
is in unchangeable equality with himself, exalted above all
becoming and development. Different from this is the way
in which God's essence manifests itself within the
actual development of the world, by which that eternal being
in its absolute unity is untouched. (``Gegensatz von Ewigkeit
und zeiterfüllender Genesis''). As we can know
the laws of space apart from what fills space, so
we can also, through the development of the a priori
forms, know the essence of God. But this does not lead,
as pantheism thinks, to an adequate and absolute
comprehension of the divine, still less to our thinking's
being identical with the divine self-knowledge.
For this knowledge of forms does not yet reach
God in himself, the fullness of the divine essence
% --- p. 155 ---
\setcounter{footnote}{0}%
\opage{155}and its infinite freedom. While our ontological thinking
is purely a priori and dialectical, the divine
thought that reveals itself in the reality of the world is absolutely
positive and concrete. The opposition between the theocentric
and the anthropocentric standpoint must be maintained;
the first is absolutely inaccessible to us. ---

What is characteristic of this attempt to get
beyond absolute idealism is that, by giving
the objective knowledge of Ideas a subjective grounding
through reflection on the essence of consciousness, one modifies
its essence and results and turns it from absolute into
relative. The question then is whether the knowledge of Ideas does not
lose its speculative and objective significance by being supported
on a reflective theory of knowledge. Against this
point criticism also turned, not only from the camp of the Hegelians,
who saw in Fichte's attempt a return
to Kant's subjective standpoint, but also from a man
who otherwise in his whole direction seemed to stand close to Fichte
and was generally placed together with him. C. H.
\emph{Weisse} showed in his polemic against Fichte, ``Das
philosophische Problem der Gegenwart'' (1842), that
the starting point Fichte had taken in his father's subjective
idealism made it impossible for him to appreciate, still less
to take up into his view, what Schelling and Hegel had
taught about the essence of the Idea. The Idea of the Absolute, as
these thinkers conceived and developed it, is the
assurance, established in an intellectual intuition or a dialectical development of concepts,
that it is not the bare subject
that is what truly is and is real, but that
this requires the subject's passing over into the objective and its
return to itself again, in short the absolute
subject's infinite potentiation by being mediated
with its objectivity. But this Idea is so far from being
gained through a reflective theory of knowledge that
the latter must, on the contrary, always presuppose it. The theory of knowledge
cannot establish the reality of the Idea, since it,
% --- p. 156 ---
\setcounter{footnote}{0}%
\opage{156}as Fichte too admits, itself rests on it. One
in any case comes to move in a circle:
the reality of knowledge is to prove the reality of the Absolute,
on which it itself depends. Out of this circle
one gets only by distinguishing between the preparatory school the individual
must pass through in order to come to the truth, and
the truth itself, and then, in developing the latter, by proceeding
from what is in itself the original.

Yet it was not this consideration that later
moved Fichte to seek another grounding of
the Absolute and another transition from the theory of knowledge
to the other parts of the system. He maintained,
on the contrary, the necessity of the theory of knowledge's going
before metaphysics (the doctrine of Ideas, the doctrine of the Absolute).
But whereas the transition had earlier been immanent,
in that the essence of consciousness itself led back to the Absolute
as its ground, he later rejected this pure
immanence as pantheistic. God was by it only
seen as the indwelling subject-object, in which the finite
knowing subject was dialectically sublated through
the infinite act of knowledge. On the other hand, the
pantheistic principle is abolished when one regards the
result of the theory of knowledge as a new problem, in whose solution
we leave the negative, immanent development. This
problem then becomes: what is the ground of the identity
of subject and object on which all knowledge depends?

In his ``speculative theology,'' in which he in
this way corrects his earlier doctrine, he states
that the latter could only lead to conceiving God as the
spiritual unity of the finite world. It is chiefly studies in natural science
and anthropology that have led
him to the conviction of the reality of the individual,
that it is not a merely finite thing that is sublated in the Absolute's
infinite process. With this there follows, then, a
new method too. If this had earlier been dialectical-ontological,
in that it sought by way of immanent speculation to
% --- p. 157 ---
\setcounter{footnote}{0}%
\opage{157}develop the idea of God, there is now laid down as foundation a line of thought
corresponding to the old cosmological proof,
in which one infers from existence as a given fact
back to God as its ground and cause.

The unity of subject and object which knowledge
presupposes now becomes only a single fact,
which, to be seen in its right light, must be ordered into the whole
fact of the world. For this shows us as what
really is not, as Hegel thought, a series of
vanishing moments, but an infinity of real
positions, of primal qualities, which constitute what is original
and abiding in all change, the substantial in the finite.
Pantheism knows only two terms, the Absolute and
the finite that becomes (the phenomenal); but if the
real development of the world in its manifoldness is not to be reduced
to a mere semblance, a third term, a
middle term, becomes necessary, the concept of the substantial
finite, the real primal position. Pantheism is thereby refuted,
as it were, from below, starting from the immediate
fact of the world. But while the explanation of the
given thus leads to the assumption of an infinite, real
manifoldness --- an assumption that contains what is justified
in Herbart's doctrine and in his opposition to
speculative idealism --- it nevertheless soon appears that
this manifoldness only subsists, and only can subsist,
in and by means of an infinite unity that pervades
them all. Each position is not something for itself, but
stands in inner interrelation with all the others. In the same
way, therefore, as the fact of knowledge leads
to the assumption of a unity that embraces and conditions
both the subject and things, it now appears that everything
existing, in its manifoldness and individual peculiarity,
must be borne and grounded by an unconditioned
unity which contains the final ground of the whole
connection of the world as well as of every relative unity
% --- p. 158 ---
\setcounter{footnote}{0}%
\opage{158}and whole within it. The unity of the world can only
be a product of the creative primal unity.

The next question then becomes how such an
absolute unity of the infinite is possible. By assuming
the absolute unity one really only moves the difficulty
one step back, from the conditioned to the unconditioned.
On this question too Fichte stresses
that its answer must be sought on the basis of and in
analogy with what is originally given. He sets up, on the
whole, as his guiding principle, that only from the reality of the world
can the idea of God be known. --- The actual connection
between single beings proves, on closer view, to
presuppose an ideal connection in which the concept of the one
is conditioned by the concept of the other and requires it as
its necessary presupposition: for the one to be,
the other must also be. However separate single beings
may be in their external reality, ideally regarded they are
nonetheless present in and for one another. This ideal
connection is a teleological connection. The relation
between aim and means is, after all, precisely this, that the former
presupposes the latter as a necessary condition, while
conversely the latter first finds in the former its complete
explanation and grounding. This ideal, teleological
connection in turn presupposes a surveying and
comprehending thought which at the same time, in its real essence, actually
posits and encloses the existing manifold.
The Absolute has been thought of as a world-soul
that acts and forms in the world with omnipresent
but unconscious reason. The ideal comprehending is
here a beholding imagination, and the creative activity
consists in this imaginative activity's rising
to express, individually producing intuiting. So
in Schelling's philosophy of nature, where everything existing
is a particular form of the Absolute's self-intuiting. A
higher stage is reached when one sees, like Hegel, that the
Absolute cannot be lost in objective intuiting,
% --- p. 159 ---
% print: derved naaes (misprint: der lost before derved)
\setcounter{footnote}{0}%
\opage{159}but must be free as universal reason, absolute spirit. Yet
it is still always only the concept of a world-spirit that is thereby reached, and the conclusion becomes only relative,
poses a new problem. For the concept of the world-spirit
is only a concentrated expression of the problem of the world,
of the fact that in the universe the highest
unity of reason and reality is realized; but the
question remains unanswered how this is possible,
what the final ground of this fact
is. Herein lies the necessity of going beyond
Hegelian philosophy\footnote[1]{Cf. above p. 21 f. Fichte's characterization of Hegelian philosophy in an earlier work.}.

To this is added yet another difficulty.
Absolute idealism cannot develop the concept of spirit
in such a way that it comes to stand free over against the unconscious.
Unconscious life, nature in the narrower sense,
therefore remains unexplained. On the other hand, the knowledge
that all immediate differences of time and space, all immediate
existing whatever, is taken up into the teleological
connection, leads to the conviction that this whole
connection must be possessed by the Absolute as the archetype of the world
in eternal clarity. Inner purposiveness
thus leads of necessity beyond
itself. The archetype of the world, the eternal thought-cosmos, the all-consciousness,
leads back in turn to a unity of self-consciousness
as its ground. And so we stand at the
last link of the long teleological proof, which in its
course has step by step supported itself on the fact of the world.
Only a divine power of will can contain
the cause of the world, since the world is not one with
the pure concept. And such a power God possesses as
the self-conscious unity that eternally posits and unites the
real primal positions.

Both in the concept of the highest principle and in the concept
% --- p. 160 ---
\setcounter{footnote}{0}%
\opage{160}of the finite, correctives have thus been shown against
the pantheistic results of abstract speculation. And
both correctives stand in inner connection. The Absolute's
concentrated unity as self-consciousness would not
be possible if a true reality were not contained
in it, and on the other hand the finite would
be only an untrue semblance if its real ground were not contained
in the divine essence. The primal positions, which
here enter as the necessary intermediate determination
between the infinite and the finite, are therefore eternal,
``uncreated through God's eternal willing of himself.'' --- The question
now is how the primal positions are more precisely related,
on the one side to God, from whose essence they spring,
on the other side to the finite world,
whose substance they constitute.

As regards the first part of this question,
there lies in it a difficulty which Weisse already,
in his polemic mentioned above, has pointed out with great
acuteness. For since the primal positions are, by
an inference from effect to cause (via causalitatis),
led back from actual existence, whose individual
foundation they constitute, to the essence of God, this
movement really comes to the same thing that Fichte himself
reproaches Schelling and Hegel with, namely seeking an
explanation of the given by moving the difficulty one
step further back. But the difficulty is thereby sharpened rather
than eased. ``How,'' asks Weisse
rightly, ``can it be reconciled with the concept of God as spirit,
as free, absolute spirit, to think any particular
finite determination whatever, or, what surely comes to
the same, an infinity of such determinations, posited
in him in an eternal and unchangeable way?'' Fichte's answer
to this question is unclear and wavering. While in
some places it seems as if he conceives the primal positions
as eternal in their absolute, monadic singleness (``primally created''
thus = uncreated) --- a notion that would
% --- p. 161 ---
\setcounter{footnote}{0}%
\opage{161}lead to a fabulous doctrine of pre-existence --- in
other places it is said that ``the primal positions exist in the
eternal nature in an archetypal, but not in an expressly
individual way,'' that they cannot in the last instance
remain in their atomistic singleness, but must
go back into the unity of the divine essence as its
eternal fundamental forces. Yet he declares the question
insoluble ``in what express way God
unites with himself these beings that are substantial for us, whether
he posits them as perennial or takes them back as
fluid moments into his unity.'' Certainly
this problem can be pressed in such a way that it becomes not merely
insoluble but also meaningless. Yet it cannot be
denied that Fichte, in taking up from Herbart the concept of the primal
positions, through which, by his own declaration
(Ontologie p. 194), it first became possible for him to take
the decisive step beyond Hegel, --- has thereby brought
something heterogeneous into his view, which can only with difficulty
be reconciled with its other elements, stemming from speculative
idealism. Therefore Erdmann has also
characterized his doctrine as eclecticism. Eclecticism
is what every standpoint becomes that is unable to let the
new and foreign elements it takes up undergo an
inner transformation in being appropriated.

The system of the primal positions constitutes for Fichte the inner
nature in God. By means of this concept he believes
he gets beyond not only the pantheistic but also
the vulgar theistic view. The justification of pantheism
lies in the truth that without the world, without
the universe, God would not be God; with God's reality
a universe is posited in which everything is enclosed
and nothing can arise or perish. But on the other
side it holds: the \emph{immediate} reality is
not the divine; God would be God even without
this given world. --- Ordinary theism teaches,
in opposition to pantheism, a miraculous creation out of
% --- p. 162 ---
\setcounter{footnote}{0}%
\opage{162}nothing, an actual production of something that before was not
but now comes into being as the product of omnipotence. Against
this view Fichte holds to his father's well-known
words, that it is the fundamental error of all false metaphysics,
and that no one has said an intelligible word about it.
In comparison with such a view pantheism would
necessarily have the advantage; it stands one step nearer
the truth, since it does not violate the idea of God
by attributing to God an absolute arbitrariness. The true
theism is the one that has taken pantheism up into itself
insofar as pantheism is justified, namely insofar as it
asserts the connection of essence between God and the world.
It can, in opposition to the vulgar, dualistic view,
be called concrete theism or, with
an expression borrowed from Krause, panentheism\footnote[1]{Cf. ``Vermischte Schriften'' I. p. 275: ``By theism we understand the quite general thought that the absolute principle of the world can in no case be thought as blind, unconscious power, as universal substance or abstract, impersonal reason, but only as a being existing in and for itself, for whose fundamental property human thinking has at its disposal no other analogy or other expression than absolute self-consciousness.''}.

This leads us to the second part of the question raised
above, the relation between the ideal world
in God and finite existence, or in other words
to the real problem of creation. That Fichte denies
a creation in the vulgar sense, which here is also
the ecclesiastical-dogmatic one, we have already seen. The problem of creation
must take the following form for him: why
do the primal positions not remain in their eternal, unchanging
rest, but enter into the limitation of finite existence
and ever-changing becoming? --- The divine
will is not absolutely producing; it is only permitting.
This concept of the permitting will becomes,
in the doctrine of creation, the intermediate determination between the
% --- p. 163 ---
\setcounter{footnote}{0}%
\opage{163}pantheistic and the dualistic conception, just as the concept of the primal positions
was in the conception of the general
relation between God and the world. --- Instead of the
inner, archetypal way in which the primal positions are in
the eternal ideal world, they are now released into isolated, mutually
exclusive existence, placed under the form of space,
time and corporeality. The problem of creation
is not a special problem; it is no more incomprehensible
than the problem of the arising of any single natural phenomenon
whatever. But the permitting will is only
one side of the matter. For the primal positions themselves
strive for independence; there is in each of them a
Promethean germ, in whose awakening creation precisely consists,
a germ which makes every finite being
something peculiar, and which drives the single being to
produce itself, while it is produced by the divine will. The omnipotence of this will is thus no
longer transcendent; but it is not on that account less
unconditioned. ---

This clear and rational development of the relation
between God and the world comes nevertheless, apart from the
difficulty touched on earlier in the concept of the primal positions,
to suffer from an unclarity which stems from the way
in which God's personality is conceived. Fichte asserts with
good reason that the concept of an absolute self-consciousness
contains no contradiction that would not also
lie in the concept of a mere unity of substance or of an
impersonal world-spirit. The problem remains the same,
namely that of the relation between infinity and unity,
but it is solved most completely and clearly by the concept of absolute
self-consciousness. It is no inadmissible anthropomorphism
to transfer the determination of self-consciousness, of personality,
to God; on the contrary, only God is a
pure I, a perfect unity with himself, while
in the finite spirit the I is only the work of reflection and abstraction,
and the unity therefore formal and empty. What
% --- p. 164 ---
\setcounter{footnote}{0}%
\opage{164}needs explanation is therefore rather how
the concept of the I or of self-consciousness also finds its
finite representatives. --- But Fichte finds in the principle
of personality, as he applies it to
God, ``a specific more'' than mere necessity;
the concept of the divine will is not exhausted,
in his conception, by its being seen as the
highest unity of freedom and necessity --- that would,
as he puts it, give only an acting, not a
deed. The divine deed is, on the contrary, to
exclude positively the necessary, the ``nature-like.''
God could therefore have refrained from creating the world, and
he can at any moment return to his inner, self-contained
self-satisfaction, without the world; if the world
nevertheless subsists, it can only be by his will.
Fichte here not only puts his rational doctrine of creation
in danger, but also comes into opposition to the whole standpoint of more recent
speculative philosophy, which he nevertheless does
not want to deny, but on the contrary to carry further. If
everything is in the last instance to depend on a choice for whose outcome
thought can find no necessary cause, what
is it then but a return to the old
creatio ex nihilo? An unconditioned and ungrounded choice was
precisely the nothing out of which God, according to the doctrine of the church, created
the world. And Fichte thereby arrives at the notion
of a God finished once and for all, whose nature is only
seen as the product of will, a God without nature. He polemicizes
against Schelling's attempt to demonstrate an unconscious
life of nature within the Godhead itself. Nature is
thus for Fichte essentially something derived, not something original\footnote[1]{``Der Geist der Natur, weil er sich nur in unbewusster Weisheit des Wirkens zu zeigen vermag, ist eben darum der nicht-absolute, nicht der Geist Gottes.'' ``Beiträge zur Charakteristik der neuern Philosophie''. 2. Aufl. p. 765.}.
% --- p. 165 ---
\setcounter{footnote}{0}%
\opage{165}The great result that speculative idealism had
reached with regard to the idea of God, although still in
one-sided form, was on the contrary that the divine essence
could not be an abstract subject, a pure spirit without nature,
but that, as the final ground of everything, it must also
be akin in essence to everything. The younger Fichte
forfeits this result through his spiritualism, which
has no eye for the significance of the life of nature with regard to
the determination of the idea of God. ---

We follow Fichte further to his indications concerning the philosophy
of history and the philosophy of religion. As single beings
pass over from the eternal primal position into finite
existence, temporal development, what is really
central in these single existences becomes the spiritual primal disposition,
in whose peculiar form individuality consists.
This is what Fichte calls genius. Through this
ideal element the single existence stands in inner connection
with the divine essence. The ground of
evil lies not in finitude itself, nor in
an incomprehensible arbitrariness in man, but in the possibility
that the finite being does not reach the end
set for it. Since, therefore, neither finitude nor
sin entails an absolute opposition to God, or
leads outside God, God's action in existence
cannot take place in an external, miraculous way either. Through
the ideal primal disposition (genius) in the finite
the Godhead works, strengthening, furthering, potentiating, toward the end of the world's development. By this divine
action the isolated and
dualistic in finite existence is sublated step by step. All truly historical
deeds are brought forth by the divine
awakening of man's genius, by the enthusiasm
that seizes human freedom, so that, in
performing its own work, it becomes conscious that its
content lies beyond itself, must have been infused into it.
At the highest stage it is the religious-ethical Idea
% --- p. 166 ---
\setcounter{footnote}{0}%
\opage{166}that inspires, and the individual thereby becomes mediator
between the Godhead and the rest of humanity, religious
genius, prophet. The highest end to be reached through
this development, borne throughout by divine forces,
is the most inward and complete unity of
self-satisfaction and devotion. All salvation consists
in the conscious perfection of one's own essence in and through
love.

Divine action is divine precisely
because it does not come forth in external, crude palpability,
but wholly clothes itself in the freedom of single beings and the
laws of finite existence. The true miracle consists
in the way in which, through this inner development,
new Ideas can suddenly come forth and geniuses be awakened
to bear them. The constant miracle is the divine
spirit's entry into the finite, whereby
the finite becomes co-creator and participant in the
process of the world's development. On the other hand, no philosophical thinking
will stoop to assume miracles proper, interruptions
or abolitions of the order of nature; for these are only the
purposeless and spiritless, the sensuous parody of
those spiritual miracles. God's action is universal, not
bound within the limit of certain events in time.

Yet here in Fichte, just as in the development
of the idea of God and of the problem of creation, there occurs a break
with the rational. Just as he there demanded ``a
specific more'' besides the unity of freedom and necessity,
so here too he is not satisfied with
the determination of the philosophy of history, somewhat indefinite indeed but
nevertheless rational in its principle, but demands,
under the evident influence of positive religion, that
the development of the world must reach its summit in a God-man,
since the unity of the divine and the
human spirit is supposed originally to be thinkable only in
one subject. ``It cannot be a multiplicity
of I's in which God originally becomes man, but
% --- p. 167 ---
\setcounter{footnote}{0}%
\opage{167}as in the whole of creation nothing wavers in indefinite outlines,
but everything is pointed to the decided peculiarity of the individual, so the highest work of creation, too, must be wholly concluded and perfected in a single fact.
Divine humanity, whereby the spirit of God is in a
universal way to make itself one with the human race,
is a wholly unclear thought, as abstract
and half-hearted as the pantheistic concept of God, to which
it wholly corresponds, both resting on taking
the abstract, the merely universal, for the highest. The God-man
alone solves the confused riddle of history, in that his
existence and its historical consequences give us the effective
guarantee that in him the divine spirit is in the midst
of us and will as surely accomplish the moral redemption
of the human race as it has first sent the God-man.'' --- Yet this is to be only a general
postulate, arising with necessity; from it there are
to follow no further consequences regarding the
way in which it is historically realized, although
Fichte does not fail to remark in passing that there
arises from it ``an inner proof of a peculiar kind'' for
the credibility of the historical accounts that treat
of the God-man.

It is easy to see that all this conflicts with
the premises from which it is derived. But this is a conflict
that necessarily arises from the unclear way in which
Fichte conceives the relation between immanence and
transcendence. It is, according to a frequently recurring
turn of phrase, immanence itself that leads to transcendence;
but that implies, too, that transcendence
is on the other hand conditioned by and enclosed by immanence.
It is in particular a most superficial line of proof
that underlies the postulate of the God-man.
Strauss's conception of the Idea is pantheistic, to be sure;
but his statement is nevertheless correct, that the Idea
never wholly exhausts itself in a single existence.
% --- p. 168 ---
\setcounter{footnote}{0}%
\opage{168}On the other hand, the way in which the Straussian conception
is to be corrected lies in Fichte's own doctrine of the
significance of individuality. The Idea of the unity of God and man
is not enough in itself, but is to be personally appropriated
and to live in the single individuals. Each one is,
as the mystics taught, to become a Christ. It is this
emphasis on the Idea's indwelling in personalities that leads
beyond the pantheistic, not the fantastic demand
for a single and absolute realization of the Idea.

Other statements suggest, moreover, that it was not
Fichte's intention by that postulate to pay homage to the positive
dogma. He names it as a great benefit of
his upbringing that he had been brought up entirely religiously,
but that the doctrines of faith proper, with
their ``mysteries'' and incomprehensibilities, remained
remote from him. ``This sparing of what is wholly superfluous
preserved me from the dangerous conflict of having to break with
the authority of faith on entering
a more mature culture. The eternal, the universally human in faith
never left me. The external, historical setting I could
confidently leave to sifting criticism (Verm.Schr.
I. p. 327). Philosophy must, he asserts against
Schelling and Weisse, wholly free itself from the
theological circle of representations. Christian theism
is to develop into universal, humanistic theism.
``For we would have it considered that not only the
fundamentally changed and enlarged cosmological world-view,
but still more the incomparably freer survey of
the essence and common origin of the great world-historical religions,
as also, finally, the deeper knowledge of
the divine in every genuine development of culture, --- that
all this also speculatively makes necessary a more enlarged and at the same time
more deeply developed concept of God'' (p. 119). ---

The writings of I. H. Fichte with which we have so far
occupied ourselves, namely ``Grundzüge zum Systeme der Philosophie''
in their three parts: theory of knowledge, ontology
% --- p. 169 ---
\setcounter{footnote}{0}%
\opage{169}and speculative theology, date from the thirties and forties.
It was then, as we have often stressed, the pantheistic
problem that was on the agenda. Despite all
his unclarities Fichte contributed much to the critique
of one-sided pantheism. With regard to the
problem that pressed forward in the fifties and sixties, materialism,
the question of the relation between the spiritual
and the material, he also has merits. --- If in
the first controversy it was dialectical and metaphysical
investigations that mattered above all, now psychology became
the philosophical science that attracted the
highest interest. Already in his first writings Fichte
appealed to experience as the supplement of speculation,
and it was the study of experience that convinced him
of the reality of the individual. He even states that
he grounds his doctrine solely on experience and on the necessary inferences
built upon it. But it has in part
already appeared in the foregoing, and will appear more clearly on
closer consideration of Fichte's anthropological and psychological
writings, that he is only with difficulty
able to place himself in an objective and unbiased relation
to the given. Just as, by his attempt to overcome
the Schelling-Hegelian speculation, he became the object
of a justified criticism from Weisse, who otherwise
agreed with him in many ways, so too
the way in which he later sought to work up the results of empirical science
in anthropological and psychological
respects met with a sharp objection from \emph{Hermann}
\emph{Lotze} (Streitschriften. 1stes Heft. 1857), who likewise,
as regards the whole general direction, declared his
agreement with him. On both points the explanation
of this peculiar relation seems to lie in the fact that
Fichte, by his own declaration mentioned above, is a
philosopher of personality and always seeks immediately to
satisfy the religious interest. He has not resignation
enough to bow to the objective powers of speculation and
% --- p. 170 ---
% print: et Uforanderligt (misprint: et doubled over the line turn (et|et))
\setcounter{footnote}{0}%
\opage{170}experience in order through them to reach the true
knowledge, which would then also have to include the true
satisfaction of the religious interest; instead this
interest intervenes unconsciously in the course of the investigations and
hinders their objective, scientific execution.

This is seen, among other things, in his polemic against atomic theory.
Modern physics has led to the necessity
of assuming that matter in the last instance consists of
discrete parts or atoms. For the exact explanation
of what changes in phenomena requires
something unchangeable, something constant. ``The point is,'' says W.
Weber\footnote[1]{In a letter to Fechner, quoted in the latter's ``Atomenlehre'' 2te Aufl. p. 88.}, ``to separate out in the causes of motion such
a constant part that the remainder is indeed changeable, but
that its changes can be thought of as dependent solely on
measurable relations of space and time. By this way one arrives
at a concept of mass with which the notion of
spatial extension is not at all necessarily connected. Consistently,
then, the size of atoms in the atomistic
view is by no means measured by spatial extension,
but by their mass, i.e. by the relation, constant for each atom,
in which force stands to the increase of
velocity.'' Atomic theory in this form is a
hypothesis at which science must necessarily
arrive. But it is clear that no dogmatic view is thereby
taught; atomic theory is no conclusion,
only what is last for physics, thus relatively last.
But the demand must be made of him who would take
the results of modern physics into his world-view
that he does not evade, but really thinks through, this
necessary presupposition. As little here as earlier
in face of the speculative conception of the Absolute is it
enough to make partial alterations and reinterpretations.

Nevertheless Fichte in his Anthropologie (1856) regards
% --- p. 171 ---
\setcounter{footnote}{0}%
\opage{171}atomic theory more as a fiction than as a necessary
scientific hypothesis. He therefore seeks to combat
it with all the weapons at his disposal,
without considering that he lacks precisely the weapon
on which alone everything here would depend, namely the exact
method. And yet he means to attack atomic theory ``at its
real center, on the very ground of natural science''!
He here becomes the idealist counterpart of
Czolbe, who, by virtue of the principle of materialism, polemicized
against modern physiology. But while Czolbe
turns against the clear, positive teaching of science, Fichte
confuses the opinion of a few individual materialist atomists
with the atomic theory of physics. Hence it comes
that he regards the latter as a doctrine that can be refuted
with speculative weapons, and he wholly overlooks the skeptical
character that every such hypothesis of a special science
possesses in metaphysical respects. It is therefore,
as will appear, more a lack of sharp apprehension
of the principle and method of empirical science than real
disagreement with its results that underlies
his polemic against atomic theory.

What compelled physicists to take refuge
in atomic theory was, in his opinion, the false
presupposition of empty space from which they
started. If they now did not want to lose the firm ground of the real,
empty space had to be thought of as filled by something else that
was just as original, and this ultimately real
and space-filling thing was then thought of as the indivisible atoms,
which thus contained the limit of divisibility
that the visible things composed of atoms did not
possess. But the presupposition of empty space he believed
he had refuted in his Ontologie, where he had
shown that space and time do not subsist independently and apart
from things, but are only forms of the real, insofar as
the real partly subsists and endures in itself, partly asserts
itself against another by positing and filling its sphere of being and
% --- p. 172 ---
\setcounter{footnote}{0}%
\opage{172}action. All quantity presupposes quality altogether
as what fills and determines, just as conversely
every quality entails its peculiar quantitative
expression. Space and time are only the forms, in themselves dependent,
of everything real; only in abstraction can
they acquire independence --- such an abstraction as,
in Fichte's opinion, underlies the physicists'
atomic theory. If one sublates this abstraction again
and instead infers back from the purely quantitative forms
to the qualitative that underlies them, then
mechanical atomic theory is thereby overthrown too. Instead of
the qualityless, mechanically indissoluble points of space
or atoms there then enter qualitative atoms, which posit
and fill their space and, by their mutual kinship
and the interaction thereby conditioned, produce the phenomenon
of relatively impenetrable bodies.

Besides mechanical atomic theory Fichte combats
the speculative-dynamic conception of matter as well.
This conception, which was founded by Kant\footnote[1]{In his ``Metaphysische Anfangsgründe der Naturwissenschaft'' (1786). On the other hand, Kant in an older treatise, ``Metaphysicæ cum geometria junctæ usus in philosophia naturali'' (1756), stood on the atomistic standpoint. Cf. \emph{Lotze} in ``Gött. gel. Anz.'' 1855. p. 1095---1099.} and later
was developed further by Schelling and Hegel, forms the counterpart
to atomic theory. While the latter from first to last
seeks to punctualize and centralize, and only from the
discrete elements rises to the total and coherent,
the dynamic conception regards matter
as a continuum whose final ground it seeks. The
fundamental peculiarities of matter, impenetrability
and the differing filling of space, are led back to two
different forces, the force of attraction and the force of repulsion, from
whose mutual relation phenomenal matter results.
Fichte objects to this construction that every
force presupposes a real subject whose property it
% --- p. 173 ---
\setcounter{footnote}{0}%
\opage{173}is, or by whose relation to another subject it can first
manifest itself as force. But this real subject of force
the speculative construction leaves
out of consideration. It overlooks that an individualizing
moment is lacking, by which the forces first become
real forces. Such a moment physics offers in its
atoms, if we only deprive them of their abstractly quantitative
character. In his qualitative atomism
Fichte thus believes he has found an expression of unity for what is
justified both in mechanical atomic theory and in
speculative dynamics. But by his misconceived
polemic against the former he has deprived himself of the real
foundation that metaphysical construction needs;
and since he first transforms and reinterprets the reality
that is to be reconciled with the demands of the ideal, the reconciliation
itself can only be ideal too. The real foundation
becomes a phenomenon without real significance. ---

The investigations in the philosophy of nature form the preparation
for the anthropological ones. The general relation
between soul and body is already determined by
qualitative atomism. As every real thing produces its
space and its time in filling them, so the soul,
the real thing disposed to consciousness, brings forth its
body. Thereby, in Fichte's opinion, an end is made
of both dualism and materialism. Bodily
being in space and time is no longer seen as something
foreign to the spirit, but proceeds with necessity
from its essence, and so it cannot have
any absolute power over the soul, whose own product
it is. But how is this producing possible?
The difficulty of seeing it rests, according to Fichte,
solely on our having grown used to conceiving the soul only
as a clearly conscious, reflecting and deliberating
being. So, in investigating the essence of man, one was bound
to stop at a duality: conscious thought
and material corporeality. But closer
% --- p. 174 ---
\setcounter{footnote}{0}%
\opage{174}consideration soon shows an intermediate determination. The body is
not mere matter. Matter is animated, formed and transformed
by the organic forces. We thus get a trinity:
spirit, organic force and bodily matter. Matter
presents no difficulty; its laws are seen through
by organic chemistry. The organic
force, on the other hand, seems to be something obscure and enigmatic. It is
this force that, by its transformation of the elementary substances,
brings forth the external, phenomenal corporeality, whose
real holding and uniting bond it is.
While the external corporeality, which is generally
called ``body,'' is never the same, but changes
incessantly in its single parts, the system of organic forces
forms, on the contrary, a fixed, inner corporeality, a pneumatic
organism, whose reflex or ``effect'' the external
body is, which ``is thrown into the changing world of matter,
much as the magnetic force prepares for itself out of iron filings
an apparently solid body, which nevertheless
splits apart on all sides when the binding power is withdrawn
from it.''

But in what relation does this force now stand to
the soul itself? It is this question above all whose answer
presupposes that one has laid aside the prejudice, nourished by centuries
of scholastic culture, that the soul
is supposed to be a purely intellectual and conscious principle. The
organic force bears in all its effects the stamp of
perfect rationality and inner purposiveness.
What it produces is no less than if a most
perfect intelligence had chosen it with free deliberation
and carried it out with artistic virtuosity. And yet
all organic processes are unconscious, hence unconsciously rational;
they bear the stamp of spirit without being spirit. The
intensive sureness in the acting of this principle and the superior
reason it manifests make it an individual
providence for the body, which acts not only in
the body's first production, when the ``primal position'' is
% --- p. 175 ---
\setcounter{footnote}{0}%
\opage{175}embodied, but also manifests itself during continued life
in all processes of nutrition, reproduction and self-healing.
An analogy to such acting
we have in the artist's involuntarily forming sense and in
the art-drive of animals. In the latter especially there is manifested
a most energetic, dreamlike activity of imagination.
And since the life-process now shows us an acting similar
to the instinctive working of animals, it is
natural to assume the same cause for it. Imagination
is precisely such an intermediate determination between consciousness and unconsciousness as we sought. ``The life-process can
be explained in all its characteristic phenomena only
as the soul's most intensive activity of imagination. It
stands deepest beneath the threshold of consciousness (deeper than
the instincts proper, which already lead beyond the
unconscious, but still in demonstrable continuity with them),
but precisely for that reason it acts most richly and most beautifully, since it
lies furthest from the sensuous-reflecting state,
which is the one most foreign to the soul's original
inwardness.'' Imagination is thus the plastic
force by which the soul forms its body out of the elements
of the material world, a doctrine for which, among others,
Stahl (c. 1700) already prepared the way, and which, in Fichte's
opinion, only prejudices have pushed back. ``Imagination
is not at all merely a spiritual faculty, i.e. an activity
belonging only to the soul's conscious sphere; it
forms quite properly an intermediate determination, is a faculty as much
unconsciously realizing as ideal in itself, and reaches
therefore, in just the same way, down into the domain of the life-processes,
of unconsciously purposive formation and maintenance of the body,
as it serves the highest Ideas
as an ensouling form.''

The soul as ``der Herr seiner Verleiblichung'' cannot
be subject to the perishability of the bodily.
The external body, with everything that belongs to it, stands
rather in an inessential relation to the spirit than
% --- p. 176 ---
\setcounter{footnote}{0}%
\opage{176}being something to be regarded as a necessary condition of
its realization and its consciousness. This sensuous
body in its actual constitution is for the spirit a
limit that acts as a hindrance to all its functions of consciousness.
There are therefore also states (clairvoyance,
magical actio in distans, etc.) in which the soul
ideally frees itself from the limits of this corporeality,
in that it beholds and acts independently of the bodily organs.
Since in these ecstatic states the horizon of ordinary
consciousness is embraced by a freer and
more comprehensive consciousness, which conversely does not exist at all
for the ordinary, sensuously conscious state,
we must have in the ecstatic state ``das Vollbewusstsein
des Geistes,'' ``das vollgeistige Dasein,'' while
the ordinary state is only a bound and halved
existence. The ideal and temporary liberation
from corporeality becomes a real one when at death
the external body is again dissolved into its elements. In
this change the soul does not pass into a purely
incorporeal existence; death consists, on the contrary, only in
the complete liberation and release of the inner, pneumatic organism.
Through its inner primal body, the system of organic
forces, the soul can thus still put
itself in connection with a surrounding world. Only a matter is
here presupposed that can serve as medium for this
acting; but such a matter modern physics, after all, offers
in the ether, which, though inaccessible to sensuous
apprehension, is nevertheless what really acts in natural processes
such as light, heat, etc.

These are the main features of Fichte's doctrine of the relation
between soul and body. He lays strong
emphasis on supporting himself everywhere on facts of experience;
but these are, to be sure, partly taken in
a very wide sense, partly interpreted in a very free
way. Only in a very unscientific sense can one
thus call phenomena such as clairvoyance etc.
% --- p. 177 ---
\setcounter{footnote}{0}%
\opage{177}facts of experience; for that they are still too little
investigated and critically tested. In any case
these phenomena belonging to ``the night side of life'' must rather
be regarded as organic-psychical cases of illness than
be granted a direct psychological significance. Precisely for
the higher mental life one cannot ascribe to them
any value, still less that one should, with Fichte,
be able to find in them the soul's ``vollgeistige Dasein.'' Only lack
of strength to master the true and real facts of
experience can lead a thinker to lay
such weight on these phenomena.

As regards Fichte's anthropology as a whole,
Lotze has strikingly shown in his polemic that he has not
achieved what he wanted (and declared in rather lofty tones
that he had achieved), namely to show, against materialism and dualism,
the true unity of soul and body.
``\emph{This} body,'' says Lotze (p. 131), ``on whose connection
with the living soul natural feeling
alone sets store, this warm blood, this elastic flesh,
the sure firmness of the bones and the tension of the sinews, all
this remains, by your own theory, wholly foreign
to the soul. An inner, invisible body you declare to be
the true one; `the other, the external appearance of it,
which is formed out of an incessant exchange of matter, one may
well go on calling body; --- it is not in truth
one, not subsisting, but the mere effect and reflex
of that inner corporeality'\dots{} No, it is this that
has always alone been called body; only in the unity of this body
with the soul can natural feeling take an interest,
for this alone it knows\dots{} To call this body
the image and effect of another is empty
words, so long as you have not shown by what force
the one is at all able to produce an effect
in the other, or according to what laws it can produce
its image in the masses that are wholly foreign and
external to it. That is the great chasm from which
% --- p. 178 ---
\setcounter{footnote}{0}%
\opage{178}you always shrink back; whatever the forming
principle of the organism may be, it can act only so
much, and in the ways, that the general law of nature permits.
This process of mechanical realization,
the seeking out of the starting points from which the single
concurring forces act, the investigation of the
special conditions, maintained in the course of nature by an unbroken tradition
and transforming themselves in fixed rhythm,
from which the unfolding, growth and propagation of creatures
proceed, --- all this is necessary tasks
for science, and I cannot convince myself
that they are solved by the notion of an ingenious, organizing
fragrance''\footnote[1]{Lotze's critique of I. H. Fichte's anthropology recalls in several respects Leibnitz's critique of Stahl's doctrine. Opera ed. Dutens. vol. II. pars 2. p. 131 ff.}. --- It is more in an unclear tendency
than by a really scientific achievement that Fichte
gets beyond the speculative philosophy of nature and psychology.
Reality is converted for him into an ideal
``fragrance,'' whose unity with the spiritual it is then not
difficult for him to show. ---

In his last more considerable work (Psychologie. 1ster
Theil. 1864) Fichte declares it to be one of the chief tasks
of his psychology to lay the ground for the transformation
of the science of religion on a psychological foundation.
What this foundation is, is seen by bringing out more
closely the way in which Fichte conceives the relation
between consciousness and the life of the spirit as a whole. --- Consciousness
is as such not productive. It produces
nothing new, but only accompanies with its light certain
real states and representations in the soul. Its light
glides like a fleeting torch over the soul's various
states, alternately illuminating them and leaving
them in darkness. But the highest in the spirit is, as we
have already seen, the real content lying behind consciousness,
% --- p. 179 ---
\setcounter{footnote}{0}%
\opage{179}which only in certain extraordinary mental states
comes forth before the spirit, as the spirit is raised to
a higher, more comprehensive consciousness. This mystical-ecstatic
character the truly religious
phenomena now possess. This Schelling has established for all time
in his Philosophy of Mythology. Only he is not
right in placing the power that in those states irresistibly
acts on and in consciousness outside the essence of man,
by regarding the various religions as
effects of and testimonies to a theogonic process.
Yet Fichte, like Schelling, distinguishes between mere
ecstasy and revelation proper. In the former
the spirit of man is raised, indeed, beyond its usual limitation,
but remains nonetheless within its own world. Such is
ordinary somnambulism, presentiments, etc.
But once the spirit has entered into this state
of spiritual beholding, it becomes involuntarily receptive also
to the reflection of another world of spirit in its consciousness;
by this ``ecstasy in the second potency'' the spirit enters
into the region of revelation and is raised above
the human stage of consciousness. This is the case
where there arises in the religions a content such that
mythological activity of imagination is no longer
a sufficient explanation. Such a content is God's
pure spiritual unity, exalted above the world. Here
the cause can be sought only in an objective influence on the
human spirit by the spirit that in revelation
bears witness to itself. The concepts of revelation and
the supernatural do not necessarily belong together. On the contrary,
God works his highest and mightiest works in
continuous connection with natural phenomena,
by heightening the general order of the world.

Since Fichte now sees in this ``supra-conscious'' content
the true essence of spirit, consciousness becomes of merely formal
significance. Originally he started from another conception,
inasmuch as he sought through the theory of knowledge a
% --- p. 180 ---
\setcounter{footnote}{0}%
\opage{180}grounding of the Absolute. Later he took this back, as we
have seen, and saw knowledge as a
fact whose ground and cause one was to seek.
In his Psychologie, finally, he goes yet another step
further. One is not to infer from consciousness to the
Absolute. Consciousness is not something independent and
subsisting for itself, but only a finite being's apprehension
and illumination of itself. As many centers of consciousness,
therefore, so many references to real, individual
spiritual beings. Only by this line of thought is philosophy
supposed to be introduced in such a way that pantheism and idealism
are excluded. Consciousness is nothing but the
most intensive effect of drive, where the being affected leads
to an inner illumination of the drive for the spirit itself.
The soul is ``ein instinktbehaftetes Triebwesen.'' The germ of
consciousness lies in the urge the drive has to seek
its object. This seeking and tracking is heightened
and clarified when the drive finds what satisfies
it; thereby it reflects itself in itself and produces
consciousness of the relation between itself and its
object. Drive is thus the intermediate determination between
the objective, the merely organic, and the subjective,
conscious sensation.

The result would thus be that consciousness is only
a predicate of a presupposed subject, not something in itself
real and substantial. That would be the complete
opposite of speculative idealism, for which precisely
thought and self-consciousness were the expression of man's
true essence. On the theological side this result has been
accepted with joy, since it abolishes the rational and necessary
character of knowledge and makes it
possible to build up the transcendent doctrines of ecclesiastical dogmatics
on a new foundation. Yet this is
surely somewhat hasty, even if one sees the matter from Fichte's
point of view. However much he regards it as an
effect of ``centuries of scholastic culture'' to
% --- p. 181 ---
\setcounter{footnote}{0}%
\opage{181}set conscious life above unconscious life, it is
by no means his intention thereby to justify
what has no really rational and psychological
grounding. Even the ``supra-conscious'' is, in
his opinion, only a higher rational. And on closer
inspection one will find that he cannot consistently
maintain drive as the original mental element.
Imagination is, after all, according to his doctrine, the original, producing
and organizing principle of the soul, and on it
drive too must depend, --- so that after all the
ideal or intellectual element becomes the original: drive
presupposes the purpose that is posited in the unconsciously
forming imagination. ---

3. \emph{Christian Hermann Weisse} (b. 1801, d. 1866)
started out from Hegel's philosophy, from which he
moved away only gradually, inasmuch as, in his own
words, like a Sibyl he kept demanding more concessions
from Hegel in return for less and less recognition.
Already at an earlier point (above p. 23f.)
we have had occasion to mention the main points
on which he departed from Hegel. It was at the summit of the
world-view, in the doctrine of absolute spirit, that
the disagreement first dawned on him. Hegel states
the relation between the highest Ideas, which make up the
content and essence of absolute spirit, in this way: the Absolute
first appears as the Idea of beauty, in art,
in the form of intuition; from there one advances
to the more inward form in representation and feeling, in religion,
until absolute inwardness and the highest
unity are reached in the Idea of science, of truth, which for
Hegel is the logical Idea, the realm of the pure categories.
To Weisse this was a complete formalism. The unity
of form and content was taken by Hegel in the
sense that the content was absorbed into the pure, abstract
essence of form; for him there could be no question of a highest
purpose, of a true and living activity in the Absolute.
% --- p. 182 ---
\setcounter{footnote}{0}%
\opage{182}The whole of existence serves here only to develop, in a dialectical,
law-governed way, the Absolute in its infinite
interplay of forms, and the perfection of the real
depends on the degree to which it can bring the Absolute's
logical forms to revelation. Against this Weisse sets up
the demand that in the Absolute, in the ultimate
ground and the highest purpose, it must be possible to point out a
real value, a good valid in itself and living, which
is to be developed and realized in existence on the basis
of those forms. Only so will feeling and intuition, too,
find rest in the Absolute, whereas they are repelled,
indeed annihilated, by pure abstraction.

The main proposition of the Hegelian system was, as has often been
emphasized, the unity of reason and reality,
but in such a way that only that was regarded as reality
which was absolutely one with reason, and reason here
means the sum total of the pure logical categories. To this proposition
Weisse gives another meaning, interpreting
it thus: ``the rational is absolute determination of form,
the unconditioned foundation and conditio sine qva
non of the real.'' The logical Idea, then, does
not exhaust the content of reality; this content is a positive
more, free reality in space and time, which Hegel
regards as something purely negative, a falling away from the Idea.
The Idea of truth in its pure universal form can then no
longer form the conclusion of the doctrine of absolute
spirit; it is, on the contrary, only the beginning, the foundation
on which a content of higher value can unfold
itself. This unfolding takes place as that universal is specified
and individualized in an infinite manifold
of forms by a producing or generative activity,
which goes on in what Weisse calls the divine
nature or the divine heart. This process
it is the task of aesthetics, as the science of the Idea of beauty,
to present, and Weisse attempted such a
presentation in his first major work ``System der
% --- p. 183 ---
% print: Deutshland (misprint: so printed)
\setcounter{footnote}{0}%
\opage{183}Aesthetik'' (1830)\footnote[1]{On the specific content of this work see \emph{Lotze}: ``Geschichte der Aesthetik in Deutshland.''}. Aesthetics forms the transition to the
highest Idea in the realm of spirit, the Idea of the good, the content of ethics
and of the philosophy of religion. The Idea of the good presupposes
the beautiful, because the good consists in the will to
communicate something real. The Idea of the good becomes empty of content
if it does not have as its presupposition a living
fullness of reality such as the Idea of the beautiful reveals
to us in the nature of the Godhead.

Hegel's fundamental error is that he hypostatized the universal
logical categories into something positively and really existing,
made the logical Idea into God. It is an error
he shares with the whole of philosophical dogmatism, which
in the necessary determinations of thought immediately sees
what really is. But what is great in Hegel is that
he was the first to conceive and carry out the idea of a complete
development and presentation of the universal and necessary
determinations that underlie all being. Absorbed
in this work, he took the very sum total of these determinations,
the logical Idea, for the only reality.
If one wants to understand Hegel rightly, one must integrate
this error. His logic must be regarded as a presentation
of the necessary forms of existence of reality, not of
positive reality itself. That this conception is
correct, Weisse finds proof in the fact that it teaches
us to see Hegel's fundamental error in the very thing that
constitutes his greatness, so that nothing
accidental and unexplained is left over.

But this fundamental error regarding the relation
between reason and reality brings other defects with it in
particulars. Reality unfolding in space and time
Hegel regarded as a
falling away from the Idea, because space and time were for him the
forms of mere externality, without essential significance
% --- p. 184 ---
\setcounter{footnote}{0}%
\opage{184}for the Absolute. A changed conception of reality
had to bring with it a changed conception
of the forms of reality as well. Weisse reproaches Hegel for not
developing space, time, and the determinations connected with them
in the logic, since after all all reality is in space
and time, and these are therefore, just as much as the pure concepts,
absolute and necessary forms of existence. Hegel has thereby
cut himself off from the possibility of an inner transition
from the logic to the philosophy of nature and brought a
veritable dualism into his view.

By this critique of Hegel Weisse had arrived, by his own
path, at a result similar to the one Schelling
voiced some years later in his preface to Cousin (see
above p. 126 f.). With this pronouncement of Schelling's he could
therefore on the whole agree. But there still remained
a thoroughgoing difference between the two thinkers,
which failed to come out more definitely only because Schelling's
later doctrine became known as a whole only much
later. Weisse adopted Schelling's separation between
negative and positive philosophy. The logical Idea, the sum total
of the pure, universal concepts of reason, then became for
him ``the negative Absolute''. But in the more precise determination
of the relation between the negative and the positive
the disagreement comes to light. In Schelling,
as we have seen, positive philosophy coincided
with the philosophy of religion, and the positive starting point
from which it proceeds could be grasped only by one who, unsatisfied
by life in the Idea, postulated the ``Idea-free'' God by a religiously
determined act of will. Schelling did not, to be sure, deny the principle of necessity, as
some of his followers did (notably the philosopher of right Stahl,
through whom Schelling's new doctrine first became known in
North Germany), who by ``positive philosophy'' understood
a tracing back of all law and all truth in existence
to an unlimited divine will; but in the last instance
negative philosophy could still acquire only an accidental
% --- p. 185 ---
\setcounter{footnote}{0}%
\opage{185}and vanishing significance if one could go immediately, through an act of will, to positive philosophy.
What Weisse approves in Schelling's later doctrine is,
in the first place, the conception of the opposition between necessity
and freedom as a real opposition, which cannot
be sublated by a formal dialectic, as in Hegel,
who in his logical Idea sees the unity of these opposites.
It is, further, the conception of necessity
as the negative foundation or prius of real
existence, in Weisse's expression: as the pure possibility of existence.
Only the free is the real, the existent,
the positive. But philosophy cannot grasp this positive
in a positive way --- by this proposition
Weisse distinguishes himself from Schelling. Negative
philosophy, which for Weisse coincides with metaphysics,
but for Schelling was one with the whole of non-religious
knowledge, therefore has essential significance here for the
positive. The logical Idea contains the root of the positive,
and only through the indwelling of the metaphysical, of the logical determinations,
does the positive become something free. Freedom
thus springs from necessity and has its ground
in it. Freedom, or the fact that the content can also not
be, does not therefore consist in an external relation to the
form of pure necessity or in an independence of the
form; rather, the content becomes free precisely and only through the form,
which in its completion is the form of freedom\footnote[1]{Grundzüge der Metaphysik (1835). Das philosophische Problem der Gegenwart (1842).}.

In his opposition to Hegel Weisse further had
a comrade-in-arms in I. H. Fichte, with whose view his own
was therefore for a time regarded as essentially one.
Yet they are not only at odds in their conception of Hegel;
in the rest of their views, too, such great
differences appear that they prompted Weisse to
compose the polemical work mentioned earlier: ``Das
% --- p. 186 ---
\setcounter{footnote}{0}%
\opage{186}philosophische Problem der Gegenwart.'' Fichte saw in
the Absolute of the Hegelian logic the concept of the world-spirit,
and believed he could accept this as relatively justified;
his critique then amounted to this, that experience, particularly
the fact of inner purposiveness, compelled one
to go a step further, since the concept of the world-spirit did not
contain the sufficient explanation. Weisse, on the other hand, saw
in Hegel's logical Idea the pure, necessary determinations of thought,
which hold for all reality but in themselves
make up only a realm of possibility. The categories
were for Fichte only determinations of a presupposed real,
of which he sought to attain certainty either through the theory
of knowledge or through experience. In this Weisse saw
the one-sidedness opposite to Hegel's. ``With you,'' he writes
to Fichte (p. 188), ``the same categories appear
as products of a dependent reflection, as
abstract shadows of a reality beyond, which we
see in Hegel, no less against their nature, hypostatized
into the sum total of what truly is''. The
right conception, in his opinion, is the one which in the universal
determinations of the possibility of existence does not yet,
to be sure, see real actuality itself, but seeks to develop
them to a point where with inner necessity they pass
over into it. Fichte came, as we have seen, to hover
between pantheism and vulgar theism, because he
would have nothing to do with any a priori or ontological grounding
of the idea of God; in consequence, as with
Schelling, the determination of freedom in God had to enter by an
uncomprehended positing.

In this endeavor of his toward a correct estimate
of the nature of metaphysical knowledge Weisse is conscious
of having a predecessor in Kant\footnote[1]{See, besides the academic address cited on p. 119, Weisse's essays in the 25th and 26th volumes of Fichte's journal: ``Ueber die transscendentale Bedeutung der Urtheilsformen und \opage{187}Schlussfiguren'' and ``Ueber den letzten Grund der Gewissheit im Denken'', as well as ``Philosophische Dogmatik'' 1ster Band (1855).}. Kant contested
% --- p. 187 ---
\setcounter{footnote}{0}%
\opage{187}the dogmatism which believed it had in the pure concepts of reason
a knowledge of things in themselves, notably a
knowledge of the infinite and unconditioned, the eternal
and necessary in things, and thus confused the necessary
basic forms of knowledge with the real essence of actual
things. Kant, by contrast, saw the original property of the understanding
as something merely formal, which a priori points us
to experience, by which alone it can reach a real
content. To reach the concept of an object the
understanding must always add something of its own to the material
given by sensation; from the combination of these two elements
the concept of the object arises. What is
added here is, in contrast to the singularity of experience,
the universal; in contrast to the matter given in experience,
the formal; and in contrast to the reality of the matter of experience,
the merely possible. But the validity of these elements of the understanding
Kant conceived only as a subjective-psychological
necessity for the human
mind; he did not see that the possibility of thought contained in the prius of the understanding
is at the same time a possibility of existence.
From this defect his view must be freed in order
to correspond to the truth. Such a liberation has taken place
through the later philosophical development (I. G. Fichte,
Schelling, Hegel), which led to a knowledge of the
universal essence of the Idea. Yet the very energy with which
the liberation from the subjective limits was carried through
led back once more, for a while, to dogmatism, in that
one now saw in the metaphysical Idea, whose objective significance
one had demonstrated, the highest real itself.
The philosophical movement begun with Kant
must thus, when the right consciousness of the nature and task of the pure
science of reason once again breaks
% --- p. 188 ---
\setcounter{footnote}{0}%
\opage{188}through, return once more to its starting point, but enriched
and strengthened by the insight gained.

In every act of thought, in all judging and inferring,
that pure possibility of thought and existence is present. It
enters as copula between representations, which thereby
arrange themselves as subject and predicate of the judgment; it is
the force which in the syllogism joins the terms inwardly
through the middle term. It has its
more precise content in the concepts which for the understanding underlie
the objects of knowledge and the knowing
consciousness itself, as also consciousness and
the objects in their mutual relation. All reality
thus rests on these presuppositions. The
natural or immediate understanding does not, to be sure, discover this; but
when doubt arises, consciousness seeks its way back to
the universal foundation, the absolute necessity, which
at the same time contains the absolute possibilities. We have, then,
here the concept of truth in its most abstract form,
purely logical truth, which constantly presupposes itself,
and of which the words of Thomas Aquinas hold: ``Qui negat
veritatem esse, concedit veritatem esse; si enim veritas
non est, verum est, veritatem non esse.'' In this
abstract form the concept of truth has its expression in the
so-called principle of contradiction, the first negative
form of the possibility of existence: everything is possible that
does not contain a contradiction. This proposition leads
us to see that pure reason already, without yet
positing or presupposing anything real, ties the possibility
of existence to certain fundamental conditions. These
conditions, which are the primal forms of all existence, make up
the content of pure necessity of thought, the pure Absolute of reason,
which philosophical speculation has at all
times regarded as its first and proper object,
without yet having succeeded in determining
it in its full purity. The Absolute thus becomes the
ultimate ground of the certainty of thought, because in it the primal possibility of the
% --- p. 189 ---
\setcounter{footnote}{0}%
\opage{189}objective is one with the highest necessity of thinking.


Even for the immediate understanding, however, this
primal possibility emerges in a more concrete form. Even if the
immediate understanding does not, like speculation, become conscious of the necessity of the categories
of reason, it nevertheless ascribes to mathematical determinations greater evidence and more perfect
certainty than to mere experience. And yet it cannot
ascribe to mathematical concepts an independent
being. It then seeks to conceive them as tied to
and dependent on what is given in experience; but neither
can this succeed, for how can something that
owes its existence to the given surpass it in certainty
and necessity? --- That is why, as Plato
and Kant have already insisted, the surest way to lead the
immediate understanding to speculative knowledge is to
start from the inquiry into the ground of mathematical
necessity. In number, space, and time
speculative reason sees fundamental presuppositions of all existence,
whereby existence falls under mathematical
laws. Metaphysics is to develop the concept of what is,
and as soon as it has freed itself from the one-sidedness of idealist
dogmatism, it must see that the concept of being
stands in an original and necessary connection
with the concepts of number, space, and time. The answer to the
question: what is being? then leads us, through the
categories by which the necessary basic forms of the filling of space and time
are designated, to a fundamental concept in
which all possibility of a content that can fill these
forms, hence all possibility of a being and becoming, is
comprised in one primal possibility. This fundamental concept,
the absolute Idea, is one with the concept of the Godhead, but
here, where we still stand on the metaphysical
standpoint, that means the concept of the possible Godhead. There
is only one truth necessary for thought, namely that
% --- p. 190 ---
\setcounter{footnote}{0}%
\opage{190}originally only God is possible, and that in his possibility
the possibility of all things is contained.

Metaphysics follows the dialectical method. While
Weisse soon came to the conclusion that the dialectical
method in Hegel's sense, as the self-movement of the concept or
of the content, could find no application
in the domain of real philosophy, where thought step
by step enters into interaction with experience, he nevertheless
believed he could still keep this method in
the metaphysical domain. Thus he writes in the preface
to his Metaphysics: ``The formal truth and material untruth of the Hegelian
philosophy, the sterling excellence of its method and the disconsolate
emptiness of its results forced themselves on my mind with equal evidence.''
Rightly, as we have seen (p. 24), both followers
and opponents of Hegel urged that it was
a self-contradiction to restrict the dialectical method,
which is the self-movement of the content, to merely formal
validity. Weisse has therefore specified his
conception of the method more precisely. By the dialectical method
he understands the very concept of philosophical knowledge,
but not specifically the form it took on in Hegel's
panlogism. Dialectic has a different character
in the metaphysical and in the real-philosophical
domain. In the former it shows itself in the power
with which every category as it were presses forward
to hold as a determination of the whole of existence,
as the sum total of all-possibility. Such is dialectic
in the history of philosophy. The struggle and the development in that history
rest precisely on this, that the categories successively put forward as determinations
of the whole are constantly displaced by
others, more comprehensive and more thoroughgoing. This is
also the course of development within metaphysical science.
The standard is the totality of reality, which
proves the inner contradiction in the inadequate categories.
In the real-philosophical domain the relation is
% --- p. 191 ---
\setcounter{footnote}{0}%
\opage{191}the reverse. Here every real given in experience is regarded
as though it included the totality of the categories, but
for that very reason it points with necessity beyond itself
to other and higher expressions of this totality. In this domain
the dialectical method becomes one with the genetic
method. For what connects categorial thought
with empirical phenomena is the law that holds for the coming into being
and development of these phenomena. By the law,
which here, as in the pure speculation of reason, expresses
the possibility of existence and the necessity of thought, the
movement on both sides is governed: the genetic movement of existence,
which is posited by the law as possible, and the movement of thought,
which is posited by the law as necessary.

In the same way as the single real-philosophical
concept relates to the existing real, the
totality of metaphysical concepts relates to concrete
reality. The course of metaphysics coincides, according to Weisse,
with what constitutes the truth in the old
so-called ontological proof of God's existence. This
proof in fact leads no further than to the concept
of the possibility of the Godhead, of God as the one who cannot
be thought of as not being. Yet the metaphysical
consideration leads at its terminal point to a sharper
determination of the essence of the Godhead. Since the absolute
Idea, the sum total of primal possibility, is a unity in the infinity
of the differences and oppositions sublated in it,
the necessary form for this existence becomes
the form of thought. The determinations of the Idea exist only in and for
thought. Now inasmuch as this thought in an eternal circuit thinks
itself in all those determinations, what eternally is is a
self-conscious primal subject, primal personality. We thus arrive
at the result that if the eternal possibility does not
remain closed within itself but passes over into reality,
this primal reality must be primal personality.
One could imagine a standpoint of abstraction
where all real being is disregarded, and where
% --- p. 192 ---
\setcounter{footnote}{0}%
\opage{192}therefore God too can be thought of as not really being.
There is not, as Hegel thinks, an immediate, dialectical
transition from concept to existence. Hegel took the
ontological proof to be what demonstrated this transition,
and therefore regarded this proof as the only genuinely
scientific one. Against him one may rightly
turn that abstraction. As soon as one posits
something as existing, on the other hand, one must also posit God's existence.
Otherwise one would entangle oneself in the contradiction
of at once positing something real and cancelling the sum total
of the conditions without which nothing can become real.
--- This is the at once dialectical and factual transition
from metaphysics to real philosophy, or from
the ontological proof to the cosmological. The two are
most inwardly united, and make up the two halves of
the same whole.

And yet this opposition between
the negative and the positive Absolute seems to set an abyss of possibility
before God's real existence, so that
he works his way forward to reality only by an act of coming-to-be or self-creation.
According to
Weisse\footnote[1]{See his article ``Gott'' in Ersch and Gruber's Encyclopedia Sect. I. vol. 75. p. 424.} there is an irreconcilable opposition between the Absolute of philosophical
speculation and the demand of religious experience
for personality, so long as one makes the
demand that God's personality is to be posited \emph{immediately}
as the Absolute, or the Absolute \emph{immediately}
as personality, instead of the Absolute of pure thought
containing only the infinite possibility of a
living, personal existence. I.H. Fichte (Verm. Schriften
I. p. 100) remarks on this that such a separation
is unthinkable if it is to hold objectively, in God's
own essence, and not be a merely theoretical distinction
of human thinking. Lotze (Geschichte der
% --- p. 193 ---
\setcounter{footnote}{0}%
\opage{193}Aesthetik p. 202, cf. already Metaphysik p. 322 f.)
derives the doctrine of the negative Absolute from the fact that, when our
thinking disregards the content of the world and immerses itself
solely in the realm of logical truth, the illusion arises
that this logical truth is something original and independent
in relation to living reality. And we
have already pointed out above (p. 120 f.) that this
point was bound to become the most difficult one for the whole tendency
to which Weisse belongs. That Weisse does not want to split
the concept of God into two parts is seen from his express declaration
(Phil. Dogm. I. p. 439) that the separation between
possibility and reality in God is only conceptual,
both being moments or ideal constituents of the
one and indivisible concept of God. And this separation is
necessary in order to see that God becomes free only through the Idea, through the
negative Absolute, becomes the absolute being,
the lord of necessity. Herein lies the real (not
merely logical) unity of freedom and necessity. Without
such an inner opposition in God's essence the concept of the living God
cannot, in Weisse's opinion, be upheld, and he puts this opposition
forward precisely in opposition to the abstract concept of God
of dogmatic philosophy, God as actus purus.

The same separation appears in another form
when Weisse, following Schelling, distinguishes
between what God is (quid Deus sit) and that God is
(quod Deus est), between God's essence and existence, in such a way
that the ontological inquiry is to give us the Qvid,
the general content of the concept of God, and the cosmological
inquiry adds to it the Qvod, the determination of existence. That this
opposition is untenable we have already seen in Schelling.
But Weisse in fact uses it only to express
the two-sided starting point that is necessary.
Only on the standpoint of absolute abstraction
would this opposition, according to Weisse, have validity.
Real philosophy (the philosophy of nature and of spirit) has,
% --- p. 194 ---
% print: Metaphysisk til (misprint: for Metaphysik)
\setcounter{footnote}{0}%
\opage{194}moreover, not merely the significance of helping knowledge to
the one determination: existence; on the contrary, with
the knowledge of existing reality there follows an
ever richer and more concrete development of the content of the idea of God;
to the first, purely rational Qvid is added a
second, qualitative Qvid, a Qvale, which grounds the aesthetic
and ethical predicates of the Godhead. Metaphysics
in and for itself cannot give us the right, full
knowledge of the idea of God; we reach such knowledge only when,
on the basis of metaphysical knowledge, we immerse
ourselves in the content of reality. Pantheism has its indisputable
scientific justification in the speculative
knowledge of a metaphysical necessity that pervades and conditions the
shapes of the real world, a necessity which, however, it is unable to separate
from the real existence of things, since it does not reach the Idea in
which this necessity, in itself only formal or negative,
finds its completion. To get beyond
pantheism it is therefore necessary to lead metaphysical
necessity back to its right significance
within the concept of God and to show that it is only a presupposition
for a fuller content. With respect both to
abstract theism and to one-sided pantheism, that
separation thus has its great significance.

The transition from the ontological to the cosmological
inquiry, from metaphysics to real philosophy,
was, as we have seen, at once a priori and a posteriori,
an inner transition in the concept at one with an experience
of what really exists. The new element which
thereby enters cannot at once find its justification.
An opposition is posited between experience
and concept which can be sublated only gradually. This
takes place, as already indicated, through the real-philosophical
sciences. Here, where we are first to inquire
what qualitative contributions these sciences can make
to the knowledge of the idea of God, the task is dialectically to
% --- p. 195 ---
% print: ogsaa ogsaa (misprint: so printed)
\setcounter{footnote}{0}%
\opage{195}develop the content of the world, as it appears in experience,
to the point where the concept of the world passes over into
the concept of God. --- Nature does not merely stand under mechanical
laws; these laws fit into a higher
teleological connection in which all its parts stand in
infinite interaction with one another and form a graded series
from apparently dead mass up to the highest
rational and spiritual beings. Experience leads to this result,
and it already lies in the essence of the pure Idea of
reason. Since this Idea is itself infinite unity,
the existence whose universal conditions it
contains must also have connection and unity as its basic forms.
The Idea of reason therefore extends as a guiding
thought from the ontological inquiry into the cosmological,
and points out the places where the teleological
connection as it were comes to a point in universal concepts,
which as predicate-determinations attach themselves, filling out
and completing it, to the metaphysical concept of the
primal subject. The cosmological inquiry then passes over,
at its summit, into the ethicological, for it is
in the spiritual domain that the ultimate purpose and
thus also the highest divine predicate-concept
must be sought. Inasmuch as thought, now guided by the scientific
consciousness of the Idea, seeks to bring its consideration
of reality into harmony with the content of this consciousness,
the very end it here seeks
to reach appears as what is at the same time the end of the whole development
of the world. The highest principle then presents itself as the knowledge
of the Idea of reason in its unity with reality,
or in other words as the Idea of truth.
But this Idea has not merely theoretical significance; it
also becomes the highest law, the moral principle. The principle of
morality consists in the finite subject's making
the purpose of its action what is the purpose of
the objective world-order. Raised above the empirical will of the finite
spirit, the concrete Idea of truth
% --- p. 196 ---
% print: og absolut Idealitet (misprint: og doubled over the line turn (og|og))
\setcounter{footnote}{0}%
\opage{196}nevertheless also shows it the objective purposes it is to make
its own. But alongside the theoretical and practical
striving for truth there runs another series of activities,
which no less have their purpose in themselves.
Here the Idea of reason and its reality in nature
and history come forward as the source of a world of forms,
which spirit, fertilized by contact with the realized
Idea, brings forth from within and enjoys in
its beholding and bringing forth. All that belongs here is comprised
under the Idea of beauty.

To these two Ideas, that of truth and that of beauty,
the real-philosophical inquiry thus leads. But
they stand in an inner opposition to each other. According to
the Idea of truth the highest purpose of the world is unity
and absolute ideality; the Idea of beauty, on the other hand, asserts
the independent value of the aesthetic states and activities of the finite spirit
and thereby posits an infinite
manifold of real world-purposes. This is the
true significance of Kant's antinomy between morality
and happiness. This contradiction is not merely subjective
but objective as well, insofar as one stops
at those two Ideas as real powers of the world and of spirit
and does not seek a higher reconciliation. That such a reconciliation
must be sought follows from the very essence of the Idea of reason, which
cannot permit the independence of such powers. They
must find their place within the idea of God. It lies
in the concept of the primal subject that it cannot posit itself
as teleological principle without thereby positing the possibility
that the finite, individual existence, too,
becomes a purpose in and for itself. Both logical and
aesthetic reality thus have their ground
in the Godhead. This consideration leads into the
ethical-religious domain, the highest and concluding sphere of experience.
The philosophy of religion, or speculative
theology, forms the summit and
end point of the philosophical system, where the unity of Idea and reality
% --- p. 197 ---
\setcounter{footnote}{0}%
\opage{197}is completed, as the highest concept of speculation meets
ethical-religious experience, which encloses within itself all
other experience. Everything here turns, then, on the more precise
conception of the essence of religious experience and of the philosophy of
religion\footnote[1]{``Philosophische Dogmatik oder Philosophie des Christenthums''. 1855---1862. 3 vols.}.

The concept of religion falls, according to Weisse, under
the general concept of experience. Experience is on the
one hand more than a mere receptivity or a purely subjective
state in the soul, but on the other hand less
than knowledge proper, or science. Long
before Kant, at the end of the last century, for the
first time gave a precise determination of the essence of experience,
it had often been declared from the side of religious consciousness
that religion is essentially experience (παθεῖν τὰ θεῖα. Dion. Areopag.). It is the mystics who have
had the clearest insight into the experiential
character of religion. The same lies in the old Church teachers'
determination of the relation between πίστις and γνῶσις.
And yet, even after the concept of experience had long
been established in other domains, it was a long time before
a similar conviction could break through with
regard to religion. The scholastic-orthodox tendency
in the Church regarded the content of religious experience as
something communicated to the human race from outside as
finished doctrines, and only by this communication, in its view,
are the subjective movement and sensation in individuals
then called forth. On the other hand,
the mystics believed that in what they experienced in the heart they immediately
had a knowledge of something objective. Only
through Kant's investigations, which in the end extended also
to the philosophy of religion, which he really founded,
did it become possible to see this domain, too, more clearly. The
purely objective doctrinal content he sets aside as merely
% --- p. 198 ---
\setcounter{footnote}{0}%
\opage{198}statutory; but on the other hand the religious does
not for him, as in the so-called religion of reason, coincide
with metaphysical knowledge; he held fast to
a new starting point for the philosophy of religion in the
ethical-religious consciousness, which for him is essentially conditioned
by the law given in man's inner being, the practical
a priori. He overthrew dogmatism not only in
the philosophical but also in the religious domain.
The latter was the more difficult task, since the religious
content, because of its character of inwardness, is
harder to distinguish from the content of pure reason.

What is given ethically-religiously stands in the same relation
to knowledge as everything else given in experience; it, too,
is bound to the laws which give the transformation of
what is subjectively lived through into an objective content of knowledge
its scientific form and significance. The laws of all
experience stand in organic connection, so that the
laws of spiritual experience have as their presupposition
the laws of physical experience and without this presupposition
can have no significance. An experience
that does not rest on fixed, unchangeable laws
cannot be distinguished from a delusion. And although
religious experience, as \emph{religious}, has its own peculiar
laws, this lawfulness nevertheless rests
on the psychological laws in general, and these
in turn on the physical. Without these presuppositions
religious experience would not be \emph{experience}. The notion
of miracle is unethical and irreligious precisely because
the ethical-religious concepts rest on the psychological
and physical foundation. Consistently carried through,
supranaturalism has also constantly, instead of leading
to living religious experience, favored blind faith in authority,
which is precisely the most complete opposite
of all real and living experience. If the immediate
religious standpoint is nevertheless the fertile soil of the notion
of miracle, this stems from the need
% --- p. 199 ---
\setcounter{footnote}{0}%
\opage{199}to find, by poetizing imagination, an expression for the
inexpressible in the facts of religious experience;
the miraculous then arises because the inner connection of natural (``ausserreligiöse'')
experience does not stand as
fact before consciousness. For religious experience
includes all other experience within itself as an ideal totality;
conversely, natural experience at its
summit passes over into the religious. Thereby all
content of experience is gathered in the center of the inwardness of spirit, where
it is grasped by philosophical knowledge through the Absolute's
form of reason, and no longer relates in an
external way to this knowledge, but really becomes
one with it.

The question now becomes that of the psychological organ
of religious experience. Schleiermacher, as is
well known, defined religion as piety and placed the essence of piety
in feeling. But he then restricted the concept of feeling
to such a degree that it is cut off from every
content that goes beyond the purely individual and subjective.
And whereas one usually regards the opposition
between weal and woe, pleasure and displeasure, as the peculiar
characteristic of the life of feeling, according to Schleiermacher
this opposition cannot in and for itself be ascribed
to religious feeling; only when the lower,
sensuous sensation enters into relation with it, either
inhibiting or promoting it, does that opposition find application
in its case as well, to that extent. The
imperfection in these determinations Weisse derives from the fact that
Schleiermacher did not compare religious feeling
with aesthetic feeling, whose kindred counterpart
it is. He would thereby have seen that the opposition
between pleasure and displeasure can also have a spiritual
significance and need not stem from sensuous
nature. Purely theoretical experience has from first
to last an objective character; subjective sensation
has value here only insofar as it passes over into
% --- p. 200 ---
\setcounter{footnote}{0}%
\opage{200}observation and sensation, for which reason feeling as such
yields no material for knowledge here. In the
aesthetic and religious domains, by contrast, the oppositions of the life of feeling
play an essential role, and on their more precise
determination the difference between the aesthetic
and the religious in turn depends. In the aesthetic domain
the feeling of pleasure and displeasure is the ruling power
and contains the principle of the natural and
spiritual forces at work there. Feeling is here, as aesthetic feeling,
in close connection with imagination, the creative power of imagination.
In the religious domain, by contrast, the feeling
of pleasure and displeasure is not the purpose, but appears partly as starting point, partly as point of passage
for an activity and a movement which do not in turn
end in a feeling, but issue in inner or outer action
or in an abiding form for the essence of the spiritual
personality. Religious experience thus belongs
within the domain of practical or moral experience,
which forms an opposition to the aesthetic on the
one side and to the theoretical on the other. The difference
between the aesthetic and the religious comes out especially
strikingly in the fact that the oppositions of feeling in
the one domain can take on the reverse of the significance
they have in the other. In the aesthetic domain
the feeling of pleasure is thus always of positive value,
the feeling of displeasure of negative value; but religiously seen, the feeling of pleasure can
take on a negative, hindering significance and, conversely,
the feeling of displeasure can act as a furtherance.

Feeling, then, does not have absolute significance in religious
respects, as Schleiermacher thought. It is not
the innermost of the life of spirit, but the revelation of an inner
hidden behind it and entirely different from its changeable
alternation. Only this inner, the ethical substance,
is the proper element of religious experience; feeling
is only the mark by which inquiring
thought recognizes the material of moral empiria. This leads
% --- p. 201 ---
\setcounter{footnote}{0}%
\opage{201}to the question what, in the last instance, is the
objective moment in religious experience. Schleiermacher
could assume such an objective moment only
inconsistently. But if religious experience belongs
within the domain of general moral experience,
it must also be placed under the concept that governs
ethics, the concept of the good. On this concept both
religious and philosophical consciousness have constantly
built the idea of God. Philosophical ethics understands
by a good a community of life of rational, personal
spirits, which are mutually bound by organic laws
into the unity of a living common body. The
inner, moral experience is then the observation of a moral
community's vital functions through the states of feeling
in the subject who belongs to the community and acts and
suffers along with it in its functions. A purely subjective religious
feeling would make such a community impossible. Like the
ethical goods, the highest good, too, the content and object of religious
experience, must consist in such a
community of life. But whereas the ethical communities
can also become the object of an outer (sensuous or
historical) experience, inner experience in the religious
domain is the only possible assurance of the
corresponding objectivity. But at this summit
the religious and the speculative, the religious
representation of the Kingdom of God and the philosophical
concept of the highest good, are then united.

The subjective moment of religious experience is feeling,
its objective moment the Idea of the good. The laws
which hold with regard to the relation between these two
moments must be the same as those which hold for
spiritual life as a whole. Since the life of spirit is not confined
to single individuals each by itself, but these,
in order to utter and express, develop and enrich their own peculiar
spiritual life, join together and form communities,
the Ideas that move the life of spirit step out of
% --- p. 202 ---
\setcounter{footnote}{0}%
\opage{202}personal inwardness onto the stage of history. It
is therefore, in the case of religious experience as well,
the historical mode of consideration built on criticism
that furnishes the organ for grasping the relation
between the subjective and the objective moment. Yet
one must not confuse with such a historical consideration, with the
historical school, which is the peculiar product of modern culture and
science, the
theological use of the concept ``historical conception'', where
a presentiment of the historical character of the positive religions
does, to be sure, lie at the bottom, but where one does not take it in
earnest, but on the contrary believes one honors the historical by
detaching it from all the laws of experience that hold in other domains.

Nothing can become religious experience by way of outer communication
unless it already was that in those from
whom the communication originally proceeded. Now when religious
experience is shaped through communication and interaction
in a community, it finds its expression in a definite,
positive form. The positive forms thus arise
historically and must be explained historically. The positive arises
at a certain point in historical development and
has its whole subsistence and activity within it. The
historical religions are different attempts to grasp
the highest good, to tie a bond that can organically
bind the human spirit to the Godhead and the
supersensible. But besides the properly religious
moment, other moments also play their part. In the
first place, the political-social drive exerts its effect and
merges immediately with religious experience, so
that the limits of the former also become those of the latter. Thereby
the historical religions become national religions. Next,
the creative power of imagination with its luxuriant
products often pushes the properly religious almost entirely
into the background. Thereby the historical religions become mythologies,
which does not, however, exclude their being
% --- p. 203 ---
\setcounter{footnote}{0}%
\opage{203}testimonies to and sources of religious experience. But the
political and aesthetic covering must be cast off for
religious experience to appear in its purity.

Every religion rests on the presupposition of a
community of life between the divine and the human
spirit. Therefore every religion lays claim to
being divine revelation and must be recognized as
such insofar as it really possesses religious import.
The concept of revelation in this generality becomes one
with religious experience, but since the latter does not appear
pure in the national religions, aesthetic, speculative, and national
elements are also introduced under the name of revelation.
Revelation acquires a narrower, more definite meaning
when its concept is separated from the
elements that are not immediately religious. In Christianity
the concept of revelation appears already in
its founder and first followers as the expression of a
liberation, an opposition to an earlier bound and darkened
state. This comes from the fact that the properly religious
content of experience dawns in its full clarity without
foreign forms. The historical Christ first grasped
clearly, and expressed in teaching, life, and death, the true idea of
the highest good, the Kingdom of God or Kingdom of Heaven, into
which man must let himself be born in a spiritual way.
Salvation, entry into this Kingdom, depends only on the conditions
that lie in the very essence of the highest good, and
is not hindered by limits such as those which in the national religions
and the mythologies were raised partly by national egoism,
partly by poetizing imagination. To that extent the form
religious experience has acquired through the historical
Christ can be designated as the redemption effected by him.
To this fundamental idea of Christianity the philosophy of religion attaches itself, in that, after having sifted the content of religious experience by historical criticism,
it becomes specifically a philosophy of Christianity. The philosophy of religion
is connected with ethics by its fundamental concept of
% --- p. 204 ---
\setcounter{footnote}{0}%
\opage{204}the highest good. But only at the point in historical
development where this fundamental concept comes clearly
before religious consciousness can a
true and inward unity of the philosophical concept and
religious experience be tied. The character of the philosophy of religion
depends on its relation to philosophical ethics on
the one side and to religious experience on the
other\footnote[1]{Weisse's essay ``Ueber Eintheilung und Gliederung des Systemes der Philosophie'' in the 27th volume of ``Zeitschr. für Phil. und phil. Kritik.'' p. 88 ff.}. A philosophy of religion therefore first becomes
possible through Christianity. But also only possible; for
on the one hand it presupposes a consciousness of the essence of philosophical
knowledge and of experience which only the
philosophical and natural-scientific development of modern times
could produce, and on the other hand the great fundamental idea of Christianity
could not permeate, in pure and unaltered form,
the world and the consciousness of the peoples.

The historical Christ, to be sure, who declared that
he would give no other sign than the sign of the prophet
Jonah (i.e. the miracle of personality and truth itself),
was in his own person raised above all mythological and supranaturalist
elements. But in his surroundings he could not
and dared not ward off the drive to poetic religious
formation of legend. The mighty spiritual upheaval, the
birth pangs of the new faith, awakened this drive anew, and it now wove
through his life and personality a web of
significant images of miracle whose character is essentially
mythological. Already for the consciousness of the first believers
such elements of the universal
substance of salvation as lie in every genuine religion were bound to be mixed
with the historical concept of the person and life of Christ.
The notion of salvation thereby took on an exclusive
form, which conflicts with the universal in the essence of salvation,
according to which salvation is tied to religious
% --- p. 205 ---
% print: og den historiske (misprint: og doubled over the line turn (og|og))
\setcounter{footnote}{0}%
\opage{205}experience as a whole and not merely to the specifically Christian.
Already the first disciples transferred the idea of the Logos,
that great result of pre-Christian speculation,
the idea of God's eternal thought or word, one in essence with himself,
to the historical Christ, and by this
application, which struck like a flash of lightning into the disciples'
souls, the ecclesiastical community was founded, built on the
ecclesiastical doctrinal concept thereby established\footnote[1]{This is one of the points where Weisse in his historical criticism departs from the Tübingen school. According to that school, the idea of Christ as Logos, as the Son of God become man, is the fruit of a long course of development which was completed only toward the end of the 2nd century after Christ, to which time the composition of the Gospel of John, which contains this result, is therefore assigned. Weisse, on the other hand, starts from the view that the undisputed Pauline epistles, too, presuppose the idea of God's becoming man, so that this idea must have developed among the apostles themselves, as they sought to make clear to themselves the content and significance of what they had lived through. As regards the Gospel of John, Weisse regards only the didactic part (the prologue and the discourses) as the work of the apostle John, while the insertion of this part into the form of a narrative presentation is only ``a not very successful attempt, testifying to a thoroughgoing unfamiliarity with the true connection of the matter, by a foreign, unknown hand.''}. What,
according to the fundamental idea of Christianity, must happen to everyone
who is to attain salvation, that he be ``born of God'', the Church
beheld in the dogma of the incarnation as an outer historical
fact, accomplished in the person of its founder. There
was no longer any difference between the ideal
and the historical Christ. In no other way could
that fundamental idea be preserved and brought into the consciousness
of the race. The need to stake out its domain positively,
notably against the bold speculations of the Gnostics, led
to a definite confession of faith and a fixed collection
of sacred scriptures. At the same time the Church teachers made use,
in their dogmatic investigations, of the
Greek philosophy, which had no clear consciousness
of the relation between speculation and experience, but
% --- p. 206 ---
\setcounter{footnote}{0}%
\opage{206}was in its essence dogmatism. All this led to a Church authority
stiffened in dogmatic forms of doctrine, within
whose narrow bounds only the untiring labor of thought of the scholastics
and the enthusiastic inwardness of the mystics pointed
back to the original sources. Not even Luther's profound spirit, fresh as nature itself, could escape the power of dogmatism
and supranaturalism; he too, but
still more those who have called themselves after him, were caught
in the bonds of faith in the letter\footnote[1]{Cf. Weisse's essay in the history of dogma: ``Die Christologie Luthers.''. (1852).}. The splitting into several smaller church communities that followed was bound to bring with it a
greater weight on the statutory, the fixed dogmatic
determinations. Theology now became only a thetic
presentation of these, which moreover lacked the grandeur
of medieval scholasticism. The essence of religious
experience became more and more unknown, the less
theologians cared to keep up with the great advances
in the domain of natural-scientific experience.
Schleiermacher was the first to assert energetically the significance of religious
experience, and at the same time historical criticism has
provided the instrument by which the content of that experience
can be scientifically sifted and made ready for
the philosophy of religion. With the development of the latter the
dogmatic forms and limits will fall, and the free knowledge
of religious truth will take the place of
ecclesiastical theology. As regards the Church's doctrines
of faith, too, the word of the Gospel will then prove
true, that the seed-corn must die to be able to bear fruit.

Historical criticism leads Weisse to a sharp
distinction between the ideal and the historical Christ.
--- He finds in Christ's teaching three main concepts, expressed
in pregnant forms: the Kingdom of Heaven, ---
the heavenly Father, --- the Son of Man.

The fundamental thought is, as already touched upon, that man's
% --- p. 207 ---
\setcounter{footnote}{0}%
\opage{207}salvation consists in living community with God and all
the spirits striving toward the same highest end. This
great thought cannot have arisen from the mere transformation
of a half-superstitious notion fixed in popular
belief, whereas one can well imagine
that in the hearts of those who were not yet able
to bear it, it could become the source of a series of fantastic
notions of that Kingdom's past and future.
It is the task of true apologetics to show that
Christ was so superior to his surroundings
and to a long series of following generations that these, in order
to hold fast to his fundamental thought, had to clothe it in
mythical forms, which soon stiffened into dogmatic forms;
every other kind of apologetics itself stems from
these errors and works to perpetuate them.

No conditions for entry into this Kingdom are set other than
those that lie in the nature of the case. Christ has
summed up these conditions in the pregnant demand
to let oneself be born anew, to let oneself be born of God,
to become a child of God. Herein the deep significance
of the designation of God as the heavenly Father is understood.

In equally close connection with the idea of rebirth
stands the expression Son of Man. It designates the sum total
of the human beings reborn out of natural humanity.
It is thus an ideal collective concept.
But it also became a designation of the historical Christ
as the personal representative of this ideal
concept. That this concept goes beyond this individual
application is seen, among other things, from the passages where Christ
distinguishes himself from ``the Son of Man'' by letting
the latter appear as a third person (especially in the intimations
of the Passion and the Judgment of the world). Only in the visions
in which the disciples believed they saw their crucified master
as risen did the ideal and the historical
Son of Man melt together in such a way that the master's person now
became the bearer of all ideal and religious import, and faith in
% --- p. 208 ---
\setcounter{footnote}{0}%
\opage{208}him as a historical person was thereby posited as a necessary
condition of salvation. The knowledge was soon lost
that the mystery of the becoming-man in all who
are reborn is essentially the same as in the One in
whom the Church beholds the fullness of the Godhead in bodily
indwelling; only the mystics sought their way back to the religious
fundamental idea.

With this there followed also the change that the consciousness
of salvation was built exclusively on the consciousness of sin.
With the exclusive conception of the concept of salvation
came the one-sided conception of it. In Christ himself
the negative moment, the opposition to ``this world'',
was wholly subordinate and enclosed in the great positive
view of the powers of the new life. Already the disciples
beheld this negative as a positive fundamental corruption.
The philosophy of religion must here, too, lead
back to the true fundamental thought. It does not see the
working of divine grace merely as an abolition of
sin, but as a progressive life-process by which
the divine ideal is realized in humanity
in the large as in the particular, a becoming-man of
the divine, the begetting and birth of a ``Son-humanity''.
The true content of the doctrine of ecclesiastical dogmatics
concerning the union of the two natures in Christ is
precisely this life-process, whose result is human nature
united with the Godhead. This Son-humanity, coming to be in the human
race, is the ideal Christ,
from whom salvation proceeds.

In the historical Christ this ideal
Christ personified himself by an act which has found its mythical
expression in the stories of his conception and baptism,
and which is best compared with the way
in which among the ancient peoples the mythological notions of the gods
and the peculiar national character came into being
by one and the same act. Such an act is, further,
% --- p. 209 ---
\setcounter{footnote}{0}%
\opage{209}faith as well, by which the individual human being in his
whole personality appropriates the highest good. Faith
is in the religious domain what self-consciousness is
in the logical and psychological domain. The object of faith
is the highest good itself, in whose subjective
appropriation salvation consists, and it is in itself independent
of all historical and dogmatic presuppositions.
Or should the faith of an Antigone, she who knows herself born
``not to hate with, but to love with'', who knows
that ``the time is longer in which she must seek to please
the powers of the beyond than that in which she
must please powers here on earth'', --- should this faith
be of less worth than the faith of Rahab the harlot\footnote[1]{Hebr. XI, 31.}? Should
the speculative certainty of a Socrates and a Plato, grounded on the Idea of the good
and the consciousness of its eternity,
be of less worth than the faith of a Barak and Gideon, a
Jephthah and Samson\footnote[2]{Hebr. XI, 32.}, which is mixed with such
dark elements of wild and cruel superstition?

What the Church has mixed up with the universal,
saving faith are the special determinations that
have validity within a narrower community of faith. The
Christian Church traces its consciousness of salvation back
to the historical Christ, who uttered the idea of the
highest good with a purity and clarity like no other.
Faith in this Christ therefore becomes here, by the nature of the case,
a condition of personal participation in the outer
church community. Yet only this faith as such, as a pointer
to the historical point of attachment for the consciousness of salvation;
a more precise dogmatic determination would
violate the universality of faith. The path the Church struck
out on when it sought to give such more precise determinations
has led to a plurality of separate churches, each
with its dogmatic confession. The true evangelical
% --- p. 210 ---
\setcounter{footnote}{0}%
\opage{210}Church will, in contrast to them, be the community that
expresses the ideal content of Christianity free from all
dogmatic restriction. Until it develops,
the state, as the bearer of the general humane and ethical
religion of reason, is a necessary integrating power over against
the dogmatic narrowness of the particular churches, in which that humane
element has been pushed back and obscured. As such
the state also represents the principle of tolerance, which
cannot have its home in any of these churches. The state thereby
becomes a means to the development of the true Church
and to its self-liberation from the hampering limits. ---

Weisse's doctrine, of which we have now given a general
survey, though in such a way that we dwelt chiefly on
its starting point and end point, seeks in its various
parts to give a progressive proof of the reality of the Godhead.
Metaphysics points out the universal foundation,
which in an inner way unlocks and specifies itself in
concrete forms, whose stages of development the real-philosophical
sciences follow, until the whole world-view
reaches its highest point in the philosophy of religion,
the most inward unity of speculation and experience. As
in our survey, so too in our assessment it is the starting point
and the end point that we shall chiefly dwell
on; it is they that give Weisse's doctrine its distinctive character.


Metaphysical knowledge is directed, according to this
doctrine, not at something real, but at the universal fundamental conditions
of all reality, the pure possibility of existence.
Liberation from the dominion of conceptual dogmatism
is thereby made possible, while at the same time the
significance of speculative conceptual development is retained. It thereby
becomes possible, further, to ascribe to experience
a significance of principle for the world-view without
lapsing into an immediate empiricism. And yet
Weisse himself, as we have seen, has not avoided the defect
that must have lain close at hand for him, namely that of giving
% --- p. 211 ---
% print: methaphysisk (misprint: so printed)
\setcounter{footnote}{0}%
\opage{211}that sum total of categories, ``the possibility of existence'',
a closed being for itself in opposition to the concrete
idea of God and the real world, like the dark fate
that in Greek mythology formed the background of
divine and human life. In part
this rests, no doubt, on an imperfection in the presentation, which
is connected with Weisse's whole theosophical tendency,
which often leads him to let concepts crystallize
for the imagination; and that fate is
also, as shown above, taken up in the further development into the
concrete forms. But on this point Weisse is nevertheless
still strongly influenced by Hegel. This shows itself notably
in his developing metaphysics purely a priori and placing
it at the head of the whole system. Even if metaphysical
knowledge is in its principle independent of experience,
it nevertheless presupposes experience in its actual development,
insofar as the human spirit only through
the critical working-through of the content of experience
is able to become conscious of the universal and necessary
truths. In this way the knowledge of reality itself acquires
metaphysical significance. Weisse has himself
made a beginning toward this insofar as, in opposition to Kant and Hegel, he has asserted the metaphysical significance of space and
time; on this point he is
a forerunner of Trendelenburg's acute logical investigations.


In the philosophy of religion Weisse has the merit
of having led beyond the one-sided emphasis of the speculative conception
on the intellectual, without on that account
dissolving the religious into pathological illusions, like Feuerbach,
or leading us out into fantastic transcendence, like Schelling.
He ties the essence of religion closely to the
psychological (feeling) and the ethical (the highest
good), and can thereby at once assert the specifically
religious in its distinctive character and preserve its rational,
humane character. The difficult point lies
% --- p. 212 ---
\setcounter{footnote}{0}%
\opage{212}here in the very relation between the subjective (psychological)
and the objective (ethical) moment. Weisse shows
us only by way of analysis that these two moments,
personal experience and the objective ideal, are contained in
the religious, but the inner relation between them, the
bond that ties the psychologically individual to the
highest, all-embracing purpose, is not demonstrated. He
has here left behind a task whose solution can be
won only through deep-going psychological, ethical, and
religious-historical investigations, but for which he has himself
in many ways provided preparatory work.

In his treatment of problems in the philosophy of religion
Weisse constantly seeks points of attachment within
earlier Church doctrine. This endeavor for organic
development has --- however much that is instructive and interesting
it also brings to light, owing to Weisse's uncommon learning
--- led him to overrate the significance
of the ecclesiastical symbols after their dogmatic
validity has been overthrown. That ecclesiastical theology
has played out its role, and that only the philosophy of religion
can take the vacant place, he harbors no
doubt. But just as in his conception of Christianity
as a whole he lays the weight on the humane and immanent
side of its essence, on its ideal content,
not on the form of transcendence under which
that content appears, a form he puts down to the account of ecclesiastical
mythology and dogmatics, so he thinks
that the philosophy of religion will become the doctrine of faith in a church of the future,
where the idea of Christianity has come into its
own against the distortions of the particular churches. He proceeds
from the assumption that the consciousness of the historical Christ was in itself
raised above the mythological and supranaturalist
standpoint, and wants, like neo-rationalism, to lead back
to the religion of Christ in the sense of the religion that
Christ himself had. But this is in any case
something that only historical investigations can
% --- p. 213 ---
\setcounter{footnote}{0}%
\opage{213}establish, and it becomes, when it is set up as fact ahead
of such an investigation, just as much a mere dogma as
the ecclesiastical doctrines. And as regards the possibility of an
ecclesiastical community built on such an ideal Christianity,
the central dogma that has constituted
and hitherto held the Church together arose precisely through the fusion
of ``the ideal'' and ``the historical Christ'',
and the possibility of an ecclesiastical community in which,
together with this dogma, all other dogmas are to be excluded
from the confession may surely be doubted as much
as that of the Platonic state. Nevertheless it is perhaps
along this path that the human spirit will gradually work
its way out of the dominion of faith in authority and of dogmatic intolerance,
and Weisse's works will in this respect, too, keep their significance. ---

Before we leave Weisse, we must still add to the general
survey of his doctrine a closer presentation
of his idea of God\footnote[1]{``Das philosophische Problem der Gegenwart''. --- ``Philosophische Dogmatik.'' Erster Band.}, with which, in turn, the main points
of his philosophy of nature\footnote[2]{``Philosophische Dogmatik. Zweiter Band.''} stand in inner connection.

The foundation of his idea of God is formed, as we have
seen, by the absolute Idea, the sum total of the
eternal conditions and possibilities of reality. This Idea is absolutely
inherent in God, present from the beginning to his
consciousness; indeed it is the necessary condition of this consciousness.
It is not the Idea itself, in its rest, absolutely indifferent to all change
and alteration in time, and in its unconditioned
equality with itself, that is God; but the
thought which has this Idea as its sole content and itself
springs from its womb, this thought is the foundation of the divine
self-consciousness, as unchangeably resting as it is restlessly and ceaselessly
moved. In general, self-consciousness and infinity are regarded
% --- p. 214 ---
\setcounter{footnote}{0}%
\opage{214}as incompatible. But the relation is, on the contrary, this, that
it is infinity that first makes self-consciousness
possible. Consciousness apprehends both itself and things
only by virtue of the infinite possibility of existence, which
to that extent is constantly its proper and immediate
object (cf. above p. 188—192). In the finite I this consciousness is, to be sure,
conditioned by a sensuous life of
representation; but it exists only from the moment
when the infinite Idea of reason becomes an object for
it, even if at first only in general and indeterminate
outlines, which only by a thinking expressly directed to them
are developed to determinateness, to the express knowledge
of the limits of the possible, which makes up the content of mathematics and
metaphysics. In the Godhead, on the other hand, it is
the infinite thought-content of the Idea itself, in its absolute inwardness,
that forms consciousness, and not, as in the finite
I, thoughts that mirror an external and foreign content.
But the divine self-consciousness cannot
be infinite unless infinite time and infinite
space are inherent in it. Every divine act of consciousness
is at the same time an act of filling space and time;
the life of divine thought is an eternal begetting and bearing,
by which the universal in the essence of the primal unity is specified
into the infinite. The new content is present in the universal
Idea in a way corresponding to that in which, according to the laws of logical
thinking, the particular and the single are present in the universal
concept, which encloses a manifold of possibilities,
which, by being realized, appear partly side by side,
partly one after another.

The principle according to which the specification of the universal
content takes place is matter in God. The scholastics already
called matter principium individuationis
formarum. This principle can theosophically be called God's
heart, his creative imagination, the expansive or
centrifugal principle in his essence. The intuition of such
a centrifugal principle lay at the root of the pantheistic stirrings of
% --- p. 215 ---
% print: Prædikat, Dette (misprint: a comma for a point)
\setcounter{footnote}{0}%
\opage{215}paganism, which were more or less
held down by the moral and intellectual powers, but
here and there nevertheless broke through in the wild offspring of unbridled
imagination, notably in the notion of a
blindly begetting and destroying world-power. In its inward
unity with the primal subject, with the absolute Idea, this
principle appears in the Christian doctrine of the eternal
Logos. That God as Father begets from eternity a
Son one with him in essence means, then, that God leads not merely an
infinite life of thought in pure reason, indifferent to every limit
and determination of time, but also a producing
life of nature, such as Giordano Bruno, among others,
has depicted with such great enthusiasm. In this inner,
infinite producing the primal subject works as teleological
principle, ordering and gathering all.

But if we held the idea of God fast at this stage,
we would not get beyond a spontaneous life of nature, guided
by a metaphysical principle, and would not reach free reality.
To thought and nature in God the will must be added as a third,
concluding element. The will is essentially
a limiting power, and one that forms by limiting,
not an immediately producing one. It therefore lies in its concept
that matter and content must be given, or rather
constantly be given anew, by the living thoughts that
spring from the same source as the will itself. The divine
nature is, so to speak, the body of the will; the
eternal forms of nature come into being as the expression of the divine
will. The supranaturalist conception errs
in regarding the divine will as natureless;
it thereby cancels the very concept of will and can only
by the salto mortale of absolute thoughtlessness bring about a
transition from the concept of God to the concept of the world.
But one errs likewise when, with Plato, Leibnitz,
and Hegel, one immediately ascribes to God's pure absolute
essence the predicate of the good. This predicate can then
acquire no real content different from the purely
% --- p. 216 ---
\setcounter{footnote}{0}%
\opage{216}metaphysical attributes. The good consists essentially
in the will to communicate something real, the possession of which is thus
presupposed in the willing subject. The ethical
has the aesthetic, the good the beautiful, the will the
rich heart and creative imagination, as its real presupposition.
The ethical concepts lie, to be sure, as possibilities
in the ideal presuppositions of the divine understanding,
which as such are also the presuppositions of the will;
but also only as possibilities, which are first realized
not merely through but in the will. By the direction given in the will
to the real, producing forces there is first called forth
the content which the ethical concepts express, namely
partly the will itself, which only now is a real will, partly
the divine determinations of nature, which are now absorbed into
and shaped by the will. This good, i.e. communicating,
will is what appears in the doctrine of the Church as
the Holy Spirit. What Church doctrine presents mythologically
as a trinity of persons are thus necessary
moments in the idea of God, which can be separated from one another
only conceptually, in order to be seen in their clarity and
distinctiveness.

God's essence does not stand closed and finished in itself.
Since the infinity of time is the element in which God's
life eternally unfolds, there must also take place an eternal
progress and development of his essence. All divine
elements are in an infinite process of elevation.
By this alone does an existence relatively independent over against God
also become possible. The divine will is
communicating; it posits itself as basis and possibility in order
that the content of the eternal nature may pass over into reality.
It thereby becomes world-matter, which is the intermediate determination
between nature in God and nature outside
God. Through the concept of nature the consciousness of the world is tied
to the consciousness of God. When we designate the content of the world
as nature, we regard it as that which moves
in a constant process, in which it alternately begets and
% --- p. 217 ---
\setcounter{footnote}{0}%
\opage{217}is begotten, produces and is produced. This process has
its principle in the eternal nature or matter in God, from which
it springs through world-matter, the divine will posited as
basis and substance, as intermediate determination.
The change which is thereby effected in
the eternal content of nature has its expression in the force of resistance,
in the law of inertia, which is therefore the characteristic mark
of world-matter\footnote[1]{When Weisse combats the modern atomic theory and theory of monads (which we shall come to know more closely in the next chapter), this rests partly on a misunderstanding similar to the one we have already pointed out in I. H. \emph{Fichte} (p. 170 f.), partly on factual mistakes (as when he assumes that the law of inertia ceases along with gravity). In his general conception he really does not stand so far from the modern theory of monads, since he, like it, sees the principle of matter in the spontaneous, individualizing force.}.

The outer, material existence which thus develops
is not forsaken by Idea or spirit. Inwardness does not,
to be sure, rule here immediately, as in the inner life of the
Godhead. The teleological now has at every point
the mechanical as its necessary foundation and must work
its way forward out of it. But there is nevertheless an indwelling
stream of ideality in the material. Between these forces,
the spiritual-teleological principle and the mechanical-material one,
a struggle for existence is waged, from which the
successive subjugation of matter results. As long as the
forming force still has to do immediately with the masses,
no life of consciousness can arise. The end does not
immediately coincide with the beginning. Development does not
even proceed in a straight line. For in the material forces there stirs
a spontaneous independence in which the possibility of evil
lies; this spontaneity can, namely, miss its end,
turn aside from the way. It is an inhumane and barbaric view
when Christian orthodoxy derives physical
evil from moral evil, from sin, and in turn conceives
sin as man's conscious guilt. On the contrary, just as
% --- p. 218 ---
\setcounter{footnote}{0}%
\opage{218}the contrary opposition always presupposes the contradictory\footnote[1]{Cf. S. \emph{Heegaard's} formal logic p. 13 f.},
so moral evil presupposes
physical evil. The blessedness of the divine nature cannot
be transferred immediately to finite nature. Such
a transfer must after all constantly presuppose, and be modified
by, the material foundation as intermediate determination.
The feeling of pleasure thereby constantly turns over into wretchedness,
and animal life is a constant alternation of pleasure
and displeasure. In the feeling of displeasure and pain is expressed the
struggle by which the living being, at the beginning of its existence
as well as at every moment of its existence,
must assert its selfhood against dead matter. Herein
lies the metaphysical necessity of physical evil.

The animal's soul is not a separate thing added to matter from
outside, but is the unlocked kernel of the material substance itself.
The animal's mental life springs from the
spirit of nature which from the first dwells in world-matter and
manifests itself in every single act of the cosmogonic
process. Through the circuit of organic processes
this spirit is held fast in the individuals of the animal kingdom in a
lasting succession of inner vital functions and vital states.
From the lowest stages of organization to the highest
there extends a common basic stock, whose offshoots are the
animal species that have died out at various times or still exist.
The assumption of such a basic stock, to which Lamarque's
and Geoffroy St. Hilaire's investigations seem to lead,
is as little in conflict with the rightly conceived concept of creation
as, e.g., the natural connection in the
propagation of the human race.

Already in animals there occurs a sporadic thinking.
But rational beings proper arise only
when consciousness can raise itself above the sporadic and
spontaneous and take possession of itself by a free act.
Self-consciousness is acquired; it is not originally given.
% --- p. 219 ---
\setcounter{footnote}{0}%
\opage{219}Only through language is such a firmness won in consciousness
that the sporadic can be pushed back and self-consciousness
develop. Only through language, therefore, does
reason become the specific character of man. Self-conscious
reason is, however, still only the formal expression
of man's essence. The highest end, the appropriation
of the real ``image of God'', is reached, as has been shown
in the foregoing, only by a free rebirth, whereby
the individual enters into what modern philosophy calls
the sphere of absolute spirit, and what Christianity calls
the Kingdom of God. With this we have come back once again to
the philosophy of religion, which we presented earlier. ---

Against this whole doctrine it has been objected\footnote[1]{\emph{Pfleiderer}: Das Wesen der Religion. 1869. p. 236.} that it
makes to pantheism and materialism the questionable concession
that spirit cannot be the original,
but must be the result of the development of the non-spiritual.
Yet this objection can be raised against
Weisse only from a standpoint that regards the divine spirit
as closed and isolated in itself. It is precisely
Weisse's merit, and it is the point by which he
especially distinguishes himself from the other representatives
of speculative theism, that he has led the concept of the divine
spirit beyond that dualistic presupposition.
It has thereby become possible for him to enter into the grand
views that are pressing forward in the natural science of our time,
and to show how the ideal and the ethical-religious,
rightly and universally conceived, are not violated
but precisely come into their own through them. Without, as it
seems, knowing Darwin himself, he has, by using
Darwin's forerunners, arrived at the problems that
Darwinism must necessarily set in motion. And
however much may be objected, on the whole and in particulars,
to his treatment of these problems,
it nevertheless undoubtedly points in a right direction.

% --- p. 220 ---
\setcounter{footnote}{0}%
\chaphere{CHAPTER FIVE.}
\opage{220}The investigators with whom we have been occupied in the preceding chapter had all developed under the
overwhelming impression of speculative idealism, and despite all their resistance had to go on retaining its stamp.
In this way the continuity in the historical development of thought was no doubt maintained; but on the other hand
the pure and full development of the new elements that are pressing forward was thereby hindered, and these still
gain only a subordinate influence on the view as a whole. One will be able to give oneself up with greater energy to
these new, formerly suppressed elements when historical continuity plays no immediate role, but the investigator from
the outset stands freer and more critical toward the thought of the preceding age.

Speculative idealism was a view of the whole. In the Idea of the Absolute it saw the essence and content of existence,
and through the development of this Idea by way of the pure, dialectical movement of thought it believed it reached
the whole fullness of existence. The particular received justification only insofar as it could find its place at one
point or another of this dialectical development. In opposition to this, weight is now laid on particular
investigations; through the analysis of the particular one seeks to approach the heights from which the great,
comprehensive syntheses become possible, whereas one regards it as illusion when thought would spin a knowledge of
reality out of itself. Hence the necessity of an intimate relation between philosophy and exact empirical science is
stressed. ``The life of science,'' says Trendelenburg already in the preface to the first edition of his ``Logische
Untersuchungen'' (1840), ``consists, like all life, in struggle, struggle as much with opinions that set themselves
% --- p. 221 ---
% print: og og (misprint: so printed)
\setcounter{footnote}{0}%
\opage{221}against it as with facts that will not bend to thought\dots{} The struggle with facts is on the whole the most
difficult; for they stand, when they are rightly apprehended, unbending, and thought must comply with them in order
to be able to subject them to itself.'' In this last sentence the peculiar opposition to speculation lies expressed.
Speculation took the victory of thought to be an immediate one, and did not see that the way to this victory often
goes through apparent submission and resignation. With the stress on facts the ideal does not lose its worth and its
reality. On the contrary, one holds fast precisely to the conviction that in the ideal all worth is contained, and
that over against it everything factual, all mere being, stands in a serving relation. ``The true and living
substance of every thing consists in reality only in the idea which it expresses; the real is in general nothing
other than the ideal.'' (Lotze. Streitschriften 1857. p. 45). The exact, punctual apprehension of reality leads of
itself to an ideal view, which has this advantage over speculative idealism, that it does not, like the latter, have
real existence outside itself and against itself, but is won precisely by penetrating it. The mechanically acting
atom itself becomes from this standpoint an ideal moment, taken up into the inner coherence of the whole. Therefore we
designate this whole school as \emph{the Idealism of Reality}.

The merit of this school, seen from the other side, consists in this, that the ideal, which in speculation fluttered
loose, as it were, to all the winds, is made punctual and determinate by being brought into true and intimate
connection with reality. But precisely in this connection, by the nature of the case, the greatest difficulty must
lie. In the last resort every philosophy must rest on a highest principle, from which all reality springs according to
an inner law of development. This last and highest task here retreats into the shadow before analysis and induction,
and on account of the
% --- p. 222 ---
% print: videnskabe-Ideal (misprint: videnskabe-|Ideal, for videnskabelige Ideal)
\setcounter{footnote}{0}%
\opage{222}surfeit that the bold deductions of the preceding age had left behind. Yet the
importance of the task is expressly acknowledged. ``The dialectical development of the Idea which every creature and
every phenomenon is destined to express in the rational whole of the world counts for me too as the only true
statement of their concept, and from it alone do I expect an exhaustive insight into all subordinate forms of
existence and into the ground of their conformity to law. But the more vividly I am penetrated by this conviction,
the less inclined I am to take the good will for the deed and to admire the cheerfulness with which I saw many a one
constantly hasten with good courage to this insoluble task. Do not believe, therefore, that this ideal of science
should be unknown to me and to many who in this respect express themselves as I do; it moves our heart too, but too
deeply for us to want to help ourselves past the seriousness of the problems with such meager answers as those we for
the most part see given.'' (Lotze. Streitschriften p. 98). In contrast to the dialectical facility with which
speculation mediated all the oppositions of existence, we here often meet with problems that are declared insoluble.
But the Absolute is held fast as a scientific and practical ideal, with an enthusiasm that is quieter, but not less
deep, than in the heroes of speculation.

1. \emph{Adolf Trendelenburg} (b. 1802, d. 1872). --- We begin with a thinker who still points definitely to historical
continuity, if not specifically to continuity with the immediately preceding period. Trendelenburg even lays a special
weight on the history of philosophy, just as he himself is one of the most gifted cultivators of this
science\footnote[1]{``Historische Beiträge zur Philosophie''. 3 vols. 1856---1867.}. It is a German prejudice, he says
% --- p. 223 ---
\setcounter{footnote}{0}%
\opage{223}(preface to the second edition of ``Logische Untersuchungen'' 1862), that every philosopher must begin on his
own account, have his own wholly peculiar principle, a mirror ground according to a special formula in order to catch
the world in it. Philosophy will not attain its old power until it gains a firm standing by growing in the same way
as the other sciences and developing continuously, without beginning anew in every head, but in such a way that it
takes up the problems historically and carries them further. The principle of philosophical knowledge has been found
and does not first need to be sought: it lies in the organic world-view founded by Plato and Aristotle, which after
them has been further developed and is gradually to be completed through deeper investigations of the fundamental
concepts and in interaction with the real sciences. --- It is the ancient, especially the Aristotelian, conception
that Trendelenburg means to carry through, with the transformations that the changed presuppositions and the whole
intervening history of knowledge make necessary. Not without reason has it been objected that he, who is so
unrelenting a critic of Hegel, passes lightly over obvious errors in Aristotle. But this attachment to the ancient,
and the close connection between thought and intuition in him that is akin to it, give his expositions their noble
plastic form; it becomes easy for him to let thought move in the world of reality without losing its original force.
He is, in Søren Kierkegaard's expression about him\footnote[1]{``Uvidenskabelig Efterskrift''. p. 79.}, ``a man who
thinks soundly and is happily educated by the Greeks''. Yet one can perhaps also derive from this a certain lack of
thoroughgoing movement of thought, a resting in intuition without a corresponding conceptual foundation, especially in
the treatment of the last, concluding problems.
% --- p. 224 ---
% print: Methaphysik (misprint: so printed)
\setcounter{footnote}{0}%
\opage{224}While most philosophical systems proceed from the whole in order to reach a knowledge of the parts
contained in it, Trendelenburg means to strike out on the opposite way. Knowledge strives to interpret the miracle of
divine creation by creating it anew in thought; but when one begins on this task in the particular, it leads further
of itself. With the same power by which things have proceeded from the eternal ground, they point back again to the
ground. When, then, the investigation of the particular leads to fixed points, a foundation is thereby reached from
which a more perfect understanding of the whole can be sought.

When knowledge has reached its completion, the universe being completely reproduced by the knowing spirit, the
opposition between philosophy and the special sciences is sublated, all sciences becoming members of the same
organism. But this is an ideal. In reality each science represents the knowledge of a special, limited sphere, and
philosophy seeks the Idea of the whole in the parts, the universal in the separate. For the individual sciences point
beyond themselves. Each of them concerns only one special being, and the question then becomes one about being as a
whole and about the ground of the special forms of being. Seen from this side, the special sciences point to
metaphysics, which concerns itself with what is as such, the universal ground of the separate objects. Every science,
further, has its peculiar method, its special procedure, by which it embraces its object and sees through to its
ground. But at the basis of all methods lies the same thought, and the question arises how the different methods
spring from the essence of thought. Seen from this side, the sciences point back to logic. The fundamental science,
the general theory of science, becomes at once logic and metaphysics (when it is designated as logic, this name is
thus taken
% --- p. 225 ---
% print: Grundsætninger. (misprint: a point for a comma)
\setcounter{footnote}{0}%
\opage{225}in an extended sense). In this fundamental science the unity of thinking and being in necessity is
demonstrated. Without thought there would be only a dead actual, nothing necessary. But without being there would
equally be no necessity. Thought, in order to reach necessity, must at every point let itself be determined and bound
by the nature of the thing, make itself one with the thing and demonstrate the law out of it. In necessity, which
makes science science, being is in a peculiar way taken up into thought, thought into being.

Therefore such standpoints are to be rejected in which thought is either entirely separated from being or made
absolutely identical with it. The first is the case in formal logic and in Herbart, the second in Hegel. --- Formal
logic distinguishes between thought and its object and means to investigate the forms of thought apart from its
content. But the essence of thought first appears in interaction with objects, just as the sense organs can be
apprehended according to their essence only by being seen in function. And formal logic, which itself defines truth
as the agreement of thought with the object, excludes itself from truth by isolating thought from the object. ---
Herbart's metaphysics is a supplementary counterpart to formal logic. He regards the content as given by an absolute
position, and thought has only the task, relying on the principles of formal logic, of removing the contradictions
with which the given is immediately burdened. But it turns out that what Herbart calls the contradiction in the
concepts of experience arises precisely only through his sharp distinction between thinking and being, and that by
this distinction he thus sins both against thought and against the object of experience.

Hegel held fast to the absolute unity of thinking and being, and believed that in his dialectical method he had a
self-movement of pure thought which was at the same time the thing's
% --- p. 226 ---
\setcounter{footnote}{0}%
\opage{226}own self-production; thought unfolding itself was to unfold the inmost essence of things. But on closer
investigation it turns out that pure thought is an illusion. a) Hegel presupposes motion from the outset and all the
way through. This is seen already in the first categories of the Logic: ``being'' and ``nothing'' are both in
themselves motionless and cannot through their unity produce moving becoming. ``Becoming'' could not \emph{become} out
of being and non-being if the representation of becoming did not go before. The ``presuppositionless'' dialectic
thus has at least one presupposition, motion. b) In motion space and time are contained, and the concept therefore
always has intuition as its foundation. Without intuition the dialectical process would be impossible. According to
Hegel it is logical negation that is to be the driving power in this process. But logical negation cannot itself
produce anything new; it can only cancel what precedes. Only with the help of intuition, which posits a real
opposition, can the movement be carried further, but then of course the immanence of the development of the concept
is broken off at every transition. c) If pure thinking were one with the thing's own essence, the dialectical method
would have to be a genetic method, one with a demonstration of the thing's own origin. But this is not so; the
dialectical thought develops, according to Hegel, a priori, without relying on experience, which alone can lead to
insight into the coming-to-be of things. We come here to the dilemma: either dialectical knowledge is independent and
a priori, but then it has experience outside itself and becomes no knowledge of the thing, or else it presupposes
empirical science, but then it is not independent and a priori.

By this critique the result is reached that the fundamental problem of the essence of science and of necessity becomes
insoluble both when one absolutely separates and when one absolutely identifies thinking and being. We must therefore
% --- p. 227 ---
\setcounter{footnote}{0}%
\opage{227}hold fast that thinking and being in their difference must stand in an inner connection, and the question
then becomes one about the possibility of this connection. To this extent one begins with an opposition, a dualism.
But the human spirit is finite and limited and lives on the influences it receives. If the human spirit were merely
free, merely independent, so that it received nothing but itself formed everything, it would no doubt be its own
master; but this solitary dominion would be like that of the bird on the desolate fields of the snow region, without
connection with the living world. Nor can one begin with the Absolute in order to derive from it the knowledge of the
finite and relative. The unconditioned no doubt contains the last ground of everything finite; but for finite beings
the finite and conditioned lies nearest, and from there they must raise themselves to the knowledge of the infinite.
As little as man can jump over his own shadow, so little can he leap over finite thinking.

We must therefore seek the elementary intermediate determination which in finite knowledge unites thinking and being.
It must be an activity; for a resting property would not be able to unite, but would remain closed in itself. This
activity must be the most general and most original, and it must be simple, not capable of being resolved into
elements. Such an activity is \emph{motion}. That motion lies at the basis of being is easily shown. Everything that
seems to rest is, regarded more deeply, nevertheless placed in motion. Astronomy has sought in vain for a fixed point
in the universe. Likewise thought is in itself motion, a motion that apprehends being, a constructive motion, which
manifests itself both in intuition and in the activity of the understanding (separating and connecting, inferring
etc.). Motion cannot be explained from anything earlier, but must be presupposed to infinity. Nor is it composite.
Space
% --- p. 228 ---
\setcounter{footnote}{0}%
\opage{228}and time, which some (Zeno, Kant and Herbart) have separated from motion, are its elements, which show it
from two sides: time is the inner measure of motion, space its outer phenomenon. --- Since motion is thus the most
general, most original and absolutely simple activity, no definition of it can be given, for a definition would of
course have to proceed from a more general and original concept.

The human spirit thus possesses an original activity, which is the counterpart of outer motion and which conditions
all apprehension. Motion is prior to experience, since it first makes experience possible. On this rests the
possibility of an a priori knowledge, and of the free, spontaneous life of the spirit, which develops with inner
necessity and springs from the spirit itself. Without such independence the spirit would not be able to behave
receptively either. Our sensuous intuition always presupposes a geometrical intuition, according to which we order
things perspectivally and at distances. The impressions serve as motives for a construction by which they are
transformed into images of external objects.

But this free, constructive motion is not absolute. A substrate is always presupposed, a bearer behind the activity.
This substrate is matter. It cannot be derived from motion, as the dynamic conception would have it. The forces of
motion are borne by an unknown thing which is no longer motion. There remains here, then, a gap in the development of
concepts, where something given in experience thrusts itself in. Motion no doubt penetrates here too and dissolves
what is fixed, but representation cannot do without the substrate, and always posits it again. Thus we come into a
circle: motion conditions everything that is, but in order that there can be motion there must be something that
moves. On the other hand, the atomistic view cannot be carried through either, but points back to the dynamic: the
atoms cannot
% --- p. 229 ---
\setcounter{footnote}{0}%
\opage{229}be thought of as absolutely inert and at rest, for then their mutual influence would become impossible. The
atoms must therefore be thought of as points of force, their proper essence consists in motion, and motion is
consequently what is original. In atomic theory too, then, matter can be thought only by means of motion. --- We
must content ourselves with the general result that insofar as the essence of matter opens itself to the spirit, this
happens only through motion.

Through constructive motion mathematics becomes possible as a pure, a priori science. From motion in and for itself
the mathematical categories spring (figure and number). Through motion, further, construction and calculation
penetrate into the material elements, whereby applied mathematics arises. As the spirit in this way immerses itself
in reality, we can derive the real categories from motion, which is common to thought and being; outer being is
blind and deaf and does not perceive its own acting and suffering, but thought, as self-consciousness, can regard
itself in motion, and through the observation of the different forms and products of motion the categories arise.

In pure mathematical intuition a figure or a number is produced. On the presupposition of the material substrate
this causal relation stands out more clearly. Causality, the efficient cause, is the first real category. And just
as figure and number arise through ideal, constructive producing, so material shape and magnitude arise through real
motion. Just as in language substantives are formed from verbs, so there springs from motion a thing, a closed
substance. Yet motion does not die away in its product: substance is itself causality, and acts in a peculiar way
through its qualities. In the living connection between the acting properties the concept of interaction manifests
itself, and with it in turn
% --- p. 230 ---
% print: Sammenhæng Ved (misprint: no point before Ved)
\setcounter{footnote}{0}%
\opage{230}the concept of force is closely connected: everything that contributes to the result in interaction must
be designated as force. --- All these categories serve in the domain of the efficient cause as a firm foundation for
our thoughts, and can be compared to the skeleton of the body. The categories of representation and the principles of
things are not two different worlds, but are one in their origin. As long as motion, as pure construction, is kept
within the mathematical domain, one can determine action and reaction in advance, and knowledge is a priori. But as
soon as matter is taken up as a presupposition, experience must decide how things act and motions meet one another.

But when we now, equipped with the categories won by considering motion in and for itself and matter (the
mathematical-physical categories), which have their typical expression in the efficient cause (causa efficiens),
step out into the world of experience, it turns out that these categories do not suffice to explain everything we
meet here. Nature offers facts such as the sprouting seed, the eye that opens to the light, the organs of locomotion
of animals corresponding to their element, where a plurality of moments are so connected that the one points to the
other, so that the explanation of the single parts can be sought only in the whole connection. With the efficient
cause, conversely, the parts explain the whole as their effect; here, on the other hand, the whole (the effect)
becomes purpose and as such cause, the end guides the beginning, the future the present. This is possible only
through the effect's being present beforehand as thought, and through this thought's being one with the efficient
cause. Through the efficient cause the purpose (the thought) attains its reality; in the purpose (the thought) the
efficient cause has its truth.

On the presupposition of the reality of purpose in nature a new relation between thinking and
% --- p. 231 ---
% print: til sin Fuldbyrdelse (misprint: til doubled over the line turn (til|til))
\setcounter{footnote}{0}%
\opage{231}being appears. Hitherto we have seen being act on thought and thereby become the cause of knowledge (causa
cognoscendi). In purpose, on the other hand, thought acts on being (as causa finalis), since it requires an
efficient cause to realize its content. The necessary presupposition here, then, is that thought itself is one with
the essence of what is; only thereby can it have power over the real. For this too the possibility lies in motion.
Just as we know outer motion only through our spirit's own motion, so too we know the purposes realized in nature
only because our own spirit is able to set itself purposes. Motion was the first a priori, purpose is the second. But
purpose presupposes motion and material being, and is therefore only a relative a priori. What compels us to give up
the blind sole dominion of the efficient cause is its own impotence. In the face of phenomena such as organic life
another way must be taken. In particular cases mistakes are no doubt possible here, but on the whole it stands firm
that light cannot be explained by darkness. The concept of purpose makes an explanation possible where the efficient
cause does not suffice. But purpose we can grasp only from the ideal side, as the thought of the whole which
determines the parts, consistently positing the means to its fulfillment. On the other hand we cannot observe the
point where thought grips force and guides it toward the goal. Just as the ancients represented Helios standing on
his chariot and steering the horses of the sun with his hand, but gave the hand no reins, these tools of human force,
so the thought of purpose steers the acting forces with invisible reins.

The inner purpose is the properly individualizing principle in the world. When force and purpose coincide in the same
subject, the self arises, and only the self is an individual in the proper sense. In the mere machine there is no such
unity. In the plant, the animal and
% --- p. 232 ---
% print: med fordres (misprint: med for men)
\setcounter{footnote}{0}%
\opage{232}man this unity shows itself at different stages. The soul is a self-realizing thought of purpose. The
mental activities are reflexive; they proceed from the self and have the self as their aim. In the feeling of pleasure
and displeasure the individual life appears as furthered or hindered; the acting of the being is no longer foreign to
the being itself. In the animal the thought of purpose is still hidden from itself; blind drive directs itself toward
the underlying purpose, and pleasure or displeasure follows from its satisfaction or disappointment. But man becomes
conscious of the purpose and, in knowledge of the universal, raises himself above immediate drive and sensation. In the
human, the organic passes over into the ethical. The ethical Idea is the inner purpose which thought makes into a task
for action and life. Purpose, which in nature is purely objective, becomes personal in man, and on this rests the
possibility that the dominion of reason can grow in the world. The ethical world has its history; it acts organically
and develops itself as an organism.

The concept of purpose thus leads us into two new spheres, the organic and the ethical. But the categories won earlier
(the mathematical and physical) are not therefore given up. They continue to lie at the basis, but are penetrated by
purpose and gathered under a new center, whereby they gain a closer determination. Thus the efficient cause now
becomes means, substance becomes machine or, with a deeper unity of matter and form, organism. Interaction is no
longer a play of blind forces, but an inner relation between the whole and the parts. Matter is no longer the dead
substance lying behind everything, but is required with necessity by thought, which needs it as its handmaid so as not
itself to become solitary and without life. Without matter, which separates and bears, there would be no hold for
thought, no individual life.
% --- p. 233 ---
\setcounter{footnote}{0}%
\opage{233}At its highest point the concept of purpose leads us to the Ideas of the true, the beautiful and the good.
The efficient cause gives us only the actual; truth reveals itself where the real whole corresponds to its purpose. In
beauty the inner purpose appears at the same time in a harmonious form for intuition. In the ethical, purpose is taken
up into knowledge and disposition.

It is, then, the same fundamental concepts that are freely constructed in the mathematical, filled out in the
physical, deepened in the organic and raised in the ethical. Thus we saw, e.g., the concept of the whole first in
abstract form as figure or number, next as material and concrete substance in nature, as organism in the living, and
finally as personality and moral organism in the ethical. Purpose at last takes up all other categories into itself
as a higher providence; for its sake they all are. From the highest standpoint one would be able, starting from the
last purpose and going backward, to arrive at all the categories. But we, who do not stand at the standpoint of the
Absolute, must start from motion (the efficient cause) and through it work our way forward to the knowledge of what is
in itself first and original, purpose.

The fundamental concepts developed so far all concern the inner nature of the thing. But in the course of thought's
labor on the work of knowledge, new fundamental concepts arise which have their root in thinking knowledge and concern
the thing's relation to it. These are the so-called modal categories, possibility, actuality and necessity. To sense
intuition things stand as something immediately actual, as phenomenon. The comprehending understanding now turns to the
phenomenon and asks about its ground. The ground is no more an absolute unity than the cause is, but a sum of
conditions that belong together. Now if only one or some of these conditions are present, the concept of possibility
arises, thought
% --- p. 234 ---
\setcounter{footnote}{0}%
\opage{234}supplying what is lacking. From possibility thought strives further to find the whole ground, the whole
system of conditions. Herewith the transition from possibility to necessity takes place. The inner, fully developed
possibility passes, in the direction of being, over into actuality, and presents itself to thought as binding
necessity. In its immediate form necessity is no doubt only the unavoidable, but through the acknowledgment of thought
this negative becomes a positive grounding. Thought is no doubt not a link in the connection of events, but the
concept of their necessity first arises when they are measured by thought. Until then they are only occurrences; now
they become necessity. The acknowledgment the spirit renders here is not weakness, as though it would rather have
something else but had to submit; it is rather its essence and its strength, since it impresses on blind existence the
higher stamp of necessity.

The question now becomes in what peculiar forms thought solves the task whose possibility has thus been demonstrated.
The forms and connections of thought must correspond to the forms of being, since the possibility of knowledge rests on
motion, which is an activity common to thinking and being. In the parallelism of the forms the original unity mirrors
itself. --- Just as substance emerges as the result of activity and is in turn active in its properties, so the
concept is the result of an original judging and develops its essence through continued judging. The concept is
substance apprehended as universal. As apprehension of substance the concept has its determinate content, as universal
its determinate extension. Content and extension are related to each other as the law to the phenomena. --- To these
moments of the concept the main kinds of judgment attach themselves: the categorical (and hypothetical) judgment is
essentially a judgment of content and aims at a generalization of the subject; the disjunctive judgment is a judgment of
extension and aims at a
% --- p. 235 ---
\setcounter{footnote}{0}%
\opage{235}specification of the subject. --- The essence of the syllogism, too, is conditioned by the two moments of the
concept. It is the movement of thought either from the law (the content of the concept) to the phenomenon (the
extension of the concept) --- the categorical syllogism; or conversely --- the disjunctive syllogism. --- With the help
of the syllogism knowledge advances synthetically and deductively. But the analytical investigation of experience is
the constant presupposition, and not even in mathematics can knowledge always advance by the way of deduction; it needs
apagogic proofs and auxiliary lines. In place of the genetic proof, which demonstrates necessity in the thing's own
coming-to-be, there steps the indirect proof, which delimits it by demonstrating the impossibility of the opposite.
Only in connection with a disjunctive judgment that sets out all the cases that are possible at all can the indirect
proof lead to a positive result. In particular the highest principles cannot be proved directly or genetically; in that
case they would of course not be principles, but would have a foreign origin. They can be established only by a ground
of knowledge: the necessity of principles consists in the impossibility of the opposite. But in this negative
expression a positive principle lies hidden, and the more comprehensive the point of view, the more compelling force it
has. For the unconditioned principle, for God, the indirect proof must be drawn not from anything particular but from
the whole of existence.

But it is only finite knowledge whose possibility and whose forms we have now investigated. When, having run through
the forms of finitude, we seek the Absolute, the earlier means no longer suffice. One might think that, just as the
mathematical-physical categories underwent a transformation in the organic and ethical domain, so the categories of
finitude could by a similar metamorphosis be developed so as to hold for the Absolute. But such an infinitizing of the
categories is impossible, and it is an illusion when one
% --- p. 236 ---
\setcounter{footnote}{0}%
\opage{236}believes one can reach by this way a constructive knowledge of the essence of God. The content becomes a
wholly different one when what is peculiar to finitude is to be negated. Logically regarded, every intermediate
determination, every analogy between the conditioned and the unconditioned, is lacking. Thought can no more apprehend the
pure and clear light in its sphere than the eye can in its own; its work is to behold the subdued rays, broken in
the play of colors --- and when it does so, it does not, after all, deny whence they come. Of God it holds:
nesciendo scitur. Only theosophical imagination sets itself above these critical scruples.


But although one cannot construct the unconditioned, every view that does not, like Kant's, regard the categories as
purely subjective forms must nevertheless find in the finite essential pointers as to how the notion of the infinite
must be formed. What is true in the finite by its necessity cannot be untrue in the infinite. The necessary in
conditioned knowledge is the fixed point to which the supposition of the unconditioned can attach itself. The necessity
which the finite spirit has seen into may perhaps, seen from the point of view of the infinite, enter into a larger
connection and itself be enlarged, but it nevertheless remains in itself what it is. The divine spirit is like a poet
of whom we know only a single work. Our world is no doubt only a fragment, but still a whole by itself, like a single
drama in a trilogy. The way in which the Absolute is apprehended rests on the whole world-view. Here in particular the
organic world-view stands over against the merely physical or mechanical\footnote[1]{Cf. the treatise ``Ueber den
letzten Unterschied philosophischer Systeme'' in ``Historische Beiträge zur Philosophie''. 2ter Band.}. The latter
sees the world from the
% --- p. 237 ---
\setcounter{footnote}{0}%
\opage{237}standpoint of the efficient cause: motion and matter are the first and the last, and purpose is only
semblance. The organic world-view, on the other hand, sees the world from the standpoint of purpose, and thought as the
last ground of everything. After what has gone before, the decision cannot be doubtful; for it has turned out that the
efficient cause was insufficient to explain the facts of experience, and that it was absorbed as an element in the
comprehensive purpose. The organic world-view reaches its completion in the Idea. The concept becomes Idea when it is
seen as determined by the higher purpose, or shows itself in the light of the unconditioned. The Idea is contained in
the essence of God, for without God there is nothing divine.


--- Trendelenburg's ``Logische Untersuchungen'', whose fundamental thoughts we have now tried to set forth, is
indisputably the most important logical-metaphysical work since Hegel's Logic. But what an opposition there is between
these two works! Hegel goes by the way of pure thought, renounces intuition and experience, and believes that in this
way, through the rhythmic self-movement of concepts, he reaches the highest truth and reality. In Trendelenburg it is
precisely the first work of thought to put on intuition as its garment, so as thus to undertake the toilsome journey
through reality. Hegel's pure concepts are in his opinion only weakened and diluted intuitions. Human thinking lives on
intuition and dies the death of starvation when it is to nourish itself on its own entrails. Philosophy therefore
presupposes the empirical sciences, and its peculiar problems arise when the knowing spirit seeks to become conscious of
its activity within these, and therefore raises the question: how is science possible? on what is necessity grounded?

Trendelenburg's logical investigations mark the
% --- p. 238 ---
\setcounter{footnote}{0}%
\opage{238}first sober self-reflection after speculation's intoxication of enthusiasm, in which the wings of thought
became so light that it believed it could soar up into the pure ether. Kant had already laid the ground for the
doctrine of pure thought by his sharp distinction between space and time as forms of intuition and the pure concepts
of the understanding. Hegel then developed in his Logic the essence of the Idea in the form of pure thought, and only
in the philosophy of nature came upon space, time and motion. Speculative idealism thereby really became, as Weisse
showed, a complete dualism. Weisse and I. H. Fichte sought to remedy this one-sidedness, but they were themselves still
too strongly influenced by speculative dialectic to break through its bounds completely. Metaphysical knowledge was not
yet put in close connection with an analytical critique of the given. At this point lies Trendelenburg's merit.

In opposition to speculation he declares that he will occupy himself only with the finite. He does not deny the
Absolute, but regards only an indirect grounding of it as possible: philosophy leads to the knowledge of the Idea; the
Idea is divine; without God there is nothing divine. Through this indirect and indeterminate apprehension of the
Absolute there is in Trendelenburg an attachment to the religious tradition, ``the handing down from generation to
generation to which the concept of God owes its proper life''. Only in the face of finitude, in which his thought moves
with Greek joy of life and harmony, does he feel a need for a determinate and closed apprehension; the Absolute is for
him the inaccessible, though all-bearing, infinity. Yet in his own exposition two points of view here come into
opposition to each other. On the one hand he maintains that the difference between the finite and the infinite is too
great for the categories that hold for the finite to be able to undergo such a change that they could become the
expression of the infinite's
% --- p. 239 ---
\setcounter{footnote}{0}%
\opage{239}essence. But just as definitely he maintains on the other hand that the essence of necessity must be the
same whether it is seen in the world of finitude or enters into the infinite connection of the Absolute. Now if the
concept of necessity forms the firm kernel of human knowledge, then all other categories have meaning only insofar as
they stand in inner connection with it and are different forms of it. And if the concept of necessity cannot lose its
meaning by being transferred from the finite to the infinite, then the separate forms under which it appears must also
have in themselves more than merely finite meaning. Thus Trendelenburg himself in principle corrects the opposition he
has set up between the finite and the infinite.

Not only in the face of the Absolute, but also in the face of the world of fact, there is in Trendelenburg a halt in
the development of concepts. He begins with mathematical construction, but in order to get beyond this sphere he must
take up the concept of matter empirically. Now after the physical categories have been developed, a new appeal to
experience becomes necessary in order to ground the concept of purpose. Trendelenburg is here evidently influenced by
Kant's critical separation between the a priori and the a posteriori, thinking and the thing in itself. He himself
admits that matter, insofar as we can apprehend and know it, is determined by and expresses itself in motion. With
what reason, then, does he allow representation ever anew to posit a dead substrate behind the content of the living
concept? Intuition and concept, which according to his own doctrine are to be in harmony, are here in unreconciled
discord. The same thing happens with the concept of purpose. It appears in opposition to the preceding categories, an
opposition which neither philosophy nor recent empirical science can allow. A real carrying-through of the development
of concepts would here have satisfied both,
% --- p. 240 ---
\setcounter{footnote}{0}%
\opage{240}the concept of purpose being brought in an inward way into connection with the earlier categories by virtue
of Trendelenburg's own fundamental thought, that it is the Idea of the whole which conditions the single parts. The
chapter on purpose is, on the side of its exposition, one of the most beautiful in Trendelenburg's ``logische
Untersuchungen'', but here he has in several respects not risen above the popular conception. ---

There still remains to consider Trendelenburg's practical philosophy, which he has developed above all in his
``Naturrecht auf dem Grunde der Ethik'' (1860). At its basis lies the organic world-view won in the ``logische
Untersuchungen'', in which thought is seen as the last cause, from whose center force is determined. The whole grounded
in an original thought completes itself in itself and in its parts. To this principle the ethical leads back with
necessity. The views that find the starting point of the moral in pleasure, self-love or self-preservation proceed from
isolated individuals and can reach only a mechanical connection between them. The particular itself first comes into
its right by being attached to the ideal, by being seen as that which is to develop according to a purpose. Herewith we
are pointed to another starting point, namely purpose, the universal. Yet this would be an empty ideal if it did not
take root in the particular through inner articulation. Both starting points come into their right when we determine
the ethical principle as the Idea of man. What is peculiarly human rests, as the ``logische Untersuchungen'' already
showed, on this, that thought, the inner purpose, which in the lower organic beings is hidden from itself or only
sensed blindly, comes in man to self-consciousness. As man knows and wills the creative thought of his own essence, the
ethical becomes the organic set free. But the single man by himself would not be able to raise himself to this
% --- p. 241 ---
\setcounter{footnote}{0}%
\opage{241}point; only in society does the essence of man manifest itself in its full extent, and only there does the
self-liberation of individuals also become possible. Society in turn develops through history, whose inner moving force
is the Idea of man. In this great historical connection, in which the single man thus stands, it is a matter on his
side of strengthening and consolidation, of widening the scope of his powers. To this is added, from the side of the
whole, the need for articulation. Both directions must coincide in order that the ethical be present; the realization
of the Idea and the self-preservation of the individual go hand in hand in the ethical.

In personal devotion to ideal reason, with the surrender of egoistic drive, the inmost element of the ethical shows
itself: the disposition by which the will becomes the good will. Through the inwardness of the disposition the ethical
has its ground in the religious. When the Idea of man is posited as what moves in the individual and in history, the
concept of the Idea, the creative thought, points back to God. The human does not stand in opposition to the divine.
Positive religion grasps the divine Idea in certain historical facts, but philosophy has a more universal character and
maintains the universally human, of which the Christian can only be a peculiarly determined direction that cannot deny
the universal foundation. The element of religion is disposition; its essence lies at the point where the notion of the
divine meets with human affect. It has its point of attachment in man's fearing and hoping, insofar as these are
determined by the notion of a superhuman being. While in true faith it is divine thoughts that call forth the affect,
in superstition --- and man is by nature superstitious --- contingent drive rules through the associations of ideas. In
the religion of spirit man is freed from superstition. Christianity is
% --- p. 242 ---
\setcounter{footnote}{0}%
\opage{242}the center of world history insofar as in it the idea of the unity of the divine and the human comes forth
for consciousness. It presupposes the organic world-view (ἐν ἀρχῆ ἦν ὁ λόγος), and the Idea of man has found its
concrete expression in the notion of the body of Christ and its members, growing through the generations. But religion
emphasizes only the disposition, inwardness. Here too philosophy is more universal; it directs its attention also to
the carrying-through of the good in the world. Here, then, what matters is not only the good disposition, but also
the truth of the action, namely that it corresponds to the ideal purpose of man, and the harmony and beauty of its
execution. The Idea of the good in the wider sense contains these three elements: the good in the disposition, the
true in the concept, the beautiful in the presentation.

The highest good is the realized Idea as moral organism, from which the single goods, duties and virtues spring. In
this lies also the source of right. A sharp opposition has wrongly been set up between the moral and the juridical.
From the same spirit from which duty arises, right too springs, which with the power of the whole preserves the outer
conditions for the realization of the ethical. Right is not merely to defend and preserve; as in organic life, so here
too preservation is a further development, the possibility of which must therefore lie in the essence of right. Duty is
the morally necessary, right the morally possible; duty is therefore of narrower, right of wider extent. But it does not
follow from this that right, the more comprehensive possibility, should not itself be morally conditioned and carry a
moral purpose within it. This purpose is the preservation of the conditions for the life and development of the
ethical Idea in society. From this the coercion bound up with right is explained; it is grounded on the power of
society as a whole. But right does not reach its goal when the individuals' way of thinking
% --- p. 243 ---
% print: Virkeliggiørelsen (misprint: so printed)
\setcounter{footnote}{0}%
\opage{243}and acting does not come to meet it, and there does not rule in them the same moral spirit that
right is to preserve against all resistance and hindrance. Right cuts away excrescences, but cannot itself heal an
inner disease. Thus Roman law remained standing, but the Roman people decayed.

The individual has no right outside society, and within society duty comes into being together with right. The modern
doctrine of innate rights is therefore one-sided, especially when it is formulated more definitely in such a way that
the individual makes a claim on society for a livelihood. The ancients, full of piety, regarded as a benefaction what
the moderns set up as a legal claim. The ancients did not set themselves in isolation over against the state. In
recent times, on the other hand, one speaks of an opposition between state and church, state and science etc., as if
the state were a special factor and not precisely the power that embraces all the elements of human life. The state
is the realization of the universal man, of the Idea of man in the individual form of a people. Every people labors to
realize the Idea of the state in proportion to its spiritual and material powers. While human activity in the
individual is only a small fragment, in the state it appears as a many-sided whole. What the guiding reason is in man,
the government is in the state; to man's single pieces of knowledge there corresponds in the state the ramified system
of the sciences, to his single skills the arts and crafts etc. ---

Trendelenburg's ethical view, by its doctrine of the individual's being absorbed in the state, recalls the ancient
conception. Yet with him this conception is not only carried back to a deeper and more many-sided grounding, but in
two essential points in particular he sets himself in opposition to it. --- In the ancient state the humane life fell
to the share of only a few, at the expense of the greater mass, who had to
% --- p. 244 ---
\setcounter{footnote}{0}%
\opage{244}work for the necessities of those few. Trendelenburg has not only acknowledged this evil, but he finds it
again at all times. What the slaves were in antiquity, the serfs were in the Middle Ages and the factory workers are in
recent times. Their conditions stand in glaring contrast to the growing culture, which constantly provides greater
possibility for human satisfaction. Yet the development of culture itself produces its own canker. It makes
necessities more manifold, widens and increases need and desire, increases the population and with it the number of
those who desire. Those who desire desire more and become more numerous. It is easier to refute the communist and
socialist theories than to heal the evils that give them their power. For this power comes from the passions that are
in motion, not from the element of truth they contain. Laws and right can do but little here; by directly opposing
the evil they strike only the effect, not the cause. The task always remains of conquering spiritually and morally
the physical force that breaks forth through overpopulation, through human labor-power's falling in value (the fate
of our time), and that threatens to dissolve society in a struggle of all against all. --- The ancient sense for
harmony and for the plastically formed whole has thus not blinded Trendelenburg's eye to the disharmonies of existence,
or led him over into an obstinate optimism.

The ancient state let the individual be so absorbed in the state that it retained no personal inwardness, no inner
infinity; the state ruled everything and everyone immediately as the adequate form of the Idea. Trendelenburg, on the
other hand, lays weight, as we have seen, on the disposition, on the inner relation of the will to the Idea, and finds
in this an attachment of the ethical to the religious. ``Man thinks in the divine a concept that has its worth in
itself. Where therefore, as in the ethical religions, he enters into an essential
% --- p. 245 ---
% print: Værd.. (misprint: two points, spaced)
\setcounter{footnote}{0}%
\opage{245}relation to the divine, he thereby acquires an inner worth himself. In religion he ceases to have a mere
market price\dots{} In the general stream of world history and under the general fate of the peoples, Christianity
labors to let the soul find its ground in itself and not to let it go without its laying hold of itself through God.
Thus man, who otherwise loses himself in the large, remains man in himself, and only through man in the particular is
the state, man in the large, man.''

But Trendelenburg turns back again to an ancient line of thought in apprehending the religious objectively and as a
direct object of the state's care; at least it lies nearer, in his case, to derive this from the after-effects of
antiquity than from those of the Middle Ages. The state, he says, cannot do without the church and its spiritual help,
the possibility of the most penetrating effect on the disposition, so as to raise individuals above egoism and the
craving for enjoyment. It is, as already mentioned, in his opinion wrong to set the state in opposition to the church;
the state needs the church, or rather the worldly side of the state needs the spiritual. From this it does not follow
that it is to attach itself to a particular religion or confession and have a state church; it is to let the different
schools unfold freely and represents over against them the Idea of humane ethics and of tolerance. But Trendelenburg
maintains that the state must see to it that children are brought up in a particular religion, even if the parents
renounce all recognized churches. ``The supposition that it should be possible for those who want it to live in the
state without religion [i.e. according to the context: without positive religion or confession] may as little be
thought, as the old lawgivers would think of the possibility of parricide.'' As Socrates was accused of not accepting
the gods the state accepted, so, then, one would now be able to imagine an accusation
% --- p. 246 ---
\setcounter{footnote}{0}%
\opage{246}of not belonging to any of the religions the state recognizes. Trendelenburg has not brought out definitely
enough the opposition between the positive and the rational-humane element in religion. His grounding of the
importance of the religious for the state really rests on the latter element; only in the conclusion is the former
inserted. And yet it must have lain near at hand for him to bring forward the great problem he has here pushed aside,
since he himself has brought out the concept of the Idea as the principle both for the knowledge of the divine and for
the knowledge of the human. A sharper apprehension of the essence of the Idea, relying on the absolute validity of the
thought of necessity, which Trendelenburg himself acknowledges, would also have led him to a sharper apprehension of
the problem of the philosophy of religion, which is of no less importance than the social problem to which
Trendelenburg points with such noble humanity, and which surely has many points of contact with it.

2. \emph{Gustav Theodor Fechner} (b. 1801) was for a
time Professor of Physics in Leipzig, but besides
works in this science he also published various poetic
and humorous writings. The free course of thought and imagination
was not crippled in him by any one-sidedness of specialism.
Soon he sought also to develop philosophically
his world-view, as it had developed
for him in interaction with his exact
studies. From earlier philosophy he has, by
his own declaration, received important impulses. It
was a work that had sprung from Schelling's views,
namely Oken's philosophy of nature, which by its titanic
boldness first led him out beyond the ordinary
conception of nature. ``The best fruit'' he has ``plucked
from a branch of Hegel's philosophy that is admittedly bent far to the side,''
inasmuch as he owes to Weisse his conception of Christianity
in the philosophy of religion. Finally,
Herbart's mathematical psychology has been to him a
% --- p. 247 ---
% print: i den totale (misprint: i doubled over the line turn (i|i))
\setcounter{footnote}{0}%
\opage{247}model, though chiefly through its shortcomings: ``I have
burnt a piece of coal from Herbart's ashes on my hearth''\footnote[1]{Ueber die physikalische und die philosophische Atomenlehre. 2te Aufl. p. XIV.}.
Yet he sets his philosophizing in decided opposition
to the earlier philosophy. Its fundamental error consists,
for Fechner, in the endeavor to find what
lies behind the phenomenon, in the presupposition
that all that we see, hear, feel, indeed even what we think,
is to be a subjective semblance, which must be removed in order that
one may reach what is really real. Neither empirical science
nor philosophy can go behind the sensuous,
the phenomenon and what one must represent to oneself in connection
with it. In this bright world thought is to remain,
and not suppose that it will win precious treasures by sinking
into the underworld, where it will find only shadows.
How, then, can the ideal interest be satisfied,
the longing that is not content with what is merely
given? Never, says Fechner, will one be able to confine
the human spirit and the human heart so
in the phenomenon that it finds itself content in this
captivity and does not strive to reach something deeper, something supersensible.
But this need is to be satisfied not
by going back to a being behind the phenomenon, but
by following the whole connection of phenomena into the
smallest particulars and the finest nuances, as well as
into the total, concluding unity. Experience does not, indeed, reach
the infinite, nor even far into the world of
finitude; but the thread that leads us within the domain
of experience we can follow out beyond it. The points of view
that hold within experience need only to
be extended, raised, potentiated, in order that with their help
one may be able to reach the higher spheres. The natural scientist,
who step by step relies on experience and
mathematics, has to do only with the outer side of nature,
% --- p. 248 ---
\setcounter{footnote}{0}%
\opage{248}i.e.\ with nature as the object of outer observation. But
by induction, analogy and rational combination
a further-reaching consideration of nature can so extend and
make infinite what is immediately given that it leads out
beyond itself. This is the only theoretical way by
which one can reach a knowledge of what lies out
beyond what is given in experience. It is a passage through
the visible world in order to find the invisible. To want to
go the opposite way, as speculative
philosophy does, is to want to begin the building of a
tower from the top instead of first providing a solid
foundation.

The first thing one must therefore make sure of is the method
and results of empirical science, in order to build further
on this foundation. What matters here, then, is partly
the limiting concepts with which exact empirical science
concludes, so that metaphysics may then in its own way
work further in the same track, and partly such
points of view as make possible an exact conception of
the relations of opposition that play an essential role in
metaphysical respects. In both respects
Fechner has produced exceedingly meritorious works.
In his work ``Ueber die physikalische und die
philosophische Atomenlehre'' (1855. Second edition 1864)
he has given precision to the concept of matter as it
presents itself from the standpoint of natural science, and has then
sought to show how this concept can enter
into a philosophical view. In his ``Elemente der Psychophysik''
(2 vols. 1860) he has given an exact doctrine
of the relation between soul and body, which is indeed independent
of further-reaching metaphysical presuppositions,
but nevertheless points of necessity to such. In
his other writings\footnote[1]{Namely: Zendavesta (3 vols. 1851). Ueber die Seelenfrage (1861). Die drei Motive des Glaubens (1863).} he seeks to carry out the great proof by analogy
% --- p. 249 ---
\setcounter{footnote}{0}%
\opage{249}and induction on which, by his conviction, a
philosophical world-view can alone be built, and he
shrinks from no consequences, even if they seem mystical
or fantastic. Often the prose style (especially in Zendavesta)
passes over into dithyrambic poetry. --- It was not Fechner's
exact investigations that led him to his
philosophical world-view; on the contrary, his
philosophical inquiry led him to set the problem by which
the foundation was laid for ``Psychophysics''\footnote[1]{Elemente der Psychophysik, II. p. 553.}. ---

Speculation, says Fechner in his atomic theory, descends
into nature like the bear into a beehive. It praises
the coherent honey, but does not heed that the
individual little beings have brought it together and have a
right to it; it takes the honey with one grab and overturns
the whole beehive. Natural science, by contrast, is like
the beekeeper, who knows that the structure in all its unity is
the work of countless individuals, carried out under the rule of a quiet and hidden
law, and who therefore tends them,
follows their flight and makes their whole life and doings
the object of his attentive observation. As
irreconcilable strife reigns between the beekeeper and the bear,
so also between natural science and speculation.
Speculation wants to fathom the essence of nature by an a priori-dialectical way
and therefore ends in a dynamic view,
which derives matter from the interaction of universal forces,
whereby every firm foundation is abolished. In
opposition to this, natural science maintains the exact
discreteness of matter. By a series of facts
from various regions of physics and chemistry (the undulatory theory,
the theory of heat, etc.) Fechner seeks to
establish that natural science cannot conceive matter
as a continuum, but must regard it as a
system of discrete parts which for it are indissoluble.
Everything thus individualizes itself down to the most particular,
% --- p. 250 ---
\setcounter{footnote}{0}%
\opage{250}but is precisely thereby joined into the firmest whole. As
the heavenly bodies move in cosmic space, so
the atoms in infinitely small space, and here as there
a universal law reigns. Modern natural science
is therefore atomistic in the sense that it holds fast to
the discreteness of matter. But it does not, like
the old atomism, undertake to fathom the absolutely ultimate.
If by atom one understands the absolutely indivisible, then
the name atomism does not fit modern
natural science, which maintains only that discreteness must
go further than the eye and the microscope reach, but by no means
denies a divisibility in indefinitum. Modern atomic theory
is therefore relative. The physicist does not assert that
the space between the atoms is empty, and that there cannot be found
between them a fine matter which has no influence
on the atomic motions he can control,
just as the ether fills cosmic space between the heavenly bodies.
Between the atoms of this matter there could then
in turn be thought an even finer ether, etc. We move
on the whole here in the domain of the relative. It is
only of importance to hold fast to such points where
exact analysis can find its point of attachment. But
this is impossible for the speculative-dynamic view,
which explains matter as the result of a continuously
acting, space-filling force. This concept is an absurdity
for physics. For physics, force is incommensurable
with matter and cannot produce it or
dialectically become one with it. All forces are concepts of relation
and therefore presuppose something that stands in these
relations. What one ascribes to a body as its determinate
force is the share with which it works toward the fulfillment
of the law. The matter that is presupposed is the atoms,
which are not measured (like force, space and time), but
counted. The motion of the atoms fills time and space.
The mass of the atom is not extension; it is, physically speaking,
one with the constant ratio in which, for this atom,
% --- p. 251 ---
\setcounter{footnote}{0}%
\opage{251}the force always stands to the increase of velocity. The concept
of mass is therefore not materialistic. On the contrary,
the fundamental conception toward which physics in the last instance
points is absolutely idealistic. The extension of matter
is = o, and the material atom is a single being, a
punctual intensity, a center of force.

Although Fechner grants the relative and skeptical
character of modern atomic theory, he nevertheless himself goes out
beyond its limits and carries atomism through absolutely.
And he does this expressly for philosophical
or metaphysical reasons. The task of metaphysics consists
for him in finding and exploring the limiting concepts of experience
in their general relation and connection,
and its method in generalizing and carrying through
what exact empirical science has found and
established, to such a conclusion as thinking
demands, the very conditions of being something ultimate
and highest forming the form for the concepts thereby found.
Metaphysics is therefore not a priori, but answers
to its name, inasmuch as it comes after physics. Metaphysical
knowledge cannot directly extend
exact knowledge; physically speaking, it always remains a hypothesis,
but it is nevertheless the conceptual conclusion of physics,
which one reaches when one carries the way that
exact knowledge goes with certainty to its end in
the Idea. Not physically, then, but metaphysically does Fechner
arrive at the absolute atomism that conceives the
material atoms as centers of force. But he arrives
at it by holding fast immediately to the exact limiting concepts
to which physics leads. He refuses
to go behind the phenomenon to find the essence of matter.
The material is the tangible, the general
foundation for the phenomena of nature; behind it nothing can
be known, but in connection with it and beyond it there is
room for abstractions and higher points of view. The concept
of substance is one with the concept of the whole
% --- p. 252 ---
\setcounter{footnote}{0}%
\opage{252}phenomenal connection, in whose lawfulness the objective
in matter consists. It is from this starting point,
then, that physics arrives at the concept of the discreteness
of matter, and no other concept of matter can
metaphysics take as its basis either. But when one
carries this concept through, all burden and weight vanish;
every impenetrable barrier that stands against
spirit is transformed into force.

Behind the phenomenon nothing can be known. But the
phenomenal itself contains an important difference,
by which it divides into two great groups. The bodily
is that which appears only for another; the mental
rests on an inner appearing for itself. We
thus get an opposition between inner and outer phenomena.
In the thing itself, however, there is no opposition, but
it arises through the two different standpoints
that can be taken. What on the inner standpoint
appears as mind appears on the outer standpoint
as body, just as the same circle seen from within
shows itself concave, seen from without convex. Whoever therefore
knows only the outer standpoint will regard everything
as bodily and deny mind; so the materialists,
who look about only in the periphery of the circle, and since they
find no center here, deny that there is
any. Conversely, the idealists want to determine the periphery
from the center; but about the same center there can,
after all, be described manifold circles. Yet the opposition
between these standpoints does not remain standing: experience
shows us a determinate and lawful relation between
the mental and the bodily, and in psychophysics we seek
an exact doctrine of this relation. In order to
be able to establish such a doctrine, a measure is needed
for the mental and the bodily motions. The difficulty
here is to find a standard for the mental, since
it cannot be measured immediately, although there can
be no doubt that it is subject to quantitative
% --- p. 253 ---
\setcounter{footnote}{0}%
\opage{253}determinations, since there is, after all, a greater or lesser
strength of sensations and drives, higher and lower
degrees of attention, of clarity of consciousness,
etc. The only way in which these mental
motions can be measured is to see them in relation to the
bodily motions to which they are of necessity
bound. If we now keep to sensation, the
most elementary mental determination, then it has
as its necessary condition partly the outer impression, partly the
organic functions by which the impression is converted into
inner sensation. But these organic functions
(the psychophysical in the narrower sense) are unknown to us
in their inner essence, and yet it is they that
are the middle term between the physical and the psychical.
Inner psychophysics, which investigates the relation between
the mental and the organic (between the psychical
and the psychophysical), must therefore be built on
outer psychophysics, which immediately investigates the relation
between sensation and impression (between the psychical
and the physical). From the latter we shall then be able to infer
back to the former.

No sensation stands suddenly at its full
height, but each runs through a whole scale from
the point where it is still imperceptible, and that often in
so short a time that it seems to be present suddenly in
its whole strength. This successive growth makes it possible
to measure it. For when one observes the increments of sensation
with the successive increases in the strength or extent
of the impression, one will have given the elements
on whose relation the measure rests. It does not, indeed, seem
easy to make sure of the equality of the increments
of sensation. But one can, after all, judge the equality of two tones
or tone intervals, and besides, Fechner gives
three methods by which at least an exact
approximation can be reached. The first is ``the method of just
noticeable differences,'' which by experiment
% --- p. 254 ---
\setcounter{footnote}{0}%
\opage{254}determines the size of difference (in weight, etc.)
that is necessary for the difference to be capable of being
sensed at all; the two others rest on the calculus of probability
and seek by its help, through
repeated experiments, to eliminate the errors that must necessarily
creep in where, as here, subjective sensation
plays so important a role.

The simplest relation would be that the increment
in sensation was directly proportional to the increase
of the impression. But this is not the case. The physiologist
E. H. Weber's experiments have instead established that
the next simplest relation obtains, namely that \emph{relatively}
equal increments in impression correspond to equal
increments in sensation --- provided, of course, that the receptivity
is the same and the conditions on the whole
normal. This is the content of what Fechner calls
Weber's law, which is only the exact expression of
what is known to everyone from daily experience.
One hears one's neighbor's voice clearly when it is quiet,
but amid loud noise one cannot hear what one
says oneself, i.e.\ cannot notice the increase of impression produced by one's own words,
because this is vanishing in
relation to the impression already present beforehand.
When a room is brightly lit, one does not notice that
one more light is added, etc.

Every impression, or difference between impressions, must
already have reached a certain finite magnitude before it can
be noticed. The zero point of sensation therefore lies above
the zero point of the impression. The zero point of sensation, where
it just begins or vanishes, Fechner calls
the threshold. As long as the impression is not strong enough
to raise the sensation above this point, it is unconscious,
and the deeper it lies below the threshold, the
higher is the degree of unconsciousness. An increment in itself imperceptible
can here turn the scale and, by joining itself to the
% --- p. 255 ---
\setcounter{footnote}{0}%
\opage{255}impression already at hand, raise the sensation
above the threshold, i.e.\ let it become conscious.

A mathematical expression of Weber's law is had
in formulas that were already set up by Dan. Bernouilli at
the beginning of the last century to express the relation
between relative and absolute value (ex conditione personæ
--- ex re ipsa). The first of these formulas, which
Fechner calls the psychophysical fundamental formula,
is: dγ = k dβ/β, where γ denotes the sensation, β the impression and k a constant dependent on the units chosen for β and γ.
This formula thus states that the increment of sensation (dγ) depends on the relation between the increment of impression (dβ) and the impression present beforehand (β).
To this there is then joined, since the relation between
the increments of sensation and of impression is the
same as that between the increments of the logarithm and of the number,
the so-called formula of measurement, whose simplest form is: γ =
log β. This formula states that sensation rises
in logarithmic relation to the impression. And if it
leads to negative values, these indicate the depth of
unconsciousness. The positive values indicate the strength of consciousness.

Weber's law, however, does not prove to be valid
in all cases. If the formula of measurement had unlimited
validity, sensation would have to grow to
infinity with the increase of the impression. But because
of the constitution of our organism, in particular of the sense organs,
the law ceases to hold for very strong impressions.
On the other hand it seems that it must have unrestricted validity
in the domain of inner psychophysics, since here,
because of the immediate relation between the psychical
and the psychophysical, there are no such grounds
for its limitation as in the relation between sensation
(the psychical) and the impression (the physical).
But more than the general conviction of the
validity of Weber's law in inner psychophysics
% --- p. 256 ---
\setcounter{footnote}{0}%
\opage{256}cannot be reached. For we still know too little of the relation
between the psychical and the psychophysical for
us to obtain here a secure point of attachment for
the analysis. While in outer psychophysics we can always
determine the value of the β of the formula of measurement, which denotes
the magnitude of the impression (weight, etc.), in
inner psychophysics, where β denotes the organic
(psychophysical) motion that conditions conscious life,
it is impossible to determine the value of this magnitude; and
since γ is, after all, the unknown magnitude that is to be found, the
calculation becomes impossible\footnote[1]{Fechner had conceived the idea of his psychophysics with Herbart's mathematical psychology before his eyes, but in such a way that he wanted to avoid the one-sidedness in the starting point that hindered its being carried through. For he saw Herbart's defect in this, that Herbart did not secure for himself a reliable measure in the material foundation of mental phenomena. Therefore he starts from sensation, whereas Herbart started from representation. But this secure foundation forsakes him precisely when he passes from outer psychophysics to inner, and the latter is therefore subject to the same shortcoming as Herbart's mathematical psychology.}. On the other hand, inner psychophysics can,
when the general lawful foundation stands
firm, lead us to general theoretical results concerning
mental life and its relation to the bodily.

As regards, first, the general conception of
the essence of the soul, Fechner declares himself decidedly
against the monadological conception, which is represented
in particular by Lotze, and which sees in the soul a punctual
unity. In opposition to this Fechner designates his
own conception as the synechological one, and sees the soul
as the comprehending principle for the bodily plurality
of parts and the succession of bodily functions.
Only in its inner appearance for itself is the mind a
unity; in outer appearance it unfolds itself in a
plurality, just as what inwardly is a single sensation
appears outwardly as a composite and
% --- p. 257 ---
\setcounter{footnote}{0}%
\opage{257}successive nerve process. On the synechological conception,
then, the motions in the mental and in the
bodily are simultaneous, parallel expressions of the same fact.

The fact of the threshold is of importance also
for inner psychophysics, by furnishing a firm foundation
for the concept of the unconscious. So in the relation
between sleep and waking. Sleep is not a mental
nothing. The life-motion of the soul does not take place only
in the waking state. This is only the life of the soul above
the threshold, while sleep is mental life below the threshold
of consciousness; if we denote the threshold by 0, then
the positive values denote the height of waking consciousness and
the negative values the depth of unconsciousness in sleep.
We have here a scale that shows a continuous
transition in the subject between the waking and the
sleeping state. And just as in outer psychophysics
the impression could have a positive value at the point
where sensation is 0 (the threshold), so the
limiting point between sleep and waking occurs not at a
zero value but at a certain positive value of the underlying
bodily activity.

This relation can now find a further application.
Between the different organisms, the continuity of consciousness
by which a plurality of
phenomena is bound together in the same mental unity shows itself interrupted. And yet
all organisms are united by universal nature into
a single system. The psychophysical activity does not, then, extend
further \emph{between} the organisms in
the same way as in them. But a cessation of psychophysical
activity means nothing other than the greatest
possible depth below the threshold. What is below the threshold
separates, as unconsciousness, the conscious beings, while
at the same time it forms the foundation for them and maintains
the physical connection between them. The synechological
conception must here, then, finally lead us to the assumption
of a universal ensoulment of nature at different
% --- p. 258 ---
% print: Feehner (misprint: so printed)
\setcounter{footnote}{0}%
\opage{258}stages, since everything, as Thales said, is full of gods, ---
to ``the view of a God omnipresent in nature,
in whom all spirits live, move and have their being, as he in
them, with intermediate stages between him and us in the souls of the heavenly
bodies, so that the created spirits are borne
just as closely knit together in them as these themselves in the
divine spirit, and as the created spirits in turn
bear their senses, and these in turn their particular sensations,
within them.''

Here, then, we stand at the limiting point where Fechner,
from the limited sphere of exact knowledge, with the
boldest flight of thought, fearing no paradoxes, sets
out on the paths of his poetic-speculative world-view.
There is, as already mentioned, no merely accidental
or external connection between his exact and his
speculative view. Not only does he himself emphasize
repeatedly that it was his speculative ideas about
the relation between the mental and the bodily that first
led him to the thought of the possibility of a psychophysics;
not only, conversely, do the psychophysical investigations lead
him back to his speculative fundamental view;
--- but the method and the guiding principles
are the same: he seeks to remain on the way of induction and
analogy even after the domain has been left
where an exact experience is possible.

In general, says Fechner, one regards only
men and animals as ensouled, nothing else, at least
on our earth. The reason why one thus restricts
the sphere of mental life and lets life be a parasite
in the realm of death is chiefly that one believes one
must raise God so high above the world that it
dare not regard itself as his body. Thus one has
come to the unhappy, peaceless faith in a natureless
God and a spiritless nature. One believes that exact
experience forbids extending ensoulment further. But
the whole question of the soul is from the very outset a matter
% --- p. 259 ---
\setcounter{footnote}{0}%
\opage{259}of faith. An exact proof rests on experience and mathematics;
but only of one's own soul does one have direct experience,
and mathematics lacks every foundation for proving
the reality of another soul. Even in the soul of my father, of my brother,
I must believe, inasmuch as I infer it by way of analogy and
induction. We begin in
reality, like Cartesius, with a cogito ergo sum, and
an unconscious inference about the other souls of men and animals
becomes a dogmatic presupposition, which one can only
blindly take for a proven fact. But
if this is how matters stand, what can then prevent us
from extending the same procedure that we unconsciously
apply to men and animals further --- in particular to
the world of plants and to the heavenly globes, provided that
valid reasons for it appear?

It is objected that plants cannot have any
soul, since they have no nerves; but do polyps, then, have
nerves? --- that the plant is not an individual closed in itself;
but the individual is only something relative, since no
living being stands absolutely isolated over against the world around it.
What prevents there corresponding, to what of
the plant we can apprehend with our senses, an
inner unity that senses itself? There is indeed so
characteristic a difference between the plant and the animal that there
must also be a peculiar difference in the
inner essence; but on the other hand the plant shows
so great an analogy with the animal organism, and there is
so continuous a transition from the animal to the plant world,
that there is no reason to fix an opposition
between them as between the ensouled and the soulless.
There must be found in the plant a peculiar mental life, even
if it perhaps stands as far below that of the animal
as this stands below that of man. --- What holds
of plants holds also of the heavenly bodies. Why
should the earth not be ensouled? One must only
not begin by setting a sharp limit between
% --- p. 260 ---
\setcounter{footnote}{0}%
\opage{260}the organic and the inorganic. The inorganic is
something lower only insofar as it is, as in physics and
chemistry, thought of as torn loose from its living and natural
connection with the organic; the two belong together in
their persistence as in their origin. Men and
animals have grown together with the earth still more inwardly
than rocks and stones. Our bodies are parts and organs
of the earth, our souls parts of the earth's soul. However unlike
us the earth is, the unlikeness means nothing other than greater
height and fullness, not the lack of what a soul needs
for the expression of its essence. By a spiritual bond
all earthly souls are united into one, as a multitude of
smaller circles within the greater, more comprehensive one. One
has often spoken of the human spirit, of the Idea in the history
of humanity. What is understood here by the soul of the earth
is really nothing else, only that the human is not
isolated from the great connection to which it belongs,
and in particular that we do not conceive the Idea as conscious in the
particular, but unconscious in the whole. The lower spirits
can find their common unity only in a higher conscious
spirit, whose particular intuitive images and motives they are.

As the lower spirits are parts of the higher, so
these finally of the highest. The proof of the latter's
existence is therefore of the same nature as the proof of
the former's. What the historical forms of faith of humanity point
to, what practical need demands, that follows
also from the only theoretical proof that is possible here.
In investigating existence in its whole
extent we find a highest, general world-law, the principle of
all particular laws, namely this: whenever and wherever
the same circumstances return, the same consequences
return too; but under other circumstances
other consequences occur. This law shows that one
being pervades all places and times, everything spiritual
and bodily. In it the highest law is one with
the highest freedom. It now remains to know
% --- p. 261 ---
\setcounter{footnote}{0}%
\opage{261}what the essence of the Godhead is like, to know it as
spirit, as infinite consciousness. Here it holds,
as with finite spirits, that only one's own consciousness
is certain of itself. With regard to the
divine spirit, too, we have no other means of knowledge
than analogy and induction. From the consideration of
our spirit we rise to the thought of the absolute spirit;
from the consideration of our body as the reflection and
bearer of our spirit we rise to the intuition of the
whole world as the reflection and bearer of the infinite
spirit. We must indeed, as we have seen, think of
a chain of higher spirits as intermediate links; but the absolute
spirit is nevertheless immediately present even in the very lowest,
just as a small circle is immediately in a greater one,
even if within the latter it is enclosed by a third circle
that lies between them both in size. As we think,
feel and intuit, the infinite spirit thinks, feels and intuits
immediately by and in us. Here there is no externality,
but absolute inwardness. In general, to be sure,
one maintains the opposition between the infinite and
the finite spirit as between finite spirits. One
sets the infinite over against, above, beyond, outside the
finite, indeed perhaps even puts up an enormous partition wall,
so that they shall not come too near each other.
But they do not lie outside each other at all. The finite
is the content of the infinite. Therefore the infinite is
not incomprehensible either; it can be grasped in its content; only
it cannot be encompassed. Therefore one need not fear
a confusion of the divine and the
human; are you God because you are in God? is the
smallest moment the same as the whole? --- Not
only the finite spirits are in God, but also the material
world. The world is not something lower than God, but
something lower in God. An outer world exists only for
finite, not for infinite consciousness. God
sees, with some of his creatures as with parts and organs
% --- p. 262 ---
\setcounter{footnote}{0}%
\opage{262}of his body, other parts and organs set over against them,
and with his higher consciousness reaches out beyond them,
as our spirit does beyond all particular sense observations.
The forces of nature that seem to us unconscious are in themselves
elements in a will that reaches out beyond our consciousness
and works through the whole world. Therefore
not every single working of a force of nature need
signify a single act of will. The will is directed at the goal
of all the workings of force, and these unconsciously render
their contribution to the total end. Thus
the spinner thinks only of her yarn, while her foot mechanically-unconsciously
treads the wheel.

As the concept of the world thus in the last instance coincides
with the concept of God, the view becomes pantheism
in the full meaning of the word. The widest conception of
the concept of God is the best founded, and the narrower
conception only an abstraction. Yet ordinary
pantheism (the Hegelian) is wrong in regarding only
single beings as conscious; all consciousness must on the contrary
finally gather itself in a highest conscious being\footnote[1]{In the same way Fechner distinguishes his view from Ørsted's in ``Aanden i Naturen'' (which had recently appeared when Fechner wrote his ``Zendavesta''). He will not, like Ørsted, call the laws of nature thoughts of nature, since the laws are the object and content of thought, not themselves thoughts; to make laws and thoughts one can, in his opinion, easily lead one to forget to ask after the real thoughts and the infinite, all-embracing consciousness from which they proceed. Therefore --- he says further --- a conscious spirit of nature does not break through in Ørsted except in the mere \emph{name} of God.}.

Since the life of the world is the life of God, this cannot
be closed in itself, but develops and unfolds itself
with and through the development of the world. God's reason is not
completely developed and conscious prior to his working;
like the artist of genius, he first beholds his work
when it is completed (``And God saw everything that he had
made, and, behold, it was very good''). The divine primal reason
% --- p. 263 ---
\setcounter{footnote}{0}%
\opage{263}did indeed rule from eternity, but it raises
itself ever higher above the material foundation by ever
more incorporating this into its realm of spirit. God is not
imperfect because he develops in this way, for his
perfection does not consist in a finished conclusion,
but in an unlimited progress. Nor does
this inward relation to the world violate God's holiness.
For the sin of the finite will cannot be ascribed to God,
any more than a single thought or impulse that arises
in our spirit can make it evil. Just as
a man is called good or evil according to the whole direction,
the higher, comprehensive will, the ruling points of view
in his spirit, so we must measure God by the
total and eternal in his essence. God's eternal will leads
out beyond evil, so that this is finally dissolved in the
victorious, all-reconciling harmony. If it is asked why
and how evil exists at all, then, no answer can
be given other than that evil exists with metaphysical necessity, not independently of God's essence, to be sure,
but independently of his will, which precisely works constantly at
its reconciliation and sublation. God feels evil with
and in the sufferer; but this does not hinder his blessedness,
since in his absolute consciousness he is raised above everything particular.

This doctrine is nothing other than a real and serious
carrying through of Christianity's sayings that
God is spirit and must be worshipped in spirit and truth, ---
that in him we are, live and move, --- sayings that in general
stand as mere words. No one has presented
God and the highest commandment as purely and holily as
Christ. The eternal in Christianity, that by which
it went beyond everything earlier, its Idea, is ``that the highest
is posited as what unites all, that the most comprehensive
is posited as what is to be united, and that the best
is posited as the highest. The Idea of the union of all from the
point of view from which alone a unity of all is possible
was first brought by Christ, with consciousness, to the consciousness of the
% --- p. 264 ---
\setcounter{footnote}{0}%
\opage{264}earthly world, and by teaching and life
he made a living beginning of the spreading
and realization of this Idea.'' --- On the other hand, the doctrine presented
is certainly not Christian if one counts
as something essential to Christianity ``the belief in
the bite of the apple in Paradise with its mystical consequences, in the
irrevocable damnation of those not elect, in miracles
against the laws of nature, in God's remoteness from his
world, and in all the unedifying things with which theologians
in general furnish Christianity, indeed
out of which they build it up.'' The Idea of Christianity must
rather be grasped in such a way that it coincides with the very
idea of religion as the most universal, humane
point of union. It is Christianity because Christ first
brought this Idea to consciousness. But by his
ideal ethical commandments he indirectly founded spiritualism
and dualism, although he himself does not loose God
from the world. It is this spiritualism that has in turn
led to setting the infinite outside the finite.
When this one-sidedness of the present world-view
is sublated, the primal ground from which both heathenism
and Christianity have developed will come forth in
all its splendor, heathenism being explained by Christianity
and Christianity rejuvenated by heathenism.
It is a renewal of the old religion of nature, but
such that the religion of spirit, far from being thereby violated
in its true and eternal essence, is precisely enriched and developed,
indeed only now comes into its full right.

Fechner declares of his doctrine that it is materialistic
insofar as it sees everything as bodily, but just as
much idealistic, because it sees everything as spirit; it is dualistic,
inasmuch as it sees spirit and bodiliness as parallel
series, but just as much a doctrine of identity, inasmuch as, with Spinoza,
it sees both these series as forms of one and
the same fundamental being. Yet it differs from Spinozism
by maintaining an interaction between the two
% --- p. 265 ---
\setcounter{footnote}{0}%
\opage{265}parallel series, whereby not only a real relation of causality
and with it an exact functional relation becomes
possible, but also, what Spinoza denies, a teleological
connection, inasmuch as the ultimate ground of the
inseparability of both series is sought in the unity of the divine
consciousness. But about the fundamental being that lies behind
those two series no question can be asked if one does not
want to go behind the phenomenon, out into the transcendent. ---

By Fechner's own statement, a study of
modern philosophy did indeed reach deep into the course of his development,
but it also led him, as we have seen,
into a peculiar opposition to it. Starting from
empirical science and animated by a childlike joy
in all that is good and beautiful in the world, he could only be satisfied
by a world-view that did extend
and heighten the life of phenomena (into which he immersed himself so sympathetically
that he could even feel the pulse of mental life
where all others denied it) but
did not explain it by pointing back to a dark
background from which it was to spring. This immediate
devotion to the phenomenal is what gives
his view its attractive freshness; but it is
also what, philosophically speaking, is the essential charge
against it. The deeper grounding comes to be lacking.
He shuns it, regarding it as a darkness that
obscures the clear light. And yet a real grounding
itself brings with it a higher, ideal clarity, by
which thought feels itself purified and set free, even if
for immediate intuition nothing seems to have been gained.
Here Steno's words hold: pulchra sunt qvæ videntur,
pulchriora qvæ sciuntur. This ideal grounding
philosophy seeks partly in metaphysics, partly in the theory of knowledge;
such a clarification of the fundamental concepts and
the fundamental laws of spirit is necessary in order that what is gained through
experience and occupation with phenomena
may acquire philosophical significance. Fechner
% --- p. 266 ---
\setcounter{footnote}{0}%
\opage{266}demands that immediate intuition be fully satisfied
also in the domain of the development of concepts, which is a
flat contradiction, since abstraction here precisely goes
out beyond what is immediately intuitable. Philosophical thinking
is, as Plato already taught, a dying away from
outer life in order to win the Idea; but this dying away is what
Fechner fears. Immediate devotion to the
phenomenal leads him, since he skips those philosophical
preliminary investigations, to a fantastic religion of nature,
which in many of its features is of high poetic
beauty and is developed with amiable naivety, and which
is at the same time grounded in its own way with a logical acuteness
that makes it harder to refute than it
seems at first glance. Opponents of the hypothesis
of the plant soul would thus do well to read Fechner's
argument in the work ``Ueber die Seelenfrage.''
But on the one hand Fechner thereby comes to lay too great
weight on something that in the last instance is not decisive,
and on the other he is also led into a kind of scholasticism,
inasmuch as with an apparatus of the understanding made up of inferences by analogy
and induction he wants to prove something that
in the end nevertheless rests on a sympathetic satisfaction of the imagination.


The more Fechner seeks to lead everything back to consciousness
as the ultimate point of support, from which all experience
also proceeds, the nearer, it seems, must it also
lie to him to sink himself into the essence of consciousness and to
seek its fundamental laws and its significance in theoretical and
practical respects. The lack of such an analysis
shows itself in the opposition that in Fechner always remains
between spirit and the phenomena. He conceives
spirit (synechologically) as that which comprehends
and binds together the phenomena; as their highest
unity, spirit is the ultimate ground of everything. But the phenomena
stand unexplained by the essence of spirit; they bob
up in it as scattered points whose origin is not
% --- p. 267 ---
% print: Philosophis (misprint: for Philosophies)
\setcounter{footnote}{0}%
\opage{267}investigated. Therefore Fechner comes to waver between an
immediate identity and an equally immediate dualism:
everything is one substance --- but this can be seen from two entirely
different standpoints. In particular, atomic theory,
by which Fechner has earned so much
and so generally recognized merit for the philosophy of nature,
stands in an entirely accidental relation to his other views.
The atoms are to be only the fixed points
for intuition in the domain of exact experience. But
this is as yet no metaphysical explanation; these fixed
points, as long as they stand unresolved, become dark
points, which stand out all the more strongly in the
sea of light in which Fechner revels. His antipathy to
the monadological view and his whole pantheistic
tendency have prevented him from himself making his atomic theory
philosophically fruitful. In Lotze and Hartmann
we shall find attempts in this direction.

Fechner wants to avoid the ``transcendent'' philosophy's
search for the dark ``thing in and for itself'' behind
the phenomenon. And yet this ``Ding an sich'' haunts
his own view, especially in the hint of a fundamental being
to which the double (psychical and physical) series
of phenomena points back. On his way of considering things, spirit,
the highest unity and connection, would indeed have
to be also that which is in and for itself. But because
of the lack of metaphysical clarification, there emerges ever
anew for his thought, which is immediately governed by intuition,
the possibility that there could, after all, lie something
behind consciousness and phenomena. Hence his constantly
repeated polemic against modern philosophy, a polemic
which thus at bottom applies to his own view, not yet
penetrated by thought.

Yet Fechner himself indicates the point where the properly
philosophical foundation that we miss in him would
have to have its place. He states (Atomenlehre. 2te Ausg.
p. 149) that when one, by way of analogy and induction,
% --- p. 268 ---
\setcounter{footnote}{0}%
\opage{268}goes out beyond experience to find the conclusion
that thought seeks, the very conditions for winning
something ultimate and highest determine the form of the concepts found, while
experience must lead to determining their content.
What he here calls ``the condition for finding something ultimate
and highest'' are the determinations of essence without which
the Absolute would not be absolute (cf.\ Weisse). Without
a prior investigation of these determinations of essence
it is impossible to apply analogy and induction
in the domain of philosophy; for otherwise one would,
as is implied in Fechner's own words, be unable to know
whether it really was the Absolute, the ultimate and highest,
that one arrived at. ---

Fechner's positive and direct merit lies in
his thorough and acute analysis of the limiting concepts
between empirical science and philosophy (in particular
in his atomic theory and psychophysics). But indirectly, too,
his spirited and enthusiastic writings have their great
significance, by maintaining the value of the intuitive in opposition
to abstract speculation and by opposing
spiritualism and materialism. ---

3. \emph{Eduard von Hartmann} (b. 1842). --- For Fechner
everything was soul and consciousness; an eternal clarity and harmony
shone for him through the whole of existence. In
the author of ``Philosophie des Unbewussten,'' a new
work that has aroused an unusual amount of attention (the first edition appeared in 1869,
the fourth in 1872), we seem, then, to have a complete
antipode to Fechner. For Hartmann makes, as already
the title of his work indicates, the Unconscious the principle
of his world-view. And yet the two have not
a few important points of contact. Not only do both lay
experience and induction at the basis of speculation
and find in the fundamental psychological determinations the key
to the understanding of the essence of existence; but precisely by
both conceiving their apparently quite opposite
% --- p. 269 ---
\setcounter{footnote}{0}%
\opage{269}principles (consciousness --- the Unconscious) in an extreme and
abstract way, they meet in the end at the same nodal point.

Fechner came upon the Unconscious in his psychophysical investigations,
showing that when
the impression did not exceed a certain limit, the sensation did
not come to consciousness. But such an unconscious
sensation is not what Hartmann means when he
philosophizes about the Unconscious. An unconscious sensation
in Fechner's sense is in his opinion no
real sensation, is something purely negative. On the other hand,
the Unconscious is for him something positive, a real representation,
that is to say: an ideal content, and a real
will, that is to say: a striving to convert this
ideal content into reality. --- The motive for setting up the
Unconscious (the unity of unconscious representation and unconscious
will) as principle Hartmann has found in earlier philosophy.
This was essentially a working through of consciousness
and its content. This vineyard is now so
thoroughly worked that the time has come to dig down
into the underground world to find the treasures it
conceals. Many riddles that previously remained unsolved
will then perhaps be seen in a clearer light. Already the earlier
philosophers have on many points prepared
the development of the concept of the Unconscious. First and foremost
Leibnitz, who expressly distinguished between unconscious
and conscious representation (perception and apperception),
maintained that the soul could well think without being
conscious of it, and showed the significance of unconscious representations
for knowledge and action. Kant showed more closely
how the human spirit applies the forms of thought (the categories)
long before it becomes conscious that it
is on them that all experience and knowledge rest. Fichte
lets the universe proceed from the unconscious producing of the subject,
which precedes every conscious reflection, and
Schelling makes this producing an activity of the
Absolute, ``the eternally Unconscious, the eternal sun in the realm
% --- p. 270 ---
% print: Form, Den (misprint: a comma for a point)
\setcounter{footnote}{0}%
\opage{270}of spirits, which hides itself in its own undarkened light.''
Hegel's logical Idea, whose life pulses in everything, is the unconscious
thought that appears only in the finite subject
under the form of consciousness. It has its counterpart in
Schopenhauer's unconscious will, which for him is the only
metaphysical principle, while representation and thought
are only vanishing phenomena of the brain. --- But not
only philosophers, natural scientists too have been led by their
investigations to emphasize the concept of the Unconscious.
Thus the physiologist Carus declares that
the key to knowledge of the essence of conscious mental life
lies in the region of the Unconscious. And the famous Helmholtz
has found, as an essential factor in sense observation,
an unconscious logical inference, which in its
essence is akin to the conscious inference, insofar
as the latter too takes place in an immediate and spontaneous way.

Now when this principle, which thus seems to be pressing
forward, is to be developed and justified, this can only
be done by a union of the inductive-empirical and the
deductive-speculative method. The former would in and of itself
give a good foundation, but does not reach the ultimate
principles; the latter starts from the principles, but does not reach
factual reality. Therefore one must start
from experience and show how, by being interpreted
rightly, it meets speculation. Hartmann therefore makes
his motto: speculative results according to the inductive-scientific
method.

First it must be emphasized that the Unconscious is the unity
of representation and will. Hegel and Herbart maintain
only representation (thought, the Idea) and regard will
only as the mutual tension of representations or of the moments of the Idea.
On the other hand Schopenhauer maintains
will as the only real and regards representation
as illusion. But without representation there is no will;
for every willing must will something, have a content, and
this content can only be given as representation.
% --- p. 271 ---
\setcounter{footnote}{0}%
\opage{271}Conversely, representation needs will in order to be real.
Will is the very transition from possibility to reality,
in itself an empty form that finds its filling in
representation, which for its part is entirely indifferent,
a mundus intelligibilis, which only through will is carried over
into real existence. This real existence is what
empirical science establishes.

--- In general one takes will in a narrower sense,
of conscious self-determination. But whether self-determination
becomes conscious or not is a circumstance
that does not concern the essence of will as such,
but depends only on whether it passes through the brain
or not. From the lower nerve centers too (spinal cord
and ganglia) real manifestations of will can proceed,
which are unconscious insofar as they are not apprehended by the
central organ (the brain) on which the consciousness of the whole individual
rests. The activities of the lower nerve centers are
more than mere reflex actions. Experience establishes
the presence of a real affect and of a sustained
consistency, which show us out beyond the purely mechanical
sphere. Thus, e.g., the two halves of an
Australian species of ant fight each other with fury; and the
headless frog shows clear traces of calculation and
fear. Comparative anatomy, moreover, leads us
to regard the brain as a collection of ganglia that
are connected by nerve fibers, and thus in its essential constitution
is of a kind with the spinal cord and the ganglionic system.
Indeed, even in entirely nerveless animals like the polyps there
appear clear traces of affect and will. If, then,
in the higher animals we have both lower and higher
nerve centers, there is nothing to prevent something
from being conscious in a lower center without being
so in the higher one, which conditions the consciousness of the whole individual,
and it thus proves necessary to extend the concept
of will out beyond the domain of conscious life.

But by this we have so far reached only the relatively
% --- p. 272 ---
% print: de nødvendige Mellemled (misprint: de doubled over the line turn (de|de))
\setcounter{footnote}{0}%
\opage{272}unconscious. That there is an absolutely unconscious will is seen,
on the other hand, when one investigates how voluntary
movement takes place. For the conscious will to
lead to real movement, an impulse must occur at a
certain point in the brain; but which point this is,
consciousness itself does not know. The only possible connection
is this, that the conscious will, e.g.\ to move
the finger, calls forth an unconscious will to move the
nerve ending in the brain that brings about the movements of the finger;
and this unconscious will then in turn presupposes an
unconscious representation of the location of that nerve ending.
Since this representation and this will, which form
the necessary intermediate links between the conscious will and
the organic movement, do not come to consciousness in
the brain, where the process takes place, they cannot
come to consciousness in other centers either, and are thus absolutely
unconscious.

In instinct, too, an unconscious will manifests
itself. Instinct cannot be a consequence of bodily
organization; for it can be different with the
same organization, and conversely. Just as little does
it rest on a brain mechanism or mental mechanism implanted once for all;
for then the means could not vary while
the purpose remains the same. Nor can instinctive actions
proceed from conscious deliberation; not only do
they occur in the lowest animals, but they are performed with
immediate certainty and presuppose knowledge of the
future and in general of things the animal cannot
know, a kind of clairvoyance. Thus, e.g., the male larvae
of certain insects, at pupation, form for themselves a cavity
that is twice as large as they themselves, so that there may
be room for the long antennae of the developed insect.
Instinct can only be conceived as a determinate willing of
a means to an end unconsciously willed. Only on this
conception do we understand how instinct can be so inwardly
one with the whole essence of the individual. This is seen
% --- p. 273 ---
\setcounter{footnote}{0}%
\opage{273}especially in the drive to self-preservation and to the propagation of the species.
When one tears the larva's web apart, it repairs
it again and keeps on doing so until it
dies of exhaustion. A bird which was constantly robbed of
an egg went on mating and laying a new egg
until it was finally (at the 29th) found dead in
the nest. Such an absolute self-sacrifice is intelligible only
when a universal power drives the individual to summon all
its intellectual and physical forces in order to serve it.
But of this purpose itself the individual does not become conscious
as long as it is governed by instinct; otherwise its
egoism would set itself up against the rule of the universal.

Reflex movements are the instinctive actions of the subordinate nerve centers.
Likewise the natural
healing power and organic new formation can also be led back
to an unconscious activity. There can only be a
difference of degree between the cause of the bird's repairing
its destroyed nest, or the spider its torn
web, and the cause of the snail's house and the bird's plumage
being renewed again after an injury. We have in
all these phenomena only means to an end that
is unconsciously represented and willed, and the difference between them
rests only on this, that the different kinds of outer impulses
call forth different kinds of reactions.

Instinct proves its power in the human
spirit, too. So in the fear of death, in shame, in sympathy,
in the sexual drive, etc. A feeling in itself conscious
can be grounded in an unconscious will and be accompanied
by unconscious representations; on this rests what is unclear
and inexpressible in the essence of feelings. Indeed, the very representations
that call forth the feeling of pleasure or displeasure
can be unconscious to us, inasmuch as they come to consciousness only
in lower nerve centers, where they awaken a feeling
that is conducted to the brain without the effecting cause
reaching it. This shows itself clearly in hypochondria
and similar states. In general, in
% --- p. 274 ---
\setcounter{footnote}{0}%
\opage{274}our acting we know only the conscious motive, but not the
way in which the will reacts against the motive and thereby
passes over into real willing. Hence it comes that we
can learn to know our own character only through practical experience,
and are often surprised by the insights
we thereby gain into our own essence. The workshop of the will
lies in hiding, and only indirectly does our gaze reach
down there. When one supposes that one is immediately conscious of
the will, one confuses it either with its
motive or its consequence (the action) or with the feelings
that accompany it. Through this absolutely unconscious character
of the will proper it becomes intelligible how
thinkers like Spinoza, Hegel and Herbart have been able
to dissolve its essence into thought and representation, while
on the other hand the most dilettantish of the modern
philosophers, Schopenhauer, like ordinary
common sense, believed himself to be immediately conscious
of the will.

In sense observation an unconscious
psychical process takes place, which leads to the intuition of the object
as an outer one; without such a process we would not
get out beyond subjective sensation, since this
in itself contains no ground for going out
beyond the sphere of the subject. It is this process that, when
the subject becomes philosophically conscious of itself, proves to
rest on a relation of causality between it and the object.
The Unconscious likewise works in the association of ideas,
inasmuch as, in reaction against the conscious interest, it lets
precisely that representation emerge which is needed, but
which consciousness by express searching would perhaps precisely
prevent itself from finding. Everything truly a priori is
posited by the Unconscious and is an object for consciousness only as a result;
for in reflecting on the content it finds,
consciousness knows a posteriori
the unconsciously active a priori. All mysticism consists
% --- p. 275 ---
\setcounter{footnote}{0}%
\opage{275}in the filling of consciousness with a content through
this content's emerging involuntarily out of the Unconscious.

We further meet the working of the Unconscious wherever
a plurality works together toward a common goal, in one
and the same direction, and where both conscious agreement and
immediate mechanism are impossible. So in the formation
of language, which is thinkable only on the presupposition
of a mass instinct of the kind we meet in a swarm of bees,
an ant state or a termite state. So in history, where
individuals through their conscious, finite purposes
unconsciously work in the service of the whole great purpose of the world.
I unconsciously will something quite other than what
my consciousness believes it exclusively wills. What is
fate or providence other than the rule of the Unconscious, of historical
instinct, in the actions of men,
as long as their conscious understanding is not mature
enough to make the purposes of history its own?

The result, then, is that it is the Unconscious that
forms and sustains the organism, guides its movements
in a purposive way and forms the intermediate link between
it and the conscious will. In instinct it gives every
being what it needs for self-preservation and
for which conscious thinking is not sufficient.
Thus it gives man instinct for the understanding of
sensuous sensation, for the formation of language and state,
etc. It sustains the species by the drive to propagation
and maternal love, ennobles them by sexual and
qualitative selection, and through historical development leads
the human race forward toward the highest perfection.


But so far we have occupied ourselves only with organic
and spiritual phenomena. Here we cannot
remain standing. Modern natural science teaches us
(cf.\ Fechner) to conceive matter as a system of
atomic forces. No single atom can
have mass, any more than one can speak of the numerical magnitude of
% --- p. 276 ---
\setcounter{footnote}{0}%
\opage{276}a unit. If, then, we keep to atomic force as the
ultimate element to which the analysis of the real leads
us, it is clear that in the striving of atomic force there can
be found determinate connection and consistency only if
the content and goal of this striving are given beforehand in an ideal way
as a logical necessity, from which the outer factual
necessity in the working of atomic force under the
different circumstances springs. What, then, is
the striving of atomic force other than a will whose content
is the unconscious representation of what is striven for?
The concept of force thus passes of necessity
over into the concept of will. Force in its usual meaning
(of the elementary forces) first becomes intelligible
through will. Thus matter is dissolved into representation
and will, and the radical opposition between spirit
and matter is sublated. The difference between them rests
only on a higher or lower form of the revelation of the same
being, the eternally Unconscious. Matter, organic
formation or instinct, and consciousness are three different
forms of activity of the Unconscious, which makes up the world's
innermost essence.

The problem then presents itself how the origin
of consciousness from the Unconscious is possible. In the
Unconscious, representation and will are one; nothing is willed
without being represented, nothing represented without being willed. Consciousness,
on the other hand, contains the possibility of the emancipation
of thought from will. Unconscious mental activity
is the original one; only in and through normal brain function
does it become consciousness. The material influence
on the unconscious spirit determines the content of representation;
but this relation can in itself also be thought of as
unconscious. The form of consciousness arises through the will's
opposition to the representation produced by the organic function.
Without opposition, no consciousness.
But consciousness does not necessarily belong to the essence
of representation; it is only a special form of it.

% --- p. 277 ---
\setcounter{footnote}{0}%
\opage{277}The rise of consciousness is the culmination of the
activity of the Unconscious. But this activity is everywhere the same;
the Unconscious reacts against the material conditions and
effects, and these different reactions denote the
different stages in the development of the forms of existence. The
Unconscious seizes life wherever it can seize
it. Thus in generation it acts on the organic
germ and thereby develops mental life as by a kind of
process of crystallization. In analogy with this the
primal generation, the origin of organic life, must also be thought: when
all conditions were present, the first organic
cell was the necessary and lawful reaction of the Unconscious.
This factor Darwin overlooks in his brilliant theory,
which keeps solely to the conditions, but must always
presuppose the inner producing force. This is acknowledged
by men like Wallace and Nägeli. In particular,
Darwin confuses varying with transition into a
higher order. If qualitative selection really had the
enormous significance Darwin attributes to it, all higher
species would have to have greater vital force than the lower; but this
is far from the case. Moreover, the struggle for
life explains the physiological, but not the morphological
changes. Qualitative selection thus indicates one of the
mechanical means the Unconscious uses, but not
the sufficient explanation of all the forms toward which
it works its way.

Everything existing is individual. We have already found in
the atom the absolute single being. The living
organism is a world of individuals. For within the organic
individual we meet, as relative individuals,
first the cells, then the various organic systems
and the subordinate nerve centers. Out beyond these
rises, then, the general consciousness of the individual. The higher individual
always has the lower as its presupposition. The concept
of individuality is therefore relative. This has already been seen by
minds like Spinoza, Leibnitz and Goethe,
% --- p. 278 ---
\setcounter{footnote}{0}%
\opage{278}and it is confirmed in full measure by modern natural science.
--- Finally everything is absorbed in the absolute, all-embracing
individual, the Unconscious. Knowledge of this
all-unity of the Unconscious one gains only by freeing
oneself from the practical instinct that holds fast to the I
and comes to a point in pure self-consciousness, and by rising
to the conviction that consciousness does not belong
to the essence, but only to the phenomenon: ``Was an mir
\emph{Wesen} ist, bin ich nicht.'' The philosophy of the Unconscious
indicates the true reconciliation of monism and individualism
by referring the former to the metaphysical, the latter to
the physical-real domain. Plurality has only phenomenal,
not metaphysical significance. Space and
time are indeed not, as Schopenhauer thought, mere forms of
the subjective intuition of the brain; they are just as much forms
of outer reality. But they hold only for
the activity of the Unconscious, not for its essence in and
for itself. The activity, not the essence, is in space and time.

Since the unconscious functions are sufficient for
the explanation of the phenomena, it is a pure hypothesis
that the Absolute should become conscious of them, while
we do not. Of the Absolute's representing we can
form no representation. Against that hypothesis stands the fact
that consciousness, as far as we know it, arises only
in and through brain and ganglia. Moreover, the
Absolute's thinking is timeless, raised above all succession; it
thinks everything in one and at once, is absolutely immersed in
intuition, and therefore cannot turn reflectively
against itself. Yet the really decisive reason
against theism lies in the reality of evil, which on the presupposition
of theism becomes entirely inexplicable. Pessimism
is decidedly right in this, that this world cannot be a
work of a conscious thought, but can only have sprung
from a blind will.

Yet this is not to say that the Unconscious is in itself
blind. On the contrary, its wisdom surpasses all possible
% --- p. 279 ---
\setcounter{footnote}{0}%
\opage{279}consciousness (in this respect it is rather to be called superconscious
than unconscious), and it is this wisdom that determines
the content of creation and of the world-process. The philosophy of the
Unconscious is therefore in principle in agreement with
speculative theism, which ascribes to God a consciousness
absolutely exalted above the limits of finite, individual
consciousness; the Absolute is absolute individuality
with will and representation. But it nevertheless prefers,
for the reasons given, to designate the Absolute by
the negative expression: the Unconscious. This would indeed
have more right to the name ``God'' than Spinoza's
substance; but God is an expression of specifically religious origin
and should therefore be avoided as much as possible in
philosophy. The negative expression has its drawbacks, to be sure,
but it has a purifying effect, and when the anthropomorphic
notions have once lost their power, it
will surely already have been replaced by a positive one.

That which wills and represents is one substance. It
is not a blind man carrying a lame one, but a sound and
whole man, who admittedly cannot see with his legs or walk
with his eyes. This substance is absolute spirit.
But its original absolute rest has been sublated in a groundless
way by the will, blind in itself, which
has taken the initiative to leave potential being.
Because of its necessary unity with representation
(the Idea), it tears this along with it and carries it
over into being. Only so, in a blind, accidental way,
can the existence of the world be comprehended. For although Schopenhauer
is not right in his absolute pessimism and could only
in a sophistical way prove that this world is
the worst possible, it is nevertheless certain that its existence
is worse than its non-being. Even if a perfect
contentment were possible, life would only be worth
as much as non-being; but there is no
real contentment. What one extols as goods are
only privative goods, that is to say, one misses them when
% --- p. 280 ---
\setcounter{footnote}{0}%
\opage{280}one does not possess them, but their possession procures
no positive pleasure. Therefore humanity discovered
also, at the close of the ancient world, that it was an
illusion to expect happiness here in the world. Even for the
most favorably placed, pain outweighs enjoyment.
But from the first stage of illusion one passed
over to the second, renouncing, with Christianity,
all happiness in this world in order to find it again in
superabundant measure in a future, supernatural life.
The power of this illusion was broken by a new return
to earthly life and a new confidence in its validity.
But, taught by experience, one no longer expects
an immediate enjoyment of happiness; one hopes for
a positive state of bliss for humanity some time
in the future and strains all one's forces to
pave the way for it. This is the third and
last stage of illusion, in which humanity will perhaps still
remain for long ages. That this stage too is
under the power of illusion is seen from the fact that the greatest possible
progress can, after all, procure only increased means and
conditions, without a positive happiness thereby becoming more
possible than it has always been. Indeed, the pessimistic
consideration will precisely gain in force through this growing
disproportion. But this is also the intention of the Unconscious
with the whole historical process, and already with the rise
of consciousness. The reality of the world dawned through the
initiative of the blind will, which representation, because of
its unity with it, was compelled to follow. If it now
succeeds in emancipating representation from will, which is
in itself alogical, but through the consequence of that initiative
has become antilogical, then redemption will be possible.
Consciousness already arises through a victory over the
blind will, which comes into opposition with itself.
And when the illusions fall, one after another,
there comes finally the point where man makes the
end of the Unconscious the end of his consciousness and
% --- p. 281 ---
\setcounter{footnote}{0}%
\opage{281}devotes himself without regard to its service for the sake of the redemption of the world.
When this burning out of egoism,
which goes hand in hand with the carrying through of pessimism, has
penetrated so far that the greater part of all spiritual
beings in the universe give themselves up to it, then
actual willing is flung back again into the nothing from which
it came, and the world-process ceases. The Unconscious
is then what it has always been; it can learn nothing and
remember nothing, and is, after the conclusion of the world-process,
what it was before its beginning. ---

It will be of interest to add to this outline of the fundamental thoughts
of ``the philosophy of the Unconscious'' a brief
indication of Hartmann's relation to the earlier thinker
to whom he stands nearest in his whole world-view,
namely Schopenhauer. He corrects the latter in his
main thought by maintaining the necessary connection
between will and representation. Schopenhauer regarded
the essence of things as will, but considered representation
only an accidental semblance. But when he takes so
decisive a determination away from the concept of will, as
we alone know it, he loses the right to extend
it out beyond the psychological sphere. Such
a warrant Hartmann, on the other hand, can establish,
since he shows how the concept of representation too can
be extended to hold for the essence of things. Thereby
he also corrects Schopenhauer's conception of thought
as a mere phenomenon of the brain; thought has not only
psychological but also metaphysical significance. This
in turn brings with it another conception of historical
development. Schopenhauer regarded it as something
accidental and meaningless, like the figures in the clouds
or the frost flowers on the windowpane. Redemption from the
wretchedness of existence, which for him too is the ultimate
purpose, every individual is to carry out by asceticism and
the abolition of the will to life. Against this Hartmann objects
that this individual liberation accomplishes no more
% --- p. 282 ---
\setcounter{footnote}{0}%
\opage{282}than what would happen through the individual's natural death;
only when the species as such is led, through historical
development, to the abolition of the will to life is
a real redemption attained, such as, according to
Christianity too, the whole creation longs for and expects
from those who have the firstfruits of the Spirit. --- The
pessimistic fundamental thought itself Hartmann has thus retained,
and his developments in this direction bear the stamp of
having been carried out at second hand

In other parts of his doctrine Hartmann agrees with
Fechner and Lotze. So in his atomic theory and monadology.
Lotze is recalled in particular by his doctrine of the atoms
as monads and of their indwelling in the Absolute. With
Fechner he has not only, as already indicated, a
point of contact, insofar as the absolute consciousness that
in Fechner is one and all is, in relation to finite
consciousness, one with absolute unconsciousness. But
he also has him as a predecessor in leading the concept of force
back to that of will and yet setting himself in opposition
to Schopenhauer (Fechner Atomenlehre, 2te Ausg. p. 132
--- 135; cf.\ above p. 262). All the more remarkable,
then, is it that he himself does not mention at all the relation of his view
to these thinkers, whose writings (Fechner's
atomic theory and the 1st part of Lotze's Mikrokosmus) were available
13 years before Hartmann's book appeared.

What is attractive in Hartmann's work rests especially
on the clear and ingenious use of the material of natural science.
Only imperfectly have we in our presentation
been able to indicate the wealth of special analyses
on which he grounds his metaphysical principles. His
philosophical talent evidently lies in the domain of the philosophy of nature,
whereas in psychological subtlety he
cannot measure himself with Schopenhauer, and in metaphysics
ends in a fabulous doctrine of potencies that recalls
Schelling's. The very principle he wants to assert,
and after which he has named his work, evidently has
% --- p. 283 ---
\setcounter{footnote}{0}%
\opage{283}its great warrant, which he himself has indicated by
showing how it emerges at once in philosophy
and in natural science. But partly through the overvaluation
to which a new principle is always exposed, partly through
after-effects from Schopenhauer, there are in his application
of this principle decisive shortcomings, which criticism
must point out.

Hartmann himself seems to forget his critique of Schopenhauer.
At the really decisive point in his
world-view, the problem of evil, in whose
treatment he, as we have seen, essentially follows
Schopenhauer, he also retains the latter's opposition
between will and representation. Like
Schopenhauer, he finds the principle of the world's development in
a blind and groundless will to life, which by taking
the initiative has dragged representation along into the movement.
Here, then, the presupposition lies at the basis that
the will isolated from representation is self-subsistent and
can move without impulse from representation. But
this conflicts both with Hartmann's own doctrine and
with the psychological observation that every drive
is determined by a representation, however
dim or unconscious this may be. On closer inspection even that ``groundless'' will also presupposes a
representation as its motive; for it wills itself, is,
in Hartmann's expression, ``velle volens,'' and thus has
its own essence as its unconscious content. If this germ of thought
is developed, the absurdity will vanish that
blind chance should be the principle of reality,
and the darkness will be dispelled into which the obstinate pessimism
lets the origin of everything go back. We do not, then, conclude
with Schopenhauer and Hartmann that, because of
the reality and overwhelming power of evil and of pain,
thought cannot be the ultimate principle of everything;
but we conclude that, since thought is necessarily the original,
the ultimate presupposition, evil can only
% --- p. 284 ---
\setcounter{footnote}{0}%
\opage{284}be a metaphysical necessity (cf.\ Weisse and Fechner),
on the foundation of which the development of life goes forward
toward its goal.

We have seen that Hartmann does not conceive the Unconscious
purely negatively, but declares himself in agreement in principle
with speculative theism, which maintains the consciousness
of the Godhead, provided only that one conceives this consciousness
as entirely raised above the limits of the finite I.
But this declaration (which, incidentally, is first found in the work's
third edition) does not seem to agree with the fact that the Unconscious
nevertheless stands in a purely negative relation to the life of nature
and of history. This life has significance only
insofar as through it the Unconscious frees itself
from itself. When the development is completed, it has gained
nothing. It thereby also comes into a negative relation
to consciousness and stands lower than it, instead
of being the ``superconscious'' that reaches out beyond
consciousness and personality and takes them up into itself according to
the inner kernel of their essence. For at the culmination of its development
conscious life possesses a wealth and
a fullness that stand in glaring opposition to the pure
nothing that the Unconscious in the last instance is and wills. Hartmann
has in his conception of the Absolute as the Unconscious
not shown dialectical caution enough\footnote[1]{Cf., by contrast, an author with whose doctrine Hartmann's has much in common, K. G. \emph{Carus}: Psyche. (Pforzheim 1837). p. 402: ``The highest divine consciousness, the consciousness of God's spirit in and for itself, is to be thought by us only as something so immeasurable, infinite and all-embracing that for so entirely conditioned a consciousness, bound to finitude, as the human, it will always in the end coincide completely with the very mystery of the Unconscious; but conversely, precisely for that reason, what we have called the divinity of the Unconscious lies only in the immeasurability and incomprehensibility of the highest divine consciousness.'' The content of this last sentence Hartmann has not allowed to come into its own; therefore his ``Unconscious'' is in reality not the Absolute either.}. There remains
always a sharp opposition between essence
% --- p. 285 ---
\setcounter{footnote}{0}%
\opage{285}and activity, essence and consciousness. Activity
is to belong only to the phenomenon; but since the forms of the phenomenon
do not concern the essence at all, the latter remains
entirely untouched by its own activity. Thus
consciousness, which is conceived as mere form, does not concern
the real kernel either: ``was an mir Wesen ist, bin
ich nicht.'' --- We end here in a dualism that through
Schopenhauer points us back to Kant and his distinction
between the phenomenon and the thing in and for itself.

It is therefore no wonder that Hartmann's attempt
to explain the origin of consciousness fails.
Through the figurative expressions, which here are certainly
unavoidable (``the will's opposition to the organically
produced representation''), what was precisely to be proved shines through
as a presupposition. For what is
one to think by the organically produced representation
that calls forth the will's opposition, other than the
first sensation that crosses the threshold of consciousness?
Before it has crossed this threshold, it has,
after all, no reason to awaken the will's counter-striving; and
when it has crossed it, this counter-striving can
indeed develop and potentiate it, but is unnecessary for
producing it.

Conscious life has, according to Hartmann, only the significance
of successively emancipating representation from the rule
of will. Humanity is finally to end in
the millennial kingdom of pessimism, which forms the transition
to complete Nirvana. But if one seeks the ultimate
cause of Hartmann's pessimism, it lies,
as with Schopenhauer, in the separation between individual
happiness and the ideal task. If happiness is
absolutely to be found outside the task that is set the individual
and the species, then the pessimistic consideration is surely
right. But it separates what cannot
and should not be separated. The warrant of the individual and of the species
rests only on their peculiar purpose and
% --- p. 286 ---
% print: at godtgjøre (misprint: at doubled over the line turn (at|at))
\setcounter{footnote}{0}%
\opage{286}task, through whose fulfillment their own essence is developed,
and only in this very development can happiness and satisfaction
be sought. What is worked out and developed along
this way cannot, then, be meaningless
and negative either; it must, however one may more precisely think
of the matter, preserve its reality, because it preserves and develops
the very essential in reality. The eudaemonological
and the ethical stand in an inner connection,
and the significance of pessimism lies in indirectly
establishing the necessity of this unity. ---

4. \emph{Rudolph Hermann Lotze} (b. 1817). ---
Like Fechner, Lotze too passed from investigations in natural science
to philosophical ones. But with him, too, these two
directions of study do not stand in an external and
accidental relation to each other; in him they join together
still more inwardly than in Fechner. It was
a lively interest in art and poetry that first led
him to philosophy. He felt himself drawn to the
great circle of thought to which Fichte, Schelling and Hegel had
opened the entrance, but he appropriated it more as
a characteristic tendency in general culture
than as a finished edifice of doctrine. A more definite influence
on his views was exercised by Weisse, under whom he
studied in Leipzig, and whom he still remembers in his latest
writings with gratitude and piety. In his
``Streitschriften'' (1857), where he gives this information
about the course of his development, he declares that he owes Weisse
not merely many influences in general,
but in particular such instruction concerning a narrower
circle of thoughts that he has since had no reason
to give it up. This ``circle of thoughts'' evidently concerns
the general metaphysical and ethical fundamental view,
the germs of which we can find in Weisse's
doctrine, but in such a way that Lotze's clear and independent
mind and his penetrating insight in natural science and the history of culture
have developed and transformed these germs
% --- p. 287 ---
\setcounter{footnote}{0}%
\opage{287}in a distinctive way. In particular, the study of medicine,
which he had first chosen as his life's work,
early drew his attention to the necessary
real connection without which even the highest and
most glorious Idea is only an empty fancy. In these two extremes,
the ideal and the system of factual conditions, his
thought took its stand, in order to seek out from them the unity
which speculation had believed it could seize
all at once by a bold grasp. At the age of 22
he was able to appear as a lecturer in both the medical
and the philosophical faculty (later he became professor
of philosophy in Göttingen). His writings go in both
directions. In the 1840s he published on the one side
a Metaphysik and a Logik, on the other side a
Pathologie and a Physiologie, besides some aesthetic treatises.
His medical writings attracted great attention,
in particular through his endeavor to seek, by an exact
mechanical route, an explanation of the phenomena
which had hitherto been derived in an obscure and mystical way
from a so-called vital force. He tries to show that
the organic functions differ from the inorganic
only by the distinctive arrangement of the elementary forces,
but not by a difference in principle in the essence of the
acting forces. His penetrating acuteness
let him see what was unprovable in much that had
hitherto stood as settled, and many therefore regarded
him as a skeptic. In philosophical respects
these works acquired significance in particular because Lotze through
them further developed his conviction of the necessity
of assuming a multiplicity of elements
in mechanical interaction as the real foundation
of reality. The way in which he expressed
this conviction led many to take him for a Herbartian;
but in the polemical work just cited he derives
it expressly from his studies in natural science. On the other hand
he refers to Leibnitz as the philosophical
% --- p. 288 ---
\setcounter{footnote}{0}%
\opage{288}model that stood nearest to him with regard to the
further way of regarding the nature of the real foundation.
--- Of his general philosophical views he declares
in the same place that they can most nearly be compared
with the fundamental conviction of the elder Fichte, that
only the Idea of the good and its content is the sufficient
ground of all being and becoming, that the world of
valuable purposes (die Welt der Werthe) contains the key
to the world of forms. Only one must not, like Fichte,
restrict what is valuable to what is strictly moral;
the tranquil blessedness of the beautiful, the holiness of the restful and
passionless mood, and the inner
consistency of the true, with all the peace and harmony it calls forth,
also belong to the ideal world. Nor must one
suppose that one can demand or point out the various parts of nature piecemeal
as the foundation of corresponding ideal
moments; the Idea does not live thus from hand to mouth,
but nature is like a great household that follows
its own distinctive laws, independent of the particular purposes,
and only as a totality shows itself bound
to the Idea.

Since Lotze's physiological writings maintained the mechanical
conception of nature, the materialists (who of course
did not read his philosophical writings) claimed him
for themselves. His ``Medicinische Psychologie
oder Physiologie der Seele'' (1852) put an end to this. Here too, it is true, he
does not go back to the ultimate presuppositions,
but he nevertheless shows that the concept of matter is not a presupposition
at which thought can stop, and that
mechanism, however necessary its concept is for
natural science, still acquires its significance only by
being the form and foundation of an inner fullness of life. Finally
Lotze's chief work ``Mikrokosmus. Ideen zur Naturgeschichte
und Geschichte der Menschheit. Versuch einer
Anthropologie'' (3 vols. 1856—1864) gives a survey of
his whole world-view. When in a philosophical
% --- p. 289 ---
\setcounter{footnote}{0}%
\opage{289}psychology he wished to set forth his doctrine of human
life, it became clear to him that the deepest forces of our nature
can be observed distinctly only where they appear
in the great context of human development. Therefore
the philosophy of history became the necessary completion
of psychology; the consideration of individual
life and of the cultural history of the human race unite
in anthropology. Supported on this rich foundation,
he then seeks at the end of the work to answer the decisive
problems of the connection of things, of God
and the world. --- In his latest writing too, Geschichte der
Aesthetik in Deutschland'' (1868), there are many contributions
to his world-view.

As an author Lotze occupies an eminent place.
A clear and precise grasp of the facts is united
in him with a poetic sympathy for the real
and with an energetic conviction of the validity of the ideal.
The exact, the poetic and the philosophical
elements fuse in a distinctive way in his
exposition. And in this his exposition is a faithful image of
his view. The ideal is for him not the Idea developing itself in abstract
forms, as in speculation, nor
the external mechanical connection, as in materialism;
but those forms as well as this connection
are only the outer shells of life's eternal kernel.
And yet they have decisive significance. One cannot
grasp the kernel immediately; the attempts to do so lead only
to grasping the shells in place of the kernel.

``Truth,'' says Lotze at the end of his Metaphysik,
``is not the first thing, but is conditioned by the fact that the
realm of the good produces it as its necessary presupposition.''
Speculation overlooked this. For science
easily takes its own domain, thinking knowledge,
for the whole of spiritual life, or at any rate for its summit.
But thought always presupposes a real content.
The essence of thought is in itself a critique of the given. By the
% --- p. 290 ---
\setcounter{footnote}{0}%
\opage{290}critical, comparative activity by which knowledge
of the law arises, the human
spirit distinguishes itself from the lower forms of soul, but not by
being absolutely absorbed in pure thought. The original spiritual
facts, such as the sensation of colors and
tones, the intuition of space, the feeling of pleasure and displeasure, and
all aesthetic and ethical valuation of what is known, are
not merely preliminary or failed attempts to think,
but make up the matter which thought cannot itself produce,
but only compare and work up. In both
ancient and modern philosophy Lotze misses
the recognition of this. The Greeks founded logic,
a magnificent work that may well be placed beside
modern natural science. But Greek thought
was inclined to confuse the logical and the
metaphysical, truths that hold and things that are.
From mythology, moreover, thinking inherited the belief in
symmetrical and rhythmical forms in the course of the world, and
this belief was stronger than the conviction that there
must after all be definite and necessary laws for the
working together of the particular elements within the whole course of the world.
Therefore the Greeks confused ideal interpretation
and causal explanation, and the latter did not come
into its own. Even in modern philosophy antiquity's
overvaluation of the Logos is at work. Realism teaches us only
what is, idealism only what ought to be; neither
in the real nor in the ideal is that brought forward by which
they are more than all reason. In the last instance
the essence of every thing consists in the Idea it is to realize,
and the true statement of a thing's concept can
therefore be won only by a dialectical development of this
Idea in its grounding by the highest Ideas. But one
cannot go back to the first principles at every detail;
they are too weighty to be applied
immediately.

% --- p. 291 ---
\setcounter{footnote}{0}%
\opage{291}There are, then, two main objections that Lotze raises
against earlier philosophy. It has not laid due
weight on the mechanical connection by which alone
the realization of the Idea is possible, and it has not
conceived the Idea in its whole fullness, as the absolute purpose,
the highest good. The highest tendency of knowledge
is teleological and includes within itself a tendency to genetic
explanation and a tendency to ideal interpretation,
which often wrongly come forward in conflict with each other.
--- From Lotze's own youthful metaphysical work we shall
here bring forward only one main thought, to which he
constantly returns in his later investigations. ---
After having investigated the forms of being (ground, cause,
purpose) he passes on to the doctrine of the phenomenon.
The concept of the phenomenon demands not merely a being that
underlies it, but also a being to
which it appears. Not merely what goes
on \emph{between} beings, but also what goes
on \emph{within} them, is a true and real occurrence.
That the reality of things consists in their revealing themselves
in the phenomenon means, then, that they are for one another.
Reality is not merely the external object; our
apprehension of it also belongs to the real, just as
sound first becomes real when the vibrations of the air, by
acting on the auditory nerves, are transformed into sensations. The
spiritual is not a superfluous addition, which at most
would have the task of mirroring a reality fully
finished in itself. The spiritual world is on the contrary
the most essential constituent of the universe; the way in which
reality is apprehended by the spirit is the most essential
part of all that happens, without which the world would not
be finished or complete. --- So too Lotze says
later in his history of aesthetics: ``There is no beauty
outside the feeling of the spirit that enjoys and admires
it; but the connection of things is so ordered that
% --- p. 292 ---
\setcounter{footnote}{0}%
\opage{292}it can awaken in the spirit forms of movement through
which that enjoyment falls to its share. --- That reality
on the whole is disposed so as to make such a coincidence
possible, --- that the connection of the real world fits
the receptivity of the spirit, --- that the connections of things
occur and can occur in forms whose impressions awaken the functions of the soul
to harmonious activity, --- this, that the world
and the spirit are thus for each other, is the great
fact that we enjoy in the feeling of beauty.''
--- This harmony between the forms of existence has its
ultimate ground in the highest purpose, and therefore metaphysics
too has its beginning not in itself but in
ethics. ---

Having described Lotze's general conception
of the essence and significance of philosophical knowledge,
we shall now with him immerse ourselves in ``Mikrokosmus''; yet
in such a way that besides the work of this name we also use
the more extensive investigations in ``Medicinische Psychologie''.
Through the problems of the philosophy of nature, of psychology
and of history he will then at last lead us
to the concluding metaphysical and ethical-religious view.


The original conception of nature saw all its
phenomena animated by divine beings, from
which all life and all activity sprang. But even after
this childlike way of regarding things had been replaced by
the scientific search for the causes of things, its
traces were preserved in the notion of plastically acting
drives as the ground of life in nature. It was
in particular the organic phenomena to which this idealistic
conception attached itself. One did not
properly account to oneself for what use such an unconsciously acting drive
was to be. For it would contain the explanation
of the purposiveness in the organic only if
it acted in accordance with a reason that weighed the circumstances;
but this is precisely excluded by its being
% --- p. 293 ---
\setcounter{footnote}{0}%
\opage{293}conceived as unconscious\footnote[1]{This remark would not, as it might perhaps seem, strike Hartmann, since his use of ``the Unconscious'', at least in principle, has precisely the mechanical connection as its presupposition. On the other hand it does strike such theories as the younger Fichte's.}. Science cannot seek
the explanation of organic life anywhere else than in
such law-governed activity as we meet throughout nature.
The vivifying power in which speculative
philosophy of nature believed is then here dissolved into a multitude of
elementary forces. The mechanical conception of nature now becomes
dominant. Its fundamental thought is that the capacity to
produce a definite effect is never contained, full and finished,
in a single atom or in the nature of a single body,
but that, just as the necessity of an activity
in general follows only from the mutual relation
between several elements, so too each element's
mode of action and the magnitude of its influence depend on
the nature of the other elements. The living is distinguished
from the lifeless not by a higher, peculiar force
by which it would be withdrawn from the general
law of mechanical connection. But the inner
arrangement of the living has such a form that the natural
forces of the elements, under the influence of the external conditions, must develop
a coherent series of effects according to
the same general laws by which elsewhere, too,
one state everywhere tends to follow upon another.
Thus the body is no longer a garment which the soul
weaves for itself, but a mechanical system of external conditions
for the life of the soul, --- the piece of the surrounding world
that stands nearest to the soul.

This mechanical conception often gains such
power over thought that thought cannot free itself from
it in domains where its validity is no longer
absolute. In regarding nature we grow accustomed to such
a degree to the mediate effects, which we explain
% --- p. 294 ---
\setcounter{footnote}{0}%
\opage{294}by the composition of particular contributions, grow accustomed to tracing
significant differences back to insignificant
changes in the magnitude and mode of connection of homogeneous elements,
that in the end we lose the understanding of everything
immediate and involuntarily presuppose an intricate machinery
for the origin and existence of everything. All our knowledge
then seems to be a cognitio circa rem, a knowledge
of the conditions under which a phenomenon comes
before us and is altered in its interaction with other
phenomena. And yet at certain points we also possess
a cognitio rei, a knowledge of the object's own
essential nature. Such a knowledge is possible where we can not
merely observe an object in its external relations,
but have it with such immediate intuitiveness that we
can, both in feeling and in representation, make ourselves one
with it, put ourselves into its state and sense
how such an existence feels in accordance with its distinctive
essence. Such a positive and immediate
intuition we possess only of what is living and active;
only this do we understand and penetrate. Matter as
such is a dark kernel for us.

The fact that we can become conscious of ourselves
presupposes such a unity in our being that
it becomes necessary to assume a distinctive mental
principle, which enters into interaction with the bodily
mechanism. Otherwise one would end in the obscure
and self-contradictory assumption of an activity without any
acting principle. But the soul is not active without
the influence and cooperation of bodily functions. In
this interaction between soul and body there lies
no greater riddle than in any other causal relation whatever.
Only one is often inclined to conceive
the principle of causality one-sidedly, dwelling solely on
the dependence of the effect on the cause. The proposition can
also be stated in the form that every cause has
its effect, which means that every element of
% --- p. 295 ---
\setcounter{footnote}{0}%
\opage{295}reality intervenes in its own distinctive way
in the whole connection, in accordance with the laws holding
in it. Our knowledge of nature determines
nothing about the inner bond between the acting elements.
With such a cognitio circa rem the standpoint is
properly occasionalistic: we seek only occasions and
conditions. This standpoint we must also hold fast
when we wish to set forth the first fundamental features of psychology
and to investigate the relation between soul and body.
Only when these fundamental features have been won shall we be able
to change our standpoint and overcome the opposition
that has shown itself to us between external mechanism
and inner life, between cognitio circa rem and cognitio
rei.

We first characterize the essence of the soul more closely, in order
then to see it in its relation to the body. --- The unity of
the spirit manifests itself in this, that it not merely
collects and comprehends its various states into
an inner whole, but also, by a critical and comparative
activity, strives to interpret the particular and
assign it its definite place within this whole,
which thereby becomes an ideal world. In this world the
spirit then finds a reflection of its own unity. Even in purely
sensuous apprehension this peculiarity shows itself.
What we apprehend by the senses we take up not merely as
an indifferent content, nor merely as that which causes
pleasure or pain, but we ascribe to it a worth of its own,
by which it fills its place in the significant connection
of phenomena. Sense perception itself already rises
above purely mechanical necessity and seeks
to subordinate its activity to an order of life that is not
immediately given in nature.

The unity that thus runs through mental life is not explained in the views that either, like materialism,
derive the mental as a product of the
bodily, or at any rate regard the unity of consciousness as
% --- p. 296 ---
\setcounter{footnote}{0}%
\opage{296}the result of the interaction of bodily and mental forces.
Insofar as one here appeals to the analogy
of the parallelogram of forces, it in fact points in
the opposite direction. For the diagonal motion arises
only when two forces act on one and the same
point; it is then this point that moves along the
diagonal. Here, therefore, a unity is expressly presupposed
from which the new force can proceed. --- But, it will be
asked, where do we find this unity in the soul? we know
it, after all, only as knowing, feeling, willing, etc.
Lotze answers to this that the being of every thing consists in
its distinctive content; outside or behind this
content it has no being. If one were to demand information
about what makes this content something that is,
that would be to ask how being comes to be.
Just as it is the essence of water to be liquid under some
circumstances, vaporous or solid under others,
so too it is the essence of the soul under certain conditions
to appear as representing, under others as
feeling, under others again as willing, and its distinctive
unity does not suffer thereby.

Modern psychology has often believed that it must maintain
the unity of the soul by tracing its various functions
back to a single one. Thus Hegel regarded all
spiritual phenomena as different, lower and higher
forms of conceptual knowledge. Herbart saw
the activity of the soul as representing; feeling and will were
for him only results of special relations of opposition
between representations. Lotze grants that representations
are for the most part the points of attachment for the feelings,
and that the drives in turn develop out of these. But
it does not follow from this, he adds, that the representations are also
the whole and full cause of the feelings and drives
they call forth. In the representation itself there lies no ground
for the distinctive function that appears in the feeling
of pleasure and displeasure. The capacity to produce this
% --- p. 297 ---
\setcounter{footnote}{0}%
\opage{297}must be sought in the soul itself and comes into activity when
the soul, through the course of the life of representation, is roused to express
itself in this way. In the same way the
utterance of the will is not produced by the feeling, but released on occasion of it
by the soul itself. The three fundamental faculties are thus graded
dispositions which mutually rouse one another to activity.
Whereas in a merely mechanical unity the result of
the interaction between two elements is completely
determined by the general laws of the course of nature and the given
circumstances, in spiritual life the nature of the soul
enters in as an acting factor that helps to condition
the result of the movement.

If the soul is such an original principle, then
it follows of itself that, however thoroughly it
is conditioned by the bodily organism, it nevertheless cannot
be immediately receptive to influences from
it, but must constantly transform these in a distinctive
way. The bodily influences do not bring the soul
finished sensations or representations, but only signals,
which it is to translate into its own language, just as
conversely the inner states of the soul are not as such
transferred into the bodily organs, but become the occasion
for these to come into function. The will knows nothing
of the means by which the arm is set in motion, and
as regards sensation, modern science has
clearly demonstrated the dissimilarity between it and the
external influence. This holds also with regard to
our apprehension of things as each occupying its own place
in space. How does this intuition of space arise?
One thing is the external extension of objects, another our
inner image of this extension. This inner image
is not itself extended. We reproduce in an ideal
way the spatial relation in our intuition; the extensive
is converted into an intensive. The impressions received spread out
in the soul into a world of space, not because there
is a real spatial relation within consciousness,
% --- p. 298 ---
\setcounter{footnote}{0}%
\opage{298}but because in the mental movements called forth by the impressions
relations are preserved that correspond to the relations
between the external objects. Every influence acquires,
on account of the point in the nervous system where it takes
place, a distinctive coloring, a characteristic secondary determination,
which can be called its local sign, because it
is this that leads the soul to refer the object
from which the impression derives to a certain place in space.
The local signs cannot be of a passive nature, but
must be thought of as a kind of reflex movement; at the same time
they must, in order that a clear, geometrical representation of space may become
possible, make up a system. But the local signs do not produce
the intuition of space itself; they serve only to specify
it. They lead the soul, which in accordance with its own nature
strives to unfold its intensive content in space,
to apply its representation of space in accordance
with the nature and mutual relations of the objects.
Originally unconscious and mechanical, spatial localization can
be developed to higher degrees by practice. ---
This doctrine of Lotze's concerning the local signs, which springs in a natural
way from his fundamental psychological view,
has won much recognition, among natural scientists too
% print: Hemholtz (misprint)
(e.g. Helmholtz).

Two questions still remain; the one
concerns the seat of the soul, the other the organs of mental life. ---
As regards the first, Lotze starts from the proposition
that every thing is where it acts. The seat of the soul
must be sought in the brain, which is the part of the body with
which it stands in immediate connection. Many, it is true,
speak of the omnipresence of the soul; but they cannot
then explain of what use the nervous system is. A
conscious sensation first comes about when the impression
is conducted from the sense organ to the brain. An
immaterial substance which lacks all extension cannot, it is true,
fill any space; but nothing prevents
its having a definite place in space, from which it
% --- p. 299 ---
\setcounter{footnote}{0}%
\opage{299}acts as from a center of force, and to which all influences
stemming from external nature must propagate themselves
in order to exist for it at all. --- The notion
of an organ of the soul Lotze considers thoughtless. The ground
of the distinctive quality of representation, feeling and drive
must be sought in the soul's own nature, and no bodily
activity is able to impart it, although such activity
is the necessary condition for the actual occurrence of those spiritual movements.
The bodily organs
work in the service of the higher activities of the spirit by
delivering and preparing the material on which
the soul is to exercise its powers; but in this too consists
their whole significance. The significance of reflex movements
for the intuition of space and for movement compels us to assume
central organs for spatial intuition and for the regulation
of movements, whereas it is meaningless
to assume central organs for the activity of the understanding,
for moral and aesthetic feeling, etc., since these
activities consist in an ideal relation between given
elements, a relation which cannot itself be explained mechanically.
The soul's dependence on
the bodily does not, however, become any less thereby. The inwardness and height
its activity is to reach are conditioned by the sense organs
and the central organs supplying it with impressions from without accurately and with sufficient
intensity. The richer and
nobler the means, the nobler the activity. The highest
spiritual activity is like the flower of the plant; but no one
will demand that the flower should have its own special
root, different from the one that bears the whole plant.

If the methodical occasionalism with whose help
this conception of the relation between soul and body
has been won were to be a real theory, it would
lead to a dualism between two worlds connected only by external, mechanical
interaction, the mental
and the bodily, each with its own distinctive laws.
The task now is to get beyond this dualism by
% --- p. 300 ---
\setcounter{footnote}{0}%
\opage{300}abolishing the methodological presuppositions that lay
at the basis of what preceded.

All mechanical connection presupposes a plurality
of elements between which it takes place. Therefore
the ultimate presupposition of modern physics is an infinite
multitude of atoms which cannot be dissolved by the
forces known to us. But metaphysical thinking cannot
stop at this presupposition. It
cannot, like physics, whose only task is to
investigate the mechanism itself, stop at the concept
of matter as a dark kernel that moves
in a bright net of relations. The notion of a being
that does not exist for itself is self-contradictory.
Such a being would then exist only for the consciousness
that apprehends it, and the one part of the world would
thus be the blind and lifeless means for the other's
end. Only to living beings does a real
being belong; all other forms of existence must be explained
from spiritual life. We must put a cognitio rei into
each of the elements which the cognitio circa rem compels
us to assume. Each of the atoms\footnote[1]{When Fechner's ``Atomenlehre'' appeared, Lotze could declare (Göttinger gelehrte Anzeiger 1855 p. 1093 ff.) that he had long held such a view as the one Fechner developed in the work named. What he misses in modern as well as in ancient atomic theory is a sufficient ground for the absolute fixity with which the atoms are nevertheless to keep their inner constitution under every conceivable interaction. This question he does not believe he can solve without assuming qualitatively different atoms, an assumption which the atomism of antiquity denied on account of its antinominalistic tendency, and which Fechner has rejected only out of an unwarranted love of simplicity.} which physics
leads us to regard as in fact indissoluble must
be thought of as consisting of essentially different primal constituents,
between which there obtains such an elective affinity
that it is not outbid by any other. With these
% --- p. 301 ---
\setcounter{footnote}{0}%
\opage{301}primal constituents one can entirely disregard all spatial extension
and regard them as supersensible beings
which, from definite points in space, by their forces
govern a certain measure of extension without in the proper
sense filling it. Their mutual position
is conditioned by their interaction, and a spatial figure would
thereby be described just as surely and definitely as if
they continuously filled its interior. To these particular
points we must think of the attractive and repulsive
forces as attached. In these centers, from which the forces
are to spring, one must think of an inner life which
in the last instance conditions the peculiarity of the force.
We, who do not stand at the center of the world, cannot, it is true,
know how the mechanical interaction of the forces
is conditioned by the inner essence of the monads; but we are led
with necessity to dissolve the concept of matter and
see it as materiality, i.e. as the form of the revelation of a supersensible
real. The life of matter is thus a
double one, an outer and an inner. Extension divisible
to infinity is only a semblance. Extension cannot
be the predicate of a single being, but designates a relation
between many beings which are not themselves extended.
The indivisible singleness of each of these beings makes it
possible to assume in them a gathering together of the external
impressions in sensation and enjoyment. Only thereby does
each single being become a participant in the purpose of the whole
life of the world, which consists in the realization of goods
that can be felt and enjoyed by the soul. We need not,
with Leibnitz, ascribe to every monad a representation
of its place in the world or a striving to develop
its powers; but we nevertheless grant him that every
single being is a mirror of the universe, inasmuch as it,
as a distinctive center, takes up in its inner sensation
what the whole connection of the world
effects at this definite place. No part of what is
is now without life and soul; only a part of what
% --- p. 302 ---
\setcounter{footnote}{0}%
\opage{302}happens, namely the movements that put the state of the one
being in connection with that of the other, winds
like an external mechanism through the fullness of what is animated
and brings it occasions and incitements to the
changing unfolding of its inner life.

The dualism between soul and body has now vanished,
not by making the soul, in materialist fashion, into
a matter, but conversely by making matter into a substance akin to
the mental. The soul is only the monad
favored by its position, primus inter pares.
The development of spiritual existence in a monadic
element depends on the influence of favorable conditions,
which present themselves in the form of bodily organization.
Only the element that is placed at the most favorable point of such an organization
receives all impressions so
ordered, and can so act back on the body's
receptive parts, that it can develop into a soul proper.
The soul now stands in interaction not with matter
as matter, but acts as a supersensible real
first of all on the other supersensible reals, which in themselves
are homogeneous with the soul but for us present the semblance
of material existence, and in the same way
the bodily acts back on the soul. There is thus
an interaction between members of like kind. But
the soul does not fuse with the atoms that lie
at the basis of the bodily organization. The living
being is a community of many single beings, which are governed
by a common law grounded in their inner nature.

It was not specifically psychological but metaphysical
motives that led us to this spiritualistic and
monadological conception. But is not everything thereby dissolved
into an infinity of single beings that stand in connection with one another
only in an outward way? Instead of
the dualism between soul and body there now seems to have
come an unreconciled opposition between an infinite
multitude of elements. In other words, the connection
% --- p. 303 ---
\setcounter{footnote}{0}%
\opage{303}between the single beings still remains unexplained. We
began with a dark point in a bright net of relations;
now we seem to have arrived at a multitude of bright
points in a dark net. --- Neither the single beings
nor the mechanical connection is the Absolute.
All single beings are only inner parts, particular determinations of content,
of an infinite being. Just as the inner unity of our soul
is mirrored in various forms of spiritual activity, which cannot be reduced to one another, but
spring each by itself and yet in necessary connection with one another
from the nature of the soul, so the
infinite substance is related to the single beings. Only that
is which has its place in the rational connection of the eternal Idea.
Mechanical interaction is the form
by which the infinite secures to every element its
distinctive activity in relation to the others. Without the
Absolute, mechanism would be inexplicable; only through
it can the gulf be overcome that would eternally separate the
particular independent elements from one another. Just as
in the soul the mutual relations of the representations can
call forth distinctive states of feeling only by acting
back on the nature of the soul itself, so in the infinite
substance every influence of a single element
is at the same time an influence of the whole infinite, and
only thereby is a corresponding effect in another element
possible.

Thus the Idea of mechanism is demonstrated in the Absolute.
What for speculative idealism coincided
with pure finitude and externality, and therefore
in the last instance stood as a negative, a not-being,
is seen here as the necessary condition for the development of the
infinite. Whereas idealism set the Idea in opposition
to mechanism, Lotze maintains that whether or not
an Idea governs the thing, all being or
becoming is possible only by virtue of mechanical causes. Whether
the Idea's command is carried out can be known only through insight into
% --- p. 304 ---
% print: Enkeltvæserne (misprint: so printed)
\setcounter{footnote}{0}%
\opage{304}the real connection of nature. The two tendencies in
knowledge, causal explanation and ideal interpretation,
can therefore subsist side by side.
But causal or mechanical explanation, as
we have seen, when subjected to a metaphysical critique, leads
back to a multiplicity of spiritual elements as
the centers of the forces. The identity of the ideal and the
real which speculation maintained thus holds pointwise
for the particular elements of the world. Each of these
is ideal, because its essence consists in the distinctive
significance it has as a special determination of content
in the infinite; it is real, for reality is the
form of existence by which something presents itself as an independent
starting point for effects which it exerts
or suffers. Spirit is then no longer, as for speculation,
a hovering being of light, but is dissolved into an
innumerable multitude of sharply bounded, radiant points.
No single element proceeds from the absolute
ground without there proceeding from it at the same time a definite
multitude of others, which together with that one fill out
reality so as to be the complete revelation
of the Absolute. In the spiritual interior of the single beings
the Absolute has the open place where, without the necessity of the mechanical
course of the world being sublated, it can exercise its
influence and call forth such manifestations of life by
which a general progress becomes possible through
the whole natural and historical development. This development
is, since the single beings are the Absolute's determinations of content,
one with the self-development of the Absolute through
a series of world-ages. But this is a thought
which human knowledge cannot carry through.
Only because speculation regarded the earthly
forms of existence as the absolute ones could it believe it
had measured out the course of the infinite's life. This presupposition,
that the creative primal ground reveals itself only on
% --- p. 305 ---
\setcounter{footnote}{0}%
\opage{305}the narrow path of earthly nature, was bound to lead to such
dialectical idylls as the speculative constructions.

Yet this is still not the most essential objection
that must be raised against the speculative view of nature
and history. It is not enough to find for each creature
and each phenomenon an Idea whose realization
they serve. The question cannot be dismissed of how much
of this Idea's significant content is present
in the being itself as a motive for its activity or
an object of its enjoyment. Otherwise we learn only
what a phenomenon is or means for us, but not
what it is in itself. Thus in the consideration
of history it is not enough to enumerate the sequence
of the Ideas which each by itself express the character and general tendency
of an age; what matters is to
know how each age itself and the individual in
it felt, how much of the Idea which we regard
at a distance really lived in men's hearts. The thoughts,
feelings and actions of individuals have for the
speculative conception worth and significance only insofar as
they happen to agree with the particular moments of development
of the universal Idea. The course of history becomes the
great slaughter-bench on which all individual happiness, all
individual life, is sacrificed so that the universal
Idea of humanity can advance in its development. And yet it is a
self-contradiction that there should be purposes whose content
and realization did not become conscious to anyone. What
is to be a good must have its existence in
the living feeling of a spiritual being; everything that is only
fact, thing, property, relation or event
is only a means, not itself a good.

The speculative conception seeks, by classification, to
demonstrate a graded series in nature from the less perfect
to the more perfect. But there is no rational reason
why the imperfect should go on subsisting beside
the perfect. One could, on the other hand, conceive
% --- p. 306 ---
\setcounter{footnote}{0}%
\opage{306}that the subordinate stage attains a distinctive
completion which it would not reach if it passed over
as a member into higher organisms. Then the simultaneous
subsistence of the less perfect and the more perfect would bring about
a greater qualitative richness in existence. Many
physical processes which we cannot sense perhaps meet
receptive organs in lower animals. So
in particular many chemical, electrical and meteorological
processes. One must assume that everything that happens is also,
at one point or another, sensed or made inward
by a psychical substance.

The mental life in the particular material atoms is only
the work of the moment. The elements of the course of nature become spiritual
events in feeling, but no coherent whole is formed
from them, and they are not preserved in any
memory. This lawless activity of the psychical movements
becomes in the plant a definite, unchangeable dream,
a musical development without freedom and mobility.
Only in the animals is organic fixity so combined
with moving manifoldness that a mental life in the
narrower sense can arise. The difference between man
and animal does not consist in special forms or elements
that are added in man. One cannot
set up the human spirit as a being of its own in opposition
to the animal soul; it would be just as impossible to fix
a boundary between these two beings as to show
how they could become one in man. What is distinctive
of man, on the other hand, is that mental
activity rises to such an intensity that,
as comparing and combining, it can enter into a
free relation to the impressions immediately received. The animal
does indeed also sense a fixed connection between the phenomena,
but it does not apprehend the necessity of this connection;
it suffers under the law, but does not become
conscious of the law. This difference is not shaken by the possibility
that one may perhaps be able to establish man's
% --- p. 307 ---
\setcounter{footnote}{0}%
\opage{307}descent from extinct animal species; for this fact
will decide nothing with regard to the estimation
of the worth of the forms of existence now actually subsisting.

Knowledge of necessity is the characteristic mark of what is distinctively
human. The presentiment of a necessary
truth and connection which lies at the basis of
all human thinking and acting can indeed lead astray through
hasty conclusion, but in itself it is right, although
it does not directly point out the means to reach the
end it sets. The attempts to set forth coherently
the sum of the necessary truths usually
fail, because it is elements of different origin
and validity that are to be united. Some truths
hold abstractly, for all reality (the principles of identity and causality),
others rest only on our experience
(the notion of the indestructibility of matter), others again
are grounded in the aesthetic and ethical ideals of our heart
(the demand for a highest unity in existence).

The life of feeling plays a great role in the development
of the higher Ideas of reason, where the question of worth
and significance takes the place of the question of acting
causes. That is why it is so difficult to convert these
Ideas into general understanding. Science must be
content when it succeeds in showing that the
clear and irrefutable principles of the understanding are nothing other than
the forms for the ideal reason's own end: the connection of
reality into the unity of a living totality.
--- Further, only the intervention of feeling can explain the essence
of self-consciousness. For when one generally
contents oneself with determining self-consciousness as
the unity of the thinking and the thought, this is a
purely formal determination which does not explain the inwardness
and the infinite worth that lie in this inner
relation to oneself. Self-consciousness is to be apprehended only
as the interpretation of a self-feeling whose original life
is not immediately heightened by the development of our
% --- p. 308 ---
\setcounter{footnote}{0}%
\opage{308}knowledge. The concept of worth is inseparably bound to the concept
of feeling, which, far more distinctively than knowledge,
designates the true nature of the spirit. Feeling is not,
as is often supposed, unreceptive to qualitative differences
which depend on the worth of that which calls forth the feeling.
Of lesser worth is that which corresponds only to a momentary
and accidental state or an individual peculiarity;
of greater worth, what harmonizes with the general
and normal fundamental features of the organization by
whose help the spirit reaches its purpose; the highest would
be what satisfies the constant mood of an ideal heart,
in which every straying from the end of development
is impossible. But beyond this one cannot go:
the thought of an infinite worth which does not establish its
reality by calling forth pleasure contains a self-contradiction.
The moral rigorism which, in its zeal to
combat egoism, would exclude the feeling of pleasure from the motives
of action was therefore unable to give anything
but purely formal maxims; the obligating
dignity of the moral law it was unable to explain. The ethical
here points to the religious; moral obligations
follow from the religious world-view: only where
man sees himself as a fellow-worker in a supersensible
world-order does he acquire positive purposes which benefit
himself and lift him above his finite
and sensuous being. The essence of the religious consists in
the conscious recognition of the unconditional validity of the moral and the holy
as against all that is
or seems to be fact. In the positive religions
the specifically religious does not appear pure; it
is obscured by cosmological speculations, through whose fusion
with that kernel religious faith can
become an inhuman and fanatical power. In Christianity too
this obscuring has entered. The fundamental thought
of original Christianity, which must not
be confused with the faith in historical
% --- p. 309 ---
\setcounter{footnote}{0}%
\opage{309}facts that the Church demands, aimed precisely at raising man
above every given state and rooting out the worship
of blind, factual being. But the Church's
demand for faith in a sacred history stands in opposition to this,
especially when this faith, through uncertain interpretation of uncertain writings,
is shaped into a cosmological system which perhaps
even (in the doctrine of the Trinity) seeks a support in the triad
of fundamental concepts in which thinking knowledge
ends, a triad which, as will appear, is not a
proof of the strength of thought, but precisely of its weakness.

In history we meet the development of the human spirit and
the realization of its content. The investigation of the essence of
the spirit has established the untenability of the
materialistic as well as of the speculative-idealistic view.
If one wished to return to the conception
that sees history as the education of the human race,
this conception, which otherwise seems to come nearest the truth,
would necessarily demand a closer determination.
For how can one speak of education when the results
of development benefit generations that have not
helped to produce them? The human race
is, after all, not a fixed unity but a succession of generations.
The goods that are attained, moreover, hardly become objects
of consciousness and enjoyment before they pass over into
being an indifferent, habitual foundation for further-reaching
wishing and striving. Moreover, in the course of development
the content of culture becomes so manifold, splits into
so many branches, that no individual can survey, let alone
appropriate, the whole. Finally, on closer
consideration it appears that the goods of culture benefit only a small
number; beneath these favored ones we find throughout
the whole of historical development a proletariat, over
whose heads culture goes its way without essentially
intervening in its conditions. The countless advances in
particulars have as yet by no means gathered into a general
% --- p. 310 ---
\setcounter{footnote}{0}%
\opage{310}advance in the happiness of life; on the contrary, life seems
more and more to become a struggle for existence.

But we cannot on that account give up the thought of
a real progress and a real purpose in
historical development. Only thereby do we find the explanation
of all enthusiastic self-sacrifice in the service of the Idea.
If it stands firm that no fact and no form
has worth in itself, but that only that is a good which
is livingly appropriated by a spiritual being, then we must also
understand by the humanity whose essence is developed through history
not a sum of individuals to whom a
certain character of kind applies, but a living and real community
in which the human spirits, otherwise separated
by the differences of time, are joined together so as to be
for one another in such a way that each occupies his distinctive
position, and all things benefit all. Only
on this presupposition does the notion of the education
of the human race, like the speculative
conception of history as the development of the Idea, acquire its significance;
for only then does it become of any worth that the Idea
is developed. Not for the sake of individuals must this demand
be made, but for the sake of the connection of the world,
which would become meaningless without its fulfillment. Only
insight into what ought to be can open to us insight
into what is. In this idealism is right, that only such
and so much exists as has its necessary
place in the highest Idea. ---

In conclusion we gather, as Lotze himself does in the last
section of ``Mikrokosmus'', the motives given in what precedes
into a concluding world-view. --- Lotze,
like Herbart, starts from the given and infers back
from it. But already with the first inference he parts
from Herbart. The latter determines being as
absolute position, i.e. absolute, simple quality. But this
is a fallacy. Being means, on the contrary, that a content
is posited in certain definite relations. The distinctive
% --- p. 311 ---
\setcounter{footnote}{0}%
\opage{311}nature of every thing can be conceived only as a definite
member of the really possible content of the world. The interaction
between these particular members is explicable only if they are
parts of an infinite substance that embraces everything. For
us things come forward each in its place in an infinite
space; but according to the doctrine of local signs, the intuition of space
is only the form under which the soul converts
the impressions into its own language. Yet space is not a
merely subjective and accidental form. To the order in space
for our intuition corresponds the way in which things
are ordered in the intellectual connection of the infinite,
and this depends on the significance and worth of
the Idea each thing expresses. The case is similar
with time. This too has no immediate
objective significance. But what appears to us
as an occurrence in time is only the revelation
of the inner relation of conditioning between the particular
elements of the infinite. The point in time and the duration of the particular
depend on its intellectual significance, i.e. on
the significance of the Idea it expresses.

But when we thus trace everything existing back
to the Idea it expresses, and see its essence expressed
in it, we seem thereby to become entangled in several difficulties.
The thing cannot be mere Idea; if it is to have
reality, it must be an actual Idea, but the Idea can
act only by being thought. We ascribe unity to the thing;
but we know unity only from our own I; only
because we become conscious of ourselves are we one. We ascribe to the thing
a distinctive acting and suffering according to the relations in
which it is placed; but only that can in truth act
and suffer which can feel weal and woe.

These difficulties in the conception of the essence of things
have occasioned various attempts at solution. ---
One has left the real in things standing as an
obscure, uncomprehended kernel, and regarded our notion
of things in general as a subjective form of knowledge
to which reality need not correspond. But
% --- p. 312 ---
% print: Bevisthed (misprint: so printed)
\setcounter{footnote}{0}%
\opage{312}one is nevertheless always compelled to ascribe to things in themselves
a being and an interaction with us; and one
then moves in a circle, at once applying
and denying the application of our conceptual forms to
the essence of things. --- There then remains only either
to deny the reality of things altogether or to seek to develop
the concept of the thing in such a way that the difficulties
vanish. The first of these ways out now leads with necessity
to the second. For idealism's proposition that
only the spiritual is real must, if reality
is not to become a pure semblance, be reversed so that it
comes to read: everything real is spirit. This proposition
has shown itself in what precedes to contain the
ultimate explanation of what experience presented to us.
It does not volatilize the concept of things, as idealism does;
the thing acquires selfhood, being-for-itself, the most general
expression of the spiritual nature, which in self-conscious
I's reaches its highest stage, without therefore being absent where
existence is not apprehended with clear consciousness but
in the form of a more obscure feeling. The difference between
this conception, spiritualistic realism, and idealism
does not consist, as one might perhaps suppose,
in the former's ascribing to things a being outside
God, whereas the latter conceives them as wholly immanent in
the Absolute. For it rests on an unclear line of thought
to suppose that what is outside God should be more
real than what is in God. Indeed, what meaning can
be attached at all to the expression ``to
be outside God?'' The difference consists rather in this, that
idealism denies the selfhood of things, whereas spiritualistic
realism, which starts from the fact that in the spirit,
in the self, we have at least \emph{one} real whose nature we know,
believes it must ascribe this spiritual
nature to all reals in different degrees.

The earlier investigations led to the assumption of
an infinite substance, of which all
% --- p. 313 ---
% print: i det physiske (misprint: i doubled over the line turn (i|i))
\setcounter{footnote}{0}%
\opage{313}things are determinations of content. But yet this infinite, the idea of God
of pantheism, cannot be the concluding concept. We would then
end with J. G. Fichte in the notion of an infinite
world-order. But order cannot be thought without something
that orders. And the whole conception of the essence of the real
that follows from what precedes, that everything real must
be for itself, have a spiritual center, must surely
find its application here too. The objections to
applying the concept of personality to God lose their
weight, since we have shown that the essence of self-consciousness does not
rest merely on the opposition between an I and a not-I,
but is only the reflected expression of the self's
original essence. On closer consideration the matter rather
stands thus, that the concept of personality
can be applied only imperfectly to \emph{finite} beings,
whose inner originality is not absolute; it is an
ideal which can be thought of as completely realized
only in the absolute being.

The highest real, the Absolute, must be absolute
unity. But here the insufficiency of human thought
shows itself. It is unable to gather the various threads of knowledge
into a complete unity.
The earlier investigations have at every point shown
us a triad of forces which cannot be traced back
to unity: 1) the general and abstract (metaphysical
and mathematical-mechanical) laws; --- 2) the real
beings themselves, in which the forces are concentrated and from
which they act according to those laws; --- 3) the purposes, the
comprehensive plan according to which the life of the world develops.
What prevents thought from letting law and force
fuse with purpose into one being is, in
the last instance, the great unsolved problem that lies
in the reality of physical and moral evil. Only
this much must thought hold fast: that the ground of law and force
must lie in purpose itself; the highest real
% --- p. 314 ---
\setcounter{footnote}{0}%
\opage{314}unfolds itself in one movement which for us splits
into three forces. ---

It will be of interest to compare Lotze's
doctrine with those of other, earlier and contemporary philosophers
with whom it has points of contact. Among older
philosophers, Leibnitz, Kant and Herbart in particular come into consideration here,
among later ones Comte, Mill and Fechner\footnote[1]{We have heard above (p. 288) in what relation Lotze places his view to the elder Fichte's. His relation to the Schelling-Hegelian speculation has been brought out at several points in the exposition.}.

Leibnitz's doctrine of monads has at the very first glance
great similarity to Lotze's view. But apart from the
difference, already mentioned earlier, in the conception of the monad's
inner essence, which Lotze himself has brought out,
the decisive disagreement consists in this, that Lotze
identifies the monad with the atom, so that a real
interaction becomes possible, only one cannot conceive
matter as a continuum. Leibnitz, on the other hand, does not
immediately ascribe any physical significance to his monads;
a real interaction between them is
impossible, and the apparent interaction is explained
by a pre-established harmony between the self-unfolding of the monads.
Whereas Leibnitz is therefore an idealist, Lotze is
a realist. --- Thereby he acquires a point of contact with
Herbart. But not only does the latter expressly distinguish
his reals from the atoms of physics; he also regards them
as absolutely simple in the sense that no
inner development, only self-preservation against external influence,
becomes possible. Lotze's real is not unchangeable,
any more than Leibnitz's monad; Lotze is a realist,
but a spiritualistic realist. In psychology he differs
from Herbart in maintaining the originality of feeling and will,
whereas Herbart derives them from definite
relations between the representations. And whereas
% --- p. 315 ---
% print: fattelse af den phænomenale (misprint: Op lost at the line turn (exakt|fattelse))
\setcounter{footnote}{0}%
\opage{315}Herbart finally sets the theoretical and the practical
over against each other as independent and wholly separate,
Lotze finds the starting point of metaphysics in the ethical-religious,
and in what ought to be the ultimate explanation
of what is. --- As regards Kant, it must suffice here
to refer to the doctrine of space and time, of
the thing in and for itself, of the good. Indeed, even as regards atomic
theory Lotze, as already touched on,
has found a point of attachment, namely in an earlier treatise
of Kant's, in which Kant set up an atomistic theory
which he later abandoned in order, in his ``Metaphysische Anfangsgründe
der Naturwissenschaft'', to found the
dynamic conception which became dominant in speculative
philosophy of nature. Lotze's whole critical-skeptical, antidogmatic
cast of mind recalls Kant's Critiques.

With modern positivism in France and England
Lotze agrees (as in part Fechner does too)
in a double respect. Like positivism he starts
from the phenomenon, and he lays the weight on an exact
apprehension of the phenomenal connection. In this
part of his doctrine he is wholly a positivist. But he
goes further and seeks to solve the problem of the essence
and ground of phenomena. This problem Comte
and Mill regard as insoluble; but since neither of them is able
to hold fast to this skepticism, they answer it nonetheless,
partly in a materialistic, partly in an idealistic direction. --- From
Fechner Lotze differs, as already remarked earlier,
in conceiving the atoms as monads with an
inner life, an analogue of conscious spiritual life. So far
Fechner does not extend his hypothesis of animation,
which is why with him the atoms stand as the ultimate presupposition
of the analysis of experience, but without connection
with the rest of his view. In the conception of mental life
itself they are in a way antipodes, inasmuch as Fechner
sees the essence of consciousness in its gathering together the particular
sensations and representations (the synechological
% --- p. 316 ---
\setcounter{footnote}{0}%
\opage{316}conception), whereas Lotze lays the weight on the concentration
in a center which consciousness presupposes
(the monadological conception). The relation between them
can be compared with that between Spinoza and Leibnitz.

If finally we consider Lotze's doctrine in relation to
post-Hegelian philosophy as a whole, we find taken up in him
as essential elements what in the other
schools appears isolated and sole-ruling. The necessity
of a causal explanation by the mechanical interaction of elementary
forces he stresses as
strongly as materialism, and in a more scientific
way. No less than the negative criticism he opposes all
philosophical and religious dogmatism and does not close
his eyes to the psychological conditionedness of knowledge. And
finally, like speculative theism, he contests
the deification of the abstract universal and maintains the metaphysical
and philosophical-religious significance of the individual
and the personal. Not without reason, then, do we
end our exposition with this spirited thinker.

Lotze's eye for the distinctiveness of things and forces,
his sense for the full content of life
and his zeal to let it come into its own over against
a dead formalism and mechanism make the treatment
of the concrete problems, where thought still moves within
individual reality, the strong side of his philosophy.
His honesty and sharp critical vision, moreover, have the effect
that he would rather openly declare a problem
insoluble than reassure himself with an apparent answer
to it (cf. above p. 222). That is why he only opens to
his reader the view to the speculative heights, but does not lead
him up there. The scattered threads of thought, in his
conviction, human knowledge cannot
gather into one. The highest ethical Idea stands for him
as that which must also be the highest real.
But here he misses the definite attachment for law
and acting force. Yet he has in part himself
% --- p. 317 ---
\setcounter{footnote}{0}%
\opage{317}called forth this lack by his conception of the essence of thought
as essentially formally reflecting; thereby
thought does not come to find itself again in its content.
Lotze maintains that the ethical Ideas must be given to thought
from feeling and only thereafter can be developed into knowledge.
But if one does not start from a standing
opposition between thought and its content, then
in the development of knowledge it is just as much the content
that is determined by thought as the reverse. Likewise
it is one-sided when Lotze declares that what ought to
be opens to us insight into what is; the reverse
standpoint too is of the greatest importance, namely that
knowledge of what really is instructs us about
the essence of purpose. No one has practiced this principle
more than Lotze himself, who has strikingly established
how wrongly one conceives the Idea when one overlooks
the real conditions of its validity and reality.
Metaphysics, then, does not merely, as Lotze says, get its principle
in ethics, but ethics must conversely have its grounding
in metaphysics; or, more definitely expressed, the
proposition that metaphysics has its beginning in ethics
can mean nothing other than that whoever is to develop
the essential determinations of all that is must then also
demonstrate the universal foundation of the ethical. This
thus makes a demand on metaphysics, but does not
immediately settle how this demand is satisfied.

When Lotze declares that the Absolute, which in itself
is eternally one movement, splits for us into three forces,
it seems to follow of itself that one cannot
derive any one of these forces from any of the
others (e.g. law and force from purpose). They must
on the contrary be traced back as far as possible to their
common source, and our knowledge strives for this by
seeing them in their mutual connections. If they really
spring from one principle, then an inner interaction
between them must also be possible and be capable of being comprehended by us at least
% --- p. 318 ---
\setcounter{footnote}{0}%
\opage{318}approximately. Ethical-religious,
metaphysical and physical knowledge (knowledge
of purpose, of law and of force) thus do not
stand in opposition to one another, but gather
from different starting points into one center.

Although, according to Lotze, what is is to be explained by
what ought to be, his own doctrine has not arisen
in this way, but is essentially a hypothesis
that has sprung from the endeavor to explain
the given connection. Rightly has the deeper philosophical grounding
been missed here\footnote[1]{\emph{Weisse} in ``Zeitschr. für Philos. und philos. Kritik.'' vol. 47 p. 279 f.}. No more
than ethics does the philosophy of nature receive its metaphysical
foundation. This is also why Lotze does not overcome
the difficulty that meets all monadological
systems, namely to show how the independent
single beings can be united with the Absolute. The monads
appear now as completely independent,
now they are wholly absorbed in the Idea, by whose grace alone they
exist. The latter alternative must accord best
with Lotze's doctrine insofar as it maintains, after all, that what is
must be subordinated to what ought to be. That is why Lotze
will not enter into the question of individual
immortality either. But if this alternative is seized, what
then becomes of the atomic theory? --- When, further, the
Absolute is to be a personal self, the question must
arise whether it must not then be a monad, like
finite beings, and how this could be
united with the essence of the Absolute.

But Lotze has not promised us any answer to these
questions. And it must be emphasized that he, who essentially
renewed Leibnitz's view, has corrected it
at decisive points, partly by conceiving the monad
% --- p. 319 ---
\setcounter{footnote}{0}%
\opage{319}as atom, whereby he has secured for himself the exact and
real points of attachment, partly by maintaining from the very first
that the monads are only determinations of content in the
infinite substance.

\chaphere{CHAPTER SIX.}


It now remains to sum up in brief strokes
the essence and significance of post-Hegelian philosophy. The
exposition we have now concluded has in many ways
shown us the opposition to the thought of the preceding period,
while at the same time it became clear how the decisive
problems here, as always, pass over from the
one generation to the next. When we
look back on the phenomena we have come to know, the relation,
on closer consideration, stands thus: speculative
philosophy in its first magnificent development led to the view
of a lofty scientific ideal, which only successively,
through faithful and persevering working up of the
given reality, can receive its firm and unshakable
foundation. This work post-Hegelian philosophy has
begun, and from this spring both its polemic
against speculation and its sympathy with it.

Speculative philosophy was essentially dogmatic
and idealistic. It was dogmatic in that it rested on
the presupposition of the unity of thought and reality,
and all its endeavors aimed only at developing
everything contained in this original unity.
Of course this presupposition was not arbitrary; it
was in turn the result of the critical philosophy's penetrating
investigations of the essence of knowledge. But
% --- p. 320 ---
\setcounter{footnote}{0}%
\opage{320}this result was nevertheless raised above every finite and
subjective grounding; one believed one had reached, once for all,
the point from which an a priori and synthetic
development of the philosophical world-view
was possible. Yet it is clear that what one can reach by this route
can only be what one already has in the presupposition
from which one starts. Now since the presupposition
was indeed the unity of thought and reality, but
reality was here conceived only as pure abstract being,
that is to say as being thought, the world-view too had
to become an absolute idealism, which left real
existence standing as an unsolved problem.

That is why reality, existence, is also the watchword
with which the field is taken against speculation from all sides.
We have found it in all the
four groups of philosophical phenomena that we have
brought forward. One thus comes to the recognition
that the presupposition of speculation must be subjected to a
new testing. To this extent post-Hegelian
philosophy resumes Kant's labors and acquires a critical character
in opposition to the dogmatism of speculation. Yet
it does not thereby give up speculation's essential yield.
It does not fall back into criticism's sharp opposition
between thought and existence, but holds fast to the unity
as the ultimate presupposition, only this cannot
be made good immediately. Therefore it maintains
that in the face of concrete reality with its differences
special methods are necessary. Philosophy cannot
immediately speculate its way into the essence of reality,
but needs empirical science as an intermediate link.
What it can appropriate by this route it must
then itself give the special grounding by which a scientific
world-view becomes possible. The Idea loses
nothing of its luster because its activity proceeds according to
definite laws and under fixed conditions. This is the
realistic character of post-Hegelian philosophy.

% --- p. 321 ---
\setcounter{footnote}{0}%
\opage{321}But it is not enough to determine the essence of this philosophy
in relation to the preceding one. To gain
a fuller insight into what it has to say, we must
state precisely the essential results to which it has led.
It follows of itself that no scientific result
is concluded in the sense that it could now stand fixed
henceforth as dogma. What is meant here by results
can only be the general direction in which
philosophical thinking points in its striving
to win a concluding world-view.

The philosophical phenomena with which we have been occupied
divide at once, as we already indicated at the beginning
(above p. 16), into two larger groups. The one
of these has an essentially negative character. The critical
consciousness dissolved speculation by surrendering to
the power of formal consistency, which in the end led it
beyond all truth and reality. In the end here not
only all philosophy but also all science as a
whole had to show itself an illusion. Even Feuerbach, that Protean figure,
who only through the ideal richness of his own spirit
gained the strength to fight against the Idea, and who constantly
pointed forward to a positive view which he
nevertheless never reached, even he ends by saying: ``meine
Philosophie ist keine Philosophie''. But merely formal
consistency is untrue. It seizes a single moment
of existence in its isolation and then rejects with
sophistical dialectic every attempt to sublate this
isolation. The first step to truth is to see that all
reality is connection, plurality in unity. --- To
the negative group modern materialism must also be
assigned. It believes, to be sure, that it has a positive
principle, namely matter, stuff; but partly the essential thing
about materialism is its polemic against the reality of every
spiritual principle, partly it has proved impossible
to determine the essence of matter in a positive way, which is why
the materialists themselves in fact give up this principle in order to
% --- p. 322 ---
\setcounter{footnote}{0}%
\opage{322}seek a grounding by way of the postulate --- and thereby
they then come to appeal to ideal motives! ---

But although the critical and materialistic school
is essentially negative, it nevertheless has for all that a
great philosophical and cultural-historical significance. It is
a preparatory school where the dogmatic and idealistic illusions
are subjected to an ordeal by fire. Only faith in authority can
reject these schools out of hand. For philosophy
they are a necessary process of purification through which
strength is won for renewed striving toward a positive
world-view. Moreover, negative criticism
and materialism serve to break the dominion of supernaturalism and dualism
in the general consciousness. It
must stand fixed once for all that only that has reality
which can establish itself as true before testing criticism,
and which does not defy the conditions of factual reality.


If we now turn to the positive philosophical phenomena
which we have treated in the 4th and 5th chapters,
we shall sum up their essential significance by
considering them from three main points of view, namely first
with regard to the relation between speculation and experience,
next with regard to the relation between spirit
and matter, and finally with regard to the relation between
philosophy and religion.

The first relation has already been discussed in general terms,
in that we have shown the realistic character of
post-Hegelian philosophy. What we are to draw
attention to here is that experience acquires not merely
negative significance, as an authority which speculation must not
overlook, but that a positive unity and
fusion of the two is striven for. Our period begins with
the belief that, beyond the domain which the logical movement of thought
can measure out, one glimpses a full reality,
the appropriation of which is to form a supplement to purely
immanent speculation. Such was the school which for a
% --- p. 323 ---
\setcounter{footnote}{0}%
\opage{323}time was called the doctrine of freedom or positive philosophy,
and which we have treated under the name of speculative
theism. Already within this school we found
an endeavor to overcome the opposition between
the two spheres and set them in an inner relation to each other.
The views we have brought together under
the name of the Idealism of Reality carry this further by
establishing that experience itself points back to speculative
presuppositions. Here is the point where the insufficiency of French
and English realism shows itself.
Like materialism, it does not carry through its
own presuppositions, and therefore ends without principle or
skeptically. ``The Idealism of Reality, on the other hand, shows us
how precisely a full recognition of the rights of experience
and an exact apprehension of its results
lead, philosophically considered, beyond mere empiricism.

Speculative philosophy saw all reality in
the Idea. But it could not derive factual, concrete
reality from the essence of the Idea. The Idealism of Reality
now goes the reverse way and seeks to show that all reality
is in the last instance of an ideal nature. And it does
this without thereby giving up its exact foothold in the
real world. In a more modest form it nevertheless points
toward a view which in boldness and
grandeur does not yield to the speculative systems.
Here, then, the way is paved for the union which philosophy,
and in it the human spirit, constantly seeks, the union
between the ideal and reality. This union is to be
demonstrated in the very kernel of reality, just as speculation
demonstrated it in the Idea. That matter in the last instance
is spirit, philosophy has often declared; but only
now has it succeeded in beginning such a carrying through
of this proposition that the material and the real are not
volatilized but preserved in their essential significance.
Only now is materialism taken up as an element in the
philosophical world-view, and precisely thereby will
% --- p. 324 ---
\setcounter{footnote}{0}%
\opage{324}this world-view be able to carry through the idealism which it cannot
give up without giving up itself.

From this it follows in turn that philosophy, less than
ever, can feel any need to bow to a faith in authority
that denies the possibility that human
thought can by its own route reach knowledge of
the divine. Philosophy bows to every
fact in the spiritual and bodily domain,
but only because every fact, in its essence and its
ground, attests the same rational principle from which
scientific inquiry itself springs. We have
seen how close a connection there is between historical
and philosophical criticism in our time; both have
developed hand in hand. Philosophy has thereby
gained a reliable foundation for its conception of the
religious. It can now place itself in a free relation to the historical
religions, equally unbiased by antipathy
and by sympathy. And the specifically religious itself is not an
element foreign to philosophy. For philosophical
thought is no longer formal and abstract; it
nourishes itself on the very substance of reality and is
conscious of its inward relation to the psychological essence of man.
Therefore it can contain all real stirrings of life
and give them their firm and conscious form. Philosophy
is itself religious, in that it leads man to
knowledge of the essence of existence and of his relation
to its ideal powers.

The idea of God of speculative philosophy was pure
spirit, which in eternal self-unfolding develops everything that is.
It was essentially abstract pantheism. From various
points of view we have seen post-Hegelian philosophy
strive toward a more concrete, more living idea of God.
When reality becomes more than a semblance, a merely
negative form of the Absolute, then the Absolute cannot
be the pure, abstract Idea either. And when everything real is
a single thing, an individual, then this must also hold
% --- p. 325 ---
\setcounter{footnote}{0}%
\opage{325}for the highest real. Therefore
speculative theism seeks to show the divine spirit
as the absolute personality: as absolute, God is one
with the Idea and stands in an immanent relation to the world;
but within this immanence an opposition makes itself
felt, in that the divine thought gathers
itself in an inner center from which the life of the world proceeds,
and which is the highest model for finite personalities.
From the pantheistic side this line of thought is met
halfway by maintaining the essential significance of difference and
opposition. Whether one now wishes to call
such a view ``immanent theism'' or ``transcendent
pantheism'', this much is clear: that here, in general
strokes, a rational concept of God has been drawn which is at once
the starting point and the end point of all philosophical
and scientific inquiry and of all truly human
striving as a whole.



\end{document}
